\documentclass[11pt,twoside]{article}
% ---------------------------------------------------------------------------
% Packages
% ---------------------------------------------------------------------------
\usepackage[T1]{fontenc}
\usepackage[utf8]{inputenc}
\usepackage{lmodern}
\usepackage[english]{babel}
\usepackage[a4paper,margin=1.05in]{geometry}

\usepackage{amsmath,amssymb,amsthm,mathtools}
\usepackage{bm}
\usepackage{booktabs}
\usepackage{longtable}
\usepackage{array}
\usepackage{multirow}
\usepackage{graphicx}
\usepackage{tikz}
\usetikzlibrary{arrows.meta,positioning,calc,decorations.pathmorphing,patterns}
\usepackage{pgfplots}
\pgfplotsset{compat=1.18}
\usepackage{caption}
\usepackage{subcaption}
\usepackage{enumitem}
\usepackage{microtype}
\usepackage[hidelinks]{hyperref}
\usepackage{cleveref}
\usepackage{siunitx}
\sisetup{per-mode=symbol, range-phrase={ to }, range-units=single}

% ---------------------------------------------------------------------------
% Theorem environments
% ---------------------------------------------------------------------------
\theoremstyle{plain}
\newtheorem{theorem}{Theorem}[section]
\newtheorem{lemma}[theorem]{Lemma}
\newtheorem{proposition}[theorem]{Proposition}
\newtheorem{corollary}[theorem]{Corollary}

\theoremstyle{definition}
\newtheorem{definition}[theorem]{Definition}
\newtheorem{assumption}[theorem]{Assumption}
\newtheorem{model}[theorem]{Model}
\newtheorem{problem}[theorem]{Problem}
\newtheorem{example}[theorem]{Example}

\theoremstyle{remark}
\newtheorem{remark}[theorem]{Remark}
\newtheorem{observation}[theorem]{Observation}

\Crefname{assumption}{Assumption}{Assumptions}
\Crefname{model}{Model}{Models}

% ---------------------------------------------------------------------------
% Macros
% ---------------------------------------------------------------------------
\newcommand{\R}{\mathbb{R}}
\newcommand{\SOthree}{\mathrm{SO}(3)}
\newcommand{\sothree}{\mathfrak{so}(3)}
\newcommand{\SEthree}{\mathrm{SE}(3)}
\newcommand{\dd}{\,\mathrm{d}}
\newcommand{\dt}[1]{\dot{#1}}
\newcommand{\ddt}[1]{\ddot{#1}}
\newcommand{\vect}[1]{\bm{#1}}
\newcommand{\mat}[1]{\bm{#1}}
\newcommand{\transp}{^{\mathsf{T}}}
\newcommand{\hatmap}[1]{\widehat{#1}}
\newcommand{\veemap}[1]{{#1}^{\vee}}
\newcommand{\tr}{\operatorname{tr}}
\newcommand{\diag}{\operatorname{diag}}
\newcommand{\rank}{\operatorname{rank}}
\newcommand{\sgn}{\operatorname{sgn}}
\DeclareMathOperator*{\argmin}{arg\,min}
\DeclareMathOperator*{\argmax}{arg\,max}

% physical symbols used throughout
\newcommand{\Adisk}{A}                 % total actuator disk area
\newcommand{\vind}{v_i}                % induced velocity
\newcommand{\Pind}{P_{\mathrm{ind}}}
\newcommand{\Pprof}{P_{\mathrm{prof}}}
\newcommand{\Pideal}{P_{\mathrm{ideal}}}
\newcommand{\Pshaft}{P_{\mathrm{sh}}}
\newcommand{\Pelec}{P_{\mathrm{el}}}
\newcommand{\Paux}{P_{\mathrm{aux}}}
\newcommand{\mfix}{m_{\mathrm{fix}}}
\newcommand{\mstr}{m_{\mathrm{str}}}
\newcommand{\mprop}{m_{\mathrm{prop}}}
\newcommand{\mbat}{m_{\mathrm{bat}}}
\newcommand{\eb}{e_b}
\newcommand{\FM}{\mathrm{FM}}
\newcommand{\Ct}{C_T}
\newcommand{\Cdz}{C_{d_0}}
\newcommand{\Vtip}{V_{\mathrm{tip}}}
\newcommand{\Mzero}{M_0}
\newcommand{\Mmax}{M_0^{\max}}
\newcommand{\ez}{\vect{e}_3}

\newcommand{\code}[1]{\texttt{\small #1}}

% Generated table rows are read inside tabular environments with the
% primitive \input: LaTeX's wrapper opens a group there and breaks the row.
\makeatletter
\newcommand{\inputrows}[1]{\@@input #1\relax}
\makeatother

% A boxed "engineering conclusion" for statements that are results rather
% than mathematics.
\newenvironment{finding}
  {\par\medskip\noindent\begin{minipage}{\linewidth}\rule{\linewidth}{0.4pt}\par\smallskip\itshape}
  {\par\smallskip\rule{\linewidth}{0.4pt}\end{minipage}\par\medskip}

% Generated by paper/make_results.py. Do not edit by hand.
\newcommand{\rGross}{1.724}
\newcommand{\rGrossG}{1724}
\newcommand{\rBattery}{288}
\newcommand{\rPower}{106.0}
\newcommand{\rPaux}{16.0}
\newcommand{\rFixed}{546}
\newcommand{\rPanel}{280}
\newcommand{\rElectronics}{181}
\newcommand{\rGimbal}{85}
\newcommand{\rFrame}{400}
\newcommand{\rArms}{81}
\newcommand{\rGuards}{155}
\newcommand{\rProp}{489}
\newcommand{\rMotors}{244}
\newcommand{\rProps}{221}
\newcommand{\rFM}{0.725}
\newcommand{\rRpm}{1036}
\newcommand{\rRpmMax}{1390}
\newcommand{\rVtip}{26.7}
\newcommand{\rVtipMax}{35.8}
\newcommand{\rCtSigma}{0.120}
\newcommand{\rReynolds}{\num{102000}}
\newcommand{\rCdzero}{0.023}
\newcommand{\rChord}{82}
\newcommand{\rSolidity}{0.212}
\newcommand{\rSpan}{1.26}
\newcommand{\rReach}{0.63}
\newcommand{\rDiam}{492}
\newcommand{\rWake}{6.0}
\newcommand{\rTipGap}{0.586}
\newcommand{\rEndReserve}{30.0}
\newcommand{\rEndEmpty}{34.1}
\newcommand{\rPackCell}{18650}
\newcommand{\rPackS}{5}
\newcommand{\rPackWh}{63.0}
\newcommand{\rSplOne}{53.8}
\newcommand{\rSplEar}{56.6}
\newcommand{\rSplFar}{54.8}
\newcommand{\rCritDist}{0.91}
\newcommand{\rBpf}{35}
\newcommand{\rAlphaRoll}{56}
\newcommand{\rAlphaPitch}{60}
\newcommand{\rAlphaYaw}{4.5}
\newcommand{\rTauRoll}{3.66}
\newcommand{\rTauYaw}{0.518}
\newcommand{\rJxx}{0.0649}
\newcommand{\rJyy}{0.0612}
\newcommand{\rJzz}{0.1141}
\newcommand{\rKR}{41.4}
\newcommand{\rKw}{12.9}
\newcommand{\rKRsug}{44.4}
\newcommand{\rKRlag}{44.4}
\newcommand{\rKRauth}{65}
\newcommand{\rPpd}{117}
\newcommand{\rScreenDeg}{16.4}
\newcommand{\rCapArcmin}{5.8}
\newcommand{\rUiScale}{278}
\newcommand{\rUiScalePref}{347}
\newcommand{\rLogicalW}{692}
\newcommand{\rLogicalH}{389}
\newcommand{\rCapArcminNear}{9.9}
\newcommand{\rLogicalWNear}{1186}
\newcommand{\rWallTwenty}{6.37}
\newcommand{\rMzeroTwenty}{0.57}
\newcommand{\rWallUse}{9}
\newcommand{\rZrotorAbove}{1.37}
\newcommand{\rZrotorBelow}{1.63}
\newcommand{\rGrossBelow}{1.720}
\newcommand{\rStandoff}{1.20}
\newcommand{\rEyeHeight}{1.50}
\newcommand{\rKappaMin}{1.000}
\newcommand{\rKappaMax}{1.055}
\newcommand{\rTurnClear}{0.61}
\newcommand{\rTurnEnv}{0.52}
\newcommand{\rTurnTrack}{64}
\newcommand{\rTurnLag}{1076}
\newcommand{\rTurnSat}{0.2}
\newcommand{\rTrackMargin}{100}
\newcommand{\rMarginAll}{all}
\newcommand{\rTrackWorst}{64}
\newcommand{\rClearWorst}{0.61}
\newcommand{\rBreezeLag}{198}
\newcommand{\rBreezeTilt}{0.65}
\newcommand{\rBreezeLean}{14}
\newcommand{\rTwindow}{16}
\newcommand{\rSettle}{3}
\newcommand{\rDt}{4}
\newcommand{\rSeed}{12345}
\newcommand{\rCoupTurnBat}{22}
\newcommand{\rCoupTurnGross}{28}
\newcommand{\rCoupTurnKappa}{1.055}
\newcommand{\rCoupTurnUp}{1.6}
\newcommand{\rCoupStandBat}{0.1}
\newcommand{\rCoupOrdLo}{1}
\newcommand{\rCoupOrdHi}{17}
\newcommand{\rCoupJlo}{0.0612}
\newcommand{\rCoupJhi}{0.0616}
\newcommand{\rCoupItLo}{3}
\newcommand{\rCoupItHi}{3}
\newcommand{\rCoupSimLo}{5}
\newcommand{\rCoupSimHi}{5}
\newcommand{\rCoupAllFeasible}{all}
\newcommand{\rEndTurn}{28.4}
\newcommand{\rCldNearKR}{41.4}
\newcommand{\rCldNearSat}{1.0}
\newcommand{\rCldNearClear}{0.56}
\newcommand{\rCldNearKp}{4.17}
\newcommand{\rCldNearInside}{47}
\newcommand{\rCldMidKR}{41.4}
\newcommand{\rCldMidSat}{0.3}
\newcommand{\rCldMidClear}{0.55}
\newcommand{\rCldMidKp}{4.17}
\newcommand{\rCldMidInside}{22}
\newcommand{\rCldFarKR}{41.4}
\newcommand{\rCldFarSat}{0.2}
\newcommand{\rCldFarClear}{0.61}
\newcommand{\rCldFarKp}{4.17}
\newcommand{\rCldFarInside}{0}
\newcommand{\rClbShortKR}{41.4}
\newcommand{\rClbShortSat}{6.5}
\newcommand{\rClbShortClear}{0.58}
\newcommand{\rClbShortKp}{4.17}
\newcommand{\rClbShortInside}{35}
\newcommand{\rClbMidKR}{26.0}
\newcommand{\rClbMidSat}{4.1}
\newcommand{\rClbMidClear}{0.55}
\newcommand{\rClbMidKp}{2.89}
\newcommand{\rClbMidInside}{24}
\newcommand{\rClbLongKR}{16.3}
\newcommand{\rClbLongSat}{8.3}
\newcommand{\rClbLongClear}{0.45}
\newcommand{\rClbLongKp}{1.81}
\newcommand{\rClbLongInside}{17}
\newcommand{\rBoomJratio}{1.72}
\newcommand{\rBoomAlphaRatio}{0.38}
\newcommand{\rBoomGross}{89}
\newcommand{\rBoomEar}{59.8}
\newcommand{\rStandEar}{56.6}
\newcommand{\rBoomTurnClear}{0.55}
\newcommand{\rNearTurnClear}{0.56}
\newcommand{\rOptSpl}{58}
\newcommand{\rOptSpan}{1.45}
\newcommand{\rOptWash}{7.5}
\newcommand{\rOptMass}{6.0}
\newcommand{\rHexGross}{1.935}
\newcommand{\rQuadGross}{1.747}
\newcommand{\rHexPower}{128.7}
\newcommand{\rQuadPower}{106.6}
\newcommand{\rHexEar}{57.6}
\newcommand{\rQuadEar}{56.5}
\newcommand{\rHexSpan}{1.23}
\newcommand{\rQuadSpan}{1.27}
\newcommand{\rHexPitch}{52}
\newcommand{\rQuadPitch}{59}
\newcommand{\rHexRoll}{43}
\newcommand{\rQuadRoll}{56}
\newcommand{\rHexYaw}{3.3}
\newcommand{\rQuadYaw}{4.4}
\newcommand{\rHexRank}{3}
\newcommand{\rHexD}{384}
\newcommand{\rQuadD}{498}
\newcommand{\rTethGross}{1.617}
\newcommand{\rTethPower}{97.9}
\newcommand{\rTethEar}{55.5}
\newcommand{\rTethArea}{0.23}
\newcommand{\rTethCurrent}{2.28}
\newcommand{\rTethMass}{37}
\newcommand{\rTethReserve}{184}
\newcommand{\rTethResLimit}{power}
\newcommand{\rTethLen}{5}
\newcommand{\rTethV}{48}
\newcommand{\rTethFeasible}{meets}
\newcommand{\rTethVbus}{8.4}
\newcommand{\rTethVdel}{45.6}
\newcommand{\rTethSizedBy}{ampacity}
\newcommand{\rTethDensity}{10.0}
\newcommand{\rTethJmax}{10}
\newcommand{\rTethAreaMin}{0.13}
\newcommand{\rTethGaugeLimited}{yes}
\newcommand{\rSubDP}{3.4e-09}
\newcommand{\rSubDTrack}{4.3e-05}
\newcommand{\rSubDClear}{3.4e-06}
\newcommand{\rCtlDP}{0.08}
\newcommand{\rCtlDTrack}{21.1}
\newcommand{\rCtlDClear}{0.03}
\newcommand{\rCtlTrackHalf}{85}
\newcommand{\rCtlTrackBase}{64}
\newcommand{\rNfSpanG}{22}
\newcommand{\rNfSpanDb}{0.22}
\newcommand{\rTlimOne}{5400}
\newcommand{\rTwallOne}{5400}
\newcommand{\rPwallAnchor}{2.83}
\newcommand{\rPwallTf}{3600}


\title{\bfseries The Flying Screen:\\
Coupled Sizing, Flight Dynamics and Control\\
of a Human-Adjacent Hovering Display}

\author{Derivation, computational treatment and design study}
\date{}

\begin{document}
\maketitle

\begin{abstract}
\noindent
A display that hovers beside a standing person is not an aircraft-sizing
problem with a payload attached. It is a coupled problem: the mass of the
machine determines the power it needs, the power determines the mass of the
battery, the battery is part of the mass, and the trajectory the controller
flies determines the power in the first place. This paper derives the whole
system and solves it. The derivation combines analytical models (momentum and
blade-element theory, thin-wall beam theory, rigid-body dynamics), empirical
fits where theory runs out (motor mass, the acoustic anchor), numerical
optimisation and simulation, and says which is which.

The paper has three parts. Part~I develops the algebraic layer. From the
actuator-disk control volume we
obtain the induced velocity and ideal hover power, add a blade-element profile
term to produce rather than assume a figure of merit, and close the mass
budget. The closure reduces to $\alpha m - \beta m^{q} = M_0$ with
$q = (3-p)/2$, where $p$ is the exponent in a disk-area scaling law
$A \sim m^{p}$. \Cref{thm:trichotomy} shows that for $q>1$ this equation admits
exactly two, one, or zero positive solutions according to the sign of
$M_0^{\max} - M_0$, that the transition is a saddle-node bifurcation, and that
the lighter root is precisely the attracting fixed point of the naive sizing
iteration. \Cref{thm:sensitivity} gives the logarithmic sensitivities of
$M_0^{\max}$ in closed form; endurance enters with exponent $-2$ and specific
energy, figure of merit and electrical efficiency each with $+2$.

Part~II derives the flight dynamics on $\SEthree$: Newton-Euler translation
and rotation with the full inertia tensor of one shared geometry, rotor thrust
and reaction torque through an allocation map, first-order rotor-speed
dynamics, a battery energy state, an orientation-dependent flat-plate model for
the screen that couples attitude to translation and back, and a gimbal. It then
builds the cascaded controller: a geometric attitude law with angular-rate
feedforward from differential flatness, a gain heuristic sized from the
machine's own control authority, which is computed by linear programming, and a
reference governor whose keep-out constraint moves with the user.

Part~III states the co-design problem as a differential-algebraic system with a
fixed-point constraint, gives conditions for the inner map to be well posed and
for the simulation-in-the-loop battery iteration to contract, describes the
engine and its numerical methods, reports verification against the closed
forms and calibration against measured hardware, and presents the resulting
design.

Four findings invert the naive expectation. First, for an indoor
human-adjacent machine the binding constraint is acoustic, not energetic.
Second, a declared sound-power level is not a pressure at one metre, and the
eight-decibel confusion between them, compounded by an uncalibrated anchor,
was worth fourteen decibels of error. Third, the standoff a user specifies is
measured to the display, while the rotors reach most of a metre back toward
them; a safety criterion evaluated at the vehicle centre is wrong by the whole
reach of the machine. Fourth, the standoff that makes the blades safe makes
ordinary text too small: at \SI{1.2}{\metre} a \num{1920}-pixel laptop panel
needs about \rUiScale\,\% interface scaling to meet a \SI{16}{\arcminute}
character height, which leaves a desktop of about
$\rLogicalW\times\rLogicalH$ logical pixels. The resolution is there; the
usable workspace is that of a tablet.
\end{abstract}

\tableofcontents
\clearpage

\section{Introduction}
\label{sec:intro}

\subsection{The question}

Consider a display that hovers beside a person while they work standing, and
follows them when they move. The engineering question is not whether a
multirotor can carry a screen; multirotors carry cameras of comparable mass
routinely. The question is whether the coupled system closes, and if it does,
whether the machine it produces is one a person would tolerate at arm's length
for half an hour.

The coupling is the whole difficulty. Write it as a cycle:
\begin{equation}
  m \;\longrightarrow\; T \;\longrightarrow\; P \;\longrightarrow\; E
    \;\longrightarrow\; \mbat \;\longrightarrow\; m ,
  \label{eq:cycle}
\end{equation}
gross mass sets the thrust required to hover, thrust sets the power through
momentum theory, power integrated over the mission sets the energy, energy
divided by specific energy sets the battery mass, and the battery is part of
the gross mass. A second cycle runs through the dynamics,
\begin{equation}
  m, J \;\longrightarrow\; \text{plant} \;\longrightarrow\; u(t)
    \;\longrightarrow\; P(t) \;\longrightarrow\; E \;\longrightarrow\; \mbat
    \;\longrightarrow\; m, J ,
  \label{eq:cycle2}
\end{equation}
because the power actually consumed depends on what the controller does, which
depends on the inertia, which depends on the mass distribution, which depends
on the battery.

Neither cycle is a mere iteration to be carried out until convergence. The
first has, under a fixed rotor geometry, a genuine bifurcation: past a
critical payload there is no fixed point at all. That is the mathematical
content of \cref{sec:fold}, and it is the reason a paper is warranted rather
than a spreadsheet.

\subsection{What this paper does}

We construct the problem in three layers and then solve them together.

\paragraph{Layer one: algebra.} Momentum theory over an actuator disk
(\cref{sec:momentum}) gives hover power as a function of thrust and disk area.
A blade-element profile term (\cref{sec:blade}) turns the figure of merit from
an assumption into a consequence of blade geometry. Substituting into a mass
budget (\cref{sec:closure}) produces a single scalar equation whose solvability
is the existence of the design. \Cref{sec:fold} analyses that equation
completely: a trichotomy theorem, the closed form of the critical payload, the
identification of the transition as a saddle-node bifurcation, and a proof that
the naive sizing iteration converges precisely to the physically meaningful
root and diverges from the other.

\paragraph{Layer two: differential equations.} Once the machine exists its mass
is constant and the flight dynamics begin (\cref{sec:dynamics}). We derive the
Newton-Euler equations on $\SEthree$, the rotor wrench and its allocation map,
first-order rotor dynamics, the battery energy state, and an
orientation-dependent aerodynamic model for the screen. The screen is what
makes this vehicle different from a camera drone: it is a large flat plate
rigidly coupled to the airframe through a gimbal, so attitude modulates
translational drag and drag applied off the centre of mass feeds back as a
pitching moment.

\paragraph{Layer three: control.} \Cref{sec:control} constructs the cascade
from position to velocity to attitude to body rate to rotor speed, states the
exponential stability result of Lee, Leok and McClamroch for the geometric
attitude law on $\SOthree$ together with the ways the implemented law departs
from it, and treats control allocation under saturation. Two items there are
specific to this vehicle. The control authority of each axis is computed as a
linear program over the rotor speed limits, and the attitude gains are bounded
by it, which matters because a rotor optimised for quietness is a rotor with
very little reaction torque and therefore very little yaw authority.
\Cref{prop:governor} states what the reference governor with its moving
keep-out constraint guarantees and what it does not; rate limiting alone,
introduced to stop the display being whipped around the user's head, will on a
curved path steer the corner-cutting reference through the user.

\paragraph{Coupling and computation.} \Cref{sec:codesign} states the full
problem as a differential-algebraic system with a fixed-point constraint and
gives conditions for the inner mass map to be well posed and for the battery
iteration to contract. \Cref{sec:engine,sec:numerics} describe the engine that
solves it and the numerical methods: a bounded root search for the closure
that reports its own failure, secant-accelerated fixed-point iteration for the
battery loop with the simulator inside it and a verification flight at the
end, fixed-step Runge-Kutta for the trajectory, and differential evolution
over the design vector.

\subsection{Contributions}

The mathematical contributions are \cref{thm:trichotomy,thm:bifurcation,%
thm:iteration,thm:sensitivity,thm:pfold,prop:twall,prop:inner}: a complete
description of the sizing equation including its bifurcation structure and the
convergence behaviour of the natural iteration; exact logarithmic sensitivities
of the critical payload; the condition on the disk-area exponent under which
the fold exists at all, and the separate statement that an endurance ceiling
survives the fold's disappearance at the constant-disk-loading boundary; and
the slope condition under which the component mass closure is uniquely
solvable. The control results (\cref{prop:governor}, the authority
computation and the gain heuristic) are stated with their conditions and are
verified by simulation, not proved.

The modelling contributions are the calibration of the acoustic and motor
submodels against measured hardware (\cref{sec:verification}), and the
near-field corrections of \cref{sec:nearfield}, which close a stated gap and,
more usefully, bound its size.

The engineering contribution is the design study of \cref{sec:results}, whose
main conclusion is that for this class of machine the binding constraint is
acoustic rather than energetic, and that stating a noise figure without stating
whether it is a sound power or a pressure, and whether it is free-field or in a
room, is worth more decibels than any design decision available.

\subsection{What this paper does not do}

No hardware was built. Every number reported is the output of a model whose
assumptions are stated where they are introduced, and \cref{sec:limitations}
collects the places where those assumptions are known to be inadequate. In
particular the rotor aerodynamics are momentum theory with a single profile
term and three near-field corrections; there is no blade-element momentum
theory, no wake model, and no treatment of the recirculating flow that a
machine of this disk loading establishes in a closed room within a minute of
taking off. The acoustic model is a broadband scaling law fitted to four
aircraft, coupled to a single-coefficient Sabine room; tonal content is
computed but not summed into the level.

\section{Notation, frames and standing assumptions}
\label{sec:notation}

\subsection{Frames}

Let $\mathcal{I}$ be an inertial frame with orthonormal basis
$\{\vect{e}_1,\vect{e}_2,\vect{e}_3\}$, $\vect{e}_3$ pointing vertically
upward, so that gravity acts along $-\vect{e}_3$. Let $\mathcal{B}$ be a frame
fixed in the airframe with origin at its centre of mass, $\vect{b}_3$ along the
nominal thrust axis and $\vect{b}_1$ pointing toward the user.

The attitude is the rotation $R \in \SOthree$ mapping body coordinates to
inertial coordinates: a vector with body coordinates $\vect{u}$ has inertial
coordinates $R\vect{u}$. Position and velocity of the centre of mass in
$\mathcal{I}$ are $\vect{r},\vect{v}\in\R^3$, and the angular velocity of
$\mathcal{B}$ relative to $\mathcal{I}$, expressed in $\mathcal{B}$, is
$\vect{\omega}\in\R^3$.

\subsection{Algebraic preliminaries}

The hat map $\hatmap{\cdot}:\R^3\to\sothree$ is the isomorphism
\begin{equation}
  \hatmap{\vect{u}}
  = \begin{bmatrix} 0 & -u_3 & u_2\\ u_3 & 0 & -u_1\\ -u_2 & u_1 & 0
    \end{bmatrix},
  \qquad
  \hatmap{\vect{u}}\vect{w} = \vect{u}\times\vect{w}
  \quad\forall\, \vect{u},\vect{w}\in\R^3,
  \label{eq:hat}
\end{equation}
with inverse the vee map $\veemap{(\cdot)}:\sothree\to\R^3$. We use repeatedly
the identities
\begin{equation}
  \hatmap{(R\vect{u})} = R\,\hatmap{\vect{u}}\,R\transp,
  \qquad
  \tr(\hatmap{\vect{u}}\transp \hatmap{\vect{w}}) = 2\,\vect{u}\transp\vect{w},
  \qquad R\in\SOthree .
  \label{eq:hatid}
\end{equation}

A unit quaternion $q = (q_0,\vect{q}_v)\in\mathbb{S}^3$ represents the same
rotation through
\begin{equation}
  R(q) = (q_0^2 - \vect{q}_v\transp\vect{q}_v)\,I
       + 2\,\vect{q}_v\vect{q}_v\transp
       + 2q_0\,\hatmap{\vect{q}_v},
  \label{eq:quat2rot}
\end{equation}
a two-to-one covering, $R(q)=R(-q)$. We integrate the quaternion rather than
$R$ because \cref{eq:quat2rot} carries no coordinate singularity, unlike any
three-parameter attitude representation.

\subsection{Symbols}

\Cref{tab:symbols} in \cref{app:nomenclature} lists every symbol. The ones used
constantly are: $m$ gross mass, $g$ gravitational acceleration, $\rho$ air
density, $N$ rotor count, $D$ rotor diameter, $A = N\pi D^2/4$ total actuator
disk area, $T$ total thrust, $\vind$ induced velocity, $\FM$ figure of merit,
$\eta$ electrical efficiency, $\eb$ usable battery specific energy, $t_f$
mission duration, $\lambda$ thrust margin, $J$ inertia tensor.

\subsection{Standing assumptions}

The following hold throughout unless explicitly suspended. They are stated
here so that later results can be read as conditional on them rather than as
claims about reality.

\begin{assumption}[Incompressible, steady, inviscid outer flow]
\label{as:flow}
The air is incompressible with constant density $\rho$, the flow through the
rotor disk is steady in the mean, and outside thin boundary layers on the
blades it is inviscid. Tip Mach numbers stay below $0.3$, which the designs of
\cref{sec:results} satisfy by a wide margin.
\end{assumption}

\begin{assumption}[Rigid airframe]
\label{as:rigid}
The airframe is a rigid body. \Cref{sec:components} sizes the arms so that the
first structural bending mode lies at least an order of magnitude above the
attitude-loop bandwidth, which is the condition under which this assumption is
self-consistent, and the engine reports the ratio so that it can be checked
rather than assumed.
\end{assumption}

\begin{assumption}[Constant mass]
\label{as:mass}
$\dot m = 0$ in flight. The battery loses energy, not mass. The relativistic
mass defect $\dot m = -P/c^2 \approx \SI{-1e-15}{\kilo\gram\per\second}$ at
$P=\SI{100}{\watt}$ is ignored.
\end{assumption}

\begin{assumption}[Uniform disk loading]
\label{as:uniform}
Within the momentum-theory layer each rotor is modelled as an actuator disk
carrying uniform loading. Non-uniformity and tip loss are absorbed into an
empirical induced-power factor $\kappa>1$ (\cref{sec:blade}).
\end{assumption}

\begin{assumption}[Rotors share thrust equally in trim]
\label{as:share}
In hover each of the $N$ rotors carries $mg/N$. Differential thrust for
attitude control is a perturbation about this trim.
\end{assumption}

\begin{assumption}[Quasi-steady rotor aerodynamics]
\label{as:quasisteady}
Rotor thrust and torque respond instantaneously to rotor speed; all actuator
lag is placed in the rotor-speed dynamics of \cref{sec:dynamics}. Inflow
dynamics, whose time constant is of order $R/\vind \approx \SI{0.1}{\second}$,
are not modelled.
\end{assumption}

\begin{assumption}[Perfect state estimation]
\label{as:estimation}
The controller has exact knowledge of $\vect{r},\vect{v},R,\vect{\omega}$ and
of the human's position. This is the assumption most obviously false in
practice and is revisited in \cref{sec:limitations}.
\end{assumption}

\begin{remark}
\Cref{as:flow,as:uniform,as:quasisteady} belong to the sizing layer;
\cref{as:rigid,as:mass,as:share,as:estimation} to the dynamic layer. The
separation matters: the algebraic results of \cref{sec:fold,sec:scaling} do not
depend on the dynamic assumptions at all, and remain valid statements about the
closure equation whatever one believes about the controller.
\end{remark}


\part{The Algebraic Layer: Existence of a Design}
\section{Momentum theory over an actuator disk}
\label{sec:momentum}

Everything downstream rests on one result: the power required to hold a mass in
the air, as a function of how much air the machine can reach. We derive it from
a control volume rather than quote it, because the exponents that appear are
the entire argument of \cref{sec:fold} and it matters that they are earned.

\subsection{The model}

\begin{model}[Rankine-Froude actuator disk]
\label{mod:disk}
The rotor is replaced by a permeable disk of area $A$ across which the static
pressure jumps by $\Delta p$ while the velocity is continuous. The flow is
axisymmetric, the fluid is at rest far upstream, and the slipstream is a
well-defined streamtube. Rotation imparted to the wake, viscous losses and the
finite number of blades are all excluded; they return in \cref{sec:blade}.
\end{model}

This is the classical model of Rankine and W.~Froude, given in the form used
here by Glauert~\cite{glauert1935} and, in the rotorcraft setting, by
Leishman~\cite[Ch.~2]{leishman2006} and Johnson~\cite[Ch.~2]{johnson1980}.

Let $\vind$ denote the axial velocity induced \emph{at the disk} and $w$ the
axial velocity in the far wake, both measured relative to the undisturbed
fluid. Let $A_\infty$ be the cross-sectional area of the far wake.
\Cref{fig:streamtube} shows the control volume and the four stations used
below.

\begin{figure}[tbp]
\centering
% Actuator-disk streamtube in hover, flow downward. The boundary radius is
% drawn from continuity, r(x)^2 u(x) = R^2 v_i, with the axial velocity of
% the vortex-cylinder solution on the axis, u/v_i = 1 + xi/sqrt(1+xi^2),
% xi = x/R positive downstream. One-dimensional, but the shape is earned.
\begin{tikzpicture}[scale=1.25,
  declare function={u(\s) = 1 + \s/sqrt(1+\s*\s);
                    trad(\s) = 1/sqrt(u(\s));}]
  \def\xa{-1.05} \def\xb{3.6}
  % streamtube boundaries (y = -xi, flow downward)
  \fill[cA!6] plot[domain=\xa:\xb, samples=80] ({ trad(\x)}, {-\x})
      -- plot[domain=\xb:\xa, samples=80] ({-trad(\x)}, {-\x}) -- cycle;
  \draw[cA, line width=0.9pt] plot[domain=\xa:\xb, samples=80] ({ trad(\x)}, {-\x});
  \draw[cA, line width=0.9pt] plot[domain=\xa:\xb, samples=80] ({-trad(\x)}, {-\x});
  % interior streamlines, scaled boundaries
  \foreach \f in {0.35, 0.7}{
    \draw[cA!55, line width=0.4pt] plot[domain=\xa:\xb, samples=60] ({ \f*trad(\x)}, {-\x});
    \draw[cA!55, line width=0.4pt] plot[domain=\xa:\xb, samples=60] ({-\f*trad(\x)}, {-\x});
  }
  % axis
  \draw[cG, dash dot, line width=0.4pt] (0,0.95) -- (0,-3.75);
  % actuator disk
  \draw[line width=2.4pt, black] (-1,0) -- (1,0);
  \node[lbl, anchor=west] at (1.28,0.1) {disk, area $A$};
  \draw[vec, cD, line width=1.2pt] (0,0.05) -- (0,0.72) node[right, lbl] {$T$};
  % control volume
  \draw[construct] (-2.45,1.0) rectangle (2.45,-3.72);
  \node[lbl, cG, anchor=north west] at (-2.42,-3.3) {control volume};
  % stations
  \foreach \y/\n in {0.92/0, 0.13/1, -0.13/2, -3.55/3}
    \node[lbl, anchor=east] at (-2.55,\y) {\n};
  \draw[cG, line width=0.3pt] (-2.45,0.13) -- (-1.05,0.13);
  \draw[cG, line width=0.3pt] (-2.45,-0.13) -- (-1.05,-0.13);
  % velocity arrows: 0 far upstream, v_i at the disk, w = 2 v_i far downstream
  \node[lbl, anchor=west, align=left] at (3.3,0.92)
      {0: $p_\infty$, at rest};
  \node[lbl, anchor=west, align=left] at (3.3,0.28)
      {1: $p_1 = p_\infty - \tfrac12\rho\vind^2$, speed $\vind$};
  \node[lbl, anchor=west, align=left] at (3.3,-0.28)
      {2: $p_2 = p_1 + \Delta p$, speed $\vind$};
  \node[lbl, anchor=west, align=left] at (3.3,-3.55)
      {3: $p_\infty$, speed $w = 2\vind$, area $A_\infty = A/2$};
  \foreach \xx in {-0.55, 0, 0.55}
    \draw[vec, cA, line width=0.8pt] (\xx,-0.35) -- ++(0,-0.5);
  \foreach \xx in {-0.35, 0, 0.35}
    \draw[vec, cA, line width=0.8pt] (\xx,-2.75) -- ++(0,-0.75);
  \node[lbl, cA, anchor=west] at (0.62,-0.62) {$\vind$};
  \node[lbl, cA, anchor=west] at (0.8,-3.1) {$w$};
\end{tikzpicture}

\caption{The actuator disk in hover. Air at rest far upstream (station 0) is
drawn into a contracting streamtube, crosses the disk (stations 1 and 2) at the
induced velocity $\vind$ with a jump $\Delta p$ in static pressure, and leaves
in a wake of speed $w$ and area $A_\infty$ at ambient pressure (station 3). The
momentum theorem is applied to the dashed control volume. The tube is drawn
from continuity, $r^2u = R^2\vind$, with the axial velocity $u$ of the
vortex-cylinder solution on the axis; it contracts to $A_\infty = A/2$, as
\cref{thm:momentum} requires.}
\label{fig:streamtube}
\end{figure}

\subsection{The three conservation laws}

\paragraph{Mass.} The streamtube carries a constant mass flow. Evaluating at
the disk,
\begin{equation}
  \dot m_{\mathrm{air}} = \rho A \vind = \rho A_\infty w .
  \label{eq:continuity}
\end{equation}

\paragraph{Momentum.} Applying the momentum theorem to a control volume
enclosing the streamtube between a station far upstream, where the fluid is at
rest and the pressure is $p_\infty$, and a station far downstream, where the
pressure has returned to $p_\infty$ and the velocity is $w$, the only external
axial force is that exerted by the disk on the fluid, equal and opposite to the
thrust:
\begin{equation}
  T = \dot m_{\mathrm{air}}\,(w - 0) = \rho A \vind w .
  \label{eq:momentum}
\end{equation}

\paragraph{Energy.} Bernoulli's equation holds along streamlines on each side
of the disk separately, but not across it, because the disk does work on the
fluid. Upstream, from the far field to the disk face,
\begin{equation}
  p_\infty = p_1 + \tfrac12 \rho \vind^2 ,
  \label{eq:bern1}
\end{equation}
and downstream, from the disk face to the far wake,
\begin{equation}
  p_2 + \tfrac12 \rho \vind^2 = p_\infty + \tfrac12 \rho w^2 .
  \label{eq:bern2}
\end{equation}
Subtracting \cref{eq:bern1} from \cref{eq:bern2} eliminates $p_\infty$ and
$\vind$:
\begin{equation}
  \Delta p \coloneqq p_2 - p_1 = \tfrac12 \rho w^2 ,
  \qquad\text{whence}\qquad
  T = A\,\Delta p = \tfrac12 \rho A w^2 .
  \label{eq:pressurejump}
\end{equation}

\subsection{The central result}

\begin{theorem}[Froude's theorem and the hover power law]
\label{thm:momentum}
Under \cref{mod:disk} the far-wake velocity is exactly twice the velocity
induced at the disk,
\begin{equation}
  w = 2\vind ,
  \label{eq:froude}
\end{equation}
the thrust is
\begin{equation}
  T = 2\rho A \vind^2 ,
  \qquad\text{equivalently}\qquad
  \vind = \sqrt{\frac{T}{2\rho A}} ,
  \label{eq:vi}
\end{equation}
the far wake contracts to half the disk area, $A_\infty = A/2$, and the power
delivered to the fluid is
\begin{equation}
  \Pideal = T\vind = \frac{T^{3/2}}{\sqrt{2\rho A}} .
  \label{eq:pideal}
\end{equation}
\end{theorem}

\begin{proof}
Equate the two expressions for thrust, \cref{eq:momentum,eq:pressurejump}:
\[
  \rho A \vind w = \tfrac12 \rho A w^2 .
\]
Since $T>0$ forces $w>0$ and $A>0$, division by $\rho A w$ is legitimate and
gives $\vind = w/2$, which is \cref{eq:froude}. Substituting $w = 2\vind$ into
\cref{eq:momentum} yields $T = 2\rho A \vind^2$, and solving for $\vind$ gives
\cref{eq:vi}; the positive root is selected because the wake is directed
downward for positive thrust. Wake contraction follows from
\cref{eq:continuity}: $A_\infty = A\vind/w = A/2$.

For the power, the rate of work done by the disk on the fluid equals the flux
of kinetic energy out of the control volume, the fluid entering at rest:
\[
  P = \tfrac12 \dot m_{\mathrm{air}} w^2
    = \tfrac12 (\rho A \vind)(2\vind)^2
    = 2\rho A \vind^3
    = (2\rho A \vind^2)\,\vind = T\vind ,
\]
so $P = T\vind$, which is also what one obtains directly as force times the
velocity of the disk relative to the fluid it acts on. Substituting
\cref{eq:vi} gives $\Pideal = T\sqrt{T/(2\rho A)} = T^{3/2}(2\rho A)^{-1/2}$.
\end{proof}

\begin{remark}
The factor of two in \cref{eq:froude} is the reason the disk sees only half the
dynamic pressure of the wake, and is the origin of the $3/2$ exponent that
drives the entire feasibility analysis. Half the velocity increment happens
upstream of the disk and half downstream; \cref{fig:axisprofile} shows both
halves.
\end{remark}

\begin{figure}[tbp]
\centering
% Axial velocity and static pressure along the rotor axis, from Bernoulli on
% each side of the disk with the vortex-cylinder axial velocity
% u/v_i = 1 + xi/sqrt(1+xi^2). Pressure in units of rho v_i^2.
\begin{tikzpicture}[declare function={u(\s) = 1 + \s/sqrt(1+\s*\s);}]
\begin{groupplot}[group style={group size=1 by 2, vertical sep=0.45cm,
                  xlabels at=edge bottom, xticklabels at=edge bottom},
    paperplot, height=4.1cm, xmin=-4, xmax=5, samples=160,
    xlabel={$x/R$, distance along the axis, positive downstream}]
  \nextgroupplot[ylabel={$u/\vind$}, ymin=0, ymax=2.25, ytick={0,1,2}]
    \addplot[cA, domain=-4:5] {u(x)};
    \addplot[construct, domain=-4:5] {2};
    \addplot[construct, domain=-4:5] {1};
    \draw[cG, dashed] (axis cs:0,0) -- (axis cs:0,2.25);
    \node[lbl, anchor=south west] at (axis cs:0.1,0.05) {disk};
    \node[lbl, anchor=north east, fill=white, inner sep=1pt] at (axis cs:4.9,1.88) {$w = 2\vind$};
    \draw[dim] (axis cs:-2.2,0) -- node[lbl, left, fill=white, inner sep=1pt]
        {half upstream} (axis cs:-2.2,1);
    \draw[dim] (axis cs:2.2,1) -- node[lbl, right, fill=white, inner sep=1pt]
        {half downstream} (axis cs:2.2,2);
  \nextgroupplot[ylabel={$(p-p_\infty)/\rho\vind^2$}, ymin=-2.25, ymax=1.9,
                 ytick={-2,-0.5,0,1.5}]
    \addplot[cA, domain=-4:0] {-0.5*u(x)^2};
    \addplot[cA, domain=0:5] {2 - 0.5*u(x)^2};
    \addplot[construct, domain=-4:5] {0};
    \draw[cD, line width=1.1pt, -{Stealth[length=2mm]}] (axis cs:0,-0.5) -- (axis cs:0,1.5);
    \node[lbl, cD, anchor=east, fill=white, inner sep=1pt] at (axis cs:-0.15,0.75)
        {$\Delta p = 2\rho\vind^2 = \tfrac12\rho w^2$};
    \node[lbl, anchor=north, fill=white, inner sep=1pt] at (axis cs:-2.4,-1.0)
        {suction ahead of the disk};
    \node[lbl, anchor=north, fill=white, inner sep=1pt] at (axis cs:3.0,-1.0)
        {recovers to $p_\infty$ in the wake};
\end{groupplot}
\end{tikzpicture}

\caption{Axial velocity and static pressure along the rotor axis, from
Bernoulli's equation applied separately on each side of the disk with the
vortex-cylinder velocity $u/\vind = 1 + \xi/\sqrt{1+\xi^2}$, $\xi = x/R$. The
velocity is continuous through the disk and rises from $0$ to $\vind$
upstream and from $\vind$ to $w = 2\vind$ downstream, the two halves of
\cref{eq:froude}. The pressure falls ahead of the disk, jumps by
$\Delta p = \tfrac12\rho w^2$ across it, \cref{eq:pressurejump}, and recovers
to $p_\infty$ in the wake.}
\label{fig:axisprofile}
\end{figure}

\begin{corollary}[Hover]
\label{cor:hover}
In hover, $T = mg$, so
\begin{equation}
  \vind = \sqrt{\frac{mg}{2\rho A}},
  \qquad
  \Pideal = \frac{(mg)^{3/2}}{\sqrt{2\rho A}} .
  \label{eq:hoverpower}
\end{equation}
At fixed disk area, $\Pideal \propto m^{3/2}$.
\end{corollary}

\subsection{Disk loading}

\begin{definition}[Disk loading]
$\mathrm{DL} \coloneqq T/A$, with units of pressure.
\end{definition}

Rewriting \cref{eq:vi,eq:pideal} in terms of disk loading isolates the
physically meaningful group:
\begin{equation}
  \vind = \sqrt{\frac{\mathrm{DL}}{2\rho}},
  \qquad
  \frac{\Pideal}{T} = \sqrt{\frac{\mathrm{DL}}{2\rho}} ,
  \qquad
  \frac{\Pideal}{m} = g\sqrt{\frac{\mathrm{DL}}{2\rho}} .
  \label{eq:dl}
\end{equation}

\begin{observation}
\label{obs:dl}
Power per unit mass in hover depends on the disk loading alone, not on mass and
area separately. A hovering machine is therefore not primarily asking for
powerful motors; it is asking for low disk loading. The quantity
$T/\Pideal = \sqrt{2\rho/\mathrm{DL}}$, the power loading, is the correct
figure of merit for a hovering vehicle and improves as disk loading falls.
\end{observation}

\Cref{fig:powerloading} places the design of \cref{sec:results} on this law.

\begin{figure}[tbp]
\centering
% Power loading against disk loading, T/P = FM sqrt(2 rho / DL), eq:dl.
\providecommand{\aThover}{4.23}
\providecommand{\aVstall}{24.71}
\providecommand{\aVdesign}{26.69}
\providecommand{\aPstall}{17.04}
\providecommand{\aPdesign}{17.56}
\providecommand{\aPind}{14.87}
\providecommand{\aCtMax}{0.14}
\providecommand{\aCtDesign}{0.12}
\providecommand{\aSigma}{0.212}
\providecommand{\aRho}{1.225}
\providecommand{\aDL}{22.2}
\providecommand{\aWake}{6.02}
\providecommand{\aKappa}{1.168}

\begin{tikzpicture}
\begin{axis}[paperplot, xmode=log, ymode=log, xmin=5, xmax=2000,
    ymin=0.02, ymax=1.3, samples=40, domain=5:2000,
    xlabel={disk loading $\mathrm{DL} = T/A$ [\si{\pascal}]},
    ylabel={power loading $T/P$ [\si{\newton\per\watt}]},
    legend pos=south west]
  % wake faster than 5 m/s: loose paper moves, requirement (R3)
  \pgfmathsetmacro{\DLfive}{2*\aRho*2.5*2.5}
  \fill[cD!7] (axis cs:\DLfive,0.02) rectangle (axis cs:2000,1.3);
  \draw[cD, dashed, line width=0.6pt] (axis cs:\DLfive,0.02) -- (axis cs:\DLfive,1.3);
  \node[lbl, cD, anchor=north west] at (axis cs:\DLfive*1.05,1.25)
      {wake $w = 2\vind > \SI{5}{\metre\per\second}$};
  \addplot[cA] {sqrt(2*\aRho/x)};
  \addlegendentry{$\FM = 1$, ideal disk}
  \addplot[cA!70, dashed] {0.75*sqrt(2*\aRho/x)};
  \addlegendentry{$\FM = 0.75$}
  \addplot[cA!45, dotted, line width=1.2pt] {0.5*sqrt(2*\aRho/x)};
  \addlegendentry{$\FM = 0.5$}
  % the design
  \pgfmathsetmacro{\TPd}{\rFM*sqrt(2*\aRho/\aDL)}
  \addplot[only marks, mark=*, mark size=2.2pt, cA, forget plot]
      coordinates {(\aDL,\TPd)};
  \node[lbl, anchor=north west, align=left, fill=white, inner sep=1.5pt]
      (dl) at (axis cs:40,0.62)
      {design: $\mathrm{DL}=\SI{\aDL}{\pascal}$,\\
       $\FM=\rFM$, $w=\SI{\aWake}{\metre\per\second}$};
  \draw[cG, line width=0.4pt] (dl.south west) -- (axis cs:\aDL*1.06,\TPd*1.04);
  % slope triangle, -1/2
  \draw[cG, line width=0.5pt] (axis cs:400,0.16) -- (axis cs:1600,0.16)
      -- (axis cs:400,0.32) -- cycle;
  \node[lbl, cG, anchor=north] at (axis cs:800,0.155) {slope $-\tfrac12$};
\end{axis}
\end{tikzpicture}

\caption{Power loading against disk loading, $T/P = \FM\sqrt{2\rho/\mathrm{DL}}$,
from \cref{eq:dl} and the definition of $\FM$ in \cref{sec:blade}. On
logarithmic axes every figure of merit gives a line of slope $-\tfrac12$:
halving the disk loading buys a factor $\sqrt2$ in thrust per watt, whatever
the rotor. The shaded region is where the wake exceeds
\SI{5}{\metre\per\second}, the speed at which loose paper begins to move
(requirement (R3) of \cref{sec:results}); the design sits just inside it, at
$\mathrm{DL} = \SI{\aDL}{\pascal}$ and $w = \SI{\rWake}{\metre\per\second}$.}
\label{fig:powerloading}
\end{figure}

This single observation determines the shape of every machine in
\cref{sec:results}: large, slowly turning rotors, because
\cref{eq:dl} says that is the only lever with any authority, and
\cref{sec:acoustics} will show that the same choice is what makes the machine
quiet.

\subsection{Axial flight and the validity of \texorpdfstring{\cref{thm:momentum}}{Theorem 3.2}}

For axial motion at climb velocity $V_c$ relative to the air, the same control
volume argument with a non-zero inflow gives
\begin{equation}
  T = 2\rho A (V_c + \vind)\vind
  \quad\Longrightarrow\quad
  \vind = -\frac{V_c}{2} + \sqrt{\left(\frac{V_c}{2}\right)^2
          + \frac{T}{2\rho A}} ,
  \label{eq:axial}
\end{equation}
which reduces to \cref{eq:vi} at $V_c=0$ and is implemented in that form. The
power is then $P = T(V_c + \vind)$.

\Cref{eq:axial} is valid in climb and in the windmill-brake state, but fails in
descent for $-2 \lesssim V_c/v_{i,h} \lesssim 0$, where $v_{i,h}$ is the hover
induced velocity. This is the vortex ring state, where
the streamtube assumption breaks down and the flow is unsteady and
recirculating~\cite[\S2.13]{leishman2006}. The machines considered here descend
slowly and the engine does not enter that regime, but it is a real limitation
of the model rather than a technicality: a fast descent to land is precisely
the manoeuvre it cannot predict.

Two further caveats are worth recording. Momentum theory gives a lower bound on
the power: it is the power required to produce the thrust by the most efficient
possible redistribution of momentum, and any real rotor does worse.
\Cref{sec:blade} quantifies by how much. Second, the theory says nothing about
how the loading is produced, so it cannot predict stall, compressibility or
Reynolds effects; those enter through the blade element layer.

\section{Blade element corrections and the figure of merit}
\label{sec:blade}

Momentum theory bounds the power from below. A real rotor pays two further
tolls: the induced flow is not uniform and the blades are finite in number, and
the blade sections have profile drag. This section derives the second from
blade element theory and absorbs the first into a coefficient, with the
consequence that the figure of merit becomes an output of blade geometry rather
than an input.

\subsection{Induced power factor}

\begin{model}[Non-ideal induced power]
\label{mod:kappa}
The induced power is written
\begin{equation}
  \Pind = \kappa\,T\vind, \qquad \kappa > 1 ,
  \label{eq:kappa}
\end{equation}
where $\kappa$ collects non-uniform inflow, tip loss and swirl. Typical values
for small rotors are $\kappa \in [1.10,1.25]$~\cite[\S2.4]{leishman2006}; we
take $\kappa = 1.15$.
\end{model}

The alternative is Prandtl's tip-loss factor $B$, with the induced velocity
computed over an effective disk of radius $BR$. The two are related by
$\kappa \approx B^{-2}$ for small losses; \cref{eq:kappa} is used because it
keeps the algebra of \cref{sec:closure} in closed form.

\subsection{Profile power}

\begin{definition}[Solidity]
For a rotor with $N_b$ blades of constant chord $c$ and radius $R$, the
solidity is the ratio of blade planform area to disk area,
\begin{equation}
  \sigma \coloneqq \frac{N_b c R}{\pi R^2} = \frac{N_b c}{\pi R}
        = \frac{N_b}{\pi\,\mathrm{AR}},
  \qquad \mathrm{AR} \coloneqq \frac{R}{c} .
  \label{eq:solidity}
\end{equation}
\end{definition}

\Cref{fig:element} fixes the notation of blade element theory.

\begin{figure}[tbp]
\centering
% Blade element theory. (a) the rotor seen from above: two blades of chord c,
% the annulus at radius r of width dr and the element it cuts from a blade.
% (b) section A-A seen from the blade tip: velocity triangle, angles and the
% resolution of lift and drag into thrust and in-plane force.
\begin{tikzpicture}[
  declare function={yt(\s) = 0.6*(0.2969*sqrt(\s) - 0.1260*\s
                     - 0.3516*\s*\s + 0.2843*\s*\s*\s - 0.1015*\s*\s*\s*\s);}]
  % ------------------------------------------------------------ panel (a)
  \begin{scope}
    \draw[cG, dashed] (0,0) circle (1.9);
    \fill[cB!10, even odd rule] (0,0) circle (1.36) circle (1.14);
    \draw[cB!60, line width=0.3pt] (0,0) circle (1.36) (0,0) circle (1.14);
    \foreach \s in {1,-1}{
      \draw[cA, fill=cA!12, line width=0.7pt]
          ({\s*0.26},-0.315) rectangle ({\s*1.9},0.315);
    }
    \fill[cB!70] (1.14,-0.315) rectangle (1.36,0.315);
    \fill[black!65] (0,0) circle (0.26);
    % section marks A-A
    \draw[line width=0.9pt] (1.25,0.45) -- (1.25,0.68);
    \node[lbl, anchor=east] at (1.2,0.66) {A};
    \draw[line width=0.9pt] (1.25,-0.45) -- (1.25,-0.68);
    \node[lbl, anchor=east] at (1.2,-0.66) {A};
    % dimensions
    \draw[dim] (0,-1.05) -- node[lbl, below, fill=white, inner sep=1pt] {$r$} (1.25,-1.05);
    \draw[dim] (1.14,1.02) -- (1.36,1.02);
    \node[lbl, above] at (1.25,1.06) {$\dd r$};
    \draw[dim] (1.72,-0.315) -- node[lbl, right] {$c$} (1.72,0.315);
    \draw[dim] (0,0) -- node[lbl, above left, inner sep=1pt] {$R$} (125:1.9);
    % rotation, clockwise seen from above
    \draw[vec, cA] (80:2.2) arc (80:35:2.2) node[right, lbl] {$\Omega$};
    \node[font=\small] at (0,-2.5) {(a) rotor seen from above};
  \end{scope}
  % ------------------------------------------------------------ panel (b)
  \begin{scope}[xshift=8.3cm, yshift=-0.2cm]
    \def\ph{12} \def\th{20}
    % plane of rotation
    \draw[construct] (-4.2,0) -- (3.6,0) node[right, lbl, align=left] {plane of\\rotation};
    % aerofoil, quarter chord at the origin, leading edge left, nose up by theta
    \begin{scope}[rotate=-\th]
      \draw[cA, fill=cA!15, line width=0.8pt]
        plot[domain=0:180, samples=70] ({3.0*((1-cos(\x))/2-0.25)}, { 3.0*yt((1-cos(\x))/2)})
        -- plot[domain=180:0, samples=70] ({3.0*((1-cos(\x))/2-0.25)}, {-3.0*yt((1-cos(\x))/2)})
        -- cycle;
      \draw[cG, dash dot, line width=0.4pt] (-3.5,0) -- (2.6,0);
    \end{scope}
    \node[lbl, cG, anchor=north] at ({-3.6*cos(\th)},{3.6*sin(\th)-0.02}) {chord};
    % direction of the resultant wind through the section
    \draw[construct] ({-3.6*cos(\ph)},{3.6*sin(\ph)}) -- (0,0);
    % angles on the upstream side
    \draw[cG] (180:1.55) arc (180:{180-\ph}:1.55);
    \node[lbl] at ({180-\ph/2}:1.82) {$\phi$};
    \draw[cG] ({180-\ph}:2.3) arc ({180-\ph}:{180-\th}:2.3);
    \node[lbl] at ({180-(\ph+\th)/2}:2.58) {$\alpha$};
    \draw[cG] (180:3.05) arc (180:{180-\th}:3.05);
    \node[lbl] at ({180-\th/2+1}:3.35) {$\theta$};
    % velocity triangle, drawn apart
    \begin{scope}[shift={(-4.1,1.95)}]
      \pgfmathsetmacro{\tp}{2.2*tan(\ph)}
      \draw[vec, cB] (0,0) -- node[lbl, above] {$U_T = \Omega r$} (2.2,0);
      \draw[vec, cB] (2.2,0) -- node[lbl, right] {$U_P = \vind$} (2.2,-\tp);
      \draw[vec, black, line width=1.0pt] (0,0) -- node[lbl, below left, inner sep=1pt] {$U$} (2.2,-\tp);
      \draw[cG] (0.75,0) arc (0:-\ph:0.75);
      \node[lbl] at (-6:1.0) {$\phi$};
    \end{scope}
    % forces at the quarter chord
    \pgfmathsetmacro{\Lx}{2.4*sin(\ph)} \pgfmathsetmacro{\Ly}{2.4*cos(\ph)}
    \pgfmathsetmacro{\Dx}{0.75*cos(\ph)} \pgfmathsetmacro{\Dy}{-0.75*sin(\ph)}
    \coordinate (F) at ({\Lx+\Dx},{\Ly+\Dy});
    \draw[construct] (\Lx,\Ly) -- (F) -- (\Dx,\Dy);
    \draw[construct] (F) -- (0,{\Ly+\Dy});
    \draw[construct] (F) -- ({\Lx+\Dx},0);
    \draw[vec, cA, line width=1.0pt] (0,0) -- (\Lx,\Ly) node[above right, lbl, inner sep=1pt] {$\dd L$};
    \draw[vec, cD, line width=1.0pt] (0,0) -- (\Dx,\Dy);
    \node[lbl, cD, anchor=north west, inner sep=1pt] at (\Dx,\Dy) {$\dd D$};
    \draw[vec, black] (0,0) -- (F);
    \draw[vec, cC!70!black, line width=1.1pt] (0,0) -- (0,{\Ly+\Dy})
        node[left, lbl] {$\dd T$};
    \draw[vec, cC!70!black, line width=1.1pt] (0,0) -- ({\Lx+\Dx},0);
    \node[lbl, cC!60!black, anchor=south west, inner sep=1pt] at ({\Lx+\Dx},0.02) {$\dd Q/r$};
    % blade motion
    \draw[vec, cG] (-0.6,-1.2) -- (-2.2,-1.2) node[left, lbl] {blade motion};
    \node[font=\small] at (-0.4,-2.3) {(b) section A-A, seen from the tip};
  \end{scope}
\end{tikzpicture}

\caption{Blade element notation. (a) A two-bladed rotor of radius $R$ and chord
$c$, so $\sigma = N_bc/(\pi R)$; the annulus at radius $r$ of width $\dd r$
cuts an element from each blade. (b) The element in section. It meets the air
at $U_T = \Omega r$ in the plane of rotation and $U_P = \vind$ through it, so
the resultant $U$ is inclined at the inflow angle $\phi = \arctan(U_P/U_T)$
and the angle of attack is $\alpha = \theta - \phi$ for a pitch $\theta$. Lift
$\dd L$ is normal to $U$ and drag $\dd D$ along it. Resolved on the rotor axes
they give the thrust $\dd T = \dd L\cos\phi - \dd D\sin\phi$ and the in-plane
force $\dd Q/r = \dd L\sin\phi + \dd D\cos\phi$. \Cref{prop:profile} keeps
the drag part of the second in the limit $\phi\to0$, where
$\dd P = \Omega r\,\dd D$.}
\label{fig:element}
\end{figure}

\begin{proposition}[Profile power of a rotor in hover]
\label{prop:profile}
Let each blade section have drag coefficient $\Cdz$, taken constant along the
span, and let the rotor turn at angular rate $\Omega$ with tip speed
$\Vtip = \Omega R$. Neglecting the inflow angle in the drag direction, the
profile power is
\begin{equation}
  \Pprof = \frac{1}{8}\,\rho\, A\, \Vtip^{3}\, \sigma\, \Cdz .
  \label{eq:profile}
\end{equation}
\end{proposition}

\begin{proof}
Consider a blade element at radius $r$ of span $\dd r$. Its sectional velocity
is $\Omega r$ to leading order, so the drag on the element is
\[
  \dd D = \tfrac12 \rho (\Omega r)^2 c\, \Cdz \dd r ,
\]
and the power absorbed is $\dd P = \Omega r \dd D$. Integrating over the blade
and multiplying by $N_b$ blades,
\[
  \Pprof = N_b \int_0^R \tfrac12 \rho \Omega^3 r^3 c\,\Cdz \dd r
         = \tfrac{1}{8}\rho\,\Omega^3 c\, N_b\, \Cdz R^4 .
\]
Substituting $c = \sigma \pi R/N_b$ from \cref{eq:solidity},
\[
  \Pprof = \tfrac{1}{8}\rho\,\Omega^3 \Cdz \,\sigma \pi R^{5}
         = \tfrac{1}{8}\rho\, \Cdz\, \sigma\, (\pi R^2)(\Omega R)^3
         = \tfrac{1}{8}\rho\, A\, \Vtip^3\, \sigma\, \Cdz . \qedhere
\]
\end{proof}

\subsection{Figure of merit}

\begin{definition}[Figure of merit]
\begin{equation}
  \FM \coloneqq \frac{\Pideal}{\Pshaft}
      = \frac{T\vind}{\kappa T\vind + \Pprof} \in (0,1) .
  \label{eq:fm}
\end{equation}
\end{definition}

\begin{proposition}[Figure of merit at fixed blade loading]
\label{prop:fmblade}
Define the blade loading coefficient
\begin{equation}
  \frac{\Ct}{\sigma} \coloneqq \frac{T}{\rho A \Vtip^2 \sigma} .
  \label{eq:ctsigma}
\end{equation}
Then
\begin{equation}
  \FM = \left[\kappa + \frac{\Cdz}{8}\,
        \frac{1}{(\Ct/\sigma)}\,\frac{\Vtip}{\vind}\right]^{-1},
  \qquad
  \frac{\Vtip}{\vind} = \sqrt{\frac{2}{\sigma\,(\Ct/\sigma)}} ,
  \label{eq:fmclosed}
\end{equation}
so that at fixed blade loading the figure of merit depends on the blade
geometry only through the group $\sigma\,(\Ct/\sigma) = \Ct$, and is
independent of rotor size and of tip speed separately.
\end{proposition}

\begin{proof}
From \cref{eq:profile,eq:ctsigma}, $\rho A \sigma = T/(\Vtip^2 (\Ct/\sigma))$
and hence
\[
  \Pprof = \tfrac18 \Vtip^3 \Cdz \cdot \frac{T}{\Vtip^2 (\Ct/\sigma)}
         = \frac{\Cdz}{8}\,\frac{T\,\Vtip}{(\Ct/\sigma)} .
\]
Dividing \cref{eq:fm} through by $T\vind$ gives the first identity. For the
second, \cref{eq:vi} gives $\vind^2 = T/(2\rho A)$ and \cref{eq:ctsigma} gives
$\Vtip^2 = T/(\rho A \sigma (\Ct/\sigma))$, so
$\Vtip^2/\vind^2 = 2/(\sigma (\Ct/\sigma))$.
\end{proof}

\begin{corollary}[Design rule]
\label{cor:designrule}
Holding blade loading at its design value and increasing solidity lowers the
tip speed required to carry a given thrust,
$\Vtip \propto \sigma^{-1/2}$, and lowers the profile power with it,
$\Pprof \propto \Vtip \propto \sigma^{-1/2}$.
\end{corollary}

\begin{proof}
Immediate from \cref{eq:ctsigma} with $T$, $\rho A$ and $\Ct/\sigma$ held
fixed, together with
\[
  \Pprof = \frac{\Cdz}{8}\,\frac{T\,\Vtip}{\Ct/\sigma} . \qedhere
\]
\end{proof}

\Cref{cor:designrule} is not obvious and is used heavily later: a fatter,
slower blade is better on \emph{both} counts, quieter by
\cref{sec:acoustics} and cheaper in profile power, until the assumptions break.
They break in two places, and both are enforced as constraints.
\Cref{fig:designrule} shows both effects at the design's thrust and the two
limits that stop them.

\begin{figure}[tbp]
\centering
% Cor. design rule: at fixed blade loading Ct/sigma the figure of merit
% rises and the tip speed falls as solidity grows, eq:fmclosed and
% eq:ctsigma, at the design's hover thrust per rotor and induced factor.
\providecommand{\aThover}{4.23}
\providecommand{\aVstall}{24.71}
\providecommand{\aVdesign}{26.69}
\providecommand{\aPstall}{17.04}
\providecommand{\aPdesign}{17.56}
\providecommand{\aPind}{14.87}
\providecommand{\aCtMax}{0.14}
\providecommand{\aCtDesign}{0.12}
\providecommand{\aSigma}{0.212}
\providecommand{\aRho}{1.225}
\providecommand{\aDL}{22.2}
\providecommand{\aWake}{6.02}
\providecommand{\aKappa}{1.168}

\begin{tikzpicture}[
  declare function={
    FMc(\s,\c) = 1/(\aKappa + (\rCdzero/(8*\c))*sqrt(2/(\s*\c)));
    Ar = pi*(\rDiam/2000)^2;
    Vt(\s,\c) = sqrt(\aThover/(\aRho*Ar*\s*\c));}]
\begin{groupplot}[group style={group size=2 by 1, horizontal sep=1.55cm},
    halfplot, xmin=0.03, xmax=0.35, domain=0.03:0.35, samples=80,
    xlabel={solidity $\sigma$}]
  \nextgroupplot[ylabel={figure of merit $\FM$}, ymin=0.3, ymax=0.8,
                 legend columns=2,
                 legend style={font=\scriptsize, at={(0.02,0.03)}, anchor=south west}]
    \fill[cG!12] (axis cs:0.05,0.3) rectangle (axis cs:0.12,0.8);
    \node[lbl, cG, anchor=north] at (axis cs:0.085,0.795) {helicopters};
    \fill[cD!10] (axis cs:0.25,0.3) rectangle (axis cs:0.35,0.8);
    \node[lbl, cD, rotate=90, anchor=center] at (axis cs:0.30,0.5) {cascade, $\sigma > 0.25$};
    \addplot[cA!40] {FMc(x,0.08)}; \addlegendentry{$\Ct/\sigma = 0.08$}
    \addplot[cA!65] {FMc(x,0.10)}; \addlegendentry{$0.10$}
    \addplot[cA]    {FMc(x,0.12)}; \addlegendentry{$0.12$, design}
    \addplot[cB, dashed] {FMc(x,0.14)}; \addlegendentry{$0.14$, stall}
    \addplot[only marks, mark=*, mark size=2pt, cA, forget plot]
        coordinates {(\aSigma,{FMc(\aSigma,0.12)})};
  \nextgroupplot[ylabel={tip speed $\Vtip$ [\si{\metre\per\second}]},
                 ymin=10, ymax=95]
    \fill[cD!10] (axis cs:0.25,10) rectangle (axis cs:0.35,95);
    \addplot[cA!40] {Vt(x,0.08)};
    \addplot[cA!65] {Vt(x,0.10)};
    \addplot[cA]    {Vt(x,0.12)};
    \addplot[cB, dashed] {Vt(x,0.14)};
    \addplot[only marks, mark=*, mark size=2pt, cA]
        coordinates {(\aSigma,{Vt(\aSigma,0.12)})};
    \node[lbl, anchor=south west, fill=white, inner sep=1pt]
        at (axis cs:\aSigma,{Vt(\aSigma,0.12)+2})
        {design, \SI{\rVtip}{\metre\per\second}};
    \node[lbl, cG, anchor=north east, fill=white, inner sep=1pt]
        at (axis cs:0.34,93) {$\Vtip\propto\sigma^{-1/2}$};
\end{groupplot}
\end{tikzpicture}

\caption{The design rule of \cref{cor:designrule} at the thrust per rotor of
the design, $T = \SI{\aThover}{\newton}$, with its induced factor
$\kappa = \aKappa$ and $\Cdz = \rCdzero$ held fixed. Left: at fixed blade
loading the figure of merit \cref{eq:fmclosed} rises with solidity, because
the profile term falls as $\sigma^{-1/2}$. Right: the tip speed that carries
the thrust falls as $\sigma^{-1/2}$, \cref{eq:ctsigma}, which is what makes
the rotor quiet (\cref{sec:acoustics}). Both gains stop at the cascade limit
$\sigma = 0.25$ of \cref{as:sigma} and at the stall limit $\Ct/\sigma = 0.14$
of \cref{as:stall}. The design ($\bullet$) sits at $\sigma = \rSolidity$ and
$\Ct/\sigma = \rCtSigma$; helicopter rotors occupy the grey band.}
\label{fig:designrule}
\end{figure}

\subsection{Two limits on solidity and blade loading}

\begin{assumption}[Blade-element independence]
\label{as:sigma}
Blade element theory treats each section as an isolated aerofoil in a uniform
inflow. This fails once the blades operate as a cascade. Helicopter rotors run
$\sigma \in [0.05,0.12]$; high-solidity propellers reach about $0.25$. We
impose $\sigma \le \sigma_{\max} = 0.25$ and report violations, because beyond
it the model flatters the rotor and one no longer has a propeller but a ducted
fan without a duct.
\end{assumption}

\begin{assumption}[Stall margin]
\label{as:stall}
Sectional lift coefficient is bounded by stall, which through
$\Ct/\sigma \approx C_L/6$ for linearly twisted blades bounds the blade
loading. We impose $\Ct/\sigma \le 0.14$ and design at $0.12$.
\end{assumption}

\begin{remark}[Fixed-pitch rotors run at constant $\Ct$]
For a fixed-pitch rotor at fixed collective, $T \propto \Omega^2$ so
$\Ct = T/(\rho A \Vtip^2)$ is invariant under changes of rotor speed. The
blade-loading constraint therefore binds identically in hover and at maximum
thrust, which is why the engine reports the same $\Ct/\sigma$ at both
conditions. Stall margin cannot be recovered by spinning faster.
\end{remark}

\subsection{Reynolds number}

Both \cref{prop:profile,prop:fmblade} take $\Cdz$ constant. A design driven
toward low tip speed by \cref{cor:designrule} operates at chord Reynolds
numbers
\begin{equation}
  \mathrm{Re} = \frac{0.7\,\Vtip\, c}{\nu} ,
  \label{eq:reynolds}
\end{equation}
evaluated at the $0.7R$ station, of order $10^5$ or below, where profile drag
rises sharply. Treating $\Cdz$ as constant would hand the optimiser a free
lunch: it would drive the tip speed arbitrarily low. We therefore use
\begin{equation}
  \Cdz(\mathrm{Re}) = \Cdz^{\mathrm{ref}}
    \min\!\left\{\left(\frac{\mathrm{Re}_{\mathrm{ref}}}{\mathrm{Re}}
    \right)^{0.2},\,3\right\},
  \qquad \mathrm{Re}_{\mathrm{ref}} = 2\times 10^5 ,
  \label{eq:recorrection}
\end{equation}
an empirical power law of the standard form for attached turbulent flow, capped
to keep it from diverging. This is the crudest submodel in the paper and is
flagged as such; its only job is to prevent an unphysical optimum.
\Cref{fig:tipspeed} shows the resulting trade for one rotor of the design.

\begin{figure}[tbp]
\centering
% One rotor of the design at its hover thrust, fixed geometry, tip speed
% swept with automatic tip speed off (engine data, gen_a.py).
\providecommand{\aThover}{4.23}
\providecommand{\aVstall}{24.71}
\providecommand{\aVdesign}{26.69}
\providecommand{\aPstall}{17.04}
\providecommand{\aPdesign}{17.56}
\providecommand{\aPind}{14.87}
\providecommand{\aCtMax}{0.14}
\providecommand{\aCtDesign}{0.12}
\providecommand{\aSigma}{0.212}
\providecommand{\aRho}{1.225}
\providecommand{\aDL}{22.2}
\providecommand{\aWake}{6.02}
\providecommand{\aKappa}{1.168}

\begin{tikzpicture}
\begin{axis}[paperplot, xmin=18, xmax=80, ymin=0, ymax=80,
    xlabel={tip speed $\Vtip$ [\si{\metre\per\second}]},
    ylabel={power per rotor [\si{\watt}]},
    legend style={at={(0.3,0.97)}, anchor=north west}]
  \fill[cD!10] (axis cs:18,0) rectangle (axis cs:\aVstall,80);
  \node[lbl, cD, rotate=90, anchor=center] at (axis cs:{(18+\aVstall)/2},50)
      {stall: $\Ct/\sigma > \aCtMax$};
  \draw[cA, dashed, line width=0.6pt] (axis cs:\aVdesign,0) -- (axis cs:\aVdesign,80);
  \addplot[black] table[x=vtip, y=Pshaft] {results/fig/a_tipspeed.dat};
  \addlegendentry{shaft, $\Pind + \Pprof$}
  \addplot[cA, dashed] table[x=vtip, y=Pind] {results/fig/a_tipspeed.dat};
  \addlegendentry{induced, $\kappa T\vind$}
  \addplot[cB] table[x=vtip, y=Pprof] {results/fig/a_tipspeed.dat};
  \addlegendentry{profile, $\tfrac18\rho A\Vtip^3\sigma\Cdz(\mathrm{Re})$}
  \addplot[cG, dotted, line width=1.2pt] table[x=vtip, y=Pideal] {results/fig/a_tipspeed.dat};
  \addlegendentry{ideal, $T\vind$}
  \addplot[only marks, mark=*, mark size=2.2pt, cD, forget plot]
      coordinates {(\aVstall,\aPstall)};
  \addplot[only marks, mark=*, mark size=2.2pt, cA, forget plot]
      coordinates {(\aVdesign,\aPdesign)};
  \node[lbl, cD, anchor=south west, fill=white, inner sep=1.2pt] (st)
      at (axis cs:18.6,6.5) {stall boundary};
  \draw[cD, line width=0.4pt] (st.north east) -- (axis cs:\aVstall,{\aPstall-0.7});
  \node[lbl, anchor=south west, fill=white, inner sep=1.2pt] (ds)
      at (axis cs:32,31) {design, $\Ct/\sigma = \aCtDesign$};
  \draw[cA, line width=0.4pt] (ds.south west) -- (axis cs:{\aVdesign+0.3},{\aPdesign+0.7});
\end{axis}
\end{tikzpicture}

\caption{One rotor of the design at its hover thrust, with the geometry fixed
and the tip speed varied (engine data). The induced power $\kappa T\vind$ does
not depend on tip speed; the profile power grows as $\Vtip^3\Cdz(\mathrm{Re})$,
slightly slower than $\Vtip^3$ at low speed because $\Cdz$ falls as the
Reynolds number rises, \cref{eq:recorrection}. The shaft power is therefore
least at the lowest tip speed the blade-loading limit admits, on the stall
boundary $\Ct/\sigma = \aCtMax$ at \SI{\aVstall}{\metre\per\second}. The
design runs at $\Ct/\sigma = \aCtDesign$ and \SI{\aVdesign}{\metre\per\second},
and pays for its stall margin with the difference between
\SI{\aPdesign}{\watt} and \SI{\aPstall}{\watt} of shaft power per rotor.}
\label{fig:tipspeed}
\end{figure}

\subsection{Electrical power}

Motor and electronic speed controller losses are lumped into a single
efficiency $\eta$, so the electrical power drawn by the propulsion system is
\begin{equation}
  \Pelec = \frac{\Pshaft}{\eta}
         = \frac{\kappa T\vind + \Pprof}{\eta}
         = \frac{\Pideal}{\FM\,\eta}
         = \frac{T^{3/2}}{\FM\,\eta\,\sqrt{2\rho A}} .
  \label{eq:pelec}
\end{equation}
The final form is the one carried into \cref{sec:closure}; the intermediate
forms are what the engine actually evaluates, so that $\FM$ is computed from
\cref{eq:fm} rather than assumed.

Adding the non-propulsive load $\Paux$ (display, computer, sensors), the total
electrical power in hover is
\begin{equation}
  P_{\mathrm{tot}}(m) = \frac{(mg)^{3/2}}{\FM\,\eta\,\sqrt{2\rho A}} + \Paux .
  \label{eq:ptot}
\end{equation}

\begin{remark}
Unlike a camera drone, this vehicle carries a payload that consumes power
continuously. $\Paux$ is not negligible: at the designs of \cref{sec:results}
it is \SIrange{14}{16}{\watt} against roughly \SI{80}{\watt} of propulsion, and
because it is constant it converts directly into battery mass, appearing in
\cref{sec:closure} as an addition to the fixed load rather than as a term
scaling with $m$.
\end{remark}

\section{The mass closure}
\label{sec:closure}

\subsection{Decomposition}

\begin{definition}[Mass budget]
The gross mass is partitioned as
\begin{equation}
  m = \mfix + \mstr + \mprop + \mbat ,
  \label{eq:budget}
\end{equation}
where $\mfix$ is the payload and avionics that do not scale with aircraft size
(display, driver electronics, computer, camera, gimbal), $\mstr$ is the
structure, $\mprop$ is motors, speed controllers and propellers, and $\mbat$ is
the battery pack.
\end{definition}

\Cref{eq:budget} is an identity. It becomes an equation, and a nontrivial one,
once the three terms on the right that depend on $m$ are made explicit. The
essential point is that they depend on $m$ in different ways, and lumping them
together, as one is tempted to do by writing a single ``structural fraction'',
destroys the structure of the problem.

\subsection{Propulsion mass: sized by peak thrust}

\begin{definition}[Thrust margin]
The propulsion system is sized so that the maximum total thrust is
\begin{equation}
  T_{\max} = \lambda\, m g, \qquad \lambda > 1 .
  \label{eq:lambda}
\end{equation}
$\lambda$ is the ratio of maximum thrust to weight; $\lambda = 2$ means the
vehicle hovers at half throttle.
\end{definition}

\begin{model}[Specific thrust]
\label{mod:sm}
Let $S_m \coloneqq T_{\max}/\mprop$ be the specific thrust of the complete
propulsion system, in \si{\newton\per\kilo\gram}. Then
\begin{equation}
  \mprop = \frac{T_{\max}}{S_m} = \frac{\lambda g}{S_m}\, m
         \eqqcolon \gamma_m\, m .
  \label{eq:mprop}
\end{equation}
\end{model}

Propulsion mass is therefore \emph{linear} in gross mass. This is worth
emphasising because it is the coupling most people worry about first, and it is
benign: it consumes a fixed fraction of the mass budget and produces a runaway
only in the degenerate case $\gamma_m \ge 1$.

\Cref{mod:sm} is a lumped model. \Cref{sec:components} replaces it with motor
mass derived from the torque the rotor demands, at which point $S_m$ becomes an
output; the engine reports both so the assumption can be audited.

\subsection{Structural mass}

\begin{model}[Structural fraction]
\label{mod:fs}
\begin{equation}
  \mstr = f_s\, m, \qquad f_s \in (0,1) .
  \label{eq:mstr}
\end{equation}
\end{model}

This is the weakest model in the algebraic layer. It is not a law; it is an
observation that across a class of similar vehicles the structure tends to a
roughly constant fraction of gross mass, typically $f_s \in [0.15,0.25]$. It is
retained because it keeps the closure in closed form, which is what makes
\cref{sec:fold} possible. \Cref{sec:components} derives $\mstr$ from beam theory
instead and the two are compared in \cref{sec:results}.

\subsection{Battery mass: sized by mission energy}

\begin{model}[Battery]
\label{mod:battery}
Let $\eb$ be the specific energy of the pack in \si{\joule\per\kilo\gram} that
the mission may use: the cell value reduced by the depth of discharge and by the
landing reserve $\varsigma$, so $\eb = e_{\mathrm{cell}}\,\mathrm{DoD}\,(1-\varsigma)$.
Where the reserve must be visible, as in \cref{sec:codesign}, it is written out
and $\eb$ denotes the value before the reserve is removed. For a mission of
duration $t_f$ flown at total electrical power $P_{\mathrm{tot}}$,
\begin{equation}
  \mbat = \frac{E_{\mathrm{mission}}}{\eb}
        = \frac{1}{\eb}\int_0^{t_f} P\big(x(t),u(t)\big)\dd t
  \;\xrightarrow[\text{hover}]{}\;
  \frac{P_{\mathrm{tot}}\,t_f}{\eb} .
  \label{eq:mbat}
\end{equation}
\end{model}

The integral form is the honest one and is what \cref{sec:codesign} uses; the
right-hand form is its evaluation under the assumption that the whole mission is
spent hovering, and is what makes the algebra tractable. The ratio between them
is measured in \cref{sec:results} and is between $1.00$ and $1.10$ depending on
the mission, which retrospectively justifies the substitution for sizing
purposes but not for endurance claims.

\subsection{The master equation}

Substituting \cref{eq:mprop,eq:mstr,eq:mbat,eq:ptot} into \cref{eq:budget}:
\begin{equation}
  m = \mfix + f_s m + \gamma_m m
      + \frac{t_f}{\eb}\left[
        \frac{(mg)^{3/2}}{\FM\,\eta\,\sqrt{2\rho A}} + \Paux \right] .
  \label{eq:closure-raw}
\end{equation}
Collecting terms gives the central equation of the algebraic layer.

\begin{definition}[Closure coefficients]
\label{def:coeffs}
\begin{align}
  \alpha &\coloneqq 1 - f_s - \frac{\lambda g}{S_m}
          = 1 - f_s - \gamma_m
          &&\text{(thrust coupling, dimensionless)}, \label{eq:alpha}\\
  \beta  &\coloneqq \frac{t_f\, g^{3/2}}{\eb\,\FM\,\eta\,\sqrt{2\rho A}}
          &&\text{(energy coupling, \si{\kilo\gram^{-1/2}})}, \label{eq:beta}\\
  \Mzero &\coloneqq \mfix + \frac{\Paux\, t_f}{\eb}
          &&\text{(effective fixed load, \si{\kilo\gram})} . \label{eq:M0}
\end{align}
\end{definition}

\begin{theorem}[The closure equation]
\label{thm:closure}
Under \cref{mod:disk,mod:kappa,mod:sm,mod:fs,mod:battery} with fixed disk area,
the gross mass of a feasible design satisfies
\begin{equation}
  \boxed{\;\alpha\, m - \beta\, m^{3/2} = \Mzero\;}
  \label{eq:closure}
\end{equation}
and conversely any positive root of \cref{eq:closure} determines a mass budget
consistent with \cref{eq:budget}.
\end{theorem}

\begin{proof}
Rearranging \cref{eq:closure-raw},
\[
  m\left(1 - f_s - \gamma_m\right)
  - \frac{t_f g^{3/2}}{\eb \FM \eta \sqrt{2\rho A}}\, m^{3/2}
  = \mfix + \frac{\Paux t_f}{\eb},
\]
where the $m^{3/2}$ term arises because
$(mg)^{3/2} = g^{3/2} m^{3/2}$. Identifying the three groups with
\cref{eq:alpha,eq:beta,eq:M0} gives \cref{eq:closure}. The converse holds
because every step is an equivalence for $m>0$.
\end{proof}

\begin{remark}[Reading the coefficients]
$\alpha$ is the fraction of gross mass left over after structure and propulsion
have taken their cut; it is a pure number, must be positive for any design to
exist, and is easy to keep positive. $\beta$ carries the mission: it is
proportional to endurance and inversely proportional to specific energy,
figure of merit, electrical efficiency and the square root of disk area. It
multiplies the nonlinear term, and it is where the design dies. $\Mzero$ is
what must be carried: note that the auxiliary power appears here, not in
$\beta$, because it is independent of aircraft mass, so a display that consumes
power converts directly into fixed load through the battery needed to run it.
\end{remark}

\begin{example}
For a 13.3-inch panel with $\mfix = \SI{0.762}{\kilo\gram}$,
$f_s = 0.20$, $\lambda = 2$, $S_m = \SI{120}{\newton\per\kilo\gram}$,
$\FM = 0.65$, $\eta = 0.75$, $\eb = \SI{900}{\kilo\joule\per\kilo\gram}$,
$t_f = \SI{3600}{\second}$ and four \SI{0.40}{\metre} rotors
($A = \SI{0.503}{\metre\squared}$), \cref{def:coeffs} gives
$\alpha = 0.6366$, $\beta = \SI{0.2271}{\kilo\gram^{-1/2}}$,
$\Mzero = \SI{0.762}{\kilo\gram}$. \Cref{sec:fold} shows that this
configuration has no solution.
\end{example}

\subsection{Generalised disk-area scaling}

\Cref{cor:hover} assumed the disk area fixed while mass varies. That is an
assumption about the design family, not a law. Suppose instead the rotors are
allowed to grow with the aircraft,
\begin{equation}
  A(m) = A_0 \left(\frac{m}{m_0}\right)^{p} ,
  \label{eq:areascaling}
\end{equation}
anchored so that $A(m_0) = A_0$ reproduces a stated rotor at a stated mass.

\begin{proposition}[Generalised power law]
\label{prop:plaw}
Under \cref{eq:areascaling} the hover propulsive power obeys
\begin{equation}
  \Pelec \propto m^{q}, \qquad q = \frac{3-p}{2} ,
  \label{eq:qexp}
\end{equation}
and the closure becomes
\begin{equation}
  \alpha m - \beta_q m^{q} = \Mzero,
  \qquad
  \beta_q = \frac{t_f\,g^{3/2}\,m_0^{p/2}}
                 {\eb\,\FM\,\eta\,\sqrt{2\rho A_0}} .
  \label{eq:closure-q}
\end{equation}
\end{proposition}

\begin{proof}
Substituting \cref{eq:areascaling} into \cref{eq:pelec} with $T=mg$,
\[
  \Pelec = \frac{(mg)^{3/2}}{\FM\eta\sqrt{2\rho A_0}}
           \left(\frac{m}{m_0}\right)^{-p/2}
         = \frac{g^{3/2} m_0^{p/2}}{\FM\eta\sqrt{2\rho A_0}}\; m^{(3-p)/2}.
\]
Multiplying by $t_f/\eb$ and inserting into \cref{eq:budget} as before gives
\cref{eq:closure-q}.
\end{proof}

Three cases matter and are treated in \cref{sec:fold,sec:scaling}:
\begin{itemize}[nosep]
  \item $p=0$: rotors fixed, $q=3/2$, the case of \cref{thm:closure};
  \item $p=1$: constant disk loading, $q=1$, the closure becomes linear;
  \item $p>1$: rotors grow faster than mass, $q<1$, the nonlinearity is
        sublinear and inverts its role entirely.
\end{itemize}

\begin{corollary}[Constant disk loading]
\label{cor:constdl}
If the rotors are sized to hold $\mathrm{DL} = mg/A$ constant, then by
\cref{eq:dl} the propulsive power is exactly linear in mass,
$\Pelec = c_P m$ with
\begin{equation}
  c_P = \frac{g}{\FM\,\eta}\sqrt{\frac{\mathrm{DL}}{2\rho}} ,
  \label{eq:cP}
\end{equation}
the battery mass fraction is $\gamma_b = c_P t_f/\eb$, and the closure
collapses to
\begin{equation}
  m = \frac{\Mzero}{1 - f_s - \gamma_m - \gamma_b} ,
  \label{eq:constdl}
\end{equation}
which has a positive solution if and only if
\begin{equation}
  f_s + \gamma_m + \gamma_b < 1 .
  \label{eq:constdlfeas}
\end{equation}
\end{corollary}

\begin{proof}
Immediate from \cref{eq:dl} with $T=mg$ and \cref{eq:closure-q} at $q=1$,
$\beta_1 = \gamma_b$.
\end{proof}

The contrast between \cref{eq:closure} and \cref{eq:constdl} is the central
structural fact of the sizing problem, and it is the subject of the next
section. In the first, feasibility fails discontinuously at a fold. In the
second, mass diverges continuously as the fractions approach unity. They are
different kinds of impossibility.

\section{The fold: existence, multiplicity and the sizing iteration}
\label{sec:fold}

This section analyses \cref{eq:closure-q} completely. The results are
elementary but they are the substance of the problem, and they are what
separates ``the design is heavy'' from ``the design does not exist''.

Throughout, fix $\alpha>0$, $\beta>0$, $\Mzero>0$ and $q>0$, and define
\begin{equation}
  F:(0,\infty)\to\R, \qquad F(m) \coloneqq \alpha m - \beta m^{q} .
  \label{eq:F}
\end{equation}
$F(m)$ is the fixed load that a machine of gross mass $m$ can support: gross
mass less what structure, propulsion and battery consume. Solving the closure
means solving $F(m)=\Mzero$.

\subsection{The trichotomy}

\begin{lemma}[Strict concavity]
\label{lem:concave}
If $q>1$ then $F$ is strictly concave on $(0,\infty)$, attains a unique
maximum at
\begin{equation}
  m^{*} = \left(\frac{\alpha}{q\beta}\right)^{\frac{1}{q-1}} ,
  \label{eq:mstar}
\end{equation}
and
\begin{equation}
  \Mmax \coloneqq F(m^{*}) = \alpha\, m^{*}\left(1 - \frac{1}{q}\right) .
  \label{eq:Mmax}
\end{equation}
Moreover $F(0^+)=0$, $F$ is strictly increasing on $(0,m^*)$, strictly
decreasing on $(m^*,\infty)$, and $F(m)\to-\infty$ as $m\to\infty$.
\end{lemma}

\begin{proof}
$F$ is smooth on $(0,\infty)$ with
\begin{equation}
  F'(m) = \alpha - q\beta m^{q-1},
  \qquad
  F''(m) = -q(q-1)\beta m^{q-2} .
  \label{eq:Fderiv}
\end{equation}
For $q>1$ and $m>0$ we have $q(q-1)\beta m^{q-2}>0$, hence $F''<0$ everywhere
on $(0,\infty)$: $F$ is strictly concave. Setting $F'(m)=0$ gives
$m^{q-1} = \alpha/(q\beta)$, and since $q-1>0$ the map $m\mapsto m^{q-1}$ is a
strictly increasing bijection of $(0,\infty)$ onto itself, so the stationary
point is unique and given by \cref{eq:mstar}. Strict concavity makes it the
global maximum. Evaluating,
\[
  F(m^*) = \alpha m^* - \beta (m^*)^{q}
         = \alpha m^* - \beta m^* (m^*)^{q-1}
         = \alpha m^* - \beta m^* \frac{\alpha}{q\beta}
         = \alpha m^*\left(1-\frac1q\right),
\]
which is \cref{eq:Mmax}. The sign of $F'$ follows from monotonicity of
$m^{q-1}$, and $F(m) = m(\alpha - \beta m^{q-1}) \to -\infty$ because the
bracket tends to $-\infty$.
\end{proof}

\begin{theorem}[Trichotomy for the closure]
\label{thm:trichotomy}
Let $q>1$, $\alpha,\beta,\Mzero>0$, and let $\Mmax$ be as in \cref{eq:Mmax}.
Then the equation $F(m)=\Mzero$ has, on $(0,\infty)$,
\begin{enumerate}[label=(\roman*),nosep]
  \item exactly two solutions $0 < m_- < m^* < m_+$ if $\Mzero < \Mmax$;
  \item exactly one solution $m = m^*$, a double root, if $\Mzero = \Mmax$;
  \item no solution if $\Mzero > \Mmax$.
\end{enumerate}
\end{theorem}

\begin{proof}
By \cref{lem:concave}, $F$ increases strictly from $F(0^+)=0$ to $F(m^*)=\Mmax$
on $(0,m^*]$ and decreases strictly from $\Mmax$ to $-\infty$ on
$[m^*,\infty)$. Both branches are continuous and strictly monotone.

(iii) If $\Mzero>\Mmax$ then $F(m)\le \Mmax < \Mzero$ for all $m>0$, so no
solution exists.

(ii) If $\Mzero=\Mmax$, then $F(m)=\Mzero$ forces $m=m^*$ by uniqueness of the
maximiser. It is a double root because $F'(m^*)=0$, so $F(m)-\Mzero$ has a zero
of order at least two; the order is exactly two since $F''(m^*)<0$.

(i) If $0<\Mzero<\Mmax$ then on $(0,m^*]$ the function $F$ is a strictly
increasing continuous bijection onto $(0,\Mmax]$, so there is exactly one
$m_-\in(0,m^*)$ with $F(m_-)=\Mzero$. On $[m^*,\infty)$, $F$ is a strictly
decreasing continuous bijection onto $(-\infty,\Mmax]$, so there is exactly one
$m_+\in(m^*,\infty)$ with $F(m_+)=\Mzero$. Strictness gives $m_-<m^*<m_+$ and
excludes further roots.
\end{proof}

\begin{corollary}[Closed forms at $q=3/2$]
\label{cor:closedform}
For the fixed-rotor case $q=3/2$,
\begin{equation}
  m^{*} = \frac{4\alpha^{2}}{9\beta^{2}},
  \qquad
  \Mmax = \frac{4\alpha^{3}}{27\beta^{2}} .
  \label{eq:closedform}
\end{equation}
\end{corollary}

\begin{proof}
Put $q=3/2$ in \cref{eq:mstar}: $1/(q-1)=2$ and $\alpha/(q\beta)
= 2\alpha/(3\beta)$, so $m^* = 4\alpha^2/(9\beta^2)$. Then \cref{eq:Mmax} gives
$\Mmax = \alpha m^*(1-2/3) = \alpha m^*/3 = 4\alpha^3/(27\beta^2)$.
\end{proof}

\begin{remark}[Which root is the design]
The physically meaningful root is $m_-$. Both roots satisfy the mass budget,
but $m_+$ is a machine that is heavier than necessary, carries a battery sized
for its own excess weight, and is, by \cref{thm:iteration} below, unstable
under the natural sizing iteration. It is a mathematical artefact of the
quadratic-like structure, not a design.
\end{remark}

\begin{remark}[Cubic reduction]
\Cref{eq:closure} at $q=3/2$ becomes a cubic under the substitution
$u=\sqrt m$: $\beta u^3 - \alpha u^2 + \Mzero = 0$. The product of its three
roots is $-\Mzero/\beta<0$, so exactly one root is negative and the other two
are the pair $\sqrt{m_\pm}$ of \cref{thm:trichotomy}, real precisely when the
discriminant is non-negative. The engine solves the equation by bracketing
rather than by radicals, for reasons given in \cref{sec:numerics}.
\end{remark}

\subsection{The transition is a saddle-node bifurcation}

Regard $\Mzero$ as a parameter and consider the one-parameter family of
equations $G(m;\Mzero) \coloneqq F(m)-\Mzero = 0$.

\begin{theorem}[Saddle-node]
\label{thm:bifurcation}
At $(m,\Mzero)=(m^*,\Mmax)$ the family $G$ undergoes a saddle-node
(fold) bifurcation: the defining conditions
\begin{equation}
  G = 0, \qquad \frac{\partial G}{\partial m} = 0
  \label{eq:sn1}
\end{equation}
hold, together with the nondegeneracy and transversality conditions
\begin{equation}
  \frac{\partial^2 G}{\partial m^2}(m^*,\Mmax) = -q(q-1)\beta (m^*)^{q-2} \ne 0,
  \qquad
  \frac{\partial G}{\partial \Mzero}(m^*,\Mmax) = -1 \ne 0 .
  \label{eq:sn2}
\end{equation}
Consequently, near the bifurcation the two roots behave as
\begin{equation}
  m_{\pm} = m^* \pm \sqrt{\frac{2(\Mmax-\Mzero)}{q(q-1)\beta (m^*)^{q-2}}}
            + O(\Mmax-\Mzero) ,
  \label{eq:sqrtlaw}
\end{equation}
so they approach one another with a square-root law and annihilate.
\end{theorem}

\begin{proof}
\Cref{eq:sn1} is \cref{lem:concave} together with the definition of $\Mmax$.
\Cref{eq:sn2} follows from \cref{eq:Fderiv} and $\partial G/\partial \Mzero=-1$.
These are exactly the hypotheses of the saddle-node normal form theorem
(see e.g.\ Guckenheimer and Holmes~\cite[\S3.4]{guckenheimer1983} or
Kuznetsov~\cite[\S3.2]{kuznetsov2004}), which asserts the existence of a smooth
change of coordinates taking $G$ to $\dot y = \mu - y^2$ near the point.
For \cref{eq:sqrtlaw}, Taylor-expand $F$ about $m^*$:
\[
  \Mzero = F(m) = \Mmax + \tfrac12 F''(m^*)(m-m^*)^2 + O((m-m^*)^3),
\]
using $F'(m^*)=0$. Solving for $m-m^*$ and inserting
$F''(m^*)=-q(q-1)\beta(m^*)^{q-2}$ gives the stated expansion.
\end{proof}

\begin{remark}[Why this matters practically]
\Cref{eq:sqrtlaw} says that near the boundary of feasibility the design mass is
an infinitely sensitive function of the requirement: $\dd m_-/\dd \Mzero \to
\infty$ as $\Mzero\uparrow\Mmax$. A design placed close to the fold is not
merely marginal; its mass is unbounded in its derivative with respect to every
assumption. This is a stronger statement than ``leave margin'' and is the
reason the engine reports the utilisation $\Mzero/\Mmax$ as a first-class
quantity.
\end{remark}

\subsection{The naive sizing iteration}

The obvious way to size such a vehicle is to guess a mass, compute what it
implies, and repeat:
\begin{equation}
  m_{k+1} = \Phi(m_k),
  \qquad
  \Phi(m) \coloneqq \Mzero + f_s m + \gamma_m m + \beta m^{q} .
  \label{eq:iteration}
\end{equation}
$\Phi$ is exactly \cref{eq:budget} read as an assignment. Its fixed points are
the roots of \cref{eq:closure-q}, since
$\Phi(m)=m \iff F(m)=\Mzero$.

\begin{theorem}[Convergence of the sizing iteration]
\label{thm:iteration}
Let $q>1$ and $\Mzero<\Mmax$, with roots $m_-<m^*<m_+$ as in
\cref{thm:trichotomy}. Then
\begin{enumerate}[label=(\roman*),nosep]
  \item $\Phi'(m) = 1 - F'(m)$, so $\Phi'(m_-)<1$ and $\Phi'(m_+)>1$;
  \item $m_-$ is an attracting fixed point of \cref{eq:iteration} with
        asymptotic linear rate $\Phi'(m_-) = 1-F'(m_-) \in (0,1)$, and $m_+$ is
        repelling;
  \item the iteration started from any $m_0 \in [0,m_+)$ converges
        monotonically upward to $m_-$ if $m_0<m_-$, and monotonically downward
        to $m_-$ if $m_0\in(m_-,m_+)$;
  \item for $m_0>m_+$ the iterates increase without bound;
  \item if $\Mzero>\Mmax$ the iterates increase without bound from every
        starting point.
\end{enumerate}
\end{theorem}

\begin{proof}
(i) Differentiating \cref{eq:iteration},
$\Phi'(m) = f_s+\gamma_m + q\beta m^{q-1}$. Using
$f_s+\gamma_m = 1-\alpha$ from \cref{eq:alpha} and
$F'(m)=\alpha - q\beta m^{q-1}$ from \cref{eq:Fderiv},
\[
  \Phi'(m) = 1-\alpha + q\beta m^{q-1} = 1 - F'(m).
\]
By \cref{lem:concave}, $F'(m_-)>0$ and $F'(m_+)<0$, giving $\Phi'(m_-)<1$ and
$\Phi'(m_+)>1$. Also $\Phi'>0$ everywhere since $f_s,\gamma_m,\beta\ge0$ and not
all zero, so $\Phi$ is strictly increasing; and $\Phi'(m_-) = 1-F'(m_-) > 0$
because $F'(m_-) \le \alpha < 1$.

(ii) Standard: a fixed point of a $C^1$ map is attracting when
$|\Phi'|<1$ and repelling when $|\Phi'|>1$, with linear rate $|\Phi'|$.

(iii) $\Phi$ is increasing and $\Phi(m)-m = \Mzero-F(m)$, which is negative
exactly on $(m_-,m_+)$ and positive on $[0,m_-)$. Hence for $m_0<m_-$ the
sequence is increasing and bounded above by $m_-$ (because $\Phi$ increasing
and $\Phi(m_-)=m_-$ imply $m_k<m_-\Rightarrow m_{k+1}<m_-$), so it converges,
necessarily to a fixed point in $[m_0,m_-]$, which must be $m_-$. The argument
on $(m_-,m_+)$ is symmetric with the sequence decreasing.

(iv) For $m_0>m_+$, $\Phi(m)-m = \Mzero - F(m) > 0$ since $F(m)<\Mzero$ there,
so the sequence increases; it cannot converge because there is no fixed point
above $m_+$; hence it diverges.

(v) If $\Mzero>\Mmax$ then $\Mzero-F(m)>0$ for every $m$, so
$m_{k+1}>m_k$ always and there is no fixed point to converge to.
\end{proof}

\begin{corollary}[Local stability criterion]
\label{cor:dmdm}
A design is a stable fixed point of the sizing iteration if and only if
\begin{equation}
  \frac{\dd}{\dd m}\Big[\mstr(m)+\mprop(m)+\mbat(m)\Big] < 1 ,
  \label{eq:dmdm}
\end{equation}
which for the model of \cref{sec:closure} reads
$f_s + \gamma_m + q\beta m^{q-1} < 1$, equivalently $F'(m)>0$, equivalently
$m<m^*$.
\end{corollary}

\Cref{cor:dmdm} is the precise form of the intuition that ``each extra kilogram
must require less than a kilogram of extra machine to carry it''. The engine
evaluates \cref{eq:dmdm} numerically at the solved design point and reports it.

\begin{remark}[Critical slowing down]
By \cref{thm:iteration}(ii) the convergence rate is $1-F'(m_-)$, and
$F'(m_-)\to 0$ as $\Mzero\uparrow\Mmax$. The iteration therefore becomes
arbitrarily slow near the fold: near-infeasible designs are not distinguished
from infeasible ones by a fixed iteration budget. This is a real trap. A
practical consequence is that a solver which reports ``did not converge in $K$
iterations'' cannot be used to decide feasibility; the analytic criterion
$\Mzero \lessgtr \Mmax$ must be used instead. The engine does exactly this and
classifies slow convergence separately from divergence.
\end{remark}

\subsection{Behaviour when \texorpdfstring{$q\le1$}{q <= 1}}

\begin{theorem}[No fold for $q\le1$]
\label{thm:pfold}
Let $\alpha,\beta,\Mzero>0$.
\begin{enumerate}[label=(\roman*),nosep]
  \item If $q=1$ and $\alpha>\beta$, then $F(m)=\Mzero$ has the unique solution
        $m=\Mzero/(\alpha-\beta)$; if $\alpha\le\beta$ there is no solution and
        $F\le 0$.
  \item If $0<q<1$, then $F(m)=\Mzero$ has exactly one solution for every
        $\Mzero>0$. In particular no critical payload exists.
\end{enumerate}
Consequently, under the disk-area scaling \cref{eq:areascaling} a fold exists
if and only if $p<1$.
\end{theorem}

\begin{proof}
(i) $F(m)=(\alpha-\beta)m$ is linear; the claim is immediate.

(ii) For $0<q<1$, $F'(m)=\alpha-q\beta m^{q-1}$ with $q-1<0$, so
$m^{q-1}\to\infty$ as $m\to0^+$ and $m^{q-1}\to0$ as $m\to\infty$. Hence
$F'(m)<0$ for small $m$ and $F'(m)>0$ for large $m$, with a single sign change
at $m_{\min}=(q\beta/\alpha)^{1/(1-q)}$. Thus $F$ decreases from $F(0^+)=0$ to
a negative minimum and then increases strictly to $+\infty$ (since
$F(m)=m(\alpha-\beta m^{q-1})$ and the bracket tends to $\alpha>0$). On
$[m_{\min},\infty)$, $F$ is a strictly increasing continuous bijection onto
$[F(m_{\min}),\infty)\supset(0,\infty)$, so there is exactly one root with
$F=\Mzero>0$, and none on $(0,m_{\min})$ where $F\le0<\Mzero$.

The final statement follows from $q=(3-p)/2$: $q>1 \iff p<1$.
\end{proof}

\begin{finding}
The catastrophe is not a property of flight. It is a property of holding the
rotor area fixed while the mass grows. Let the disk grow at least as fast as
the mass and the fold disappears entirely, replaced by the far gentler
condition \cref{eq:constdlfeas}. What is then unbounded is not the mass but the
rotor diameter, so the binding constraint migrates from energy to geometry,
noise and safety. It does not go away. Nor does the absence of a fold mean the
absence of a ceiling: at $p=1$ exactly the longest mission that closes is still
finite (\cref{prop:twall}); only for $p>1$ does it disappear.
\end{finding}

\section{Scaling laws and sensitivity}
\label{sec:scaling}

\Cref{cor:closedform} gives the critical payload in closed form. Because it is
a product of powers of the underlying parameters, its logarithmic derivatives
are exact rational numbers, and those numbers are the design guidance.

\subsection{The explicit critical payload}

\begin{proposition}[Explicit form of $\Mmax$]
\label{prop:mmaxexplicit}
For $q=3/2$,
\begin{equation}
  \Mmax
  = \frac{4\alpha^{3}}{27\beta^{2}}
  = \frac{8\,\alpha^{3}\,\rho\, \Adisk\, \eb^{2}\,\FM^{2}\,\eta^{2}}
         {27\, g^{3}\, t_f^{2}} .
  \label{eq:mmaxexplicit}
\end{equation}
\end{proposition}

\begin{proof}
From \cref{eq:beta}, $\beta^2 = t_f^2 g^3/(\eb^2\FM^2\eta^2\cdot 2\rho A)$.
Hence $1/\beta^2 = 2\rho A \eb^2\FM^2\eta^2/(t_f^2 g^3)$ and
$\Mmax = (4\alpha^3/27)\cdot 2\rho A\eb^2\FM^2\eta^2/(t_f^2g^3)$, which is
\cref{eq:mmaxexplicit}.
\end{proof}

\subsection{Exact logarithmic sensitivities}

\begin{definition}
For a positive quantity $Y$ depending on a positive parameter $x$, the
logarithmic sensitivity is
$S_x \coloneqq \partial \ln Y/\partial \ln x$, so that a relative change
$\delta x/x$ produces $S_x\,\delta x/x$ to first order.
\end{definition}

\begin{theorem}[Sensitivities of the critical payload]
\label{thm:sensitivity}
With $\Mmax$ as in \cref{eq:mmaxexplicit} and
$\gamma_m = \lambda g/S_m$, $\alpha = 1-f_s-\gamma_m$,
\begin{align}
  S_{t_f} &= -2, &
  S_{\eb} &= +2, &
  S_{\FM} &= +2, &
  S_{\eta} &= +2, \label{eq:sens1}\\
  S_{\Adisk} &= +1, &
  S_{D} &= +2, &
  S_{N} &= +1, &
  S_{\rho} &= +1, \label{eq:sens2}
\end{align}
and for the parameters entering through $\alpha$,
\begin{align}
  S_{f_s} &= -\frac{3 f_s}{\alpha}, &
  S_{\lambda} &= -\frac{3\gamma_m}{\alpha}, &
  S_{S_m} &= +\frac{3\gamma_m}{\alpha}, &
  S_{g} &= -3 - \frac{3\gamma_m}{\alpha} . \label{eq:sens3}
\end{align}
\end{theorem}

\begin{proof}
Take logarithms of \cref{eq:mmaxexplicit}:
\[
  \ln \Mmax = \ln\tfrac{8}{27} + 3\ln\alpha + \ln\rho + \ln A
              + 2\ln \eb + 2\ln\FM + 2\ln\eta - 3\ln g - 2\ln t_f .
\]
For any parameter $x$ not appearing in $\alpha$, differentiation is immediate
and gives \cref{eq:sens1,eq:sens2}; the entries for $D$ and $N$ follow from
$A = N\pi D^2/4$, so $\partial\ln A/\partial\ln D = 2$ and
$\partial\ln A/\partial \ln N = 1$.

For $f_s$: $\partial\alpha/\partial f_s = -1$, so
$\partial \ln\Mmax/\partial\ln f_s = 3 f_s\,\partial\ln\alpha/\partial f_s
= -3f_s/\alpha$.

For $\lambda$: $\alpha = 1-f_s-\lambda g/S_m$, so
$\partial\alpha/\partial\lambda = -g/S_m$ and
$\partial\ln\Mmax/\partial\ln\lambda
= 3\lambda(-g/S_m)/\alpha = -3\gamma_m/\alpha$. The result for $S_m$ follows
identically with the opposite sign since $\gamma_m\propto S_m^{-1}$.

For $g$: gravity appears both explicitly, as $g^{-3}$, and inside $\alpha$
through $\gamma_m = \lambda g/S_m$. Hence
\[
  S_g = -3 + 3\frac{g}{\alpha}\frac{\partial \alpha}{\partial g}
      = -3 + \frac{3g}{\alpha}\left(-\frac{\lambda}{S_m}\right)
      = -3 - \frac{3\gamma_m}{\alpha}. \qedhere
\]
\end{proof}

\begin{corollary}[The inverse-square endurance law]
\label{cor:tsquared}
At fixed technology and geometry, $\Mmax \propto t_f^{-2}$. Multiplying the
required endurance by $k$ divides the supportable fixed load by $k^2$.
\end{corollary}

\Cref{fig:walltf} shows the law for the lumped model of the design of
\cref{sec:results}, and what it implies once the auxiliary load is counted.

\begin{figure}[tbp]
\centering
% Lumped fixed-area closure of the design of sec17 (gen_b.wall_tf):
% the critical load M0max(t_f) falls as t_f^{-2}; the load the mission demands,
% M0(t_f) = m_fix + P_aux t_f / e_b, rises. They cross at the endurance wall.
\begin{tikzpicture}
\begin{axis}[paperplot, height=6.2cm, xmode=log, ymode=log,
  xlabel={mission duration $t_f$ [\si{\minute}]},
  ylabel={load [\si{\kilo\gram}]},
  xmin=5, xmax=480, ymin=0.01, ymax=150, legend pos=north east,
  xtick={5,10,20,30,60,120,240,480},
  xticklabels={5,10,20,30,60,120,240,480}]
  \addplot[cA, name path=wall] table[x=tmin, y=Mmax] {results/fig/b_walltf.dat};
  \addlegendentry{critical load $\Mmax(t_f)\propto t_f^{-2}$}
  \addplot[cB, name path=need] table[x=tmin, y=M0] {results/fig/b_walltf.dat};
  \addlegendentry{required load $\Mzero(t_f)=\mfix+\Paux t_f/\eb$}
  \path[name intersections={of=wall and need, name=w}];
  \fill[black] (w-1) circle (2.2pt);
  \draw[construct] (w-1) -- (w-1 |- {axis cs:5,0.01});
  \node[lbl, anchor=south west, align=left] at (w-1) {endurance wall\\of the lumped\\fixed-area closure};
  % slope triangle
  \draw[cG] (axis cs:10,0.5) -- (axis cs:40,0.5) -- (axis cs:10,8) -- cycle;
  \node[lbl, cG, anchor=north] at (axis cs:20,0.5) {$\times4$};
  \node[lbl, cG, anchor=east] at (axis cs:10,2) {$\times16$};
  \node[lbl, cG, anchor=west] at (axis cs:10.5,1.2) {slope $-2$};
  \draw[cG, densely dotted, line width=0.7pt] (axis cs:30,0.01) -- (axis cs:30,150);
  \node[lbl, cG, anchor=west] at (axis cs:30,0.02) {design $t_f$};
\end{axis}
\end{tikzpicture}

\caption{The inverse-square law of \cref{cor:tsquared} for the lumped
fixed-area model of the design of \cref{sec:results}. The critical load
$\Mmax$ falls with slope $-2$ on logarithmic axes. The load the mission itself
demands rises, because the battery that runs the auxiliary load,
$\Paux t_f/\eb$, is part of $\Mzero$. The two cross at the endurance wall of
this closure, \SI{66}{\minute}, against the design's \SI{30}{\minute} session;
beyond the crossing no mass closes.}
\label{fig:walltf}
\end{figure}

\begin{example}[Verification against the engine]
For the parameter set of \cref{sec:results} the engine computes the
sensitivities numerically by central differences in $\ln$-space and returns
$S_{t_f}=-2.000$, $S_{\eb}=S_{\FM}=S_{\eta}=+2.000$, $S_D=+2.000$, $S_N=+1.000$,
$S_\rho=+1.000$, $S_{f_s}=-0.943$, $S_{S_m}=+0.770$, $S_\lambda=-0.770$ and
$S_g=-3.770$. With $\gamma_m = 2\times 9.80665/120 = 0.16344$ and
$\alpha=0.63656$, \cref{eq:sens3} predicts
$S_{f_s} = -3(0.20)/0.63656 = -0.9426$,
$S_{S_m} = 3(0.16344)/0.63656 = 0.7703$ and
$S_g = -3-0.7703 = -3.7703$. The agreement to the printed precision is a
verification of both the closed form and its implementation. See
\cref{sec:verification}.
\end{example}

\begin{remark}[How to read \cref{thm:sensitivity}]
Three groups of parameters behave quite differently.

The \emph{squared} levers, $\eb$, $\FM$, $\eta$ and $t_f$, are the ones with
real authority, and their sign structure is unforgiving: a 20\% better cell
buys 44\% more payload, once; a doubling of endurance costs three quarters of
it. The figure of merit and the electrical efficiency are bounded above by
unity, so at $\FM=0.73$ and $\eta=0.78$ there remains at most a factor
$(1/0.73)^2(1/0.78)^2 \approx 3.1$ in that direction, ever.

The \emph{linear} levers are disk area and air density. Disk area is the only
one of the four that the designer controls freely, and because $D$ enters
squared it is in practice the strongest available lever.

The \emph{cubic-in-$\alpha$} levers, $f_s$, $\lambda$, $S_m$, have modest
sensitivity as long as $\alpha$ is comfortably positive, and become violent as
$\alpha\to0$. They are not where the design is won or lost unless one has been
careless.
\end{remark}

\Cref{fig:sensitivity} collects the sensitivities as the engine computes them.

\begin{figure}[tbp]
\centering
% Logarithmic sensitivities of M0max computed by the engine
% (sizing.sensitivities at the default parameters, central differences in
% log space). Order written by gen_b.sensitivities.
\begin{tikzpicture}
\begin{axis}[paperplot, height=6.6cm, width=0.8\linewidth,
  xbar, bar width=7.5pt, bar shift=0pt, y dir=reverse,
  ytick={0,...,10},
  yticklabels={$t_f$,$\eb$,$\FM$,$\eta$,$D$,$N$,$\rho$,$f_s$,$\lambda$,$S_m$,$g$},
  xmin=-4.6, xmax=3.1, ymin=-0.6, ymax=10.6,
  xlabel={$S_x=\partial\ln\Mmax/\partial\ln x$},
  grid=none, xmajorgrids, xtick={-4,-3,...,3},
  nodes near coords, point meta=x,
  every node near coord/.append style={font=\scriptsize,
    /pgf/number format/fixed, /pgf/number format/precision=2,
    /pgf/number format/fixed zerofill}]
  \addplot[fill=cA!65, draw=cA] table[x=S, y=i, restrict x to domain=0:10]
    {results/fig/b_sens.dat};
  \addplot[fill=cD!55, draw=cD] table[x=S, y=i, restrict x to domain=-10:0]
    {results/fig/b_sens.dat};
  \draw[cG] (axis cs:0,-0.6) -- (axis cs:0,10.6);
  \draw[cG, dashed] (axis cs:-4.6,3.5) -- (axis cs:3.1,3.5);
  \draw[cG, dashed] (axis cs:-4.6,6.5) -- (axis cs:3.1,6.5);
  \node[lbl, cG, anchor=west] at (axis cs:-4.5,1.5) {squared};
  \node[lbl, cG, anchor=west] at (axis cs:-4.5,5.0) {disk area, air};
  \node[lbl, cG, anchor=west, align=left] at (axis cs:-4.5,7.6) {through $\alpha$};
\end{axis}
\end{tikzpicture}

\caption{Logarithmic sensitivities of $\Mmax$ at the engine's default
parameters, computed by central differences in $\ln$-space. The top group are
the squared levers of \cref{eq:sens1}; the middle group are disk area, entering
as $D^2$ and $N$, and air density, \cref{eq:sens2}; the bottom group act
through $\alpha$ and are \cref{eq:sens3}, here with $\gamma_m/\alpha=0.257$.
Every bar agrees with \cref{thm:sensitivity} to three decimals. A bar of length
two means a \SI{10}{\percent} change in the parameter moves the critical load
by about \SI{20}{\percent}.}
\label{fig:sensitivity}
\end{figure}

\subsection{The exponent, the constants, and where the wall sits}

A claim one meets in discussions of this problem is that ``the exponents kill
you, not the constants''. \Cref{thm:sensitivity} lets us be precise about the
sense in which that is and is not true.

\begin{observation}
\label{obs:expconst}
The exponent $q$ determines whether a wall exists at all
(\cref{thm:pfold}: it exists iff $q>1$) and the asymptotic behaviour of the
closure. The constants $\alpha$ and $\beta$ determine where the wall sits,
through \cref{eq:closedform}, and their influence is not weak: $\Mmax$ is
cubic in $\alpha$ and inverse-quadratic in $\beta$. The correct statement is
that the exponent creates the wall and the constants locate it.
\end{observation}

For finite engineering at finite scales the location is what one actually
cares about. \Cref{eq:mmaxexplicit} says that location moves by a factor of two
for a 41\% improvement in cell energy density, or a 41\% improvement in figure
of merit, or a doubling of disk area, or a 29\% reduction in endurance.

\subsection{The endurance wall across the exponent}

Combining \cref{thm:pfold} with \cref{prop:plaw} gives the full picture, but
only once two statements that are easily conflated are separated: whether the
closure has a fold, and whether there is a longest mission it can close. They
are different, and they part company at $p=1$.

Write the closure with its dependence on duration explicit. Under
\cref{eq:areascaling} the propulsive power is $c_P m^q$ with
$c_P = c_0 m_0^{p/2}$ and $c_0 = g^{3/2}/(\FM\,\eta\sqrt{2\rho A_0})$, so a mass
$m$ closes a mission of duration $t$ if and only if
$\alpha m - (t/\eb)(c_P m^q + \Paux) = \mfix$, that is if and only if
$t = \Theta_p(m)$ with
\begin{equation}
  \Theta_p(m) \coloneqq
  \frac{\eb\,(\alpha m - \mfix)}{c_P\, m^{q} + \Paux} .
  \label{eq:theta}
\end{equation}
Define the endurance wall at fixed $\mfix$, $\Paux$ and anchor $(m_0,A_0)$ as
$t_f^{\max}(p) \coloneqq \sup_{m>0}\Theta_p(m)$.

\begin{proposition}[The endurance wall]
\label{prop:twall}
Let $\alpha>0$ and $\mfix>0$.
\begin{enumerate}[label=(\roman*),nosep]
  \item If $p<1$, $\Theta_p$ attains its maximum at a finite mass
        $\hat m(p) > \mfix/\alpha$, and $t_f^{\max}(p)=\Theta_p(\hat m)$ is
        finite and attained.
  \item If $p=1$, $\Theta_1$ is strictly increasing and
        \begin{equation}
          t_f^{\max}(1) = t_{\lim} = \frac{\alpha\,\eb}{c_P} ,
          \label{eq:tlim}
        \end{equation}
        which is finite and not attained: for every $t_f<t_{\lim}$ the unique
        mass is finite, and it diverges as $t_f\uparrow t_{\lim}$. Here $\eb$
        is net of the landing reserve, as in \cref{mod:battery}.
  \item If $p>1$, $\Theta_p(m)\to\infty$ as $m\to\infty$ and there is no wall.
  \item Where $\hat m$ is differentiable in $p$,
        \begin{equation}
          \frac{\dd\ln t_f^{\max}}{\dd p}
          = \frac{\vartheta}{2}\,\ln\frac{\hat m}{m_0} ,
          \qquad
          \vartheta = \frac{c_P\hat m^{q}}{c_P \hat m^{q}+\Paux}\in(0,1],
          \label{eq:wallslope}
        \end{equation}
        so the wall rises with $p$ if and only if the mass at the wall exceeds
        the anchor mass. It is not monotone in $p$ for every anchor.
\end{enumerate}
\end{proposition}

\begin{proof}
$\Theta_p$ is continuous, negative for $m<\mfix/\alpha$ and positive beyond.
(i) For $q>1$ the numerator grows like $m$ and the denominator like $m^q$, so
$\Theta_p(m)\to0$ as $m\to\infty$; a continuous function that is positive on
$(\mfix/\alpha,\infty)$ and tends to $0$ at both ends of that interval attains
its supremum. (ii) At $q=1$,
$\Theta_1'(m)$ has the sign of
$\alpha(c_P m+\Paux) - c_P(\alpha m-\mfix) = \alpha\Paux + c_P\mfix>0$, and
$\Theta_1(m)\to\alpha\eb/c_P$. The mass closing duration $t$ is
$m = (\mfix + \Paux t/\eb)/(\alpha - c_P t/\eb)$, positive and finite exactly
when $t<t_{\lim}$; with $\Paux=0$ this is \cref{eq:constdlfeas}. (iii) For
$q<1$ the numerator outgrows the denominator. (iv) By the envelope theorem,
$\partial_m\Theta_p(\hat m)=0$ gives
$\dd\ln t_f^{\max}/\dd p = \partial_p\ln\Theta_p(\hat m)
= -\partial_p \ln(c_P \hat m^q + \Paux)$, and with
$\partial_p\ln c_P = \tfrac12\ln m_0$ and $\partial_p q = -\tfrac12$ this is
$-\vartheta(\tfrac12\ln m_0 - \tfrac12\ln\hat m)$.
\end{proof}

\begin{remark}
With $\Paux=0$ the maximiser is explicit: $\hat m = q\,\mfix/[\alpha(q-1)]$,
which tends to infinity as $q\downarrow1$, reconciling (i) with (ii). If the
anchor is the fold mass of the fixed-area closure at the nominal duration $t_f$,
then $\sqrt{m_0} = 2\alpha/(3\beta)$ and \cref{eq:tlim} gives exactly
$t_{\lim} = \tfrac32 t_f$, independently of $\Paux$.
\end{remark}

\Cref{fig:theta} draws $\Theta_p$ for five exponents. The maxima of (i), the
unattained ceiling of (ii) and the unbounded growth of (iii) are all visible,
and the maxima lie to the right of the anchor, which by (iv) is why the wall
rises with $p$ here.

\begin{figure}[tbp]
\centering
% Theta_p(m) / t_f against m / m0 at the engine defaults, anchor m0 = fold
% mass of the fixed-area closure (gen_b.theta). For p < 1 the curve has an
% interior maximum (marked); at p = 1 it rises to alpha e_b / c_P = (3/2) t_f
% without reaching it; for p > 1 it grows without bound.
\begin{tikzpicture}
\begin{axis}[paperplot, height=6.4cm, xmode=log,
  xlabel={gross mass relative to the anchor, $m/m_0$},
  ylabel={longest closed mission $\Theta_p(m)/t_f$},
  xmin=0.3, xmax=300, ymin=0, ymax=2.1, unbounded coords=jump,
  legend style={at={(0.5,1.03)}, anchor=south, font=\scriptsize},
  legend columns=4]
  \addplot[cA] table[x=x, y=p0] {results/fig/b_theta.dat};
  \addlegendentry{$p=0$}
  \addplot[cC] table[x=x, y=p1] {results/fig/b_theta.dat};
  \addlegendentry{$p=0.5$}
  \addplot[cE] table[x=x, y=p2] {results/fig/b_theta.dat};
  \addlegendentry{$p=0.75$}
  \addplot[cB] table[x=x, y=p3] {results/fig/b_theta.dat};
  \addlegendentry{$p=1$}
  \addplot[cD] table[x=x, y=p4] {results/fig/b_theta.dat};
  \addlegendentry{$p=1.25$}
  \addplot[cB, dashed, line width=0.7pt] table[x=x, y=t] {results/fig/b_theta_lim.dat};
  \addlegendentry{$\alpha\eb/(c_Pt_f)=\tfrac32$}
  \addplot[only marks, mark=*, mark size=2pt, black] table[x=x, y=t]
    {results/fig/b_theta_peaks.dat};
  \addlegendentry{$\hat m(p)$, $t_f^{\max}(p)$}
  \draw[construct] (axis cs:1,0) -- (axis cs:1,2.1);
  \node[lbl, cG, anchor=south west, rotate=90] at (axis cs:0.97,0.05) {anchor $m_0$};
  \draw[cG, densely dotted] (axis cs:0.3,1) -- (axis cs:300,1);
  \node[lbl, cG, anchor=south east] at (axis cs:290,1.0) {nominal $t_f$};
\end{axis}
\end{tikzpicture}

\caption{The longest mission a mass $m$ can close, $\Theta_p(m)$ of
\cref{eq:theta}, relative to the nominal $t_f$, at the engine's default
parameters with the anchor $m_0$ at the fold mass of the fixed-area closure.
For $p<1$ each curve has an interior maximum $(\hat m,t_f^{\max})$ (dots), and
these lie to the right of the anchor. At $p=1$ the curve rises towards
$\alpha\eb/c_P=\tfrac32t_f$ (dashed) without reaching it; at $p=1.25$ it grows
without bound. No fold and no ceiling are different statements: $p=1$ has the
first property and not the second.}
\label{fig:theta}
\end{figure}

\Cref{tab:pwall} records the computed values and \cref{fig:pwall} draws them
on a finer grid. The default parameters do not close at fixed area at the
nominal hour, so the common anchor is the fold mass, and the $p=1$ entry is
$\tfrac32 t_f$ as the remark predicts.

\begin{figure}[tbp]
\centering
% Endurance wall against the disk-area exponent (gen_b.pwall), same
% parameters and anchor as the table. (a) t_max / t_f: the curve is the
% direct maximisation of Theta_p, the crosses the bisection of
% sweep.endurance_wall; the open circle at p = 1 is the unattained
% supremum 3/2. (b) the wall mass relative to the anchor; the wall rises
% with p exactly where this exceeds one.
\begin{tikzpicture}
\begin{groupplot}[group style={group size=2 by 1, horizontal sep=1.6cm},
  halfplot, height=5.8cm, width=0.46\linewidth,
  xlabel={disk-area exponent $p$}, xmin=-0.25, xmax=1.5]
\nextgroupplot[ylabel={$t_f^{\max}/t_f$}, ymin=0.8, ymax=1.75,
  title={(a) the endurance wall}, unbounded coords=discard]
  \fill[cD!8] (axis cs:1,0.8) rectangle (axis cs:1.5,1.75);
  \node[lbl, cD, align=center] at (axis cs:1.25,1.25) {no wall:\\$\Theta_p\to\infty$};
  \addplot[cA] table[x=p, y=tdirect] {results/fig/b_pwall.dat};
  \addplot[only marks, mark=x, mark size=2.2pt, cB] table[x=p, y=t]
    {results/fig/b_pwall.dat};
  \addplot[only marks, mark=o, mark size=2.6pt, cA, line width=1pt] table[x=p, y=t]
    {results/fig/b_pwall_lim.dat};
  \draw[cB, dashed] (axis cs:0.55,1.5) -- (axis cs:1,1.5);
  \node[lbl, anchor=south east] at (axis cs:0.98,1.52) {$t_{\lim}=\tfrac32t_f$};
\nextgroupplot[ylabel={$\hat m/m_0$}, ymode=log, ymin=0.5, ymax=2000,
  title={(b) the mass at the wall}, xmax=1.0]
  \addplot[cA] table[x=p, y=x] {results/fig/b_pwall.dat};
  \draw[cD, dashed] (axis cs:-0.25,1) -- (axis cs:1,1);
  \node[lbl, cD, anchor=south east] at (axis cs:0.95,1.05) {$\hat m=m_0$};
  \node[lbl, anchor=north west, align=left] at (axis cs:-0.2,900)
    {wall rises with $p$\\where $\hat m>m_0$};
\end{groupplot}
\end{tikzpicture}

\caption{The endurance wall of \cref{prop:twall} on a fine grid of exponents,
at the parameters and anchor of \cref{tab:pwall}. (a) $t_f^{\max}/t_f$ from
direct maximisation of \cref{eq:theta} (line) and from the bisection of
\code{sweep.endurance\_wall} (crosses), which agree wherever the bisection
runs; within about $0.01$ of $p=1$ the bisection overflows in the fold mass
$(\alpha/q\beta)^{1/(q-1)}$ and only the direct value is shown. The circle at
$p=1$ is the unattained supremum $\tfrac32t_f$ of \cref{eq:tlim}; for $p>1$
there is no wall. (b) The mass at the wall exceeds the anchor for every $p<1$
and diverges as $p\uparrow1$, so by \cref{eq:wallslope} the wall rises
monotonically on this family.}
\label{fig:pwall}
\end{figure}

\begin{table}[t]
\centering
\caption{Endurance wall against the disk-area exponent at the engine's default
parameters ($t_f=\SI{\rPwallTf}{\second}$), anchored at the fold mass
$m_0=\SI{\rPwallAnchor}{\kilo\gram}$ of the fixed-area closure. Computed by bisection on
$t_f$ with \code{sweep.endurance\_wall} and checked against direct maximisation
of \cref{eq:theta}. The last column is $\hat m/m_0$; it exceeds one throughout,
so by \cref{eq:wallslope} the wall rises with $p$ for this anchor. At $p=1$ there
is no fold but there is still a ceiling, the unattained supremum
\cref{eq:tlim}.}
\label{tab:pwall}
\begin{tabular}{@{}rrlrr@{}}
\toprule
$p$ & $q=(3-p)/2$ & fold exists & $t_f^{\max}$ & $\hat m/m_0$ \\
\midrule
\inputrows{results/tab_pwall_rows}
\end{tabular}
\end{table}

\section{First-principles component models}
\label{sec:components}

\Cref{mod:sm,mod:fs} introduced two constants, $S_m$ and $f_s$, that carry a
great deal of weight and are not laws. This section replaces both with models
derived from geometry and materials, at the cost of losing the closed form. The
engine keeps both layers and reports the implied values of $S_m$, $f_s$ and
$\FM$ so that the lumped assumptions can be audited against the detailed model.

\subsection{Arms: a cantilevered thin-walled tube}

\begin{model}[Rotor arm]
\label{mod:arm}
Each of the $N$ arms is a thin-walled circular tube of mean radius $r$, wall
thickness $t\ll r$, length $L$, elastic modulus $E$, density $\rho_m$ and
allowable stress $\sigma_{\mathrm{all}}$, built in at the hub and carrying the
rotor as a point load at the tip.
\end{model}

\paragraph{Length.} Adjacent rotor disks must not overlap. With $N$ hubs on a
circle of radius $L$, the centre-to-centre spacing of adjacent hubs is
$2L\sin(\pi/N)$, so non-overlap requires $2L\sin(\pi/N)\ge 2R$, i.e.
\begin{equation}
  L = \frac{k_{\mathrm{gap}}\,R}{\sin(\pi/N)}, \qquad k_{\mathrm{gap}}\ge1 ,
  \label{eq:armlength}
\end{equation}
with $k_{\mathrm{gap}}$ a clearance factor. The tip-to-tip span is
$2(L+R)$ and the \emph{reach}, the distance from the vehicle centre to the
nearest blade tip, is
\begin{equation}
  \mathcal{R} \coloneqq L + R .
  \label{eq:reach}
\end{equation}
\Cref{eq:reach} is trivial and turns out to be one of the most consequential
quantities in the paper; see \cref{sec:results}.

\paragraph{Section properties.} For a thin-walled tube the second moment of
area about a diameter is
\begin{equation}
  I = \pi r^3 t + O(t^2), \qquad
  \text{mass per unit length } \mu = 2\pi r t \rho_m .
  \label{eq:tube}
\end{equation}

\paragraph{Strength.} Under the limit load $F = \mathrm{SF}\,\lambda m g/N$ at
the tip, the root bending moment is $M=FL$ and the extreme-fibre stress is
$\sigma = Mr/I$. Requiring $\sigma\le\sigma_{\mathrm{all}}$ gives
\begin{equation}
  t \ge t_{\sigma} = \frac{F L}{\pi r^2 \sigma_{\mathrm{all}}} .
  \label{eq:tstress}
\end{equation}

\paragraph{Stiffness.} Euler-Bernoulli theory gives the tip deflection of a
cantilever under a tip load as $\delta = FL^3/(3EI)$. Imposing
$\delta\le L/\Lambda$ for a slenderness limit $\Lambda$ (we use $\Lambda=100$)
gives
\begin{equation}
  t \ge t_{\delta} = \frac{\Lambda F L^{2}}{3\pi E r^{3}} .
  \label{eq:tstiff}
\end{equation}

\paragraph{Dynamics.} The fundamental bending frequency of a uniform cantilever
of mass per unit length $\mu$ is
\begin{equation}
  f_1 = \frac{1.875^2}{2\pi}\sqrt{\frac{EI}{\mu_{\mathrm{eff}} L^4}},
  \qquad
  \mu_{\mathrm{eff}} = \mu + \frac{3 m_{\mathrm{tip}}}{L} ,
  \label{eq:freq}
\end{equation}
where $1.875$ is the first root of $\cos\zeta\cosh\zeta+1=0$ and the correction
to $\mu$ is the Rayleigh lumped-mass approximation for a tip mass
$m_{\mathrm{tip}}$ (motor plus propeller). The wall thickness actually used is
\begin{equation}
  t = \max\{t_\sigma,\; t_\delta,\; t_{\min}\} ,
  \label{eq:wall}
\end{equation}
and the engine records which of the three governs, because the answer is
informative: for the designs of \cref{sec:results} it is stiffness or minimum
gauge, never strength, with a working stress of \SI{44}{\mega\pascal} against
an allowable of \SI{350}{\mega\pascal}.

\begin{remark}[Consistency with \cref{as:rigid}]
\Cref{eq:freq} is what makes the rigid-body assumption checkable. With
$f_1\approx\SI{32}{\hertz}$ and an attitude loop at
$\sqrt{K_R}/2\pi\approx\SI{1.0}{\hertz}$ (\cref{sec:control}) the separation is
a factor of $32$, comfortably in the regime where treating the airframe as rigid
does not corrupt the control design.
\end{remark}

\paragraph{Frame total.} Guards are rings around each disk, so their mass scales
with circumference, not area:
\begin{equation}
  m_{\mathrm{guard}} = k_g\, N\, 2\pi R .
  \label{eq:guards}
\end{equation}
The frame is then
\begin{equation}
  \mstr = N\mu L\,(1+k_{\mathrm{hub}}) + m_{\mathrm{guard}}
          + m_{\mathrm{boom}} + m_{\mathrm{fixed}} ,
  \label{eq:frame}
\end{equation}
with $k_{\mathrm{hub}}$ accounting for the centre plate, motor mounts and
fasteners, $m_{\mathrm{boom}}$ the optional display cantilever of
\cref{sec:results}, and $m_{\mathrm{fixed}}$ landing gear and wiring.

\subsection{Motors: sized by torque, with a scale effect}

The rotor demands a torque $Q_{\max} = \Pshaft^{\max}/\Omega_{\max}$ at the
sizing condition $T_{\max}=\lambda mg/N$ per rotor. A motor's continuous torque
is set by the shear stress its airgap can sustain over its rotor surface,
$Q \sim \tau_{\mathrm{airgap}}\,(\pi D_m L_m)\,(D_m/2)$, which for
geometrically similar machines gives $Q\propto D_m^3 \propto m_{\mathrm{mot}}$,
i.e.\ constant torque density. Real small outrunners do better than that as
they grow, because the ratio of cooling surface to loss volume improves and
because the winding fill factor rises. We therefore write

\begin{model}[Motor mass]
\label{mod:motor}
\begin{equation}
  Q = k_\tau \left(\frac{m_{\mathrm{mot}}}{m_{\mathrm{ref}}}\right)^{\!e}
      m_{\mathrm{mot}}
  \quad\Longleftrightarrow\quad
  m_{\mathrm{mot}}
  = \left(\frac{Q\, m_{\mathrm{ref}}^{e}}{k_\tau}\right)^{\!\frac{1}{1+e}} ,
  \label{eq:motormass}
\end{equation}
with $m_{\mathrm{ref}}=\SI{0.1}{\kilo\gram}$, $e\approx0.30$ and $k_\tau$
calibrated in \cref{sec:verification}.
\end{model}

Because the motor is sized by the burst condition $T_{\max}$, the peak rather
than continuous torque is the right basis. \Cref{eq:motormass} is sublinear in
$Q$, which matters: at $e=0.3$ a tenfold increase in torque costs a factor
$10^{1/1.3}\approx5.9$ in mass, not ten.

Speed controllers are sized by power, $m_{\mathrm{esc}} = \Pelec^{\max}/k_P$,
and propellers by an empirical areal law $m_{\mathrm{prop}} = k_p D^{2.6}
(N_b/2)$ fitted to carbon propellers.

\begin{remark}[The electrical chain that mass does not see]
\Cref{eq:motormass} determines the size of the motor but not whether it can be
driven. A rotor turning at $\Omega$ requires a back-EMF $\Omega/K_V$ below the
pack voltage with margin for the current-resistance drop, and the motor
constant satisfies $K_t\,[\si{\newton\metre\per\ampere}] = 9.5493/K_V$
$[\si{rpm\per\volt}]$. Sizing $K_V$ so that maximum rotor speed is reached at
the \emph{minimum} pack voltage, not the nominal or full one, is what keeps the
machine flyable at the bottom of a discharge. This constraint is decoupled from
the mass closure and is treated separately in the engine; it is what forces the
low-$K_V$, custom-wound motors of \cref{sec:results}.
\end{remark}

\subsection{Battery: two constraints, not one}

\Cref{mod:battery} sizes the pack by energy. A pack must also deliver power.
Writing $p_b$ for the usable specific power and $\varsigma$ for the state of
charge held in reserve,
\begin{equation}
  \mbat = \max\left\{
    \underbrace{\frac{E_{\mathrm{mission}}}{\eb\,(1-\varsigma)}}
      _{\text{energy}},\;
    \underbrace{\frac{\Pelec^{\max}}{p_b}}_{\text{power}}
  \right\} .
  \label{eq:batmax}
\end{equation}
Which term binds depends on the mission: for short sessions the power term
dominates and endurance is effectively free, while for long sessions the energy
term takes over. The engine reports which is active, because a design that is
power-limited responds to entirely different levers than one that is
energy-limited.

High-energy lithium-ion cells are not high-power cells. Sizing against the
burst condition $T_{\max}$ requires the burst rating; the continuous hover draw
must then be checked separately against the continuous rating. Both checks are
implemented.

\subsection{Cost of the first-principles layer}

Replacing \cref{mod:sm,mod:fs} by \cref{mod:arm,mod:motor} removes the closed
form: $\mstr(m)$ and $\mprop(m)$ are no longer proportional to $m$, so
\cref{eq:closure} is replaced by the implicit equation
\begin{equation}
  \mathcal{F}(m) \coloneqq m - \Big[\mfix + \mstr(m) + \mprop(m) + \mbat(m)\Big]
  = 0 ,
  \label{eq:fpclosure}
\end{equation}
solved numerically. For the designs studied $\mathcal{F}$ has the qualitative
shape of \cref{sec:fold}: negative for small $m$, positive on the branch where a
design exists, negative again beyond the fold. That shape is observed, not
proved. \Cref{lem:concave} is a statement about the lumped residual; it does
not transfer to $\mathcal{F}$, which contains clipped quantities, maxima
(\cref{eq:batmax}) and empirical fits. Near the fold the positive stretch can
also be very narrow, and \cref{sec:numerics} describes how the engine searches
for it and what it reports when it does not find it. A failed search is not a
proof that no mass closes.

\section{Acoustics}
\label{sec:acoustics}

For a machine that hovers beside a person's head, the acoustic model is not a
secondary consideration; \cref{sec:results} finds it to be the binding
constraint. It is also the place where the most error was found, so it is
developed carefully.

\subsection{Why the sixth power}

\begin{model}[Broadband loading noise]
\label{mod:acoustic}
At the scales and Reynolds numbers considered, rotor noise is dominated by
broadband loading noise: unsteady lift on the blade sections, arising from
inflow turbulence and boundary-layer vortex shedding, radiating as a
distribution of acoustic dipoles.
\end{model}

Curle's extension of Lighthill's acoustic analogy to flows bounded by solid
surfaces~\cite{lighthill1952,curle1955} shows that the surface term is dipole
in character, and dimensional analysis of the dipole source strength gives the
acoustic power radiated by a body of characteristic area $A_b$ in a flow of
characteristic velocity $U$:
\begin{equation}
  W_{\mathrm{ac}} \;\sim\; \frac{\rho\, A_b\, U^{6}}{c_0^{3}} ,
  \label{eq:dipole}
\end{equation}
with $c_0$ the speed of sound. This is the classical $U^6$ dipole law, against
$U^8$ for the quadrupole (free jet) case; see Ffowcs Williams and
Hawkings~\cite{fwh1969} for the rotating-surface generalisation, and
Brooks, Pope and Marcolini~\cite{brooks1989} for the airfoil self-noise
mechanisms this lumps together.

Taking $10\log_{10}$ of \cref{eq:dipole} with $U=\Vtip$ gives a
$60\log_{10}\Vtip$ dependence. Loading enters through the sectional lift, for
which we take an empirical $10\log_{10}T$, and $N$ statistically independent
rotors add incoherently, contributing $10\log_{10}N$. Hence

\begin{model}[Source model]
\label{mod:spl}
\begin{equation}
  L_p(\SI{1}{\metre})
  = C + 60\log_{10}\!\frac{\Vtip}{\Vtip^{\mathrm{ref}}}
      + 10\log_{10}\!\frac{T_{\mathrm{rotor}}}{T^{\mathrm{ref}}}
      + 10\log_{10} N ,
  \label{eq:splmodel}
\end{equation}
with $\Vtip^{\mathrm{ref}}=\SI{100}{\metre\per\second}$,
$T^{\mathrm{ref}}=\SI{5}{\newton}$ and $C$ an anchor calibrated against
measured aircraft in \cref{sec:verification}.
\end{model}

\begin{observation}[The design lever]
\label{obs:sixthpower}
Halving tip speed reduces the level by $60\log_{10}2 = \SI{18.1}{\decibel}$.
No other parameter in the model has comparable authority. Together with
\cref{cor:designrule}, which says that raising solidity lowers the tip speed
needed to carry a given thrust, this determines the entire shape of the
machine: large, fat-bladed, slowly turning rotors.
\end{observation}

\begin{figure}[tbp]
\centering
% Generated by paper/figures/gen_c.py. Do not edit by hand.
\def\cM{1.724}
\def\cMbat{0.288}
\def\cL{0.3827}
\def\cR{0.2460}
\def\cReach{0.629}
\def\cSpan{1.257}
\def\cKgap{1.100}
\def\cGap{0.0492}
\def\cSoverD{0.100}
\def\cKint{1.0252}
\def\cRcp{0.130}
\def\cCoMz{0.0275}
\def\cScreenOffz{0.1025}
\def\cScreenW{0.345}
\def\cScreenH{0.194}
\def\cArmr{8.0}
\def\cWall{0.66}
\def\cWallStress{0.083}
\def\cWallStiff{0.660}
\def\cWallMin{0.60}
\def\cDriver{stiffness}
\def\cSigmaWork{44}
\def\cSigmaAllow{350}
\def\cFbend{33}
\def\cFlimit{15.2}
\def\cZrotor{1.37}
\def\cZceil{1.03}
\def\cZoverR{5.57}
\def\cZcOverR{4.19}
\def\cKground{1.0020}
\def\cKceiling{1.0072}
\def\cKc{2.0}
\def\cCeilPct{0.7}
\def\cGroundPct{0.20}
\def\cKcSpanG{13}
\def\cCeiling{2.4}
\def\cBpf{34.5}
\def\cRpm{1036}
\def\cNb{2}
\def\cAbpf{-37.5}
\def\cLwa{61.8}
\def\cRroom{20.8}
\def\cRc{0.91}
\def\cStandoff{1.20}
\def\cSplEar{56.6}
\def\cSplOne{53.8}
\def\cSplFar{54.8}
\def\cTrotor{4.23}
\def\cVtip{26.7}
\def\cAnchor{82.9}
\def\cVi{3.01}
\def\cTfCross{21}
\def\cTf{30}
\def\cQmax{0.289}
\def\cMmot{61}
\def\cKtau{5.5}
\def\cEmot{0.30}
\def\cRoomL{5}
\def\cRoomW{4}
\def\cAlphaRoom{0.20}
%
% The source model at the design's thrust per rotor: 60 dB per decade of tip
% speed, with the design and a camera-drone tip speed marked.
\begin{tikzpicture}
  \begin{semilogxaxis}[paperplot, height=5.6cm,
      xlabel={tip speed $\Vtip$ [\si{\metre\per\second}]},
      ylabel={$L_p(\SI{1}{\metre})$ [\si{\decibel A}]},
      xmin=10, xmax=150, ymin=25, ymax=95,
      xtick={10,20,50,100,150}, xticklabels={10,20,50,100,150},
      legend pos=north west]
    \addplot[cA, domain=10:150, samples=60]
      {\cAnchor + 60*log10(x/100) + 10*log10(\cTrotor/5) + 10*log10(4)};
    \addlegendentry{\cref{eq:splmodel}, $T_{\mathrm{rotor}}=\SI{\cTrotor}{\newton}$, $N=4$}
    % comfort bands
    \draw[cC, dashed, line width=0.6pt] (axis cs:10,45) -- (axis cs:150,45);
    \node[lbl, cC, anchor=south east] at (axis cs:150,45) {quiet office};
    \draw[cD, dashed, line width=0.6pt] (axis cs:10,60) -- (axis cs:150,60);
    \node[lbl, cD, anchor=south east] at (axis cs:150,60) {speech};
    % halving the tip speed
    \pgfmathsetmacro{\Ldes}{\cAnchor + 60*log10(\cVtip/100) + 10*log10(\cTrotor/5) + 10*log10(4)}
    \pgfmathsetmacro{\Vtwo}{2*\cVtip}
    \pgfmathsetmacro{\Ltwo}{\Ldes + 60*log10(2)}
    \draw[cG, line width=0.5pt] (axis cs:\cVtip,\Ldes) -- (axis cs:\Vtwo,\Ldes) -- (axis cs:\Vtwo,\Ltwo);
    \node[lbl, cG, anchor=west] at (axis cs:\Vtwo,0.5*\Ldes+0.5*\Ltwo) {\SI{18.1}{\decibel}};
    \node[lbl, cG, anchor=north] at (axis cs:1.41*\cVtip,\Ldes) {$\times2$};
    \addplot[only marks, mark=*, mark size=2.3pt, cA] coordinates {(\cVtip,\Ldes)};
    \node[lbl, cA, anchor=north west] at (axis cs:\cVtip,\Ldes-1) {design};
    \pgfmathsetmacro{\Ldrone}{\cAnchor + 60*log10(68/100) + 10*log10(\cTrotor/5) + 10*log10(4)}
    \addplot[only marks, mark=square*, mark size=2.2pt, cB] coordinates {(68,\Ldrone)};
    \node[lbl, cB, anchor=south east, align=right] at (axis cs:66,\Ldrone+1)
      {camera drone,\\same thrust};
  \end{semilogxaxis}
\end{tikzpicture}

\caption{The source model \cref{eq:splmodel} at the design's thrust per rotor.
On a logarithmic tip-speed axis the level is a straight line of
\SI{60}{\decibel} per decade, so halving the tip speed buys \SI{18.1}{\decibel}
wherever one starts. The design turns at \SI{\cVtip}{\metre\per\second};
carrying the same thrust at a camera drone's tip speed would cost about
\SI{24}{\decibel}, the difference between a quiet office and a raised voice.}
\label{fig:sixth}
\end{figure}

\Cref{fig:sixth} draws \cref{eq:splmodel} at the design's thrust.

\subsection{Sound power against sound pressure}

\begin{definition}
The sound power level $L_W$ of a source is
$L_W = 10\log_{10}(W/W_0)$, $W_0=\SI{1e-12}{\watt}$, a property of the source
alone. The sound pressure level $L_p$ at a point is
$L_p = 20\log_{10}(p_{\mathrm{rms}}/p_0)$, $p_0=\SI{20}{\micro\pascal}$, and
depends on the environment.
\end{definition}

\begin{proposition}[Free-field conversion]
\label{prop:lwlp}
For a point source of directivity factor $Q$ radiating into free space above a
rigid plane, the mean-square pressure at distance $r$ satisfies
\begin{equation}
  L_p(r) = L_W + 10\log_{10}\!\left(\frac{Q}{4\pi r^2}\right) ,
  \label{eq:lwlp}
\end{equation}
so for a source above a reflecting plane, $Q=2$, and at $r=\SI{1}{\metre}$,
\begin{equation}
  L_p(\SI{1}{\metre}) = L_W - 10\log_{10}(2\pi) = L_W - \SI{7.98}{\decibel} .
  \label{eq:eightdb}
\end{equation}
\end{proposition}

\begin{proof}
Acoustic intensity at distance $r$ for a source of power $W$ radiating with
directivity $Q$ is $I = QW/(4\pi r^2)$. In the far field of a plane wave
$I = p_{\mathrm{rms}}^2/(\rho c_0)$, and the reference quantities are chosen so
that $10\log_{10}(I/I_0) = L_p$ with $I_0 = p_0^2/(\rho c_0)$ to within
\SI{0.1}{\decibel} at standard conditions. Taking logarithms gives
\cref{eq:lwlp}; substituting $Q=2$, $r=1$ gives \cref{eq:eightdb}.
\end{proof}

\Cref{fig:hemisphere} is the geometry behind $Q=2$.

\begin{figure}[tbp]
\centering
% A source on a reflecting floor radiates into a hemisphere of area 2 pi r^2,
% which is the Q = 2 of eq:lwlp and the source of the 8 dB between a declared
% sound power and a pressure at one metre.
\begin{tikzpicture}[>={Stealth[length=1.8mm]}, font=\footnotesize]
  \fill[pattern=north east lines, pattern color=cG] (-3.2,-0.25) rectangle (3.2,0);
  \draw[line width=0.8pt] (-3.2,0) -- (3.2,0);
  \node[anchor=north east, cG] at (3.2,-0.28) {reflecting plane};
  \foreach \r/\op in {1.0/70, 1.8/50, 2.6/32}
    \draw[cA!\op!white, line width=0.9pt] (\r,0) arc (0:180:\r);
  \foreach \ang in {20,55,90,125,160}
    \draw[vec, cA] (\ang:0.22) -- (\ang:0.82);
  \fill[cD] (0,0) circle (0.07);
  \node[anchor=north east, cD] at (-0.1,-0.28) {source of power $W$};
  \draw[dim, black] (0,0) -- (35:2.6) node[pos=0.66, above left=-2pt] {$r$};
  \node[anchor=south, cA] at (0,2.62) {wavefront, area $2\pi r^{2}$};
  \node[anchor=west, align=left] at (3.2,1.55)
    {$I=\dfrac{W}{2\pi r^{2}}=\dfrac{QW}{4\pi r^{2}},\quad Q=2$\\[6pt]
     $L_p(r)=L_W+10\log_{10}\dfrac{Q}{4\pi r^{2}}$\\[6pt]
     $L_p(\SI{1}{\metre})=L_W-10\log_{10}2\pi$\\[2pt]
     $\phantom{L_p(\SI{1}{\metre})}=L_W-\SI{7.98}{\decibel}$};
\end{tikzpicture}

\caption{Sound power against sound pressure, \cref{prop:lwlp}. A source on a
reflecting plane spreads its power $W$ over hemispherical wavefronts of area
$2\pi r^2$, which is the directivity $Q=2$ of \cref{eq:lwlp}. At one metre that
area is $2\pi\,\si{\metre\squared}$, and $10\log_{10}2\pi=\SI{7.98}{\decibel}$
is the whole difference between a declared sound power level and the pressure
level a listener at one metre hears.}
\label{fig:hemisphere}
\end{figure}

\begin{remark}[A trap worth naming]
European regulation 2019/945 requires manufacturers of unmanned aircraft to
declare a \emph{sound power} level $L_{WA}$. It is routinely misread as a
pressure at one metre. By \cref{prop:lwlp} that misreading is worth
\SI{8}{\decibel}, and in the wrong direction: it makes published aircraft sound
louder than they are, and so makes a design calibrated against the misreading
look quieter than it is. \Cref{sec:verification} records the consequences.
\end{remark}

\subsection{The room}

Free field is the wrong model indoors. The classical treatment
(Sabine; see Beranek~\cite{beranek1988} or Kuttruff~\cite{kuttruff2009}) writes
the steady-state level as the sum of a direct field, which falls with distance,
and a diffuse reverberant field, which does not.

\begin{definition}[Room constant]
For a room of total interior surface area $S$ and average absorption
coefficient $\bar\alpha$,
\begin{equation}
  \mathcal{R}_{\mathrm{room}} \coloneqq \frac{S\,\bar\alpha}{1-\bar\alpha} .
  \label{eq:roomconstant}
\end{equation}
\end{definition}

\begin{proposition}[Steady-state level in a room]
\label{prop:room}
\begin{equation}
  L_p(r) = L_W + 10\log_{10}\!\left(
    \frac{Q}{4\pi r^{2}} + \frac{4}{\mathcal{R}_{\mathrm{room}}}\right) .
  \label{eq:roomlevel}
\end{equation}
The two contributions are equal at the critical distance
\begin{equation}
  r_c = \sqrt{\frac{Q\,\mathcal{R}_{\mathrm{room}}}{16\pi}} .
  \label{eq:critdist}
\end{equation}
\end{proposition}

\begin{proof}[Sketch]
The direct term is \cref{eq:lwlp}. For the reverberant term, in the diffuse
field approximation the energy density is uniform and in steady state the power
absorbed equals the power injected: $W = \tfrac14 \bar\alpha S\, c_0\, E$ with
$E$ the energy density, whence $E = 4W/(c_0 S\bar\alpha)$; correcting for the
fact that the first reflection has already occurred replaces $S\bar\alpha$ by
$\mathcal{R}_{\mathrm{room}}$, and $p_{\mathrm{rms}}^2 = \rho c_0^2 E$ gives the $4/\mathcal{R}_{\mathrm{room}}$
term. Setting the two terms equal and solving for $r$ gives
\cref{eq:critdist}.
\end{proof}

\begin{finding}
\Cref{prop:room} has a consequence that no free-field figure can express:
beyond $r_c$, moving away from the machine buys nothing. The level is then set
by how absorbent the room is, not by how far away one stands. A flying display
does not make a quiet place near itself and a loud place far away; it raises
the level of the entire room to within a few decibels of the level at the user's
own ear.
\end{finding}

For an ordinary furnished room, $\bar\alpha\approx0.20$ and $r_c$ is under a
metre, which is smaller than the standoff at which such a machine is useful.
The user therefore sits close to the boundary between the two regimes and
receives both contributions (\cref{fig:room}).

\begin{figure}[tbp]
\centering
% Generated by paper/figures/gen_c.py. Do not edit by hand.
\def\cM{1.724}
\def\cMbat{0.288}
\def\cL{0.3827}
\def\cR{0.2460}
\def\cReach{0.629}
\def\cSpan{1.257}
\def\cKgap{1.100}
\def\cGap{0.0492}
\def\cSoverD{0.100}
\def\cKint{1.0252}
\def\cRcp{0.130}
\def\cCoMz{0.0275}
\def\cScreenOffz{0.1025}
\def\cScreenW{0.345}
\def\cScreenH{0.194}
\def\cArmr{8.0}
\def\cWall{0.66}
\def\cWallStress{0.083}
\def\cWallStiff{0.660}
\def\cWallMin{0.60}
\def\cDriver{stiffness}
\def\cSigmaWork{44}
\def\cSigmaAllow{350}
\def\cFbend{33}
\def\cFlimit{15.2}
\def\cZrotor{1.37}
\def\cZceil{1.03}
\def\cZoverR{5.57}
\def\cZcOverR{4.19}
\def\cKground{1.0020}
\def\cKceiling{1.0072}
\def\cKc{2.0}
\def\cCeilPct{0.7}
\def\cGroundPct{0.20}
\def\cKcSpanG{13}
\def\cCeiling{2.4}
\def\cBpf{34.5}
\def\cRpm{1036}
\def\cNb{2}
\def\cAbpf{-37.5}
\def\cLwa{61.8}
\def\cRroom{20.8}
\def\cRc{0.91}
\def\cStandoff{1.20}
\def\cSplEar{56.6}
\def\cSplOne{53.8}
\def\cSplFar{54.8}
\def\cTrotor{4.23}
\def\cVtip{26.7}
\def\cAnchor{82.9}
\def\cVi{3.01}
\def\cTfCross{21}
\def\cTf{30}
\def\cQmax{0.289}
\def\cMmot{61}
\def\cKtau{5.5}
\def\cEmot{0.30}
\def\cRoomL{5}
\def\cRoomW{4}
\def\cAlphaRoom{0.20}
%
% Eq:roomlevel for the design in the reference room: direct field, diffuse
% reverberant field, their sum, the critical distance and the user's ear.
\begin{tikzpicture}
  \begin{semilogxaxis}[paperplot, height=5.8cm,
      xlabel={distance from the machine $r$ [\si{\metre}]},
      ylabel={$L_p$ [\si{\decibel A}]},
      xmin=0.1, xmax=8, ymin=40, ymax=78,
      xtick={0.1,0.2,0.5,1,2,5}, xticklabels={0.1,0.2,0.5,1,2,5},
      legend pos=north east]
    \addplot[cB, dashed] table[x=r, y=direct] {results/fig/c_room.dat};
    \addlegendentry{direct, $L_W+10\log_{10}(Q/4\pi r^{2})$}
    \addplot[cE, dashed] table[x=r, y=reverb] {results/fig/c_room.dat};
    \addlegendentry{reverberant, $L_W+10\log_{10}(4/\mathcal{R}_{\mathrm{room}})$}
    \addplot[cA, line width=1.4pt] table[x=r, y=total] {results/fig/c_room.dat};
    \addlegendentry{total, \cref{eq:roomlevel}}
    \draw[construct] (axis cs:\cRc,40) -- (axis cs:\cRc,78);
    \node[lbl, cG, anchor=south east] at (axis cs:0.88,40.8) {$r_c=\SI{\cRc}{\metre}$};
    \addplot[only marks, mark=*, mark size=2.3pt, cA] coordinates {(\cStandoff,\cSplEar)};
    \node[lbl, cA, anchor=south west] at (axis cs:\cStandoff,\cSplEar+0.4)
      {user's ear, \SI{\cSplEar}{\decibel A}};
    \addplot[only marks, mark=o, mark size=2.3pt, cB] coordinates {(1,\cSplOne)};
    \draw[cB, line width=0.4pt] (axis cs:0.98,\cSplOne-0.3) -- (axis cs:0.82,51.2);
    \node[lbl, cB, anchor=north east] at (axis cs:0.84,51.4)
      {free field at \SI{1}{\metre}: \SI{\cSplOne}{\decibel A}};
    \node[lbl, cE, anchor=south, align=center] at (axis cs:4.2,54.9)
      {rest of the room};
  \end{semilogxaxis}
\end{tikzpicture}

\caption{\Cref{eq:roomlevel} for the design in the reference
$\cRoomL\times\cRoomW\times\SI{\cCeiling}{\metre}$ room with
$\bar\alpha=\cAlphaRoom$. The direct field falls by \SI{6}{\decibel} per
doubling of distance, the reverberant field does not fall at all, and the two
are equal at the critical distance $r_c=\SI{\cRc}{\metre}$ of
\cref{eq:critdist}. The user at \SI{\cStandoff}{\metre} is just past it and
hears \SI{\cSplEar}{\decibel A}, \SI{2.8}{\decibel} above the free-field figure
at one metre; the rest of the room settles at the reverberant level of about
\SI{\cSplFar}{\decibel A}.}
\label{fig:room}
\end{figure}

\subsection{Weighting and tonal content}

The A-weighting of IEC 61672 is
\begin{align}
  R_A(f) &= \frac{12194^2 f^4}
    {\big(f^2+20.6^2\big)\sqrt{\big(f^2+107.7^2\big)\big(f^2+737.9^2\big)}
     \big(f^2+12194^2\big)}, \nonumber\\
  A(f) &= 20\log_{10}R_A(f) + 2.00 .
  \label{eq:aweight}
\end{align}

The blade passage frequency is
\begin{equation}
  f_{\mathrm{BPF}} = \frac{N_b\,\Omega}{2\pi} = \frac{N_b\,\mathrm{rpm}}{60} .
  \label{eq:bpf}
\end{equation}

\begin{remark}[What the broadband number omits]
\Cref{eq:splmodel} predicts a broadband level. What makes a rotor
\emph{annoying} rather than merely loud is tonal content at
$f_{\mathrm{BPF}}$ and its harmonics. For the designs of \cref{sec:results},
$f_{\mathrm{BPF}}\approx\SI{\rBpf}{\hertz}$, which \cref{eq:aweight} discounts by
about \SI{38}{\decibel} (\cref{fig:aweight}). That is a genuine perceptual benefit of turning slowly
and it is also a warning: the acoustic energy has not gone anywhere. It has
moved to low frequency, where rooms absorb poorly, where Sabine's diffuse-field
assumption fails because the modal density is low, and where humans report a
throb rather than a hiss. Furthermore several rotors at slightly different
speeds beat at the difference frequency, which is perceptually worse than a
steady tone. None of this is in the A-weighted broadband figure, and it is the
most likely way for a built prototype to be unacceptable at a level the model
calls acceptable.
\end{remark}

The Schroeder frequency, above which a room behaves diffusely,
\begin{equation}
  f_S \approx 2000\sqrt{\frac{T_{60}}{V}} ,
  \label{eq:schroeder}
\end{equation}
is of order \SI{200}{\hertz} for a domestic room, which is precisely where the
first several blade harmonics lie. \Cref{prop:room} is therefore being applied
slightly outside its domain of validity for exactly the components that matter
most perceptually. This is stated as a limitation rather than repaired.

\begin{figure}[tbp]
\centering
% Generated by paper/figures/gen_c.py. Do not edit by hand.
\def\cM{1.724}
\def\cMbat{0.288}
\def\cL{0.3827}
\def\cR{0.2460}
\def\cReach{0.629}
\def\cSpan{1.257}
\def\cKgap{1.100}
\def\cGap{0.0492}
\def\cSoverD{0.100}
\def\cKint{1.0252}
\def\cRcp{0.130}
\def\cCoMz{0.0275}
\def\cScreenOffz{0.1025}
\def\cScreenW{0.345}
\def\cScreenH{0.194}
\def\cArmr{8.0}
\def\cWall{0.66}
\def\cWallStress{0.083}
\def\cWallStiff{0.660}
\def\cWallMin{0.60}
\def\cDriver{stiffness}
\def\cSigmaWork{44}
\def\cSigmaAllow{350}
\def\cFbend{33}
\def\cFlimit{15.2}
\def\cZrotor{1.37}
\def\cZceil{1.03}
\def\cZoverR{5.57}
\def\cZcOverR{4.19}
\def\cKground{1.0020}
\def\cKceiling{1.0072}
\def\cKc{2.0}
\def\cCeilPct{0.7}
\def\cGroundPct{0.20}
\def\cKcSpanG{13}
\def\cCeiling{2.4}
\def\cBpf{34.5}
\def\cRpm{1036}
\def\cNb{2}
\def\cAbpf{-37.5}
\def\cLwa{61.8}
\def\cRroom{20.8}
\def\cRc{0.91}
\def\cStandoff{1.20}
\def\cSplEar{56.6}
\def\cSplOne{53.8}
\def\cSplFar{54.8}
\def\cTrotor{4.23}
\def\cVtip{26.7}
\def\cAnchor{82.9}
\def\cVi{3.01}
\def\cTfCross{21}
\def\cTf{30}
\def\cQmax{0.289}
\def\cMmot{61}
\def\cKtau{5.5}
\def\cEmot{0.30}
\def\cRoomL{5}
\def\cRoomW{4}
\def\cAlphaRoom{0.20}
%
% The A-weighting of eq:aweight with the design's blade passage frequency
% and its harmonics, and the band below the Schroeder frequency where the
% diffuse-field room model does not hold.
\begin{tikzpicture}
  \begin{semilogxaxis}[paperplot, height=5.6cm,
      xlabel={frequency $f$ [\si{\hertz}]},
      ylabel={$A(f)$ [\si{\decibel}]},
      xmin=10, xmax=20000, ymin=-72, ymax=8,
      xtick={10,20,50,100,200,500,1000,2000,5000,10000,20000},
      xticklabels={10,,50,100,,500,1k,,5k,10k,},
      legend pos=south east]
    \fill[cD!8] (axis cs:10,-72) rectangle (axis cs:200,8);
    \node[lbl, cD, align=center, anchor=north] at (axis cs:45,6)
      {room not diffuse,\\$f<f_S\approx\SI{200}{\hertz}$};
    \addplot[cA, line width=1.2pt] table[x=f, y=A] {results/fig/c_aweight.dat};
    \addlegendentry{A-weighting, IEC 61672}
    \addplot[only marks, mark=*, mark size=2.2pt, cB] table[x=f, y=A] {results/fig/c_harm.dat};
    \addlegendentry{$k\,f_{\mathrm{BPF}}$, $k=1,\dots,6$}
    \node[lbl, cB, anchor=north west, align=left] at (axis cs:\cBpf,\cAbpf-1)
      {$f_{\mathrm{BPF}}=\SI{\cBpf}{\hertz}$\\$A=\SI{\cAbpf}{\decibel}$};
  \end{semilogxaxis}
\end{tikzpicture}

\caption{The A-weighting \cref{eq:aweight} with the design's blade passage
frequency and its first five harmonics. The fundamental at
\SI{\cBpf}{\hertz} is weighted down by \SI{37.5}{\decibel}, and every harmonic
up to the sixth lies in the shaded band below the Schroeder frequency
\cref{eq:schroeder}, where the diffuse-field model of \cref{prop:room} does not
hold. The A-weighted broadband figure is therefore least informative about
exactly the tonal components a listener notices.}
\label{fig:aweight}
\end{figure}

\section{Near-field aerodynamic corrections}
\label{sec:nearfield}

\Cref{mod:disk} treats an isolated rotor in an unbounded fluid at rest. A
machine hovering at head height in a room violates all three conditions: there
is a floor below, a ceiling above, other rotors alongside, and after a short
time the air is not at rest but circulating. This section quantifies the first
three and measures the fourth.

\subsection{Ground effect}

When the wake is constrained by a surface a distance $z$ below the rotor, the
streamtube cannot contract freely, the induced velocity at the disk falls, and
the rotor produces more thrust for the same power.

\begin{model}[Cheeseman-Bennett]
\label{mod:ige}
Modelling the ground by an image source, Cheeseman and
Bennett~\cite{cheeseman1955} obtain, at constant power,
\begin{equation}
  \frac{T_{\mathrm{IGE}}}{T_{\mathrm{OGE}}}\bigg|_{P}
  = \frac{1}{1-\left(\dfrac{R}{4z}\right)^{2}} ,
  \label{eq:ige}
\end{equation}
valid for $z/R \gtrsim 0.5$; equivalently, at constant thrust the induced power
is reduced by the same factor.
\end{model}

The dependence is inverse-square in $z/R$, which decides the magnitude
immediately.

\begin{example}
For $R=\SI{0.246}{\metre}$ and a rotor plane $z=\SI{\rZrotorAbove}{\metre}$
above the floor, the adopted layout of \cref{sec:results},
$(R/4z)^2 \approx 2\times10^{-3}$ and the effect is $0.2\%$. Ground effect
is negligible for a small rotor at chest height, and would only matter within
about one radius of a surface.
\end{example}

\subsection{Ceiling effect}

A ceiling above the rotor restricts the \emph{inflow} rather than the wake. The
low-pressure region above the disk is bounded, the rotor is drawn toward the
surface, and thrust rises at constant power. The functional form is taken to be
the same as \cref{eq:ige} with a separate coefficient,
\begin{equation}
  \frac{T_{\mathrm{ICE}}}{T_{\mathrm{OGE}}}\bigg|_{P}
  = \frac{1}{1-k_c\left(\dfrac{R}{4z_c}\right)^{2}} ,
  \label{eq:ice}
\end{equation}
with $z_c$ the distance to the ceiling and $k_c>1$ because at equal spacing
ceiling effect is the stronger of the two. \Cref{fig:image} shows the image
construction for both surfaces.

\begin{figure}[tbp]
\centering
% Generated by paper/figures/gen_c.py. Do not edit by hand.
\def\cM{1.724}
\def\cMbat{0.288}
\def\cL{0.3827}
\def\cR{0.2460}
\def\cReach{0.629}
\def\cSpan{1.257}
\def\cKgap{1.100}
\def\cGap{0.0492}
\def\cSoverD{0.100}
\def\cKint{1.0252}
\def\cRcp{0.130}
\def\cCoMz{0.0275}
\def\cScreenOffz{0.1025}
\def\cScreenW{0.345}
\def\cScreenH{0.194}
\def\cArmr{8.0}
\def\cWall{0.66}
\def\cWallStress{0.083}
\def\cWallStiff{0.660}
\def\cWallMin{0.60}
\def\cDriver{stiffness}
\def\cSigmaWork{44}
\def\cSigmaAllow{350}
\def\cFbend{33}
\def\cFlimit{15.2}
\def\cZrotor{1.37}
\def\cZceil{1.03}
\def\cZoverR{5.57}
\def\cZcOverR{4.19}
\def\cKground{1.0020}
\def\cKceiling{1.0072}
\def\cKc{2.0}
\def\cCeilPct{0.7}
\def\cGroundPct{0.20}
\def\cKcSpanG{13}
\def\cCeiling{2.4}
\def\cBpf{34.5}
\def\cRpm{1036}
\def\cNb{2}
\def\cAbpf{-37.5}
\def\cLwa{61.8}
\def\cRroom{20.8}
\def\cRc{0.91}
\def\cStandoff{1.20}
\def\cSplEar{56.6}
\def\cSplOne{53.8}
\def\cSplFar{54.8}
\def\cTrotor{4.23}
\def\cVtip{26.7}
\def\cAnchor{82.9}
\def\cVi{3.01}
\def\cTfCross{21}
\def\cTf{30}
\def\cQmax{0.289}
\def\cMmot{61}
\def\cKtau{5.5}
\def\cEmot{0.30}
\def\cRoomL{5}
\def\cRoomW{4}
\def\cAlphaRoom{0.20}
%
% Image-rotor construction for ground and ceiling proximity, drawn to scale
% for the design's rotor in the reference room. The mirrored wake makes the
% floor a plane of symmetry, across which no air passes.
\begin{tikzpicture}[scale=1.5, >={Stealth[length=1.6mm]}, font=\footnotesize]
  \pgfmathsetmacro{\zr}{\cZrotor}
  \pgfmathsetmacro{\H}{\cCeiling}
  \pgfmathsetmacro{\rr}{\cR}
  \pgfmathsetmacro{\zi}{2*\H-\zr}
  % floor and ceiling
  \fill[pattern=north east lines, pattern color=cG] (-1.6,-0.1) rectangle (1.6,0);
  \draw[line width=0.8pt] (-1.6,0) -- (1.6,0);
  \node[anchor=south west] at (-1.6,0.02) {floor};
  \fill[pattern=north east lines, pattern color=cG] (-1.6,\H) rectangle (1.6,\H+0.1);
  \draw[line width=0.8pt] (-1.6,\H) -- (1.6,\H);
  \node[anchor=north west] at (-1.6,\H-0.02) {ceiling};
  % the rotor, its inflow and its wake turning along the floor
  \draw[cA, line width=1.6pt] (-\rr,\zr) -- (\rr,\zr);
  \foreach \s in {-1,1} {
    \draw[cA, line width=0.6pt, ->]
      (\s*\rr,\zr) .. controls (\s*0.8*\rr,\zr-0.35) .. (\s*0.72*\rr,0.9)
      .. controls (\s*0.7*\rr,0.25) and (\s*0.6,0.07) .. (\s*1.25,0.05);
    \draw[cA, line width=0.6pt, <-]
      (\s*1.25*\rr,\zr+0.02) .. controls (\s*1.35*\rr,\zr+0.4) .. (\s*2.0*\rr,\zr+0.72);
    % the ground image and its mirrored wake
    \draw[cA!45, dashed, line width=0.6pt, ->]
      (\s*\rr,-\zr) .. controls (\s*0.8*\rr,-\zr+0.35) .. (\s*0.72*\rr,-0.9)
      .. controls (\s*0.7*\rr,-0.25) and (\s*0.6,-0.07) .. (\s*1.25,-0.05);
    % the ceiling image and its mirrored inflow
    \draw[cE!55, dashed, line width=0.6pt, <-]
      (\s*1.25*\rr,\zi-0.02) .. controls (\s*1.35*\rr,\zi-0.4) .. (\s*2.0*\rr,\zi-0.72);
  }
  \draw[cA!45, dashed, line width=1.6pt] (-\rr,-\zr) -- (\rr,-\zr);
  \draw[cE!55, dashed, line width=1.6pt] (-\rr,\zi) -- (\rr,\zi);
  \node[cA, anchor=east] at (-\rr-0.06,\zr) {rotor, radius $R$};
  \node[cA!70, anchor=east] at (-\rr-0.06,-\zr) {ground image};
  \node[cE!80, anchor=east] at (-\rr-0.06,\zi) {ceiling image};
  % dimensions on the right
  \draw[construct] (\rr+0.05,\zr) -- (1.8,\zr);
  \draw[construct] (\rr+0.05,-\zr) -- (2.1,-\zr);
  \draw[dim, black] (1.55,0) -- (1.55,\zr) node[midway, right] {$z$};
  \draw[dim, black] (1.55,\zr) -- (1.55,\H) node[midway, right] {$z_c$};
  \draw[dim, cG] (2.0,-\zr) -- (2.0,\zr) node[midway, right] {$2z$};
  \node[align=right, anchor=north east] at (-1.68,\H-0.3)
    {$z/R=\cZoverR$\\ $z_c/R=\cZcOverR$};
\end{tikzpicture}

\caption{The image construction behind \cref{eq:ige,eq:ice}, drawn to scale for
the design's rotor in the reference room. A mirror rotor below the floor, with
its wake mirrored, makes the floor a plane of symmetry across which no air
passes, which is the boundary condition the floor imposes; a mirror above the
ceiling does the same for the inflow. The images sit $2z$ and $2z_c$ away, many
rotor radii, and the corrections fall as the square of $R/4z$, which is why both
are small here.}
\label{fig:image}
\end{figure}

\begin{remark}[Honesty about $k_c$]
Published values for ceiling proximity scatter considerably across geometries
and measurement techniques~\cite{sanchezcuevas2017,conyers2019}. $k_c$ is the
least certain constant in this paper. It is exposed as a user parameter and
the engine reports the resulting factor separately so that its influence can be
bounded rather than believed. For the design here it reduces induced power by
\SI{\cCeilPct}{\percent} at $k_c=\cKc$, and the entire range $k_c\in[0.5,6]$
moves the design mass by \SI{\cKcSpanG}{\gram} (\cref{fig:surface}).
\end{remark}

\begin{figure}[tbp]
\centering
% Generated by paper/figures/gen_c.py. Do not edit by hand.
\def\cM{1.724}
\def\cMbat{0.288}
\def\cL{0.3827}
\def\cR{0.2460}
\def\cReach{0.629}
\def\cSpan{1.257}
\def\cKgap{1.100}
\def\cGap{0.0492}
\def\cSoverD{0.100}
\def\cKint{1.0252}
\def\cRcp{0.130}
\def\cCoMz{0.0275}
\def\cScreenOffz{0.1025}
\def\cScreenW{0.345}
\def\cScreenH{0.194}
\def\cArmr{8.0}
\def\cWall{0.66}
\def\cWallStress{0.083}
\def\cWallStiff{0.660}
\def\cWallMin{0.60}
\def\cDriver{stiffness}
\def\cSigmaWork{44}
\def\cSigmaAllow{350}
\def\cFbend{33}
\def\cFlimit{15.2}
\def\cZrotor{1.37}
\def\cZceil{1.03}
\def\cZoverR{5.57}
\def\cZcOverR{4.19}
\def\cKground{1.0020}
\def\cKceiling{1.0072}
\def\cKc{2.0}
\def\cCeilPct{0.7}
\def\cGroundPct{0.20}
\def\cKcSpanG{13}
\def\cCeiling{2.4}
\def\cBpf{34.5}
\def\cRpm{1036}
\def\cNb{2}
\def\cAbpf{-37.5}
\def\cLwa{61.8}
\def\cRroom{20.8}
\def\cRc{0.91}
\def\cStandoff{1.20}
\def\cSplEar{56.6}
\def\cSplOne{53.8}
\def\cSplFar{54.8}
\def\cTrotor{4.23}
\def\cVtip{26.7}
\def\cAnchor{82.9}
\def\cVi{3.01}
\def\cTfCross{21}
\def\cTf{30}
\def\cQmax{0.289}
\def\cMmot{61}
\def\cKtau{5.5}
\def\cEmot{0.30}
\def\cRoomL{5}
\def\cRoomW{4}
\def\cAlphaRoom{0.20}
%
% Fractional reduction in induced power from ground and ceiling effect,
% 1 - 1/k = k_s (R/4z)^2, against distance over radius: a slope of -2 on
% log axes, with the design's two distances marked and the spread of k_c.
\begin{tikzpicture}
  \begin{loglogaxis}[paperplot, height=5.6cm,
      xlabel={distance to the surface over rotor radius, $z/R$},
      ylabel={induced power reduction [\si{\percent}]},
      xmin=0.5, xmax=12, ymin=0.03, ymax=60,
      xtick={0.5,1,2,5,10}, xticklabels={0.5,1,2,5,10},
      ytick={0.1,1,10}, yticklabels={0.1,1,10},
      legend pos=north east]
    \addplot[name path=hi, draw=none, domain=0.5:12, samples=40, forget plot]
      {100*6/(16*x^2)};
    \addplot[name path=lo, draw=none, domain=0.5:12, samples=40, forget plot]
      {100*0.5/(16*x^2)};
    \addplot[cE!18, forget plot] fill between[of=hi and lo];
    \addplot[cB, domain=0.5:12, samples=40] {100/(16*x^2)};
    \addlegendentry{ground, $(R/4z)^{2}$}
    \addplot[cE, domain=0.5:12, samples=40] {100*\cKc/(16*x^2)};
    \addlegendentry{ceiling, $k_c(R/4z_c)^{2}$, $k_c=\cKc$}
    \addlegendimage{area legend, fill=cE!18, draw=none}
    \addlegendentry{ceiling, $k_c\in[0.5,6]$}
    \addplot[only marks, mark=*, mark size=2.2pt, cB] coordinates {(\cZoverR,\cGroundPct)};
    \node[lbl, cB, anchor=north east] at (axis cs:\cZoverR,\cGroundPct) {floor, \SI{\cGroundPct}{\percent}};
    \addplot[only marks, mark=*, mark size=2.2pt, cE] coordinates {(\cZcOverR,\cCeilPct)};
    \node[lbl, cE, anchor=south west] at (axis cs:\cZcOverR,\cCeilPct) {ceiling, \SI{\cCeilPct}{\percent}};
  \end{loglogaxis}
\end{tikzpicture}

\caption{Fractional reduction of induced power from \cref{eq:ige,eq:ice},
$1-1/k=k_s(R/4z)^2$ with $k_s=1$ for the floor and $k_s=k_c$ for the ceiling,
against distance over rotor radius. On logarithmic axes both are lines of slope
$-2$; the band spans the published range $k_c\in[0.5,6]$. At the design's two
distances the floor is worth \SI{\cGroundPct}{\percent} and the ceiling
\SI{\cCeilPct}{\percent}, and even the upper edge of the band stays below
\SI{3}{\percent}.}
\label{fig:surface}
\end{figure}

\subsection{Rotor-to-rotor interference}

Adjacent rotors each operate partly in the other's induced flow. The penalty is
negligible when the tip-to-tip gap exceeds roughly a quarter diameter and grows
as the disks approach.

\begin{model}[Interference factor]
\label{mod:interference}
With $s$ the tip-to-tip gap and $D$ the diameter,
\begin{equation}
  \kappa_{\mathrm{int}}(s/D) =
  \begin{cases}
    1, & s/D \ge 0.25,\\[2pt]
    1 + \varepsilon\left(1 - \dfrac{s/D}{0.25}\right)^{2},
      & 0 \le s/D < 0.25,
  \end{cases}
  \qquad \varepsilon = 0.07 ,
  \label{eq:interference}
\end{equation}
applied as a multiplier on induced power. The quadratic interpolation is chosen
so that the penalty and its derivative vanish at the matching point, consistent
with an induced-flow interaction that decays smoothly.
\end{model}

For a quadrotor with $k_{\mathrm{gap}}=1.10$ in \cref{eq:armlength}, the
adjacent hub spacing is $2L\sin(\pi/4)=2.2R$ and the tip-to-tip gap is
$0.2R = 0.1D$, giving $\kappa_{\mathrm{int}} = 1.025$ (\cref{fig:interference}).

\begin{figure}[tbp]
\centering
% Generated by paper/figures/gen_c.py. Do not edit by hand.
\def\cM{1.724}
\def\cMbat{0.288}
\def\cL{0.3827}
\def\cR{0.2460}
\def\cReach{0.629}
\def\cSpan{1.257}
\def\cKgap{1.100}
\def\cGap{0.0492}
\def\cSoverD{0.100}
\def\cKint{1.0252}
\def\cRcp{0.130}
\def\cCoMz{0.0275}
\def\cScreenOffz{0.1025}
\def\cScreenW{0.345}
\def\cScreenH{0.194}
\def\cArmr{8.0}
\def\cWall{0.66}
\def\cWallStress{0.083}
\def\cWallStiff{0.660}
\def\cWallMin{0.60}
\def\cDriver{stiffness}
\def\cSigmaWork{44}
\def\cSigmaAllow{350}
\def\cFbend{33}
\def\cFlimit{15.2}
\def\cZrotor{1.37}
\def\cZceil{1.03}
\def\cZoverR{5.57}
\def\cZcOverR{4.19}
\def\cKground{1.0020}
\def\cKceiling{1.0072}
\def\cKc{2.0}
\def\cCeilPct{0.7}
\def\cGroundPct{0.20}
\def\cKcSpanG{13}
\def\cCeiling{2.4}
\def\cBpf{34.5}
\def\cRpm{1036}
\def\cNb{2}
\def\cAbpf{-37.5}
\def\cLwa{61.8}
\def\cRroom{20.8}
\def\cRc{0.91}
\def\cStandoff{1.20}
\def\cSplEar{56.6}
\def\cSplOne{53.8}
\def\cSplFar{54.8}
\def\cTrotor{4.23}
\def\cVtip{26.7}
\def\cAnchor{82.9}
\def\cVi{3.01}
\def\cTfCross{21}
\def\cTf{30}
\def\cQmax{0.289}
\def\cMmot{61}
\def\cKtau{5.5}
\def\cEmot{0.30}
\def\cRoomL{5}
\def\cRoomW{4}
\def\cAlphaRoom{0.20}
%
% The interference factor of eq:interference against the tip gap over
% diameter, with the geometry of two adjacent disks and the design point.
\begin{tikzpicture}
  \begin{axis}[halfplot, width=0.56\linewidth, height=5.2cm,
      xlabel={tip-to-tip gap over diameter, $s/D$},
      ylabel={$\kappa_{\mathrm{int}}$},
      xmin=0, xmax=0.4, ymin=0.995, ymax=1.08,
      yticklabel style={/pgf/number format/fixed, /pgf/number format/precision=2}]
    \addplot[cA, domain=0:0.25, samples=60] {1 + 0.07*(1 - x/0.25)^2};
    \addplot[cA, domain=0.25:0.4, samples=2] {1};
    \draw[construct] (axis cs:0.25,0.995) -- (axis cs:0.25,1.08);
    \node[lbl, cG, anchor=south west] at (axis cs:0.25,1.035) {$s/D=0.25$};
    \node[lbl, cA, anchor=south west] at (axis cs:0.005,1.07) {$1+\varepsilon$};
    \addplot[only marks, mark=*, mark size=2.3pt, cA] coordinates {(\cSoverD,\cKint)};
    \node[lbl, cA, anchor=south west] at (axis cs:\cSoverD,\cKint) {design};
  \end{axis}
  % geometry of two adjacent rotors
  \begin{scope}[shift={(11.4,2.0)}, scale=0.78, every node/.style={scale=1}]
    \pgfmathsetmacro{\rr}{1.0}
    \pgfmathsetmacro{\ss}{0.42}
    \fill[cA!10] (-\rr-0.5*\ss,0) circle (\rr);
    \draw[cA, line width=0.8pt] (-\rr-0.5*\ss,0) circle (\rr);
    \fill[cA!10] (\rr+0.5*\ss,0) circle (\rr);
    \draw[cA, line width=0.8pt] (\rr+0.5*\ss,0) circle (\rr);
    \draw[dim, black] (-0.5*\ss,0) -- (0.5*\ss,0);
    \node[lbl, anchor=south] at (0,0.08) {$s$};
    \draw[dim, black] (-2*\rr-0.5*\ss,-1.25) -- (-0.5*\ss,-1.25) node[midway, below, lbl] {$D$};
    \draw[construct] (-1.0-0.5*\ss,0) -- (1.0+0.5*\ss,0);
    \draw[dim, cG] (-\rr-0.5*\ss,1.25) -- (\rr+0.5*\ss,1.25)
      node[midway, above, lbl, cG] {$2L\sin(\pi/N)$};
    \node[lbl, align=center, anchor=north] at (0,-1.7) {$s=2L\sin(\pi/N)-2R$};
  \end{scope}
\end{tikzpicture}

\caption{The interference factor \cref{eq:interference} against tip gap over
diameter, with the geometry that defines the gap. The penalty and its slope
vanish together at $s/D=0.25$, so the model joins the isolated-rotor case
smoothly. The design, with $k_{\mathrm{gap}}=\cKgap$ in \cref{eq:armlength},
sits at $s/D=\cSoverD$ and pays $\kappa_{\mathrm{int}}=\cKint$.}
\label{fig:interference}
\end{figure}

\subsection{Net effect and why it is worth stating}

The three corrections enter the induced power as
\begin{equation}
  \Pind = \frac{\kappa\,\kappa_{\mathrm{int}}}
               {k_{\mathrm{ground}}\,k_{\mathrm{ceiling}}}\;T\vind ,
  \label{eq:kappatotal}
\end{equation}
and they partly cancel. \Cref{tab:nearfield} gives the measured influence on a
converged design.

\begin{table}[t]
\centering
\caption{Effect of the near-field corrections on the converged design of
\cref{sec:results}; the four rows differ only in the corrections applied.
$\kappa_{\mathrm{tot}}$ is the induced power factor actually used,
\cref{eq:kappatotal}. The whole span of the four combinations is
\SI{\rNfSpanDb}{\decibel} and \SI{\rNfSpanG}{\gram}.}
\label{tab:nearfield}
\begin{tabular}{@{}lrrrr@{}}
\toprule
Model & $\kappa_{\mathrm{tot}}$ & Gross mass & Hover power & $L_p(\SI{1}{\metre})$ \\
\midrule
\inputrows{results/tab_nearfield_rows}
\end{tabular}
\end{table}

\begin{finding}
The corrections were worth adding not because they changed the design but
because they bounded their own importance. A caveat that says ``ground and
ceiling effect are not modelled'' is unfalsifiable; a table that says they are
worth a quarter of a decibel is a result. The value of a model is often in the
size of the term it lets you discard.
\end{finding}

\subsection{Recirculation: the correction that is not made}

\Cref{mod:disk} assumes the fluid far upstream is at rest. In a closed room
this is false, and the timescale over which it becomes false is computable.

\begin{definition}[Room exchange time]
The volumetric flow through the disks is $\dot{V} = A\vind$. For a room of
volume $V_{\mathrm{room}}$,
\begin{equation}
  t_{\mathrm{ex}} \coloneqq \frac{V_{\mathrm{room}}}{A\vind}
  \label{eq:exchange}
\end{equation}
is the time for a volume of air equal to the room's to pass through the rotor
disks once.
\end{definition}

\begin{example}
For the design of \cref{sec:results}, $A=\SI{0.760}{\metre\squared}$ and
$\vind=\SI{3.0}{\metre\per\second}$, so in a
$5\times4\times\SI{2.4}{\metre}$ room
$\dot V = \SI{2.28}{\metre\cubed\per\second}$,
$V_{\mathrm{room}}=\SI{48}{\metre\cubed}$ and
$t_{\mathrm{ex}} = \SI{21}{\second}$. The mean return velocity distributed over
the floor plan is $\dot V/(5\times4) = \SI{0.11}{\metre\per\second}$.
\end{example}

\Cref{fig:recirc} shows the circulation this sets up in the room.

\begin{figure}[tbp]
\centering
% Generated by paper/figures/gen_c.py. Do not edit by hand.
\def\cM{1.724}
\def\cMbat{0.288}
\def\cL{0.3827}
\def\cR{0.2460}
\def\cReach{0.629}
\def\cSpan{1.257}
\def\cKgap{1.100}
\def\cGap{0.0492}
\def\cSoverD{0.100}
\def\cKint{1.0252}
\def\cRcp{0.130}
\def\cCoMz{0.0275}
\def\cScreenOffz{0.1025}
\def\cScreenW{0.345}
\def\cScreenH{0.194}
\def\cArmr{8.0}
\def\cWall{0.66}
\def\cWallStress{0.083}
\def\cWallStiff{0.660}
\def\cWallMin{0.60}
\def\cDriver{stiffness}
\def\cSigmaWork{44}
\def\cSigmaAllow{350}
\def\cFbend{33}
\def\cFlimit{15.2}
\def\cZrotor{1.37}
\def\cZceil{1.03}
\def\cZoverR{5.57}
\def\cZcOverR{4.19}
\def\cKground{1.0020}
\def\cKceiling{1.0072}
\def\cKc{2.0}
\def\cCeilPct{0.7}
\def\cGroundPct{0.20}
\def\cKcSpanG{13}
\def\cCeiling{2.4}
\def\cBpf{34.5}
\def\cRpm{1036}
\def\cNb{2}
\def\cAbpf{-37.5}
\def\cLwa{61.8}
\def\cRroom{20.8}
\def\cRc{0.91}
\def\cStandoff{1.20}
\def\cSplEar{56.6}
\def\cSplOne{53.8}
\def\cSplFar{54.8}
\def\cTrotor{4.23}
\def\cVtip{26.7}
\def\cAnchor{82.9}
\def\cVi{3.01}
\def\cTfCross{21}
\def\cTf{30}
\def\cQmax{0.289}
\def\cMmot{61}
\def\cKtau{5.5}
\def\cEmot{0.30}
\def\cRoomL{5}
\def\cRoomW{4}
\def\cAlphaRoom{0.20}
%
% Section through the reference room: the rotor wake strikes the floor,
% spreads, rises along the walls and returns across the ceiling into the
% rotor inflow, so the inflow stops being quiescent after t_ex.
\begin{tikzpicture}[scale=2.3, >={Stealth[length=1.6mm]}, font=\footnotesize]
  \pgfmathsetmacro{\W}{\cRoomL}
  \pgfmathsetmacro{\H}{\cCeiling}
  \pgfmathsetmacro{\xm}{2.1}
  \pgfmathsetmacro{\zr}{\cZrotor}
  \draw[line width=0.9pt] (0,0) rectangle (\W,\H);
  \fill[pattern=north east lines, pattern color=cG] (0,-0.08) rectangle (\W,0);
  % the machine and the user
  \draw[cA, line width=1.6pt] (\xm-0.5,\zr) -- (\xm+0.5,\zr);
  \fill[cE!40] (\xm-0.01,\zr+0.03) rectangle (\xm+0.01,\zr+0.23);
  \draw[black!70, line width=0.8pt] (\xm+1.2,0) -- (\xm+1.2,1.25);
  \fill[black!70] (\xm+1.2,1.4) circle (0.1);
  \node[anchor=west] at (\xm+1.3,1.4) {user};
  % downwash column
  \foreach \dx in {-0.35,0,0.35}
    \draw[cA, line width=0.9pt, ->] (\xm+\dx,\zr-0.03) -- (\xm+0.85*\dx,0.25);
  % recirculating streamlines
  \draw[cA!80, line width=0.7pt, ->] (\xm-0.3,0.22) .. controls (\xm-1.2,0.05) and (0.15,0.05) .. (0.12,0.8)
     .. controls (0.1,\H-0.2) and (0.8,\H-0.1) .. (\xm-1.0,\H-0.2)
     .. controls (\xm-0.6,\H-0.3) and (\xm-0.4,\zr+0.5) .. (\xm-0.3,\zr+0.06);
  \draw[cA!80, line width=0.7pt, ->] (\xm+0.3,0.22) .. controls (\xm+1.2,0.05) and (\W-0.15,0.05) .. (\W-0.12,0.8)
     .. controls (\W-0.1,\H-0.2) and (\W-0.9,\H-0.1) .. (\xm+1.1,\H-0.2)
     .. controls (\xm+0.6,\H-0.3) and (\xm+0.4,\zr+0.5) .. (\xm+0.3,\zr+0.06);
  \draw[cA!55, line width=0.6pt, ->] (\xm-0.2,0.35) .. controls (\xm-0.8,0.25) and (0.6,0.35) .. (0.55,1.0)
     .. controls (0.55,\H-0.5) and (\xm-0.9,\H-0.55) .. (\xm-0.15,\zr+0.1);
  \draw[cA!55, line width=0.6pt, ->] (\xm+0.2,0.35) .. controls (\xm+0.8,0.25) and (\W-0.6,0.35) .. (\W-0.55,1.0)
     .. controls (\W-0.55,\H-0.5) and (\xm+0.9,\H-0.55) .. (\xm+0.15,\zr+0.1);
  % labels
  \node[anchor=north, align=center] at (0.5*\W,-0.12)
    {$\dot V=A\vind\approx\SI{2.3}{\metre\cubed\per\second}$,\quad
     $t_{\mathrm{ex}}=V_{\mathrm{room}}/\dot V\approx\SI{21}{\second}$,\quad
     mean return velocity $\dot V/(\text{floor area})\approx\SI{0.1}{\metre\per\second}$};
  \draw[dim, black] (0,\H+0.12) -- (\W,\H+0.12) node[midway, above] {\SI{\cRoomL}{\metre}};
  \draw[dim, black] (\W+0.12,0) -- (\W+0.12,\H) node[midway, right] {\SI{\cCeiling}{\metre}};
  \node[cA, anchor=south] at (\xm,\zr+0.25) {machine};
\end{tikzpicture}

\caption{Recirculation in the reference room, schematic. The wake strikes the
floor, spreads, rises along the walls and returns across the ceiling into the
rotor inflow. At $\dot V=A\vind$ a volume of air equal to the room's passes
through the disks every $t_{\mathrm{ex}}$ of \cref{eq:exchange}, so within the
first minute of a session the machine is drawing in air it has already
accelerated, and the quiescent inflow of \cref{mod:disk} no longer holds.}
\label{fig:recirc}
\end{figure}

\begin{finding}
Within one minute of a thirty-minute session the machine is flying in air it
has already accelerated. The quiescent-inflow assumption underlying every
result of \cref{sec:momentum} has quietly ceased to hold. Correcting this is
not a matter of a coefficient: it requires solving the room flow, which is a
computational fluid dynamics problem with a moving boundary. It is the largest
outstanding gap in the model and is recorded as such in
\cref{sec:limitations}. A secondary consequence is a persistent draught of
order \SI{0.1}{\metre\per\second} throughout the room, far below the
\SI{\rWake}{\metre\per\second} wake velocity but present continuously and over the
whole floor area.
\end{finding}


\part{The Differential Layer: Flight}
\section{Flight dynamics}
\label{sec:dynamics}

Once the machine exists its mass is fixed (\cref{as:mass}) and the problem
becomes one of differential equations. This section derives them.

\subsection{Configuration and kinematics}

The configuration space is $\SEthree = \R^3\rtimes\SOthree$, with state
$(\vect{r},R)$ and velocities $(\vect{v},\vect{\omega})$.

\begin{proposition}[Attitude kinematics]
\label{prop:kinematics}
\begin{equation}
  \dt{R} = R\,\hatmap{\vect{\omega}} .
  \label{eq:Rdot}
\end{equation}
In terms of the unit quaternion $q=(q_0,\vect{q}_v)$ representing $R$,
\begin{equation}
  \dt{q} = \tfrac12\, q \otimes (0,\vect{\omega})
         = \tfrac12
           \begin{bmatrix}
             -\vect{q}_v\transp \\
             q_0 I + \hatmap{\vect{q}_v}
           \end{bmatrix}\vect{\omega}
         \eqqcolon \tfrac12\,\Omega(\vect{\omega})\,q .
  \label{eq:qdot}
\end{equation}
\end{proposition}

\begin{proof}
Let $\vect{u}$ be a vector fixed in $\mathcal{B}$, so its inertial coordinates
are $R\vect{u}$ and its inertial time derivative is
$\vect{\omega}_{\mathcal I}\times R\vect{u}$ where
$\vect{\omega}_{\mathcal I}=R\vect{\omega}$. Then
$\dt{R}\vect{u} = \hatmap{(R\vect{\omega})}R\vect{u}$ for all $\vect{u}$, so
$\dt{R} = \hatmap{(R\vect{\omega})}R = R\hatmap{\vect{\omega}}R\transp R
= R\hatmap{\vect{\omega}}$ using \cref{eq:hatid}. \Cref{eq:qdot} follows by
differentiating \cref{eq:quat2rot} and using quaternion multiplication;
see~\cite[\S1.4]{markley2014}.
\end{proof}

\begin{remark}
\Cref{eq:qdot} preserves $\|q\|$ exactly, since
$q\transp\dt q = \tfrac12 q\transp\Omega q = 0$ by antisymmetry. Numerically,
integration error accumulates and the engine renormalises after each step; see
\cref{sec:numerics}.
\end{remark}

\subsection{Newton-Euler equations}

\begin{theorem}[Equations of motion]
\label{thm:newtoneuler}
For a rigid body of mass $m$ and body-frame inertia tensor $J$ subject to a
total external force $\vect{F}$ expressed in $\mathcal I$ and a total external
moment $\vect{\tau}$ about the centre of mass expressed in $\mathcal B$,
\begin{align}
  \dt{\vect{r}} &= \vect{v}, \label{eq:rdot}\\
  m\,\dt{\vect{v}} &= \vect{F}, \label{eq:vdot}\\
  J\,\dt{\vect{\omega}} + \vect{\omega}\times J\vect{\omega}
    &= \vect{\tau} . \label{eq:wdot}
\end{align}
\end{theorem}

\begin{proof}
\Cref{eq:rdot,eq:vdot} are Newton's second law in an inertial frame. For
\cref{eq:wdot}, the angular momentum about the centre of mass expressed in
$\mathcal I$ is $\vect{h}_{\mathcal I} = RJ\vect{\omega}$ and
$\dt{\vect h}_{\mathcal I} = \vect{\tau}_{\mathcal I} = R\vect\tau$.
Differentiating using \cref{eq:Rdot},
\[
  \dt{\vect h}_{\mathcal I}
  = \dt R J\vect\omega + RJ\dt{\vect\omega}
  = R\hatmap{\vect\omega}J\vect\omega + RJ\dt{\vect\omega}
  = R\big(\vect\omega\times J\vect\omega + J\dt{\vect\omega}\big).
\]
Equating and cancelling $R$ gives \cref{eq:wdot}.
\end{proof}

The forces are
\begin{equation}
  \vect{F} = T\,R\ez - mg\,\ez + \vect{F}_{\mathrm{aero}}
             + \vect{F}_{\mathrm{dist}} ,
  \label{eq:forces}
\end{equation}
with $T$ the total rotor thrust along the body $\vect b_3$ axis.

\subsection{Rotor wrench and the allocation map}

\begin{model}[Quasi-steady rotor]
\label{mod:rotorcoeff}
Rotor $i$ turning at $\Omega_i$ produces thrust and reaction torque
\begin{equation}
  T_i = k_T\Omega_i^2, \qquad Q_i = k_Q\Omega_i^2 ,
  \label{eq:ktkq}
\end{equation}
with $k_T,k_Q$ determined by \cref{sec:momentum,sec:blade} rather than fitted:
$k_T = T/\Omega^2$ and $k_Q = \Pshaft/\Omega^3$ evaluated at the design point.
\end{model}

Let rotor $i$ sit at body position $\vect{d}_i = (x_i,y_i,0)$ and have spin
sense $s_i\in\{+1,-1\}$. Its contribution to the body moment is
$\vect{d}_i\times T_i\ez$ from thrust and $-s_iQ_i\ez$ from reaction torque, so
\begin{equation}
  \vect\tau = \sum_{i=1}^{N}
    \begin{bmatrix} y_i\,k_T \\ -x_i\,k_T \\ -s_i\,k_Q \end{bmatrix}\Omega_i^2 ,
  \qquad
  T = \sum_{i=1}^{N} k_T\Omega_i^2 .
  \label{eq:wrench}
\end{equation}

\begin{definition}[Allocation matrix]
Collecting \cref{eq:wrench} with $\vect{u} = (\Omega_1^2,\dots,\Omega_N^2)$,
\begin{equation}
  \begin{bmatrix} T \\ \vect\tau \end{bmatrix}
  = \mat{M}\,\vect{u},
  \qquad
  \mat{M} =
  \begin{bmatrix}
    k_T & \cdots & k_T\\
    y_1 k_T & \cdots & y_N k_T\\
    -x_1 k_T & \cdots & -x_N k_T\\
    -s_1 k_Q & \cdots & -s_N k_Q
  \end{bmatrix} \in \R^{4\times N} .
  \label{eq:alloc}
\end{equation}
\end{definition}

\begin{proposition}[Controllability of the wrench]
\label{prop:alloc}
The set of achievable wrenches is $\mat{M}\R_{\ge0}^N$. Full four-axis
authority requires $\rank \mat M = 4$, which for hubs equispaced on a circle of
radius $L$ with alternating spin senses holds whenever $N\ge4$ is even. For odd
$N$, $\sum_i s_i\ne0$ and the hover trim $\Omega_i\equiv\Omega_h$ produces a
nonzero yaw moment, so the vehicle cannot trim without unequal thrusts.
\end{proposition}

\begin{proof}
The rank statement is a direct computation on \cref{eq:alloc}: the first row is
constant, rows two and three span the $(y_i,-x_i)$ pattern which has rank two
for $N\ge3$ equispaced hubs, and row four is independent of the others exactly
when the $s_i$ are not all equal. For odd $N$ with alternating signs,
$\sum_i s_i = \pm1$, so at equal $\Omega_i$ the yaw component of
$\mat M\vect u$ is $\mp k_Q\Omega_h^2 \ne 0$.
\end{proof}

\begin{remark}
\Cref{prop:alloc} is why the engine restricts $N$ to even values. A tricopter
resolves the trim problem with a tilting motor, which is a different actuation
model and is outside the scope of \cref{eq:alloc}.
\end{remark}

\subsection{Motor dynamics}

Rotor speed does not follow command instantaneously. The simplest adequate
model is first order,
\begin{equation}
  \tau_m\,\dt{\Omega}_i = \Omega_i^{\mathrm{cmd}} - \Omega_i ,
  \label{eq:motorlag}
\end{equation}
with $\tau_m$ of order \SIrange{20}{60}{\milli\second}. The underlying
electromechanical model is
\begin{equation}
  L\dt{\imath}_i = V_i - R_m \imath_i - k_e\Omega_i,
  \qquad
  J_r\dt{\Omega}_i = k_\tau \imath_i - k_Q\Omega_i^2 ,
  \label{eq:bldc}
\end{equation}
whose electrical time constant $L/R_m$ is typically two orders of magnitude
faster than the mechanical one, so that singular perturbation reduces
\cref{eq:bldc} to a first-order system in $\Omega_i$; \cref{eq:motorlag} is its
linearisation about the trim point. We use \cref{eq:motorlag} and record the
simplification.

\subsection{Battery state}

\begin{equation}
  \dt{E} = -P, \qquad
  P = \frac{1}{\eta}\sum_{i=1}^{N} k_P\,\Omega_i^{3} + \Paux ,
  \qquad
  k_P = \frac{\Pshaft}{\Omega^3}\bigg|_{\text{design}} ,
  \label{eq:edot}
\end{equation}
or in normalised form $\dt{s} = -P/E_{\max}$ with $s = E/E_{\max}$ the state of
charge. The cubic dependence follows from $\Pshaft = Q\Omega$ with
$Q\propto\Omega^2$. \Cref{eq:edot} is the coupling that makes control effort
cost endurance, and it is what \cref{sec:codesign} closes.

\Cref{eq:edot} is calibrated at hover and has no inflow term. At fixed rotor
speed it returns the same power in hover and in axial translation, so it does
not contain the climb correction of \cref{eq:axial}, which the component model
of \cref{sec:momentum} does. It is a near-hover approximation, and so are the
mission energies computed with it: they are appropriate for following a person
indoors at walking speed and vertical rates of a few tenths of a metre per
second, and they are not a model of climbs, descents or aggressive manoeuvres.
Adding a climb term to the power alone, while thrust keeps its hover
calibration, would make the two inconsistent; a consistent treatment couples an
axial and oblique inflow model to both.

\subsection{The screen as an aerodynamic surface}

This is what distinguishes the vehicle from a camera drone. The display is a
large flat plate, rigidly attached through a gimbal, and therefore both
generates orientation-dependent drag and applies that drag at a point offset
from the centre of mass.

\begin{model}[Flat plate with orientation]
\label{mod:screen}
Let $\vect{n}$ be the outward normal of the screen in inertial coordinates and
$\vect{v}_{\mathrm{rel}} = \vect{v} - \vect{v}_{\mathrm{air}}$. Decompose
$\vect{v}_{\mathrm{rel}} = v_n\vect{n} + \vect{v}_t$ with
$v_n = \vect{n}\transp\vect{v}_{\mathrm{rel}}$. Then
\begin{equation}
  \vect{F}_{\mathrm{screen}}
  = -\tfrac12\rho A_s\Big[
      C_{D,n}\,|v_n|\,v_n\,\vect{n}
      + C_{D,t}\,\|\vect{v}_t\|\,\vect{v}_t \Big] ,
  \label{eq:screendrag}
\end{equation}
with $C_{D,n}\approx1.17$ for a normal flat plate and $C_{D,t}\ll C_{D,n}$
skin friction.
\end{model}

The moment about the centre of mass follows from the offset $\vect{d}_s$ of the
screen's centre of pressure in body coordinates:
\begin{equation}
  \vect\tau_{\mathrm{screen}} = R\transp\big[(R\vect{d}_s)\times
    \vect{F}_{\mathrm{screen}}\big] .
  \label{eq:screenmoment}
\end{equation}

\begin{remark}[The sign of $\vect{d}_s$ matters]
If the display is carried \emph{above} the rotor plane, $d_{s,3}>0$; if it
hangs below, $d_{s,3}<0$. In forward flight the drag acts rearward, so a screen
above the centre of mass produces a nose-up moment that opposes the forward
tilt, a weathercock-stable arrangement, while a screen below produces a
nose-down moment that reinforces it. The magnitude is identical; only the sign
changes, and it changes the closed-loop behaviour measurably.
\end{remark}

\begin{remark}[Why the screen is nearly transparent to the rotors]
The screen is vertical and the rotor inflow and wake are vertical, so the plate
is edge-on to the rotor flow and blocks a projected area of order
(thickness $\times$ width), some \SI{2e-3}{\metre\squared} against a disk area
of \SI{0.76}{\metre\squared}. Placing it above or below the rotor plane is
therefore almost free aerodynamically. What it changes is geometry relative to
the user, which \cref{sec:results} shows to be the decisive consideration.
\end{remark}

\subsection{Gimbal}

The display is carried on a gimbal whose job is to hold it upright while the
airframe tilts to accelerate. Modelling the pitch and roll axes as
independently actuated single-degree-of-freedom systems,
\begin{equation}
  J_g\ddt{\theta}_{g,j} + b_g\dt{\theta}_{g,j}
    = \tau_{g,j} - \tau_{\mathrm{wind},j},
  \qquad
  |\tau_{g,j}| \le \tau_g^{\max},
  \quad j\in\{\text{pitch},\text{roll}\} .
  \label{eq:gimbal}
\end{equation}
The gimbal motor torque reacts on the airframe and is included in $\vect\tau$.

\begin{definition}[Readability error]
Let $\vect{u}_s$ be the unit vector along the top edge of the display in
inertial coordinates. The readability error is
\begin{equation}
  \theta_{\mathrm{read}} = \arccos\big(\ez\transp\vect{u}_s\big) ,
  \label{eq:readability}
\end{equation}
the angle between the top of the image and vertical.
\end{definition}

\begin{remark}
\Cref{eq:readability} is deliberately not the angle between the screen normal
and horizontal. Gimbal roll leaves the normal unchanged while rotating the
image in its own plane, so a normal-based metric is blind to exactly the error a
reader notices most. Measuring the up-vector captures both axes.
\end{remark}

\subsection{The complete system}

Assembling, the state is
\begin{equation}
  x = \big(\vect{r},\ \vect{v},\ q,\ \vect{\omega},\
        \Omega_1,\dots,\Omega_N,\ E,\
        \vect\theta_g,\ \dt{\vect\theta}_g\big)
  \in \R^{3}\times\R^{3}\times\mathbb{S}^{3}\times\R^{3}\times\R^{N}\times\R
      \times\R^{2}\times\R^{2} ,
  \label{eq:state}
\end{equation}
of dimension $18+N$, and the dynamics are
\begin{equation}
  \dt{x} = f(x,u,w;\,p),
  \label{eq:ode}
\end{equation}
with control $u = (\Omega^{\mathrm{cmd}}_{1:N}, \vect\tau_g)$, disturbance $w$
comprising wind and human motion, and design parameters $p$.

Crucially, \cref{eq:ode} is not the whole problem. The parameters $p$ are
themselves constrained by the algebraic closure of \cref{sec:closure}, giving a
differential-algebraic system
\begin{equation}
  \dt{x} = f(x,u,w;p),
  \qquad
  0 = h(x,u,p) ,
  \label{eq:dae}
\end{equation}
whose algebraic part contains \cref{eq:budget} with the mission energy computed
from a solution of the differential part. \Cref{sec:codesign} makes this
precise.

\section{Control}
\label{sec:control}

\subsection{Why a cascade}

The vehicle is underactuated: four independent rotor speeds control six
degrees of freedom. Horizontal motion is available only by tilting the thrust
vector, so the causal chain is
\begin{equation}
  \tau_y \;\to\; \ddt\theta \;\to\; \theta \;\to\; \ddt x \;\to\; x ,
  \label{eq:chain}
\end{equation}
four integrations from moment to position. Near hover, linearising
\cref{eq:vdot,eq:wdot} about $R=I$, $T=mg$ gives the classical decoupling
\begin{equation}
  \ddt{x}\approx g\theta,\quad
  \ddt{y}\approx -g\phi,\quad
  \ddt{z}\approx \frac{\delta T}{m},\quad
  \ddt\phi = \frac{\tau_x}{J_x},\quad
  \ddt\theta = \frac{\tau_y}{J_y},\quad
  \ddt\psi = \frac{\tau_z}{J_z} .
  \label{eq:linhover}
\end{equation}
The attitude loop is two integrations from its input, the position loop four.
Since $\sqrt{K_R}\gg\sqrt{K_p}$ in any sensible design, time-scale separation
justifies the cascade

\begin{center}
position $\to$ velocity $\to$ attitude $\to$ body rate $\to$ rotor speed,
\end{center}

each level treating the level below as instantaneous. This is standard
practice~\cite{mahony2012}, and \cref{eq:linhover} is its justification.
\Cref{fig:cascade} shows the cascade as the engine implements it, and
\cref{fig:bandwidth} the separation of its time scales that the argument
needs.

\begin{figure}[tbp]
\centering
% The cascade of the engine's controller, from the user to the rigid body.
\begin{tikzpicture}[
    font=\footnotesize,
    box/.style={block, text width=1.95cm, minimum height=3.2em, inner sep=2pt},
    ctl/.style={box, fill=cA!7},
    phys/.style={box, fill=cB!10},
    sig/.style={font=\scriptsize, midway},
  ]
  % top row: from the user to the thrust command
  \node[phys] (user) at (0,0)      {user\\ $\vect c(t),\ \psi_h(t)$};
  \node[box]  (raw)  at (2.95,0)   {ideal display\\ trajectory};
  \node[ctl]  (gov)  at (5.9,0)    {reference\\ governor};
  \node[ctl]  (pos)  at (8.85,0)   {position loop\\ \labelcref{eq:poscmd}};
  \node[ctl]  (des)  at (11.8,0)   {desired attitude\\ \labelcref{eq:Rd,eq:omegad}};
  % bottom row: from the attitude loop back to the plant
  \node[ctl]  (att)  at (11.8,-2.6) {attitude loop\\ \labelcref{eq:leeused}};
  \node[ctl]  (mix)  at (8.85,-2.6) {mixer $\mat M^{+}$,\\ saturation};
  \node[phys] (mot)  at (5.9,-2.6)  {motors, lag\\ \labelcref{eq:motorlag}};
  \node[phys] (body) at (2.95,-2.6) {rigid body,\\ rotors, drag};
  \node[ctl]  (gim)  at (0,-2.6)    {gimbal PD,\\ screen upright};

  \draw[flow] (user) -- node[sig, above] {$\vect r_h$} (raw);
  \draw[flow] (raw)  -- node[sig, above] {$\vect r_{\mathrm{raw}}$} (gov);
  \draw[flow] (gov)  -- node[sig, above, align=center] {$\vect r_g,\vect v_g$\\ $\vect a_g,\vect j_g$} (pos);
  \draw[flow] (pos)  -- node[sig, above] {$\vect F_c$} (des);
  \draw[flow] (des)  -- node[sig, right, align=left] {$R_d$\\ $\vect\omega_d$} (att);
  \draw[flow] (att)  -- node[sig, above] {$\vect\tau$} (mix);
  \draw[flow] (mix)  -- node[sig, above] {$\Omega^{\mathrm{cmd}}$} (mot);
  \draw[flow] (mot)  -- node[sig, above] {$\Omega$} (body);
  % thrust magnitude goes to the mixer directly
  \draw[flow] ($(des.south)+(-0.7,0)$) -- ++(0,-0.62)
        -| node[sig, pos=0.25, above] {$T_c=\|\vect F_c\|$} ($(mix.north)+(0.55,0)$);
  % head position enters the keep-out constraint
  \draw[flow, cD] (user.north) -- ++(0,0.55)
        -| node[sig, pos=0.25, above] {head $\vect c,\dt{\vect c},\ddt{\vect c}$: keep-out \labelcref{eq:barrier}} (gov.north);
  % heading reference
  \draw[flow] ($(gov.north)+(0.65,0)$) -- ++(0,0.28)
        -| node[sig, pos=0.25, above] {$\psi_g,\ \dt\psi_g,\ \ddt\psi_g$} (des.north);
  % feedback of translational state
  \draw[flow, cE] ($(body.north)+(0.4,0)$) -- ++(0,0.5)
        -| node[sig, pos=0.2, below] {$\vect r,\ \vect v$} ($(pos.south)+(-0.45,0)$);
  % feedback of rotational state
  \draw[flow, cE] (body.south) -- ++(0,-0.8)
        -| node[sig, pos=0.25, below] {$R,\ \vect\omega$} (att.south);
  \draw[flow, cE] (body.west) -- node[sig, above] {$R$} (gim.east);
  \draw[flow] (gim.south) -- ++(0,-0.5)
        -| node[sig, pos=0.3, below] {$\vect\tau_g$} ($(body.south)+(-0.65,0)$);

  % the sampled controller
  \begin{scope}[on background layer]
    \node[draw=cA, dashed, rounded corners=3pt, inner xsep=6pt, inner ysep=9pt,
          fit=(gov)(pos)(des)(att)(mix)] (ctlbox) {};
  \end{scope}
  \node[font=\scriptsize, anchor=north west, align=left] at (-1.05,-4.3)
       {\tikz\fill[cA!7,draw=black] (0,0) rectangle (0.3,0.18); computed\quad
        \tikz\fill[cB!10,draw=black] (0,0) rectangle (0.3,0.18); physical\quad
        \tikz\draw[cA, dashed] (0,0) rectangle (0.3,0.18); controller, sampled at
        $1/h=\SI{250}{\hertz}$, output held between samples};
\end{tikzpicture}

\caption{The cascade as the engine implements it. Blue boxes are computed at
each sample of the controller, orange boxes are physical. The reference
governor turns the ideal display trajectory into one the aircraft can fly and
holds it outside the keep-out sphere about the user's head; the position loop
turns the governed reference into a force $\vect F_c$; the thrust magnitude,
the desired attitude $R_d$ and its rate $\vect\omega_d$ follow from that force
and its derivative; the attitude loop and the mixer turn them into rotor speed
commands. The translational state closes the outer loop, the rotational state
the inner one. The gimbal loop runs beside the cascade and keeps the screen
upright while the airframe tilts.}
\label{fig:cascade}
\end{figure}

\begin{remark}[Flatness]
The quadrotor with these dynamics is differentially flat with flat outputs
$(x,y,z,\psi)$~\cite{mellinger2011}: every state and input is an algebraic
function of the flat outputs and finitely many derivatives. This is what makes
feedforward from a desired trajectory possible and is why the controller below
uses $\ddt{\vect r}_d$ directly rather than relying on feedback alone.
\end{remark}

\subsection{Position loop}

Given a reference $(\vect r_d,\dt{\vect r}_d,\ddt{\vect r}_d)$, define errors
$\vect e_p = \vect r_d - \vect r$, $\vect e_v = \dt{\vect r}_d - \vect v$, and
command
\begin{equation}
  \vect a_c = \ddt{\vect r}_d + K_p\vect e_p + K_d\vect e_v
              + K_i\!\int\!\vect e_p\,\dd t ,
  \qquad
  \vect F_c = m(\vect a_c + g\ez) .
  \label{eq:poscmd}
\end{equation}
The desired thrust magnitude and thrust direction are
\begin{equation}
  T_c = \|\vect F_c\|,
  \qquad
  \vect b_3^d = \frac{\vect F_c}{\|\vect F_c\|} ,
  \label{eq:thrustdir}
\end{equation}
and with a desired yaw $\psi_d$ the desired attitude is completed by
\begin{equation}
  \vect b_2^d = \frac{\vect b_3^d \times \vect b_1^{\psi}}
                     {\|\vect b_3^d \times \vect b_1^{\psi}\|},
  \quad
  \vect b_1^d = \vect b_2^d\times\vect b_3^d,
  \quad
  R_d = [\,\vect b_1^d\ \vect b_2^d\ \vect b_3^d\,],
  \label{eq:Rd}
\end{equation}
where $\vect b_1^{\psi}=(\cos\psi_d,\sin\psi_d,0)$. \Cref{eq:Rd} is singular
only when $\vect b_3^d\parallel\vect b_1^{\psi}$, i.e.\ at \SI{90}{\degree} of
tilt, which the thrust limit precludes; the engine falls back to an alternative
reference axis if it occurs.

\subsection{Geometric attitude control on \texorpdfstring{$\SOthree$}{SO(3)}}

Attitude error is measured on the group rather than in coordinates, which
avoids both the singularities of Euler angles and the unwinding of naive
quaternion feedback.

\begin{definition}[Attitude error]
\begin{equation}
  \vect e_R = \tfrac12\big(R_d\transp R - R\transp R_d\big)^{\vee},
  \qquad
  \vect e_\omega = \vect\omega - R\transp R_d\,\vect\omega_d ,
  \label{eq:attitudeerror}
\end{equation}
with configuration error function
$\Psi(R,R_d) = \tfrac12\tr\!\big(I - R_d\transp R\big) \in [0,2]$.
\end{definition}

\begin{model}[Geometric attitude law]
\label{mod:leecontrol}
\begin{equation}
  \vect\tau = -K_R\,\vect e_R - K_\omega\,\vect e_\omega
              + \vect\omega\times J\vect\omega
              - J\big(\hatmap{\vect\omega}R\transp R_d\vect\omega_d
                      - R\transp R_d\dt{\vect\omega}_d\big) .
  \label{eq:leecontrol}
\end{equation}
\end{model}

\begin{theorem}[Exponential attitude stability; Lee, Leok and
McClamroch~\cite{lee2010}, Proposition~1]
\label{thm:lee}
Consider \cref{eq:wdot} under \cref{eq:leecontrol} with positive scalar gains
$k_R,k_\omega$ in place of $K_R,K_\omega$ and a smooth reference
$(R_d,\vect\omega_d,\dt{\vect\omega}_d)$. The equilibrium
$(\vect e_R,\vect e_\omega)=(0,0)$ is exponentially stable, and its region of
attraction contains the set
\begin{equation}
  \Psi(R(0),R_d(0)) < 2,
  \qquad
  \|\vect e_\omega(0)\|^2 < \frac{2}{\lambda_{\max}(J)}\,k_R\big(2-\Psi(R(0),R_d(0))\big) .
  \label{eq:roa}
\end{equation}
\end{theorem}

The excluded set is $\Psi=2$, where $R_d\transp R$ is a rotation by $\pi$. For
the unweighted $\Psi$ used here every such rotation is a critical point: if
$R_d\transp R = 2\vect u\vect u\transp - I$ for a unit vector $\vect u$, the
matrix is symmetric and $\vect e_R=0$. The critical set of $\Psi$ is therefore
the identity together with a two-dimensional family of attitudes, not three
isolated points. \Cref{thm:lee} is a statement on the sublevel set
\cref{eq:roa}, and no almost-global claim is made here. \Cref{fig:psi} shows
the mechanism along any single axis of rotation.

\begin{figure}[tbp]
\centering
% Configuration error Psi and the attitude error e_R against the rotation
% angle theta between R and R_d: Psi = 1 - cos theta, |e_R| = sin theta.
\begin{tikzpicture}
  \begin{axis}[paperplot, height=5.2cm,
      xmin=0, xmax=3.1416, ymin=0, ymax=2.25,
      xtick={0, 0.7854, 1.5708, 2.3562, 3.1416},
      xticklabels={$0$, $\pi/4$, $\pi/2$, $3\pi/4$, $\pi$},
      xlabel={rotation angle $\theta$ of $R_d\transp R$ about any axis $\vect u$},
      ylabel={dimensionless},
      legend pos=north west, domain=0:3.1416, samples=120]
    \addplot[name path=psi, cA] {1 - cos(deg(x))};
    \addlegendentry{$\Psi = \tfrac12\tr(I-R_d\transp R) = 1-\cos\theta$}
    \addplot[name path=er, cB] {sin(deg(x))};
    \addlegendentry{$\|\vect e_R\| = \sin\theta$}
    \path[name path=axisbot] (axis cs:0,0) -- (axis cs:3.1416,0);
    \addplot[cB!12] fill between[of=er and axisbot, soft clip={domain=1.5708:3.1416}];
    \draw[construct] (axis cs:1.5708,0) -- (axis cs:1.5708,2.25);
    \draw[construct, cD] (axis cs:0,2) -- (axis cs:3.1416,2);
    \node[lbl, cD, anchor=south east, align=right] at (axis cs:3.05,2.0)
      {$\Psi=2$: excluded by \cref{eq:roa}};
    \fill[cD] (axis cs:3.1416,0) circle (2.2pt);
    \fill[cD] (axis cs:3.1416,2) circle (2.2pt);
    \node[lbl, align=center, anchor=south] at (axis cs:2.2,0.08)
      {restoring term $K_R\vect e_R$\\ weakens past $\pi/2$};
    \draw[lbl, cD, -{Stealth[length=1.6mm]}] (axis cs:2.8,0.95) -- (axis cs:3.09,0.07)
      node[pos=0, above left=-2pt, align=center] {critical set:\\ $\vect e_R=0$, $\Psi=2$};
  \end{axis}
\end{tikzpicture}

\caption{The attitude error for a rotation of $R_d\transp R$ by the angle
$\theta$ about any unit axis $\vect u$, for which $\tr(R_d\transp R)=1+2\cos\theta$
gives $\Psi=1-\cos\theta$ and $\vect e_R=\sin\theta\,\vect u$. The configuration
error grows monotonically to its maximum $2$ at $\theta=\pi$, but the restoring
term, proportional to $\vect e_R$, peaks at $\pi/2$ and decays to zero at $\pi$
(shaded). Every rotation by $\pi$ is therefore a critical point of $\Psi$ at
which the proportional feedback vanishes, which is why \cref{eq:roa} excludes
$\Psi=2$ and why no continuous feedback of this kind is global on $\SOthree$.}
\label{fig:psi}
\end{figure}

The engine implements
\begin{equation}
  \vect\tau = J\big(-K_R\vect e_R - K_\omega\vect e_\omega\big)
              + \vect\omega\times J\vect\omega
              - J\big(\hatmap{\vect\omega}R\transp R_d\vect\omega_d
                      - R\transp R_d\dt{\vect\omega}{}_d^{\psi}\big),
  \qquad
  \vect e_\omega = \vect\omega - R\transp R_d\vect\omega_d ,
  \label{eq:leeused}
\end{equation}
with the desired angular velocity taken from differential
flatness~\cite{mellinger2011}. With $\vect j_d$ the reference jerk and $T_c$ the
commanded thrust,
\begin{equation}
  \vect h_\omega = \frac{m}{T_c}\Big(\vect j_d
     - (\vect b_3^d\cdot\vect j_d)\,\vect b_3^d\Big),
  \quad
  \vect\omega_d = \begin{bmatrix}
     -\vect h_\omega\cdot\vect b_2^d\\ \vect h_\omega\cdot\vect b_1^d\\
     \dt\psi_d\,(\vect b_3^d\cdot\ez) \end{bmatrix},
  \quad
  \dt{\vect\omega}{}_d^{\psi} = \begin{bmatrix} 0\\0\\
     \ddt\psi_d\,(\vect b_3^d\cdot\ez)\end{bmatrix} .
  \label{eq:omegad}
\end{equation}
\Cref{eq:leeused} differs from the law of \cref{thm:lee} in three ways. The
gains $JK_R$ and $JK_\omega$ are matrices, which coincide with scalar gains only
when $J$ is a multiple of the identity. The tilt part of $\dt{\vect\omega}_d$,
which needs the reference snap, is omitted. And the plant has motor lag and
rotor saturation, which the theorem does not. \Cref{thm:lee} therefore
motivates \cref{eq:leeused} and does not certify it: the behaviour of the
implemented loop is established by simulation for the trajectories reported,
not by proof. Normalising by $J$ makes $K_R$ carry units of
\si{\per\second\squared} and gives it a direct reading: the angular
acceleration commanded per radian of attitude error.

\begin{remark}[Why the rate feedforward matters]
An earlier version of the engine set $\vect\omega_d=0$. On a curved path the
thrust direction rotates at the turn rate of the path (\cref{fig:thrustrot}),
and a loop with no rate
feedforward follows it with a standing attitude error of about
$K_\omega\|\vect\omega_d\|/K_R$. On the turnaround mission of
\cref{sec:results} that version reached a peak tracking error of
\SI{0.27}{\metre} and a rotor disk passed \SI{0.34}{\metre} from the user's
eye. With \cref{eq:omegad} the same mission gives \SI{\rTurnTrack}{\milli\metre}
and \SI{\rTurnClear}{\metre}.
\end{remark}

\begin{figure}[tbp]
\centering
% Why a curved path needs angular-rate feedforward. Left: top view of a turn
% of radius rho at speed V; the horizontal part of b3^d points at the centre
% and turns with the path at Omega = V/rho. Right: b3^d sweeps a cone of
% half-angle theta_t = atan(V^2/(rho g)) about e3, so its derivative has
% magnitude Omega sin(theta_t); omega_d supplies that rate.
\begin{tikzpicture}[font=\footnotesize]
  % ---- left: top view ----
  \begin{scope}
    \def\rr{1.9}
    \draw[construct] (0,0) circle (\rr);
    \fill[cG] (0,0) circle (1.2pt);
    \foreach \a in {90, 30, -30, -90} {
      \coordinate (P) at (\a:\rr);
      \draw[cA, line width=0.6pt] ($(P)+(\a+45:0.16)$) -- ($(P)+(\a+225:0.16)$)
           ($(P)+(\a-45:0.16)$) -- ($(P)+(\a+135:0.16)$);
      \fill[cA] (P) circle (1.3pt);
      \draw[vec, cA] (P) -- ($(P)!0.36!(0,0)$);
      \draw[vec, cB] (P) -- ($(P)+(\a-90:0.62)$);
    }
    \node[cB, anchor=north west] at ($(-90:\rr)+(-180:0.62)+(0.0,-0.02)$) {$\vect v$};
    \node[cA, align=left, anchor=west] at ($(90:\rr)!0.4!(0,0)+(0.1,0.05)$)
         {horizontal part\\ of $\vect b_3^d$};
    % the motion is clockwise (velocity at angle a - 90 degrees)
    \draw[-{Stealth[length=2mm]}, cE, line width=0.8pt]
         (260:0.8) arc (260:190:0.8);
    \node[cE, anchor=east] at (228:0.85) {$\Omega=V/\rho$};
    \draw[dim] (0,0) -- (145:\rr) node[pos=0.6, above, sloped] {$\rho$};
    \node[anchor=north] at (0,-\rr-0.4) {(a) top view of a turn};
  \end{scope}
  % ---- right: the cone swept by b3^d ----
  \begin{scope}[xshift=6.6cm, yshift=-1.7cm]
    \def\L{3.2}
    \def\th{26}
    \pgfmathsetmacro{\rc}{\L*sin(\th)}
    \pgfmathsetmacro{\hc}{\L*cos(\th)}
    \coordinate (O) at (0,0);
    \coordinate (Z) at (0,1);
    \coordinate (S) at (\rc,\hc);
    \draw[vec, cG] (O) -- (0,\L+0.5) node[above] {$\ez$};
    \draw[construct] (0,\hc) ellipse ({\rc} and {0.3*\rc});
    \draw[construct] (O) -- (-\rc,\hc);
    \draw[construct] (O) -- (S);
    \coordinate (B1) at ({\rc*cos(-110)},{\hc+0.3*\rc*sin(-110)});
    \coordinate (B2) at ({\rc*cos(-55)},{\hc+0.3*\rc*sin(-55)});
    \draw[vec, cA!50] (O) -- (B1) node[pos=0.62, left] {$\vect b_3^d(t)$};
    \draw[vec, cA] (O) -- (B2) node[pos=1, below right=-1pt] {$\vect b_3^d(t+\delta t)$};
    \draw[-{Stealth[length=2mm]}, cE, line width=0.8pt, domain=-120:-60, samples=24]
         plot ({1.15*\rc*cos(\x)},{\hc+0.345*\rc*sin(\x)});
    \node[cE, anchor=west, align=left] at ({1.15*\rc},{\hc+0.35})
         {$\|\dt{\vect b}{}_3^d\| = \Omega\sin\theta_t$,\\
          supplied by $\vect\omega_d$ of \labelcref{eq:omegad}};
    \pic[draw, cG, angle radius=1.2cm, "$\theta_t$", angle eccentricity=1.3]
         {angle=S--O--Z};
    \node[anchor=north, align=center] at (0.6,-0.3)
         {(b) $\vect b_3^d$ sweeps a cone about $\ez$,\\ $\tan\theta_t = V^2/(\rho g)$};
  \end{scope}
\end{tikzpicture}

\caption{Why a curved path needs angular-rate feedforward. (a) On a turn of
radius $\rho$ flown at speed $V$ the commanded force has a horizontal part
pointing at the centre of the turn, so the horizontal part of $\vect b_3^d$
turns with the path at $\Omega=V/\rho$. (b) $\vect b_3^d$ therefore sweeps a cone
of half-angle $\theta_t$ about $\ez$, with $\tan\theta_t=V^2/(\rho g)$, and changes
at the rate $\Omega\sin\theta_t$ even in a perfectly steady turn. With
$\vect\omega_d=0$ the attitude loop sees that rate as a persistent error and
settles at a lag of about $K_\omega\|\vect\omega_d\|/K_R$; \cref{eq:omegad}
supplies the rate instead.}
\label{fig:thrustrot}
\end{figure}

\subsection{Control authority and gain selection}

\Cref{thm:lee} says nothing about saturation, and saturation is what actually
limits this vehicle.

\begin{definition}[Torque and angular acceleration authority]
Holding total thrust at $mg$ and the other two moments at zero, the authority
about body axis $j$ is the smaller of the two one-sided maxima
\begin{equation}
  \tau_j^{\max} = \min\{\tau_j^{+},\tau_j^{-}\},
  \qquad
  \tau_j^{\pm} = \max\Big\{\pm\tau :
     \mat M\vect u = (mg,\ \tau\vect e_j),\
     0\le u_i\le\Omega_{\max}^2\Big\} ,
  \label{eq:authority}
\end{equation}
a linear program in $\vect u=(\Omega_1^2,\dots,\Omega_N^2)$ that the engine
solves for every layout, with $\Omega_{\max}^2 = \lambda\,\Omega_h^2$ and
$\Omega_h^2 = mg/(Nk_T)$. For the symmetric quadrotor with alternating spins
it has the closed form (\cref{app:proofs})
\begin{equation}
  \tau_j^{\max} = \Delta\sum_{i=1}^{4}\big|\mat M_{j+1,i}\big|,
  \qquad
  \Delta = \min\{\Omega_h^2,\ \Omega_{\max}^2-\Omega_h^2\} ,
  \label{eq:authquad}
\end{equation}
the differential being bounded by whichever is smaller, the room to speed two
rotors up or the room to slow the other two down. The angular acceleration
authority is $\alpha_j^{\max} = (J^{-1})_{jj}\,\tau_j^{\max}$, which is
$\tau_j^{\max}/J_{jj}$ only when $J$ is diagonal.
\end{definition}

For the design of \cref{sec:results}, \cref{eq:authority} gives
$\tau_{\mathrm{roll}}^{\max}=\tau_{\mathrm{pitch}}^{\max}=\SI{\rTauRoll}{\newton\metre}$
and $\tau_{\mathrm{yaw}}^{\max}=\SI{\rTauYaw}{\newton\metre}$, i.e.\
$\alpha^{\max}_{\mathrm{roll}/\mathrm{pitch}/\mathrm{yaw}} =
\rAlphaRoll/\rAlphaPitch/\rAlphaYaw\ \si{\radian\per\second\squared}$. An earlier
version of \cref{eq:authority} divided the sum by $N$, on the reasoning that
the thrust constraint ``shares'' the differential among the rotors; it does
not, the constraint is met by pairing increments with decrements, and the
error was a factor of four. It was found by independent review and is
recorded in \cref{sec:verification}. \Cref{fig:attainable} draws the whole
attainable set, of which \cref{eq:authority} takes the axis intercepts.

\begin{figure}[tbp]
\centering
% Attainable torque sets at thrust mg and zero yaw moment, computed from the
% engine by one linear program per direction (the support function).
\providecommand{\fdQuadRoll}{3.66}
\providecommand{\fdQuadPitch}{3.66}
\providecommand{\fdQuadOld}{0.91}
\providecommand{\fdHexRoll}{6.13}
\providecommand{\fdHexPitch}{7.07}
\providecommand{\fdHexGross}{2.565}
\providecommand{\fdQuadGross}{1.724}
\providecommand{\fdKeepRadius}{1.180}
\providecommand{\fdKeepSlice}{1.174}
\providecommand{\fdEnvRadius}{1.080}
\providecommand{\fdTurnT}{5.65}
\providecommand{\fdTurnHeadX}{-1.440}
\providecommand{\fdTurnHeadY}{0.000}
\providecommand{\fdTurnPeriod}{7.54}
\providecommand{\fdStandoff}{1.20}
\providecommand{\fdRootLo}{4.5678}
\providecommand{\fdRootHi}{4.6173}
\providecommand{\fdPeakM}{4.5925}
\providecommand{\fdPeakF}{0.03}
\providecommand{\fdGridLo}{4.460}
\providecommand{\fdGridHi}{6.021}
\providecommand{\fdGridLoF}{-0.001}
\providecommand{\fdGridHiF}{-0.083}
\providecommand{\fdIslandPct}{1.1}
\providecommand{\fdPlainRate}{0.263}
\providecommand{\fdConvWindow}{8}
\providecommand{\fdHlMotorRe}{-0.0889}
\providecommand{\fdHlMotorIm}{0.0000}
\providecommand{\fdHlAttRe}{-0.0240}
\providecommand{\fdHlAttIm}{0.0000}
\providecommand{\fdHlPosRe}{-0.0082}
\providecommand{\fdHlPosIm}{0.0004}
\providecommand{\fdHlGovRe}{-0.0160}
\providecommand{\fdHlGovIm}{0.0000}
\providecommand{\fdWpos}{2.04}
\providecommand{\fdWgov}{4.00}
\providecommand{\fdWatt}{6.43}
\providecommand{\fdWmot}{22.2}
\providecommand{\fdWarm}{206}
\providecommand{\fdWnyq}{785}
\providecommand{\fdRatioPA}{3.2}
\providecommand{\fdRatioAM}{3.5}
\providecommand{\fdKR}{41.4}
\providecommand{\fdKp}{4.17}
\providecommand{\fdTauMms}{45}
\providecommand{\fdH}{4}
\providecommand{\fdMotorQ}{0.289}
\providecommand{\fdMotorM}{0.0609}

\begin{tikzpicture}
  \begin{axis}[paperplot, width=0.62\linewidth, height=0.62\linewidth,
      axis equal image,
      xmin=-8, xmax=8, ymin=-8, ymax=8,
      xlabel={$\tau_x$ (roll) [\si{\newton\metre}]},
      ylabel={$\tau_y$ (pitch) [\si{\newton\metre}]},
      legend style={at={(0.5,1.02)}, anchor=south, legend columns=2},
    ]
    \addplot[cA, fill=cA!14] table[x=tx, y=ty] {results/fig/d_attain_quad.dat};
    \addlegendentry{quadrotor, the design}
    \addplot[cB, dashed, line width=1.1pt] table[x=tx, y=ty] {results/fig/d_attain_hex.dat};
    \addlegendentry{hexacopter, same payload and rotors}
    \addplot[only marks, mark=square*, mark size=1.9pt, cB, forget plot]
      coordinates {(\fdHexRoll,0) (-\fdHexRoll,0) (0,\fdHexPitch) (0,-\fdHexPitch)};
    % authority of the quadrotor on each axis
    \addplot[only marks, mark=*, mark size=2.2pt, cA, forget plot]
      coordinates {(\fdQuadRoll,0) (-\fdQuadRoll,0) (0,\fdQuadPitch) (0,-\fdQuadPitch)};
    \node[lbl, cA, anchor=south west] at (axis cs:\fdQuadRoll,0.2)
      {$\tau^{\max}_{\mathrm{roll}}$};
    \node[lbl, cA, anchor=south west] at (axis cs:0.2,\fdQuadPitch)
      {$\tau^{\max}_{\mathrm{pitch}}$};
    % the earlier estimate with the factor 1/N
    \addplot[only marks, mark=x, mark size=3.2pt, cD, line width=1pt, forget plot]
      coordinates {(\fdQuadOld,0) (-\fdQuadOld,0) (0,\fdQuadOld) (0,-\fdQuadOld)};
    \node[lbl, cD, anchor=north west, align=left] at (axis cs:0.15,-\fdQuadOld)
      {earlier estimate\\ with $1/N$: \SI{\fdQuadOld}{\newton\metre}};
    \draw[construct] (axis cs:-8,0) -- (axis cs:8,0);
    \draw[construct] (axis cs:0,-8) -- (axis cs:0,8);
  \end{axis}
\end{tikzpicture}

\caption{Attainable torque sets in the roll and pitch plane with thrust held
at $mg$ and the yaw moment at zero, computed from the engine through their
support function: one linear program of the form \cref{eq:authority} for each
of \num{240} directions. The quadrotor's set is a square turned by \ang{45}
whose axis intercepts are
$\tau^{\max}_{\mathrm{roll}}=\tau^{\max}_{\mathrm{pitch}}=\SI{\fdQuadRoll}{\newton\metre}$,
the balanced differential of \cref{eq:authquad}; the crosses mark the earlier
estimate with the factor $1/N$, a quarter of it. The hexacopter of
\cref{tab:hex}, re-optimised under the same limits (gross
\SI{\fdHexGross}{\kilo\gram}), has a hexagonal set that reaches further in
pitch than in roll: its roll
differential also produces a yaw moment, which the zero-yaw constraint must
cancel, while its pitch differential does not. Its authority in the sense of
\cref{eq:authority} is the smaller intercept, roll.}
\label{fig:attainable}
\end{figure}

\begin{model}[Gain heuristic]
\label{mod:gains}
With $e_{\max}$ the largest attitude error the loop is expected to correct and
$\zeta$ a damping ratio, the engine sets
\begin{equation}
  K_R = \min\left\{\frac{\alpha^{\max}}{2\,e_{\max}},\
                   \Big(\frac{0.3}{\tau_m}\Big)^{2}\right\},
  \qquad
  \alpha^{\max} \coloneqq \min\{\alpha_{\mathrm{roll}}^{\max},
                                \alpha_{\mathrm{pitch}}^{\max}\},
  \qquad
  K_\omega = 2\zeta\sqrt{K_R} ,
  \label{eq:kr}
\end{equation}
and, one level up,
$K_p = \min\{a_{\max}/e_{p,\max},\ K_R/9\}$, $K_d = 2\zeta\sqrt{K_p}$, with
$a_{\max}$ the acceleration cap of the reference governor. The second term
keeps the position bandwidth $\sqrt{K_p}$ at least a factor three below the
attitude bandwidth $\sqrt{K_R}$, which is the time-scale separation the
cascade of \cref{eq:chain} rests on; a machine whose attitude authority has
been cut, by a display on a boom, is destabilised by position gains that were
fine without the boom.
\end{model}

The first bound is the linearised saturation argument: with $\vect e_\omega$
small the commanded angular acceleration is $K_R\|\vect e_R\|\le K_R e_{\max}$,
and half the single-axis authority is reserved for the rate term, for
simultaneous demands on the other axes and for yaw. The second bound keeps the
attitude bandwidth $\sqrt{K_R}$ below $0.3/\tau_m$, so that the first-order
motor lag of \cref{eq:motorlag} contributes less than about \ang{17} of phase
at crossover. $K_\omega$ follows from the linearised closed loop
$\ddt e + K_\omega\dt e + K_R e = 0$. None of this is a saturation certificate:
it neglects the rate feedback term, simultaneous torque demands, the yaw axis
and the actuator dynamics, and \cref{sec:results} verifies by simulation that
the resulting loop stays out of saturation on the missions flown. For the
design the two bounds are \SI{\rKRauth}{\per\second\squared} and
\SI{\rKRlag}{\per\second\squared}, the lag bound binds, and the value used is
$K_R=\SI{\rKR}{\per\second\squared}$, $K_\omega=\SI{\rKw}{\per\second}$.

\begin{figure}[htbp]
\centering
% The frequency ladder of the cascade: each loop is a factor of about three
% or more slower than the one it relies on, and the sampling rate is two
% orders of magnitude above the fastest loop.
\providecommand{\fdQuadRoll}{3.66}
\providecommand{\fdQuadPitch}{3.66}
\providecommand{\fdQuadOld}{0.91}
\providecommand{\fdHexRoll}{6.13}
\providecommand{\fdHexPitch}{7.07}
\providecommand{\fdHexGross}{2.565}
\providecommand{\fdQuadGross}{1.724}
\providecommand{\fdKeepRadius}{1.180}
\providecommand{\fdKeepSlice}{1.174}
\providecommand{\fdEnvRadius}{1.080}
\providecommand{\fdTurnT}{5.65}
\providecommand{\fdTurnHeadX}{-1.440}
\providecommand{\fdTurnHeadY}{0.000}
\providecommand{\fdTurnPeriod}{7.54}
\providecommand{\fdStandoff}{1.20}
\providecommand{\fdRootLo}{4.5678}
\providecommand{\fdRootHi}{4.6173}
\providecommand{\fdPeakM}{4.5925}
\providecommand{\fdPeakF}{0.03}
\providecommand{\fdGridLo}{4.460}
\providecommand{\fdGridHi}{6.021}
\providecommand{\fdGridLoF}{-0.001}
\providecommand{\fdGridHiF}{-0.083}
\providecommand{\fdIslandPct}{1.1}
\providecommand{\fdPlainRate}{0.263}
\providecommand{\fdConvWindow}{8}
\providecommand{\fdHlMotorRe}{-0.0889}
\providecommand{\fdHlMotorIm}{0.0000}
\providecommand{\fdHlAttRe}{-0.0240}
\providecommand{\fdHlAttIm}{0.0000}
\providecommand{\fdHlPosRe}{-0.0082}
\providecommand{\fdHlPosIm}{0.0004}
\providecommand{\fdHlGovRe}{-0.0160}
\providecommand{\fdHlGovIm}{0.0000}
\providecommand{\fdWpos}{2.04}
\providecommand{\fdWgov}{4.00}
\providecommand{\fdWatt}{6.43}
\providecommand{\fdWmot}{22.2}
\providecommand{\fdWarm}{206}
\providecommand{\fdWnyq}{785}
\providecommand{\fdRatioPA}{3.2}
\providecommand{\fdRatioAM}{3.5}
\providecommand{\fdKR}{41.4}
\providecommand{\fdKp}{4.17}
\providecommand{\fdTauMms}{45}
\providecommand{\fdH}{4}
\providecommand{\fdMotorQ}{0.289}
\providecommand{\fdMotorM}{0.0609}

\begin{tikzpicture}
  \begin{axis}[
      width=\linewidth, height=4.6cm, xmode=log,
      xmin=1, xmax=2000, ymin=-1.9, ymax=2.2,
      axis y line=none, axis x line*=bottom,
      xlabel={angular frequency [\si{\radian\per\second}]},
      tick label style={font=\footnotesize}, label style={font=\small},
      clip=false]
    \draw[cG] (axis cs:1,0) -- (axis cs:2000,0);
    \addplot[only marks, mark=*, mark size=2.4pt, cA] coordinates
      {(\fdWpos,0) (\fdWatt,0)};
    \addplot[only marks, mark=square*, mark size=2.2pt, cB] coordinates
      {(\fdWgov,0)};
    \addplot[only marks, mark=triangle*, mark size=2.8pt, cE] coordinates
      {(\fdWmot,0) (\fdWarm,0)};
    \addplot[only marks, mark=diamond*, mark size=2.8pt, cG] coordinates
      {(\fdWnyq,0)};
    \node[lbl, anchor=north, align=center] at (axis cs:\fdWpos,-0.3)
      {position\\ $\sqrt{K_p}=\fdWpos$};
    \node[lbl, anchor=south, align=center, cB] at (axis cs:\fdWgov,0.3)
      {governor\\ $\sqrt{k_p}=4$};
    \node[lbl, anchor=north, align=center] at (axis cs:\fdWatt,-0.3)
      {attitude\\ $\sqrt{K_R}=\fdWatt$};
    \node[lbl, anchor=south, align=center, cE] at (axis cs:\fdWmot,0.3)
      {motor\\ $1/\tau_m=\fdWmot$};
    \node[lbl, anchor=north, align=center, cE] at (axis cs:\fdWarm,-0.3)
      {first arm bending\\ mode, $2\pi f_1=\fdWarm$};
    \node[lbl, anchor=south, align=center, cG] at (axis cs:\fdWnyq,0.3)
      {sampling\\ $\pi/h=\fdWnyq$};
    \draw[dim] (axis cs:\fdWpos,1.55) -- (axis cs:\fdWatt,1.55)
      node[midway, above, font=\scriptsize] {$\times\fdRatioPA$};
    \draw[dim] (axis cs:\fdWatt,1.55) -- (axis cs:\fdWmot,1.55)
      node[midway, above, font=\scriptsize] {$\times\fdRatioAM$};
    \draw[construct] (axis cs:\fdWpos,0.1) -- (axis cs:\fdWpos,1.55);
    \draw[construct] (axis cs:\fdWatt,0.1) -- (axis cs:\fdWatt,1.55);
    \draw[construct] (axis cs:\fdWmot,1.2) -- (axis cs:\fdWmot,1.55);
  \end{axis}
\end{tikzpicture}

\caption{The frequency ladder of the design's cascade on a logarithmic axis.
Each loop is at least a factor of three slower than the one it relies on:
position at $\sqrt{K_p}$, attitude at $\sqrt{K_R}$, the motors at $1/\tau_m$.
The governor's filter at $\sqrt{k_p}$ shapes the reference and closes no loop
on the vehicle. The first bending mode of the arms and the sampling rate of
the controller lie one and two orders of magnitude above the fastest loop,
which is the condition under which the airframe may be treated as rigid and
the sampled controller as continuous.}
\label{fig:bandwidth}
\end{figure}

\begin{remark}[Yaw is the weak axis]
By \cref{eq:alloc} the yaw column involves $k_Q$ rather than $k_T L$, and
$k_Q/k_T \sim \Cdz$-scale quantities. A rotor optimised for quietness by
\cref{obs:sixthpower} turns slowly and has small reaction torque, so
$\alpha_{\mathrm{yaw}}^{\max}$ is an order of magnitude below the roll and
pitch values: \SI{\rAlphaYaw}{\radian\per\second\squared} against
\SI{\rAlphaRoll}{\radian\per\second\squared} for the design. A quiet hovering
display therefore cannot turn to face its user quickly. The engine responds by
governing the heading reference as it governs position, with a rate cap of
\SI{1.5}{\radian\per\second} and an acceleration cap of half the yaw
authority; the architectural response, a yaw axis on the gimbal where a
\SI{30}{\gram} motor solves what the airframe cannot, is not modelled.
\end{remark}

\subsection{Control allocation under saturation}

Given a desired wrench $\vect\chi_d = (T_c,\vect\tau)$ the allocation problem is
\begin{equation}
  \text{find } \vect u\in[0,\Omega_{\max}^2]^N
  \quad\text{with}\quad \mat M\vect u = \vect\chi_d .
  \label{eq:allocprob}
\end{equation}
When $\vect\chi_d$ lies outside $\mat M[0,\Omega_{\max}^2]^N$ no exact solution
exists and a priority must be chosen. We solve the unconstrained least-norm
problem $\vect u = \mat M^{+}\vect\chi_d$ and, on violation, iterate
\begin{equation}
  \vect u \leftarrow \Pi_{[0,\Omega_{\max}^2]}(\vect u),
  \qquad
  \vect u \leftarrow \vect u + \mat M^{+}
    \big(\vect\chi_d - \mat M\vect u\big)\odot \vect\nu ,
  \label{eq:allociter}
\end{equation}
with weights $\vect\nu$ that de-prioritise the thrust channel. The
engineering content is the priority: \emph{attitude is preserved and thrust is
sacrificed}, because losing attitude loses the vehicle whereas losing thrust
loses altitude, which is recoverable.

\subsection{The reference governor}

Tracking a person's face rigidly is the wrong specification, and this is a
statement about kinematics rather than taste.

\begin{proposition}[The cost of a rigid standoff during a turn]
\label{prop:turn}
Let the user rotate their heading by $\pi$ in time $\Delta t$ while the display
is required to remain at standoff $d$ ahead of them. Then the display must
traverse a semicircle of radius $d$, and its mean speed and peak centripetal
acceleration satisfy
\begin{equation}
  \bar v = \frac{\pi d}{\Delta t},
  \qquad
  a_c = \frac{\bar v^2}{d} = \frac{\pi^2 d}{\Delta t^2} .
  \label{eq:turncost}
\end{equation}
\end{proposition}

For $d=\SI{0.7}{\metre}$ and $\Delta t = \SI{1}{\second}$ this is
\SI{2.2}{\metre\per\second} and \SI{6.9}{\metre\per\second\squared}, and the
simulated demand including the translation is over
\SI{10}{\metre\per\second\squared}. The engine confirms that increasing the
thrust margin makes the outcome \emph{worse}, because the controller then tilts
harder and the gimbal cannot keep up.

\begin{model}[Rate-limited reference with a moving keep-out]
\label{mod:governor}
Let $(\vect r_{\mathrm{raw}},\vect v_{\mathrm{raw}},\vect a_{\mathrm{raw}},
\vect j_{\mathrm{raw}})$ be the ideal display trajectory and its derivatives,
and let $\vect c(t)$ be the user's head with velocity $\dt{\vect c}$ and
acceleration $\ddt{\vect c}$. The governed reference is the state
$(\vect r_g,\vect v_g)$ of a second-order filter with acceleration feedforward,
\begin{equation}
  \vect a^{\star} = \vect a_{\mathrm{raw}}
    + k_p(\vect r_{\mathrm{raw}}-\vect r_g)
    + k_d(\vect v_{\mathrm{raw}}-\vect v_g) ,
  \label{eq:gov1}
\end{equation}
whose acceleration is projected, by alternating projections, onto the
intersection of three convex sets: the acceleration ball
$\|\vect a\|\le a_{\max}$, the next-step speed ball
$\|\vect v_g+\vect a\,\dd t\|\le v_{\max}$, and the keep-out half-space
$\vect n\transp\vect a\ge\ell$. With
$\varrho = d_{\mathrm{safe}}+\mathcal{R}+\varepsilon$ the keep-out radius,
$\vect n = (\vect r_g-\vect c)/\|\vect r_g-\vect c\|$,
$h = \|\vect r_g-\vect c\|-\varrho$, $\dt h = \vect n\transp(\vect v_g-\dt{\vect c})$
and $\|\vect v_t\|^2 = \|\vect v_g-\dt{\vect c}\|^2-\dt h^2$,
\begin{equation}
  \ell = \max\left\{
    \vect n\transp\ddt{\vect c} - \frac{\|\vect v_t\|^2}{\|\vect r_g-\vect c\|}
      - 2\gamma\,\dt h - \gamma^2 h ,\;
    \vect n\transp\ddt{\vect c} - \frac{2\,(h + \dt h\,\dd t)}{\dd t^{2}}
  \right\},
  \qquad \gamma = 4 .
  \label{eq:barrier}
\end{equation}
The state then advances by
$\vect r_g\leftarrow\vect r_g+\vect v_g\,\dd t+\tfrac12\vect a\,\dd t^2$,
$\vect v_g\leftarrow\vect v_g+\vect a\,\dd t$. There is no position projection.
If the three sets have empty intersection the step is flagged infeasible, the
acceleration is scaled back to $a_{\max}$, and the violation is reported. The
heading reference is governed separately by the same filter in one dimension
with a rate cap and an acceleration cap.
\end{model}

\begin{proposition}[What the keep-out guarantees]
\label{prop:governor}
\begin{enumerate}[label=(\roman*),nosep]
  \item The caps alone do not keep the reference out of the user. On a curved
        raw path the acceleration cap makes the governed reference cut the
        corner, and since the raw path lies at distance $d$ from the user by
        construction, the governed reference can approach the user
        arbitrarily closely as $a_{\max}$ decreases.
  \item In continuous time, and while the first term of \cref{eq:barrier} is
        the active constraint, the set
        $\{h\ge0,\ \dt h+\gamma h\ge0\}$ is forward invariant: a reference
        that starts outside the keep-out sphere with an admissible closing
        rate never enters it, whatever the user does.
  \item The second term of \cref{eq:barrier} is the one-step condition
        $h_{k+1}\ge0$ under constant $\ddt{\vect c}$ over the step, linearised
        in the tangent plane. It is what the discrete update enforces.
  \item The vehicle, not the reference, is what must stay clear. If
        $\sup_t\|\vect r(t)-\vect r_g(t)\|\le\varepsilon$ on a mission, then on
        that mission the vehicle's centre stays at least
        $d_{\mathrm{safe}}+\mathcal{R}$ from the head, and hence every blade tip
        at least $d_{\mathrm{safe}}$. The hypothesis is checked per run and
        reported; it is not proved.
\end{enumerate}
\end{proposition}

\begin{proof}
(i) In the limit of small $a_{\max}$ the filter output is a convex combination
of past raw positions, and a chord of a circular arc lies inside the arc, so
the reference lies inside the sphere of radius $d$ about the user.
(ii) Differentiating $h$ twice along the governed motion, with
$\dt{\vect n} = \vect v_t/\|\vect r_g-\vect c\|$ in the notation above,
\[
  \ddt h = \vect n\transp(\vect a-\ddt{\vect c})
           + \frac{\|\vect v_t\|^2}{\|\vect r_g-\vect c\|} ,
\]
so the constraint $\vect n\transp\vect a\ge\ell_1$ with $\ell_1$ the first
term of \cref{eq:barrier} is exactly $\ddt h\ge-2\gamma\dt h-\gamma^2 h$. Put
$\psi=\dt h+\gamma h$; then $\dt\psi = \ddt h+\gamma\dt h\ge-\gamma\psi$, so
$\psi(t)\ge\psi(0)e^{-\gamma t}\ge0$, and then
$\dt h\ge-\gamma h$ gives $h(t)\ge h(0)e^{-\gamma t}\ge0$. This is the
second-order barrier construction of Ames et al.~\cite{ames2017}, with both
poles at $-\gamma$ to match the filter's own $k_p=\gamma^2$, $k_d=2\gamma$.
(iii) Expand $\|\vect r_g+\vect v_g\dd t+\tfrac12\vect a\,\dd t^2 -
\vect c-\dt{\vect c}\dd t-\tfrac12\ddt{\vect c}\dd t^2\|$ to first order in the
tangent plane about $\vect n$. (iv) Triangle inequality:
$\|\vect r-\vect c\|\ge\|\vect r_g-\vect c\|-\|\vect r-\vect r_g\|
\ge\varrho-\varepsilon = d_{\mathrm{safe}}+\mathcal{R}$, and the reach bound
\cref{eq:reach}.
\end{proof}

\Cref{fig:barrier} draws part (ii) in the phase plane of $h$.

\begin{figure}[tbp]
\centering
% Phase plane of the keep-out barrier h = ||r_g - c|| - rho. The admissible
% set {h >= 0, hdot + gamma h >= 0} (shaded) is forward invariant under the
% extremal dynamics hddot = -2 gamma hdot - gamma^2 h that the constraint
% allows; a start with hdot + gamma h < 0 is not protected.
% Closed form with psi0 = hdot0 + gamma h0:
%   h(t)    = (h0 + psi0 t) exp(-gamma t)
%   hdot(t) = (hdot0 - gamma psi0 t) exp(-gamma t)
\begin{tikzpicture}
  \pgfmathsetmacro{\gam}{4}
  \begin{axis}[paperplot, height=6.0cm,
      xmin=-0.15, xmax=1.0, ymin=-3.4, ymax=0.9,
      xlabel={$h = \|\vect r_g-\vect c\|-\varrho$ [\si{\metre}]},
      ylabel={$\dt h$ [\si{\metre\per\second}]},
      legend pos=south east]
    % admissible set
    \fill[cC!12] (axis cs:0,0) -- (axis cs:0,0.9) -- (axis cs:1.0,0.9)
                 -- (axis cs:1.0,-4.0) -- cycle;
    \draw[cC!70!black, line width=0.9pt] (axis cs:0,0.9) -- (axis cs:0,0)
         -- (axis cs:0.85,-3.4);
    \node[lbl, cC!60!black, anchor=north east] at (axis cs:0.64,-2.62)
         {$\dt h+\gamma h=0$};
    \node[lbl, cC!60!black, align=center] at (axis cs:0.62,0.45)
         {admissible set\\ $\{h\ge0,\ \dt h+\gamma h\ge0\}$};
    \fill[cD!10] (axis cs:-0.15,-3.4) rectangle (axis cs:0,0.9);
    \node[lbl, cD, rotate=90] at (axis cs:-0.075,-1.25) {inside the keep-out sphere};
    % extremal trajectories from admissible starts
    \foreach \hz/\hd in {0.5/-1.5, 0.8/-3.0, 0.2/-0.6, 0.9/0.5} {
      \addplot[cA, domain=0:2.5, samples=90, variable=\t,
               postaction={decorate}, decoration={markings,
               mark=at position 0.25 with {\arrow{Stealth[length=2mm]}}}]
        ({(\hz + (\hd+\gam*\hz)*\t)*exp(-\gam*\t)},
         {(\hd - \gam*(\hd+\gam*\hz)*\t)*exp(-\gam*\t)});
      \addplot[only marks, mark=*, mark size=1.6pt, cA] coordinates {(\hz,\hd)};
    }
    % an inadmissible start: closing too fast for the distance left
    \addplot[cD, dashed, domain=0:1.2, samples=80, variable=\t,
             postaction={decorate}, decoration={markings,
             mark=at position 0.3 with {\arrow{Stealth[length=2mm]}}}]
      ({(0.3 + (-2.5+\gam*0.3)*\t)*exp(-\gam*\t)},
       {(-2.5 - \gam*(-2.5+\gam*0.3)*\t)*exp(-\gam*\t)});
    \fill[cD] (axis cs:0.3,-2.5) circle (1.6pt);
    \node[lbl, cD, anchor=west] at (axis cs:0.31,-2.62) {$\dt h_0+\gamma h_0<0$};
  \end{axis}
\end{tikzpicture}

\caption{The phase plane of the keep-out barrier $h=\|\vect r_g-\vect c\|-\varrho$
with $\gamma=4$. The shaded set $\{h\ge0,\ \dt h+\gamma h\ge0\}$ is forward
invariant under the slowest recession that the first term of
\cref{eq:barrier} permits, $\ddt h=-2\gamma\dt h-\gamma^2h$, whose solutions
$h(t)=(h_0+\psi_0t)\,e^{-\gamma t}$ with $\psi_0=\dt h_0+\gamma h_0$ are drawn
from four admissible starts: each reaches the origin without crossing $h=0$.
A start that closes faster than $\gamma h_0$ (dashed) is not protected and
enters the sphere. This is why the governor maintains the invariant set and
not the half-space $h\ge0$ alone.}
\label{fig:barrier}
\end{figure}

Three things the proposition does not say. The invariance in (ii) is a
statement about the continuous-time reference; the engine enforces (iii) at
update instants, and between them the reference moves under a constant
acceleration, so (ii) is the limit the discrete scheme approaches, not what it
executes. The guarantee lapses whenever the step is infeasible, which is why
that is reported as a violation rather than absorbed. And $\vect c$ is the
true head position: an estimator with latency or dropouts moves the sphere
away from the head, and no estimator is modelled (\cref{as:estimation}).

\Cref{fig:turnaround} shows the governor on the most demanding mission of
\cref{sec:results}, the turnaround.

\begin{figure}[tbp]
\centering
% The turnaround flown by the engine. (a) Top view around one reversal:
% the user's head, the ideal display trajectory (a semicircle about the
% head), the governed reference and the vehicle, with the keep-out circle at
% the reversal. (b) Distances from the head over the whole window.
\providecommand{\fdQuadRoll}{3.66}
\providecommand{\fdQuadPitch}{3.66}
\providecommand{\fdQuadOld}{0.91}
\providecommand{\fdHexRoll}{6.13}
\providecommand{\fdHexPitch}{7.07}
\providecommand{\fdHexGross}{2.565}
\providecommand{\fdQuadGross}{1.724}
\providecommand{\fdKeepRadius}{1.180}
\providecommand{\fdKeepSlice}{1.174}
\providecommand{\fdEnvRadius}{1.080}
\providecommand{\fdTurnT}{5.65}
\providecommand{\fdTurnHeadX}{-1.440}
\providecommand{\fdTurnHeadY}{0.000}
\providecommand{\fdTurnPeriod}{7.54}
\providecommand{\fdStandoff}{1.20}
\providecommand{\fdRootLo}{4.5678}
\providecommand{\fdRootHi}{4.6173}
\providecommand{\fdPeakM}{4.5925}
\providecommand{\fdPeakF}{0.03}
\providecommand{\fdGridLo}{4.460}
\providecommand{\fdGridHi}{6.021}
\providecommand{\fdGridLoF}{-0.001}
\providecommand{\fdGridHiF}{-0.083}
\providecommand{\fdIslandPct}{1.1}
\providecommand{\fdPlainRate}{0.263}
\providecommand{\fdConvWindow}{8}
\providecommand{\fdHlMotorRe}{-0.0889}
\providecommand{\fdHlMotorIm}{0.0000}
\providecommand{\fdHlAttRe}{-0.0240}
\providecommand{\fdHlAttIm}{0.0000}
\providecommand{\fdHlPosRe}{-0.0082}
\providecommand{\fdHlPosIm}{0.0004}
\providecommand{\fdHlGovRe}{-0.0160}
\providecommand{\fdHlGovIm}{0.0000}
\providecommand{\fdWpos}{2.04}
\providecommand{\fdWgov}{4.00}
\providecommand{\fdWatt}{6.43}
\providecommand{\fdWmot}{22.2}
\providecommand{\fdWarm}{206}
\providecommand{\fdWnyq}{785}
\providecommand{\fdRatioPA}{3.2}
\providecommand{\fdRatioAM}{3.5}
\providecommand{\fdKR}{41.4}
\providecommand{\fdKp}{4.17}
\providecommand{\fdTauMms}{45}
\providecommand{\fdH}{4}
\providecommand{\fdMotorQ}{0.289}
\providecommand{\fdMotorM}{0.0609}

\begin{tikzpicture}
  \pgfmathsetmacro{\ta}{\fdTurnT-1.5}
  \pgfmathsetmacro{\tb}{\fdTurnT+2.2}
  \begin{axis}[paperplot, name=top, width=0.47\linewidth,
      axis equal image, xmin=-3.0, xmax=2.3, ymin=-1.35, ymax=1.65,
      xlabel={$x$ [\si{\metre}]}, ylabel={$y$ [\si{\metre}]},
      title={(a) top view, $t\in[\pgfmathprintnumber[fixed,precision=1]{\ta},
             \pgfmathprintnumber[fixed,precision=1]{\tb}]\,\si{\second}$},
      window/.style={restrict expr to domain={\thisrow{t}}{\ta:\tb},
                     unbounded coords=jump}]
    \addplot[cE, line width=2pt, window] table[x=hx, y=hy] {results/fig/d_turn.dat};
    \addplot[cG, dashed, window] table[x=rawx, y=rawy] {results/fig/d_turn.dat};
    \addplot[cB, window] table[x=govx, y=govy] {results/fig/d_turn.dat};
    \addplot[cA, window, postaction={decorate}, decoration={markings,
             mark=between positions 0.12 and 0.95 step 0.2 with {\arrow{Stealth[length=2.2mm]}}}]
      table[x=vehx, y=vehy] {results/fig/d_turn.dat};
    \addplot[cD, dashed, domain=0:360, samples=120, line width=0.7pt]
      ({\fdTurnHeadX + \fdKeepSlice*cos(x)}, {\fdTurnHeadY + \fdKeepSlice*sin(x)});
    \addplot[only marks, mark=*, mark size=2.2pt, cE] coordinates {(\fdTurnHeadX,\fdTurnHeadY)};
    \node[lbl, cE, anchor=north] at (axis cs:\fdTurnHeadX,-0.08) {head at the reversal};
  \end{axis}
  \begin{axis}[paperplot, name=dist, at={(top.east)}, anchor=west, xshift=1.7cm,
      width=0.47\linewidth, height=4.9cm,
      xlabel={$t$ [\si{\second}]}, ylabel={distance from the head [\si{\metre}]},
      title={(b) distances over the window},
      xmin=0, xmax=16, ymin=1.0, ymax=1.35,
      legend to name=turnleg, legend columns=3,
      legend style={/tikz/every even column/.append style={column sep=6pt}}]
    \addlegendimage{cE, line width=2pt}
    \addlegendentry{head}
    \addplot[cG, dashed] table[x=t, y=draw] {results/fig/d_turn.dat};
    \addlegendentry{ideal display trajectory}
    \addplot[cB] table[x=t, y=dgov] {results/fig/d_turn.dat};
    \addlegendentry{governed reference $\vect r_g$}
    \addplot[cA] table[x=t, y=dveh] {results/fig/d_turn.dat};
    \addlegendentry{vehicle $\vect r$}
    \addplot[cD, dashed, domain=0:16] {\fdKeepRadius};
    \addlegendentry{keep-out radius $\varrho=d_{\mathrm{safe}}+\mathcal{R}+\varepsilon$}
    \addplot[cC, dotted, line width=1.3pt, domain=0:16] {\fdEnvRadius};
    \addlegendentry{$d_{\mathrm{safe}}+\mathcal{R}$}
    \draw[construct] (axis cs:\fdTurnT,1.0) -- (axis cs:\fdTurnT,1.35);
  \end{axis}
  \node[anchor=north] at ({$(top.outer west)!0.5!(dist.outer east)$} |- dist.outer south)
       {\pgfplotslegendfromname{turnleg}};
\end{tikzpicture}

\caption{The turnaround flown by the engine with the governor of
\cref{mod:governor} at $v_{\max}=\SI{2.5}{\metre\per\second}$,
$a_{\max}=\SI{2.5}{\metre\per\second\squared}$ and
$\varepsilon=\SI{0.10}{\metre}$. (a) Top view around one reversal of the user.
The ideal display trajectory is a semicircle of radius
$d=\SI{\fdStandoff}{\metre}$ about the head, which \cref{prop:turn} shows is
not flyable; the governed reference yields, swings wide and catches up after
the turn, and the vehicle follows it, the arrows giving the direction of
travel. The dashed circle is the keep-out sphere at the reversal, cut at the
vehicle's altitude. (b) Distances from the head over the whole window, the
vertical line marking the reversal shown in (a). The governed reference does
not enter $\varrho$, as \cref{prop:governor}(ii) and (iii) require; the
vehicle falls below $\varrho$ by at most its tracking error and stays above
$d_{\mathrm{safe}}+\mathcal{R}$, which is \cref{prop:governor}(iv) with the
tracking bound measured on this run. The peaks are the display lagging the
user after each reversal.}
\label{fig:turnaround}
\end{figure}

\begin{remark}[The governor's bandwidth must be high]
The gains $k_p,k_d$ in \cref{eq:gov1} must place the filter's natural frequency
well above the frequency content of the mission; otherwise the governor adds a
steady lag on every curved path, which is a different and worse behaviour than
yielding during a turn. With $k_p=16$, $k_d=8$ the filter is critically damped
at \SI{4}{\radian\per\second}, and with the feedforward $\vect a_{\mathrm{raw}}$
in \cref{eq:gov1} it is exactly transparent whenever no cap binds. Without the
feedforward it lags every acceleration by $\vect a_{\mathrm{raw}}/k_p$, about
\SI{5}{\centi\metre} at \SI{0.8}{\metre\per\second\squared}, which was enough
to push the reference into the keep-out sphere a third of the time on a mission
with no turn in it.
\end{remark}

\begin{definition}[Two error measures]
With a governor in the loop, two distinct errors must be reported:
\begin{equation}
  e_{\mathrm{track}} = \|\vect r_g - \vect r\|
  \quad\text{(control performance)},
  \qquad
  e_{\mathrm{lag}} = \|\vect r_{\mathrm{raw}} - \vect r\|
  \quad\text{(what the user sees)} .
  \label{eq:twoerrors}
\end{equation}
Reporting only $e_{\mathrm{track}}$ would flatter the design, since the governor
is free to make its own reference easy to follow. The peak of
$e_{\mathrm{track}}$ over the whole window, start-up included, is the quantity
that must stay below $\varepsilon$ for \cref{prop:governor}(iv).
\end{definition}

\subsection{Hover linearisation and controllability}

Linearising \cref{eq:ode} about the hover trim by central differences gives
$\dt{\delta x} = \mat A\,\delta x + \mat B\,\delta u$ on the $12+N$ rigid-body
and rotor states, after removing the quaternion scalar component which is not
a free coordinate. The trim is computed, not assumed: with the display on a
boom the centre of mass is off the rotor centroid and hover needs unequal
rotor speeds. Controllability is assessed by the rank of the reachable
subspace, computed by an orthogonal iteration
$Q_{k+1} = \operatorname{range}[\mat B\ \ Q_k\ \ \mat A Q_k]$ on the matrices
scaled by their 2-norms with a singular-value tolerance of $10^{-9}$, which
avoids forming the ill-conditioned powers of the Kalman matrix. Controllability
requires the rank to equal the state dimension; for the design of
\cref{sec:results} it is $16$ of $16$. The eigenvalues cluster at the origin,
reflecting the six open-loop integrators of \cref{eq:chain}, which is
precisely why the cascade rather than a single loop is required.


\part{Coupling, Computation and Results}
\section{The coupled co-design problem}
\label{sec:codesign}

\subsection{Statement}

The algebraic layer of \cref{sec:closure} assumed the whole mission is spent
hovering. The differential layer of \cref{sec:dynamics} assumed the mass was
known. Neither is true independently. The honest statement couples them.

\begin{problem}[Co-design]
\label{prob:codesign}
Find a design vector $p$ and a control law $u(\cdot)$ such that
\begin{align}
  \dt{x} &= f(x,u,w;p), \qquad x(0)=x_0(p), \label{eq:cd1}\\
  m &= \mfix + \mstr(p) + \mprop(p) + \mbat, \label{eq:cd2}\\
  \eb(1-\varsigma)\,\mbat &\ge \int_0^{t_f} P\big(x(t),u(t);p\big)\dd t ,
  \qquad
  p_b\,\mbat \ge \max_{t\in[0,t_f]} P\big(x(t),u(t);p\big), \label{eq:cd3}
\end{align}
where in this section $\eb$ denotes the usable specific energy before the
landing reserve $\varsigma$ is removed and $p_b$ the pack's specific power,
subject to the path constraints
\begin{equation}
  T(t)\le\lambda mg,\quad
  \Omega_i(t)\le\Omega_{\max},\quad
  \ell(t)\ge d_{\mathrm{safe}},\quad
  \theta_{\mathrm{read}}(t)\le\theta_{\max},
  \label{eq:cd4}
\end{equation}
with $\ell(t)$ the distance from the user's eye to the nearest rotor disk,
minimising a scalar objective
\begin{equation}
  \mathcal{J}(p,u) = w_m m + w_P \bar P + w_N L_p + w_s\,\mathrm{span}
                     + w_e\, e_{\mathrm{lag}} .
  \label{eq:cd5}
\end{equation}
\end{problem}

\Cref{eq:cd3} is the coupling. The right-hand side depends on the trajectory,
which depends on the plant, which depends on $m$, which contains $\mbat$, which
is the left-hand side. The first inequality keeps the reserve: starting from
$E(0)=\eb\mbat$, the state of charge at $t_f$ is at least $\varsigma$. For a
minimally sized continuous battery one of the two inequalities is an equality,
\begin{equation}
  \mbat = \max\left\{\frac{1}{\eb(1-\varsigma)}\int_0^{t_f}P\,\dd t,\;
                     \frac{1}{p_b}\max_t P\right\} ;
  \label{eq:batreq}
\end{equation}
a discrete pack, built from whole cells, need not satisfy either with equality.

\subsection{Reduction to a fixed point}

\Cref{eq:cd2,eq:cd3} can be reduced to a scalar fixed-point problem, which is
what makes the system tractable.

\begin{definition}[Inner and outer maps]
For a given battery mass $\mu$, let
\begin{equation}
  \mathcal{G}(\mu) \coloneqq \text{the solution } m \text{ of }
  m = \mfix + \mstr(m) + \mprop(m) + \mu
  \label{eq:innermap}
\end{equation}
be the gross mass consistent with that battery, and let
\begin{equation}
  \mathcal{H}(\mu) \coloneqq
  \max\left\{
  \frac{1}{\eb(1-\varsigma)}\int_0^{t_f}
    P\big(x(t;\mathcal{G}(\mu),\mu),u(t)\big)\dd t,\;
  \frac{1}{p_b}\max_t P\big(x(t;\mathcal{G}(\mu),\mu),u(t)\big)\right\}
  \label{eq:outermap}
\end{equation}
be the battery mass that the resulting mission demands, by \cref{eq:batreq}.
\end{definition}

\begin{proposition}[Well-posedness of the inner map]
\label{prop:inner}
Let $S(m)\coloneqq\mstr(m)+\mprop(m)$ map $[0,\infty)$ into $[0,\infty)$ (masses
are non-negative) and satisfy the Lipschitz bound
\begin{equation}
  |S(m_2)-S(m_1)| \le L\,|m_2-m_1|
  \quad\text{for all } m_1,m_2\ge0,
  \qquad L<1 .
  \label{eq:innercond}
\end{equation}
Then for every $\mu\ge0$, \cref{eq:innermap} has exactly one solution
$\mathcal{G}(\mu)$, the iteration $m\leftarrow\mfix+S(m)+\mu$ converges to it
from any start at rate $L$, and $\mathcal{G}$ is Lipschitz with constant
$1/(1-L)$. If $S$ is non-decreasing, so is $\mathcal{G}$. If $S$ is
differentiable,
\begin{equation}
  \mathcal{G}'(\mu) = \frac{1}{1-S'\big(\mathcal{G}(\mu)\big)} .
  \label{eq:Gprime}
\end{equation}
\end{proposition}

\begin{proof}
The map $\Phi(m)=\mfix+S(m)+\mu$ sends $[0,\infty)$ into itself because
$S\ge0$, and is a contraction with constant $L$ by \cref{eq:innercond}, so the
Banach fixed-point
theorem~\cite[Ch.~5]{ortega1970} gives existence, uniqueness and convergence of
the iteration. For $\mu_1,\mu_2$ with solutions $m_1,m_2$,
$m_2-m_1 = S(m_2)-S(m_1)+\mu_2-\mu_1$, so
$|m_2-m_1|\le L|m_2-m_1|+|\mu_2-\mu_1|$, which is the Lipschitz bound. If $S$ is
non-decreasing and $\mu_2>\mu_1$ but $m_2<m_1$, then
$m_2-m_1-\big(S(m_2)-S(m_1)\big)\le-(1-L)(m_1-m_2)<0$, contradicting
$m_2-m_1-\big(S(m_2)-S(m_1)\big)=\mu_2-\mu_1>0$. Finally
$\Xi(m,\mu)=\Phi(m)-m$ has $\partial_m\Xi = S'-1\le L-1<0$, so the implicit
function theorem applies and gives \cref{eq:Gprime}.
\end{proof}

\begin{remark}[Why a slope bound and not a growth bound]
Continuity, monotonicity and $\limsup S(m)/m<1$ give existence but not
uniqueness. The smooth, increasing, bounded
$S(m)=4/[1+e^{-4(m-3)}]$ with $\mfix+\mu=1$ makes $m=1+S(m)$ hold at
$m\approx1.0013$, $3$ and $4.9987$. The slope condition is what excludes this.
For the lumped models $S'=f_s+\gamma_m$. For the component models of
\cref{sec:components}, evaluated on the design of \cref{sec:results}, $S'$ lies
between $0.13$ and $0.23$ over \SIrange{0.8}{50}{\kilo\gram}; the test suite
asserts $0\le S'<0.9$ on that interval. That is a check on one design over a
bounded interval, not a proof for every parameter set.
\end{remark}

\begin{theorem}[Existence of a closed design]
\label{thm:codesignfixed}
The system \cref{eq:cd2,eq:cd3} is equivalent to the scalar fixed-point equation
\begin{equation}
  \mu = \mathcal{H}(\mu) .
  \label{eq:fixedpoint}
\end{equation}
If $\mathcal{H}$ is continuous and maps a compact interval $[0,\mu_{\max}]$ into
itself, a solution exists by Brouwer's theorem. If in addition $\mathcal{H}$ is
a contraction, $\mathrm{Lip}(\mathcal{H}) = L < 1$, the solution is unique and
the iteration $\mu_{k+1}=\mathcal{H}(\mu_k)$ converges geometrically with rate
$L$ from any starting point.
\end{theorem}

\begin{proof}
Equivalence is by construction. The existence and uniqueness statements are
Brouwer's fixed-point theorem and the Banach fixed-point theorem
respectively~\cite[Ch.~5]{ortega1970}.
\end{proof}

\begin{proposition}[Contraction criterion]
\label{prop:contraction}
Suppose the mission power is well approximated by a mission-dependent multiple
of the hover power, $\bar P = \varkappa\,P_{\mathrm{hover}}(m)$ with $\varkappa$
independent of $m$, and that the energy term attains the maximum in
\cref{eq:outermap} in a neighbourhood of $\mu$. Then
\begin{equation}
  \mathcal{H}'(\mu)
  = \frac{\varkappa\,t_f}{\eb(1-\varsigma)}\,
    \frac{\dd P_{\mathrm{hover}}}{\dd m}\,\mathcal{G}'(\mu),
  \label{eq:Hprime}
\end{equation}
and in the fixed-rotor case with the lumped models, where
$\mathcal{G}'=1/(1-f_s-\gamma_m)=1/\alpha$,
\begin{equation}
  \mathcal{H}'(\mu) = \frac{\varkappa}{\alpha}\cdot
    \frac{3}{2}\,\beta\,\sqrt{m} ,
  \label{eq:Hprime2}
\end{equation}
with $\beta$ as in \cref{eq:beta}, whose specific energy is net of the reserve,
$\eb(1-\varsigma)$ in the notation of this section. The iteration therefore
contracts precisely when $\tfrac32\varkappa\beta\sqrt m < \alpha$, which at
$\varkappa=1$ is the condition $F'(m)>0$ of \cref{cor:dmdm}, i.e.\ $m<m^*$. On
the power-limited branch $\mathcal{H}' = p_b^{-1}(\dd P_{\mathrm{peak}}/\dd m)
\mathcal{G}'$ instead, and where the two branches exchange the maximum
$\mathcal{H}$ has a kink.
\end{proposition}

\begin{proof}
Differentiate the energy term of \cref{eq:outermap} by the chain rule. From
\cref{eq:beta,eq:ptot},
$P_{\mathrm{hover}}(m) = \beta\,\eb(1-\varsigma)\, m^{3/2}/t_f + \Paux$, so
$\dd P_{\mathrm{hover}}/\dd m = \tfrac32\beta\,\eb(1-\varsigma)\, m^{1/2}/t_f$.
Substituting into \cref{eq:Hprime} with $\mathcal{G}'=1/\alpha$, the factors
$t_f$ and $\eb(1-\varsigma)$ cancel and \cref{eq:Hprime2} remains.
\end{proof}

\begin{remark}[What the overhead does, and does not, do]
\Cref{prop:contraction} connects the algebraic and the simulated pictures: the
condition for the simulation-in-the-loop iteration to converge is the
light-branch condition with $\beta$ multiplied by $\varkappa$. It is tempting to
say that overhead ``degrades $\alpha$''. That is a reading of one inequality,
not an identity. Multiplying the whole power by $\varkappa$ turns the closure
into
\begin{equation}
  \alpha m - \varkappa\beta\, m^{3/2}
  = \mfix + \varkappa\,\frac{\Paux t_f}{\eb},
  \label{eq:kappaclosure}
\end{equation}
which scales the energy coupling and the auxiliary part of $\Mzero$ together,
and by \cref{eq:closedform} lowers the critical load by the factor
$1/\varkappa^2$. A change of $\alpha$ would do neither.
\end{remark}

\subsection{Residual formulation}

Rather than iterate \cref{eq:fixedpoint} naively, the engine treats the pair
$y = (m,\mbat)$ and drives the residual
\begin{equation}
  \vect{\mathcal{R}}(y) =
  \begin{bmatrix}
    m - \big(\mfix + \mstr(m) + \mprop(m) + \mbat\big) \\[2pt]
    \mbat - \max\left\{\dfrac{E_{\mathrm{mission}}(m,\mbat)}{\eb(1-\varsigma)},\;
                       \dfrac{P_{\mathrm{peak}}(m,\mbat)}{p_b}\right\}
  \end{bmatrix}
  = \vect 0 .
  \label{eq:residual}
\end{equation}
The first component is algebraic and cheap; the second requires a trajectory
simulation. The engine solves $\mathcal{R}_1=0$ at every outer step by the
iteration of \cref{prop:inner}, reducing the problem to the scalar
$\mathcal{R}_2=0$ at one simulation per update. When $|\mathcal{R}_2|$ falls
below tolerance the battery is rounded up to the last requirement and the
design is flown again; the run is declared converged only when that
verification flight requires no more battery than it carries. The energy margin
$\eb(1-\varsigma)\mbat-E_{\mathrm{mission}}$, the power margin
$p_b\mbat-P_{\mathrm{peak}}$ and the path constraints \cref{eq:cd4} of the final
flight are then reported separately. Closing the mass and energy balance and
satisfying the path constraints are different outcomes: a converged battery
loop can still be an unacceptable flight.

\subsection{The outer optimisation}

With the closure solved, \cref{eq:cd5} becomes a function of the design vector
alone, and \cref{prob:codesign} reduces to a constrained finite-dimensional
program
\begin{equation}
  \min_{p\in\mathcal{P}} \ \mathcal{J}(p)
  \quad\text{s.t.}\quad
  g_j(p)\le0,\ j=1,\dots,J ,
  \label{eq:outerprog}
\end{equation}
with $\mathcal{P}$ a box and $g_j$ the constraints of \cref{eq:cd4} together
with the modelling-validity constraints of
\cref{as:sigma,as:stall,eq:reynolds}. The objective is non-convex, cheap to
evaluate on the algebraic layer and expensive on the simulated one, and has
large infeasible regions where the closure has no solution at all. This
combination dictates the choice of a population-based method with exterior
penalties, discussed in \cref{sec:numerics}.

\begin{remark}[Why not adjoint optimal control]
\Cref{prob:codesign} in full generality is a control co-design problem and
could be attacked with adjoint-based methods over the trajectory. We do not,
for two reasons. First, the objective is dominated by steady hover, for which
the trajectory contributes a scalar $\varkappa$ near unity; the sensitivity of
the design to the control law is small compared with its sensitivity to the
rotor geometry. Second, and more importantly, the constraints that actually
bind, acoustic and geometric, are algebraic functions of $p$ and do not involve
$u(\cdot)$ at all. Spending effort on trajectory optimisation would refine the
part of the problem that is not deciding the answer.
\end{remark}

\section{The engine}
\label{sec:engine}

\subsection{Design principles}

The implementation follows three rules, each of which was adopted because
violating it caused a specific error during development.

\paragraph{One source of truth for parameters.} Every assumption is declared
once, in a schema carrying its symbol, unit, admissible range, grouping and
help text. Validation, the command line interface and the browser controls are
all generated from that schema. Nothing about the model is written down twice.
The alternative, a model that reads parameters and a user interface that
mirrors them, guarantees eventual divergence between the two.

\paragraph{Analytic and numerical layers coexist.} The closed forms of
\cref{sec:fold,sec:scaling} are implemented alongside the numerical solvers
that supersede them, and the test suite asserts agreement. The closed forms are
not documentation; they are the oracle against which the code is verified. This
is what caught the errors reported in \cref{sec:verification}.

\paragraph{Assumptions are outputs wherever possible.} The lumped constants
$f_s$, $S_m$ and $\FM$ remain as inputs to the algebraic layer, but the
first-principles layer computes them and the engine reports both side by side.
A model whose assumptions can be audited against its own detailed submodels is
much harder to fool oneself with.

\subsection{Module structure}

\begin{table}[t]
\centering
\caption{Modules and the sections of this paper they implement.}
\label{tab:modules}
\small
\begin{tabular}{@{}llp{6.6cm}@{}}
\toprule
Module & Sections & Responsibility \\
\midrule
\code{params.py}      & \ref{sec:notation}      & parameter schema, derived quantities, named presets \\
\code{sizing.py}      & \ref{sec:momentum} to \ref{sec:scaling} & momentum theory, closure, fold, sensitivities, fixed-point trace \\
\code{components.py}  & \ref{sec:blade}, \ref{sec:components}, \ref{sec:nearfield} & rotor blade model, beam-sized arms, motor and propeller mass, inertia, near-field corrections \\
\code{acoustics.py}   & \ref{sec:acoustics}     & source model, sound power, room field, A-weighting, blade passage \\
\code{electrical.py}  & \ref{sec:components}    & pack selection, $K_V$ sizing, currents, C-rate checks \\
\code{dynamics.py}    & \ref{sec:dynamics}      & state, forces, moments, allocation, hover linearisation \\
\code{control.py}     & \ref{sec:control}       & cascaded controller, authority, gain selection \\
\code{mission.py}     & \ref{sec:dynamics}, \ref{sec:control} & human trajectories, wind, reference governor with keep-out \\
\code{simulate.py}    & \ref{sec:dynamics}      & fixed-step integrator, telemetry, mission energy \\
\code{closure.py}     & \ref{sec:codesign}      & residual formulation, secant-accelerated battery loop, verification flight, margins \\
\code{numerics.py}    & \ref{sec:numerics}      & bounded mass-root search with explicit failure status, controllability rank \\
\code{optimize.py}    & \ref{sec:codesign}      & penalised objective, differential evolution, trade studies \\
\code{tether.py}      & \ref{sec:results}       & the closure with energy off-board \\
\code{sweep.py}       & \ref{sec:scaling}       & feasibility maps, exponent sweeps, endurance walls \\
\code{payloads.py}    & \ref{sec:results}       & display catalogue, readability, bill of materials \\
\bottomrule
\end{tabular}
\end{table}

\Cref{fig:modules} shows how the modules depend on one another, and
\cref{fig:pipeline} traces one evaluation of the coupled problem through them:
the order of operations that the rest of this paper derives piece by piece.

\begin{figure}[tbp]
\centering
% Import graph of the engine. An arrow A -> B means that module A imports B.
% Every module imports params.py, drawn as the base; the three interfaces at
% the top use most of the engine and their individual edges are omitted.
\begin{tikzpicture}[
    font=\footnotesize, y=1.05cm,
    mod/.style={draw, rounded corners=2pt, fill=cA!6, minimum width=1.9cm,
                minimum height=1.5em, font=\ttfamily\footnotesize},
    ifc/.style={mod, fill=cB!12},
    dep/.style={-{Stealth[length=1.8mm]}, cG, line width=0.55pt},
  ]
  % base
  \node[draw, fill=cG!12, minimum width=13.8cm, minimum height=1.5em,
        font=\ttfamily\footnotesize] (params) at (6.3,0)
        {params.py: schema, presets, derived quantities};
  % layer 1
  \node[mod] (numerics)  at (1.0,1)  {numerics};
  \node[mod] (acoustics) at (6.3,1)  {acoustics};
  \node[mod] (payloads)  at (8.6,1)  {payloads};
  \node[mod] (sizing)    at (11.6,1) {sizing};
  % layer 2
  \node[mod] (components) at (3.5,2) {components};
  % layer 3
  \node[mod] (electrical) at (0.8,3)  {electrical};
  \node[mod] (dynamics)   at (3.5,3)  {dynamics};
  \node[mod] (mission)    at (6.3,3)  {mission};
  \node[mod] (sweep)      at (10.0,3) {sweep};
  % layer 4
  \node[mod] (tether)     at (0.8,4) {tether};
  \node[mod] (control)    at (3.5,4) {control};
  % layer 5
  \node[mod] (simulate)   at (5.2,5) {simulate};
  % layer 6
  \node[mod] (closure)    at (2.8,6) {closure};
  \node[mod] (optimize)   at (7.6,6) {optimize};
  % interfaces
  \node[ifc] (api) at (2.6,7.7) {api};
  \node[ifc] (cli) at (6.3,7.7) {cli};
  \node[ifc, minimum width=2.6cm] (mk) at (10.4,7.7) {make\_results};
  \node (tool)  at (2.6,8.8)  {browser tool, \code{web/}};
  \node (term)  at (6.3,8.8)  {command line};
  \node (paper) at (10.4,8.8) {this paper};
  \draw[flow, cB] (api) -- (tool);
  \draw[flow, cB] (cli) -- (term);
  \draw[flow, cB] (mk) -- (paper);

  % engine edges: importer -> imported
  \draw[dep] (components.south west) -- (numerics.north east);
  \draw[dep] (dynamics.south west) -- (numerics.north);
  \draw[dep] (dynamics) -- (components);
  \draw[dep] (electrical.south east) -- (components.north west);
  \draw[dep] (mission.south west) -- (components.north east);
  \draw[dep] (sweep.south west) -- ($(components.east)+(0,0.12)$);
  \draw[dep] (sweep.south east) -- (sizing.north west);
  \draw[dep] (tether) -- (electrical);
  \draw[dep] (tether.south east) -- ($(components.north west)+(0.25,0)$);
  \draw[dep] (tether.west) -- ++(-0.55,0) |- (-0.35,0.52) -| (sizing.south);
  \draw[dep] (control) -- (dynamics);
  \draw[dep] (simulate.south west) -- (control.east);
  \draw[dep] (simulate.south) -- (dynamics.north east);
  \draw[dep] (simulate.south east) -- (mission.north);
  \draw[dep] (closure) -- (simulate);
  \draw[dep] ($(closure.south)+(-0.6,0)$) |- (components.west);
  \draw[dep] (closure.north) -- ++(0,0.25) -| ($(sizing.north)+(0.6,0)$);
  \draw[dep] (optimize) -- (simulate);
  \draw[dep] (optimize.south) |- ($(components.east)+(0,-0.12)$);

  % the interfaces use the whole engine
  \draw[decorate, decoration={brace, amplitude=5pt, mirror}, cB]
       ($(api.south west)+(-0.2,-0.12)$) -- ($(mk.south east)+(0.2,-0.12)$)
       node[midway, below=5pt, font=\scriptsize, fill=white]
       {the interfaces call across the whole engine};
\end{tikzpicture}

\caption{Import graph of the engine: an arrow from A to B means that module A
imports B. Every module also imports \code{params.py}, drawn as the base. The
layers read from the bottom up as the layers of the paper: the algebraic and
component models, the dynamics and control, the simulation, and the two
modules that close the coupled problem, \code{closure.py} and
\code{optimize.py}. \code{acoustics.py} and \code{payloads.py} import nothing
but the parameters. The three interfaces at the top call across the whole
engine and their individual edges are omitted.}
\label{fig:modules}
\end{figure}

\begin{figure}[tbp]
\centering
% How the coupled problem is solved: a bounded root search sizes the hover
% design, the battery loop wraps an inner algebraic fixed point and one
% simulation per update, a verification flight closes it, and the results
% flow into this paper.
\begin{tikzpicture}[
    font=\footnotesize,
    box/.style={block, text width=2.8cm, minimum height=3.3em, inner sep=3pt},
    stp/.style={box, fill=cA!7},
    sim/.style={box, fill=cB!12},
    res/.style={box, fill=cC!10},
    lab/.style={font=\scriptsize},
  ]
  % top: hover sizing
  \node[box] (p)    at (0,0)    {parameter set $p$,\\ \code{params.py}};
  \node[stp] (root) at (3.9,0)  {component closure\\ $\mathcal{F}(m)=0$ by a\\ bounded root search};
  \node[box] (hov)  at (7.8,0)  {hover design\\ $m,\ \mbat^{(0)}$};
  \draw[flow] (p) -- (root);
  \draw[flow] (root) -- (hov);

  % middle: the battery loop
  \node[stp] (inner) at (0,-2.6)   {inner map $\mathcal{G}(\mu)$:\\
       $m=\mfix+S(m)+\mu$\\ by contraction};
  \node[stp] (veh)   at (3.9,-2.6) {vehicle: $J$, $\mat M$,\\ $k_T,k_Q,k_P$, trim};
  \node[sim] (fly)   at (7.8,-2.6) {simulate a window:\\ RK4, $h=\SI{4}{\milli\second}$,\\ controller, governor};
  \node[stp] (req)   at (7.8,-5.1) {requirement\\ $\max\{E/\eb(1{-}\varsigma),$\\ $P_{\mathrm{peak}}/p_b\}$};
  \node[stp] (sec)   at (3.9,-5.1) {secant step\\ on $\mathcal{R}_2$, guarded,\\ \labelcref{eq:secant}};
  \node[box] (test)  at (0,-5.1)   {$|\mathcal{R}_2|<\mathrm{tol}$\,?};
  \draw[flow] (hov.south) -- ++(0,-0.55)
       -| node[lab, pos=0.3, below] {$\mu_0=\mbat^{(0)}$} (inner.north);
  \draw[flow] (inner) -- node[lab, above] {$m$} (veh);
  \draw[flow] (veh) -- (fly);
  \draw[flow] (fly) -- node[lab, right] {$\bar P,\ P_{\mathrm{peak}}$} (req);
  \draw[flow] (req) -- node[lab, above] {$\mathcal{H}(\mu)$} (sec);
  \draw[flow] (sec) -- node[lab, above] {$\mathcal{R}_2$} (test);
  \draw[flow] (test.north) -- node[lab, left] {no} (inner.south);

  \begin{scope}[on background layer]
    \node[draw=cA, dashed, rounded corners=3pt, inner sep=7pt,
          fit=(inner)(veh)(fly)(req)(sec)(test)] (loop) {};
  \end{scope}
  \node[lab, cA, anchor=north east] at (loop.south east)
       {battery loop, \code{closure.py}: one simulation per update};

  % bottom: verification and output
  \node[sim] (ver) at (0,-7.9)    {verification flight\\ at the rounded battery};
  \node[res] (con) at (3.9,-7.9)  {margins and path\\ constraints \labelcref{eq:cd4}};
  \node[res] (res) at (7.8,-7.9)  {\code{make\_results.py}:\\ results, macros,\\ tables, figure data};
  \node[res] (pap) at (11.5,-7.9) {this paper};
  \node[res] (api) at (11.5,-2.6) {API and\\ browser tool};
  \draw[flow] (test.south) -- node[lab, left] {yes} (ver.north);
  \draw[flow] (ver) -- (con);
  \draw[flow] (con) -- (res);
  \draw[flow] (res) -- (pap);
  \draw[flow, cB] (p.north) -- ++(0,0.4) -| (api.north);
  \draw[flow, cB, dashed] (ver.west) -- ++(-0.35,0) |- ($(ver.south)+(0,-0.45)$)
       node[lab, pos=0.75, below] {round up, fly again} -- (ver.south);
\end{tikzpicture}

\caption{One evaluation of the coupled problem as the engine carries it out.
A bounded root search on the component residual \cref{eq:fpclosure} gives the
hover design, which seeds the battery loop. Each update of the loop solves the
inner mass map by contraction (\cref{prop:inner}), rebuilds the vehicle, flies
one simulation window and turns its mean and peak power into a battery
requirement by \cref{eq:batreq}; a guarded secant step then moves the battery.
When the residual is within tolerance the battery is rounded up and flown
again, and the loop ends only on a verification flight that needs no more
battery than it carries. The results pass through \code{make\_results.py} into
this paper, and the same engine serves the API and the browser tool.}
\label{fig:pipeline}
\end{figure}

\subsection{Interfaces}

Three front ends share the engine: a command line for scripted studies, an HTTP
API, and a browser tool that generates its controls from the schema and
recomputes on every parameter change. The tool exists because the questions
that matter here, ``what does this assumption cost'', are answered by moving one
number and watching six others, which is a different activity from running a
batch job.

The browser tool renders the assembly in three dimensions from the same solved
dimensions used by the closure: arm length and wall thickness from
\cref{eq:armlength,eq:wall}, rotor diameter and chord from
\cref{sec:blade}, motor body size from \cref{eq:motormass}, and the individual
battery cells from the pack selection of \cref{sec:components}. The check that
this is a drawing of the machine rather than a picture beside it is that the
drawn parts sum to the closure's gross mass.

\subsection{Reproducing the results}

Every number in \cref{sec:results} is produced by one script,
\code{paper/make\_results.py}, which writes the numbers together with the
configuration that produced them (preset, payload, window, step, settling
interval, seed, governor limits, solver tolerances and optimiser settings) to
\code{paper/results/results.json}, and emits the table rows and macros that the
\LaTeX\ sources include. No figure in the design study is typed by hand.
\code{paper/build.py} then compiles this document.

\section{Numerical methods}
\label{sec:numerics}

\subsection{Solving the closure}

\Cref{eq:closure-q} could be solved by radicals at $q=3/2$ or by Newton's
method in general. Neither is used, for reasons that are structural rather than
matters of taste.

Newton's method applied to $F(m)-\Mzero$ has a critical defect here: by
\cref{lem:concave}, $F'(m^*)=0$, so the Newton step is unbounded near the fold,
which is precisely the region of interest. Worse, from a starting point on the
wrong side it converges to $m_+$, the unstable branch, without any indication
that it has done so.

The engine instead exploits the shape established in \cref{lem:concave}:

\begin{enumerate}[nosep]
  \item Compute $m^*$ and $\Mmax$ in closed form from
        \cref{eq:mstar,eq:Mmax}. Feasibility is decided \emph{analytically} by
        comparing $\Mzero$ with $\Mmax$, never by whether an iteration
        converged.
  \item If feasible, bracket $m_-$ on $(0,m^*]$, where $F$ is a strictly
        increasing bijection, and bisect. Bracket $m_+$ on $[m^*,\infty)$ by
        geometric expansion and bisect.
\end{enumerate}

Bisection converges linearly with rate $1/2$ and requires
$\lceil\log_2((b-a)/\epsilon)\rceil$ evaluations, about $40$ for double
precision. It is slower than Newton and it cannot fail. Given that a single
evaluation of $F$ costs microseconds, and given that the failure modes of
Newton here are silent rather than loud, this is the correct trade.

\begin{remark}[Feasibility is analytic, not numerical]
This is the single most important implementation decision in the sizing layer,
and it follows from the critical slowing-down remark after
\cref{thm:iteration}. A solver that reports non-convergence after $K$ iterations
cannot distinguish a near-feasible design from an infeasible one, because
$\mathrm{rate} = 1-F'(m_-)\to1$ as $\Mzero\uparrow\Mmax$. Deciding feasibility
by $\Mzero\lessgtr\Mmax$ removes the ambiguity entirely.
\end{remark}

For the component closure \cref{eq:fpclosure} no closed form exists, and the
shape of $\mathcal{F}$ is observed rather than proved (\cref{sec:components}),
so the search cannot borrow the analytic feasibility test. The engine samples
$\mathcal{F}$ on a geometric grid of ratio $1.35$ from \SI{e-3}{\kilo\gram}
upward, stopping at the first non-negative sample. Every sampled local maximum
is then refined by bounded minimisation of $-\mathcal{F}$ in $\ln m$, because a
positive stretch of $\mathcal{F}$ narrower than one grid interval can sit
between two negative samples. The first bracket with a sign change, in
increasing mass, is closed by Brent's method to a relative tolerance of
$10^{-12}$, which selects the light branch. The search reports one of four
outcomes: a root, a near-tangent root resolved only to residual tolerance, a
non-finite residual, or an exhausted search. The last is reported as a failure
of a bounded search and is never described as infeasibility. The case that
motivated the refinement is kept as a regression test: the design of
\cref{sec:results} at $t_f=\SI{4280.0502}{\second}$ has a positive island of
width \SI{5}{\percent} in mass, with roots at \SI{4.5678}{\kilo\gram} and
\SI{4.6173}{\kilo\gram}, which the grid alone straddled and reported as no
mass closing.

\subsection{The battery loop with the simulator inside it}

\Cref{eq:fixedpoint} is solved by iteration, but naive iteration converges only
at the linear rate $\mathcal{H}'(\mu)$ of \cref{prop:contraction}, which near
the fold approaches unity. Since each evaluation of $\mathcal{H}$ costs a full
trajectory simulation, this is expensive.

The engine applies a secant step to the residual
$\mathrm{gap}(\mu) = \mathcal{H}(\mu)-\mu$:
\begin{equation}
  \mu_{k+1} = \mu_k - \mathrm{gap}(\mu_k)\,
    \frac{\mu_k-\mu_{k-1}}{\mathrm{gap}(\mu_k)-\mathrm{gap}(\mu_{k-1})} ,
  \label{eq:secant}
\end{equation}
which uses only quantities already computed and therefore costs no additional
simulations.

\begin{proposition}[Convergence order]
\label{prop:secant}
If $\mathrm{gap}\in C^2$ near a simple root $\mu^\star$ with
$\mathrm{gap}'(\mu^\star)\ne0$, the secant iteration \cref{eq:secant} converges
locally with order $\varphi=(1+\sqrt5)/2\approx1.618$, against order $1$ for
the fixed-point iteration.
\end{proposition}

\begin{proof}
Classical; see Ortega and Rheinboldt~\cite[\S8.2]{ortega1970}. The error
satisfies $|\epsilon_{k+1}|\approx C|\epsilon_k||\epsilon_{k-1}|$, whose
characteristic exponent is the golden ratio.
\end{proof}

The step is guarded: it is accepted only when it moves in the same direction as
the gap and by no more than six times its magnitude, otherwise the plain
fixed-point step is taken. Because $\mathcal{H}$ takes the larger of an
energy and a power requirement (\cref{eq:outermap}) it has a kink where the
two exchange, and \cref{prop:secant} does not apply across it; the guard is
what keeps the iteration sane there. When the gap is within tolerance the
battery is rounded up to the last requirement and the design is flown once
more, and the loop reports convergence only if that flight requires no more
battery than it carries. On every mission of \cref{sec:results} this took
three updates and two verification flights, five simulations in all, against
eight or more for plain iteration.

\subsection{Trajectory integration}

\Cref{eq:ode} is integrated with classical fourth-order Runge--Kutta at fixed
step. Fixed step rather than adaptive, because the cost must be predictable for
an interactive tool and because the controller is evaluated once per step,
making the scheme effectively a sampled-data system whose sample rate should
not wander.

Local truncation error is $O(h^5)$ and global error $O(h^4)$. The step is
chosen from the fastest closed-loop mode. With $\tau_m=\SI{45}{\milli\second}$,
$K_R\approx\SI{41}{\per\second\squared}$ and hence
$\omega_n=\sqrt{K_R}\approx\SI{6.4}{\radian\per\second}$, the fastest pole is
the motor at $1/\tau_m\approx\SI{22}{\radian\per\second}$; at
$h=\SI{4}{\milli\second}$ the product $h\lambda\approx0.09$ sits deep inside the
RK4 stability region, whose real-axis extent is $|h\lambda|\le2.79$.

The quaternion is renormalised after each step, $q\leftarrow q/\|q\|$. This
does not cost the order. The exact flow preserves $\|q\|=1$, so the RK4 step
returns $\tilde q$ with $\|\tilde q\| = 1+O(h^5)$; the map
$\tilde q\mapsto\tilde q/\|\tilde q\|$ is smooth near $\mathbb{S}^3$ with
$\tilde q/\|\tilde q\| = \tilde q + O(\|\tilde q\|-1)$, so the projection
perturbs each step by $O(h^5)$, the same order as the local truncation error,
and the global order $O(h^4)$ is retained.

Two different things are being integrated and the distinction matters. The
controller is evaluated once per step and its output held, so RK4 integrates
the plant under a piecewise-constant input at fourth order. The closed loop is
therefore a sampled-data system at \SI{250}{\hertz}, which is what a digital
autopilot is; the $O(h)$ difference between it and a continuous-feedback
idealisation is a property of the system being simulated, not an integration
error. A convergence check with the controller rate held fixed and the
integration substeps refined was not run and is listed in
\cref{sec:limitations}.

\begin{remark}[Windowed mission energy]
Integrating a full \SI{1800}{\second} mission at $h=\SI{4}{\milli\second}$ would
require $4.5\times10^5$ steps per evaluation, which is incompatible with the
outer loop. The engine instead integrates a representative window of
\SIrange{12}{20}{\second}, measures the mean electrical power after a settling
interval, and extrapolates:
$E_{\mathrm{mission}} = \bar P\,t_f$. This is exact for a periodic mission whose
window spans a whole number of periods and is otherwise an approximation whose
error is the difference between the window mean and the mission mean. It is
recorded as a limitation.
\end{remark}

\begin{remark}[Initial conditions]
The initial state is taken to be the reference position \emph{and the reference
velocity}. Starting at rest beside a person already walking is not a gentler
initial condition but a worse one: on a curved path the display sits in the
person's path, so they close on it while it accelerates. This artefact
contaminated every mission statistic until it was found; see
\cref{sec:verification}.
\end{remark}

\subsection{Linearisation}

$\mat A$ and $\mat B$ are formed by central differences,
\begin{equation}
  \mat A_{:,j} \approx \frac{f(x_0+h\vect e_j)-f(x_0-h\vect e_j)}{2h},
  \qquad h = \varepsilon\max\{1,|x_{0,j}|\},
  \label{eq:fd}
\end{equation}
with $\varepsilon=10^{-6}$. Central differences have truncation error $O(h^2)$
against $O(h)$ for one-sided, and the total error is minimised near
$h\sim\epsilon_{\mathrm{mach}}^{1/3}\approx6\times10^{-6}$, so the choice is
near-optimal. Controllability is assessed by the reachable-subspace iteration
of \cref{sec:control} with a singular-value tolerance of $10^{-9}$ on the
norm-scaled matrices, and requires full rank.

\subsection{The outer optimisation}

\Cref{eq:outerprog} is solved by differential evolution~\cite{storn1997}, a
population-based stochastic method, with constraints handled by an exterior
quadratic penalty
\begin{equation}
  \tilde{\mathcal{J}}(p) = \mathcal{J}(p)
    + \varpi\sum_j \left(\frac{\max\{0,g_j(p)\}}{s_j}\right)^{2} ,
  \label{eq:penalty}
\end{equation}
with per-constraint scales $s_j$ chosen so that a violation of engineering
significance contributes order unity.

Three properties of the problem dictate this choice. The objective is
non-convex and its landscape contains large regions where the closure has no
solution at all, so gradient methods have nowhere to start. Several design
variables are integer (rotor count, blade count), so smoothness fails. And an
evaluation of the algebraic objective costs milliseconds, so a population of a
few hundred is affordable.

The penalty is exterior rather than a barrier precisely because the infeasible
region must remain traversable: a barrier method cannot cross from one feasible
island to another through infeasible territory, and here the islands are
separated by the fold.

\subsection{Numerical hygiene}

Three details are worth recording because each caused a visible defect.

\emph{Non-finite values in transport.} \texttt{NaN} and $\pm\infty$ are not
representable in JSON. Infeasible cells in a sweep, and quantities such as
$\Mmax$ when no fold exists, are genuinely undefined rather than zero. The API
maps every non-finite float to \texttt{null} and the front end renders it as a
dash, so an undefined quantity is never silently displayed as $0$.

\emph{Solidity and Reynolds guards.} Without \cref{as:sigma} and
\cref{eq:recorrection} the optimiser drives solidity to $0.55$ and tip speed
toward zero, exploiting the fact that \cref{prop:profile} has no notion of
blade interference and \cref{prop:fmblade} none of Reynolds number. Both
constraints exist to keep the optimiser inside the domain where the model is a
model of something.

\emph{Boundary comparisons.} Constraints active exactly at their bound were
being reported as violations because of floating-point comparison. All
constraint checks carry a relative tolerance of $2\times10^{-3}$, chosen so
that a design sitting on a bound reads as satisfying it.

\section{Verification, calibration and errors found}
\label{sec:verification}

A model of this kind is only as good as the evidence that it computes what it
claims. This section records the verification strategy, the calibration against
measured hardware, and, because they are instructive, the errors that the
strategy caught.

\subsection{Verification against the closed forms}

The closed forms of \cref{sec:fold,sec:scaling} exist in the implementation
alongside the numerical solvers and serve as an oracle. The test suite
comprises 66 assertions; those bearing on the theory are:

\begin{itemize}[nosep]
  \item \emph{Momentum theory.} $T = 2\rho A\vind^2$ holds identically to
        machine precision; with $\Cdz\to0$, $\kappa\to1$ and the near-field
        corrections disabled, the computed $\FM\to1$ to within $10^{-5}$,
        confirming \cref{thm:momentum,eq:fm}.
  \item \emph{Coefficients.} $\alpha$ and $\beta$ from \cref{def:coeffs} agree
        with hand computation to four figures.
  \item \emph{Fold.} $m^*$ and $\Mmax$ agree with \cref{eq:closedform} to
        $10^{-12}$ relative; both roots satisfy \cref{eq:closure} to $10^{-8}$;
        $m_-<m^*<m_+$ as \cref{thm:trichotomy} requires; at
        $\Mzero=\Mmax-10^{-9}$ the two roots agree to $10^{-3}$ relative,
        exhibiting the square-root merging of \cref{eq:sqrtlaw}; and for
        $\Mzero>\Mmax$ no solution is returned.
  \item \emph{Iteration.} The naive iteration converges to the analytic $m_-$
        to $10^{-4}$ relative from a distant start, and is classified as
        diverging when $\Mzero>\Mmax$, confirming \cref{thm:iteration}.
  \item \emph{Scaling.} $\Mmax(2t_f) = \Mmax(t_f)/4$ to $10^{-9}$, and the
        numerically computed logarithmic sensitivities reproduce
        \cref{thm:sensitivity} to three decimal places, including the
        non-obvious $S_g=-3.770$ which depends on $\gamma_m/\alpha$.
  \item \emph{Exponent.} The endurance wall is finite for $p=0$ and absent for
        $p=1.4$, confirming \cref{thm:pfold}.
  \item \emph{Dynamics.} With zero thrust the vertical acceleration is exactly
        $-g$; at the hover trim the allocation returns $T=mg$ and
        $\vect\tau=\vect0$ to $10^{-9}$; $\dt E$ equals the negative of the
        computed electrical power identically; rotations satisfy
        $RR\transp=I$ and $\det R=1$ to $10^{-10}$ through the quaternion
        round trip; and the hover linearisation is controllable at full rank.
  \item \emph{Review regressions.} The near-fold root pair of
        \cref{sec:numerics} is found; the quadrotor authority equals the
        balanced differential \cref{eq:authquad} on every axis; the slab
        inertia is assigned to the panel axes; a forward boom produces a
        product of inertia and a trimmed hover with unequal rotor speeds;
        telemetry at $t$ compares the state at $t$ with the references at $t$;
        the governor keeps a reference out of a head that walks toward it;
        the battery loop's reported battery survives its verification flight;
        and the inner-map slope of \cref{prop:inner} lies in $[0,0.9)$ over
        \SIrange{0.8}{50}{\kilo\gram} for the shipped component models.
\end{itemize}

\subsection{Calibration of the acoustic anchor, and a fourteen-decibel error}

\Cref{mod:spl} contains an anchor $C$ which was initially assigned by
judgement, at $C=68$. Calibration against published declarations exposed this
as badly wrong.

Manufacturers declare a sound \emph{power} level $L_{WA}$ under EU regulation
2019/945. Converting to pressure at \SI{1}{\metre} by \cref{eq:eightdb} and
inverting \cref{eq:splmodel} gives $C$ for each aircraft.

\begin{table}[t]
\centering
\caption{Calibration of the acoustic anchor. Hover rotor speeds are
manufacturer figures; tip speed and thrust per rotor follow from geometry and
mass. The four aircraft span a factor of $3.6$ in mass and agree on $C$ to
\SI{1.5}{\decibel}.}
\label{tab:acoustic}
\begin{tabular}{@{}lrrrrrr@{}}
\toprule
Aircraft & $m$ & $T/\text{rotor}$ & $\Vtip$ & $L_{WA}$ & $L_p(\SI{1}{\metre})$ & $C$ \\
\midrule
DJI Mini 3      & \SI{0.249}{\kilo\gram} & \SI{0.61}{\newton} & \SI{71.6}{\metre\per\second} & 79.0 & 71.0 & 82.8 \\
DJI Air 2S      & \SI{0.595}{\kilo\gram} & \SI{1.46}{\newton} & \SI{67.1}{\metre\per\second} & 82.0 & 74.0 & 83.8 \\
DJI Mavic 3     & \SI{0.895}{\kilo\gram} & \SI{2.19}{\newton} & \SI{68.8}{\metre\per\second} & 83.0 & 75.0 & 82.3 \\
DJI Mavic 2 Pro & \SI{0.907}{\kilo\gram} & \SI{2.22}{\newton} & \SI{67.9}{\metre\per\second} & 83.0 & 75.0 & 82.6 \\
\midrule
\multicolumn{6}{r}{mean} & \textbf{82.9} \\
\bottomrule
\end{tabular}
\end{table}

\Cref{tab:acoustic} gives $C = 82.9$ with a spread of \SI{1.5}{\decibel}. The
consistency across a $3.6\times$ mass range is strong evidence that the
functional form \cref{eq:splmodel} is right, and that the original anchor was
\SI{14}{\decibel} optimistic. \Cref{fig:acoustic} shows the four values against
the one first assigned.

\begin{figure}[tbp]
\centering
% Calibration of the acoustic anchor C: each aircraft's declared sound power
% inverted through eq. (splmodel), against the value first assigned.
\begin{tikzpicture}
  \begin{axis}[paperplot, height=5.4cm,
      xmin=0.15, xmax=1.0, ymin=64, ymax=88,
      xlabel={aircraft mass [\si{\kilo\gram}]},
      ylabel={anchor $C$ [\si{\decibel}]},
      legend style={at={(0.98,0.36)}, anchor=east}]
    \fill[cA!10] (axis cs:0.15,82.3) rectangle (axis cs:1.0,83.8);
    \addplot[only marks, mark=*, mark size=2.6pt, cA] coordinates
      {(0.249,82.8) (0.595,83.8) (0.895,82.3) (0.907,82.6)};
    \addlegendentry{calibrated from $L_{WA}$}
    \addplot[cA, dashed, domain=0.15:1.0] {82.9};
    \addlegendentry{mean, $C=82.9$}
    \addplot[cD, line width=1.1pt, domain=0.15:1.0] {68};
    \addlegendentry{initial judgement, $C=68$}
    \node[lbl, anchor=south] at (axis cs:0.249,83.0) {Mini 3};
    \node[lbl, anchor=south] at (axis cs:0.595,84.0) {Air 2S};
    \node[lbl, anchor=north east] at (axis cs:0.88,82.1) {Mavic 3};
    \node[lbl, anchor=south west] at (axis cs:0.91,82.8) {Mavic 2};
    \draw[dim, cD] (axis cs:0.42,68) -- (axis cs:0.42,82.9)
      node[midway, right, font=\footnotesize, cD, align=left]
      {the anchor was\\ this far optimistic};
  \end{axis}
\end{tikzpicture}

\caption{The acoustic anchor $C$ recovered from the declared sound power of
four aircraft through \cref{eq:eightdb,eq:splmodel}, against aircraft mass.
The four values lie in the narrow band shaded across a factor of $3.6$ in
mass, about the mean $C=82.9$. The value first assigned by judgement, $C=68$,
lies below every one of them by the length of the arrow, an error larger than
the effect of any design choice in \cref{sec:results}.}
\label{fig:acoustic}
\end{figure}

\begin{finding}
Fourteen decibels is not a refinement. It is the difference between a machine
one can think beside and one that must leave the room, and it exceeds the
entire range spanned by every design decision available in the study. The error
had two components: an uncalibrated anchor, and the eight-decibel confusion of
\cref{prop:lwlp} between a declared sound power and a pressure at one metre. It
was found only because the model was made to reproduce measured hardware.
\end{finding}

The test suite now asserts that \cref{eq:splmodel} reproduces the declared
levels of three of these aircraft to within \SI{1.5}{\decibel}.

\subsection{Calibration of the motor model}

$k_\tau$ in \cref{eq:motormass} was likewise assigned by judgement at
\SI{3.0}{\newton\metre\per\kilo\gram} and found to be pessimistic by roughly a
factor of two.

\begin{table}[t]
\centering
\caption{Motor mass model against catalogue peak torque. The model is
conservative at the large end, which is the safe direction for a design study;
the designs here use motors of \SIrange{50}{60}{\gram} where the error is
small.}
\label{tab:motor}
\begin{tabular}{@{}lrrrr@{}}
\toprule
Motor & Actual mass & Peak torque & Model mass & Error \\
\midrule
T-Motor MN2806  & \SI{66}{\gram}  & \SI{0.30}{\newton\metre} & \SI{63}{\gram}  & $-5\%$ \\
T-Motor MN3110  & \SI{96}{\gram}  & \SI{0.63}{\newton\metre} & \SI{111}{\gram} & $+16\%$ \\
T-Motor MN505-S & \SI{190}{\gram} & \SI{2.03}{\newton\metre} & \SI{273}{\gram} & $+44\%$ \\
\bottomrule
\end{tabular}
\end{table}

With $k_\tau = \SI{5.5}{\newton\metre\per\kilo\gram}$ at
$m_{\mathrm{ref}}=\SI{0.1}{\kilo\gram}$ and $e=0.30$, \cref{tab:motor} results.
The recalibration mattered structurally, not just numerically: with the
pessimistic constant the motor mass penalty overwhelmed the induced-power
benefit of a larger rotor, and the optimum diameter sat on the lower bound of
the search box. After recalibration a genuine interior optimum appears, which
is the physically expected behaviour, and the test suite now asserts its
existence. \Cref{fig:motorcal} shows the calibrated law against the catalogue.

\begin{figure}[tbp]
\centering
% The motor mass model m = (Q m_ref^e / k_tau)^{1/(1+e)} with k_tau = 5.5
% N m/kg, m_ref = 0.1 kg, e = 0.30, against catalogue motors; the connectors
% are the model errors of tab:motor.
\providecommand{\fdQuadRoll}{3.66}
\providecommand{\fdQuadPitch}{3.66}
\providecommand{\fdQuadOld}{0.91}
\providecommand{\fdHexRoll}{6.13}
\providecommand{\fdHexPitch}{7.07}
\providecommand{\fdHexGross}{2.565}
\providecommand{\fdQuadGross}{1.724}
\providecommand{\fdKeepRadius}{1.180}
\providecommand{\fdKeepSlice}{1.174}
\providecommand{\fdEnvRadius}{1.080}
\providecommand{\fdTurnT}{5.65}
\providecommand{\fdTurnHeadX}{-1.440}
\providecommand{\fdTurnHeadY}{0.000}
\providecommand{\fdTurnPeriod}{7.54}
\providecommand{\fdStandoff}{1.20}
\providecommand{\fdRootLo}{4.5678}
\providecommand{\fdRootHi}{4.6173}
\providecommand{\fdPeakM}{4.5925}
\providecommand{\fdPeakF}{0.03}
\providecommand{\fdGridLo}{4.460}
\providecommand{\fdGridHi}{6.021}
\providecommand{\fdGridLoF}{-0.001}
\providecommand{\fdGridHiF}{-0.083}
\providecommand{\fdIslandPct}{1.1}
\providecommand{\fdPlainRate}{0.263}
\providecommand{\fdConvWindow}{8}
\providecommand{\fdHlMotorRe}{-0.0889}
\providecommand{\fdHlMotorIm}{0.0000}
\providecommand{\fdHlAttRe}{-0.0240}
\providecommand{\fdHlAttIm}{0.0000}
\providecommand{\fdHlPosRe}{-0.0082}
\providecommand{\fdHlPosIm}{0.0004}
\providecommand{\fdHlGovRe}{-0.0160}
\providecommand{\fdHlGovIm}{0.0000}
\providecommand{\fdWpos}{2.04}
\providecommand{\fdWgov}{4.00}
\providecommand{\fdWatt}{6.43}
\providecommand{\fdWmot}{22.2}
\providecommand{\fdWarm}{206}
\providecommand{\fdWnyq}{785}
\providecommand{\fdRatioPA}{3.2}
\providecommand{\fdRatioAM}{3.5}
\providecommand{\fdKR}{41.4}
\providecommand{\fdKp}{4.17}
\providecommand{\fdTauMms}{45}
\providecommand{\fdH}{4}
\providecommand{\fdMotorQ}{0.289}
\providecommand{\fdMotorM}{0.0609}

\begin{tikzpicture}
  \begin{axis}[paperplot, height=5.6cm, xmode=log, ymode=log,
      xmin=0.1, xmax=3.0, ymin=20, ymax=400,
      xlabel={peak torque $Q$ [\si{\newton\metre}]},
      ylabel={motor mass [\si{\gram}]},
      xtick={0.1,0.2,0.5,1,2,3}, xticklabels={0.1,0.2,0.5,1,2,3},
      ytick={20,50,100,200,400}, yticklabels={20,50,100,200,400},
      legend pos=north west]
    \addplot[cA, domain=0.1:3.0, samples=60]
      {1000*(x*0.1^0.3/5.5)^(1/1.3)};
    \addlegendentry{model \labelcref{eq:motormass}, $k_\tau=\SI{5.5}{\newton\metre\per\kilo\gram}$, $e=0.30$}
    \addplot[only marks, mark=square*, mark size=2.4pt, cB] coordinates
      {(0.30,66) (0.63,96) (2.03,190)};
    \addlegendentry{catalogue motors, T-Motor}
    \addplot[only marks, mark=*, mark size=2.8pt, cC] coordinates
      {(\fdMotorQ,1000*\fdMotorM)};
    \addlegendentry{the design's motor}
    \draw[cD, line width=0.8pt] (axis cs:0.30,66) -- (axis cs:0.30,63);
    \draw[cD, line width=0.8pt] (axis cs:0.63,96) -- (axis cs:0.63,111);
    \draw[cD, line width=0.8pt] (axis cs:2.03,190) -- (axis cs:2.03,273);
    \node[lbl, anchor=north west] at (axis cs:0.31,64) {MN2806, $-5\%$};
    \node[lbl, anchor=north west] at (axis cs:0.65,94) {MN3110, $+16\%$};
    \node[lbl, anchor=north] at (axis cs:2.03,178) {MN505-S, $+44\%$};
  \end{axis}
\end{tikzpicture}

\caption{The calibrated motor-mass model \cref{eq:motormass} against catalogue
motors, on logarithmic axes on which the model is a straight line of slope
$1/(1+e)=0.77$. The red connectors are the model errors of \cref{tab:motor}:
small at the size of motor this design uses (the green point, the design's
\SI{61}{\gram} motor at its torque for maximum thrust) and conservative at the
large end.}
\label{fig:motorcal}
\end{figure}

\subsection{Errors found by the model itself}

Three defects were found not by inspection but because a computed quantity was
inconsistent with a hand calculation.

\paragraph{The tip-speed design rule.} Holding tip speed fixed while growing
the rotor made power \emph{increase} with diameter, contradicting
\cref{obs:dl}. The cause was that at fixed tip speed the profile power
\cref{eq:profile} scales with disk area while the induced power falls only as
$A^{-1/2}$. The resolution is \cref{cor:designrule}: a rotor is designed at
constant blade loading, not constant tip speed, and with $\Vtip$ derived from
\cref{eq:ctsigma} the expected monotonicity is recovered.

\paragraph{Safety measured to the wrong point.} The simulator's minimum-distance
check measured from the vehicle centre of mass to the user. By \cref{eq:reach}
the blades extend $\mathcal{R}$ from that point, and for these designs
$\mathcal{R}=\SI{0.63}{\metre}$, so the check was optimistic by the entire reach
of the machine. Recomputing to the nearest blade tip showed the nominal
configuration violating its own safety distance by a factor of three. This is
discussed in \cref{sec:results}.

\paragraph{A startup transient masquerading as a failure.} The simulation began
with the vehicle at rest at the reference position while the human was already
at full walking speed. On a curved path the display sits directly in the
person's path, so the separation collapsed from \SI{1.100}{\metre} to
\SI{0.603}{\metre} within \SI{0.62}{\second} before recovering. Initialising the
vehicle and the governor with the reference velocity, as noted in
\cref{sec:numerics}, removed it; the minimum envelope gap then agrees with the
static tip gap of \cref{eq:clearance} to within a centimetre, the difference
being the conservative sphere against the actual blade tips. The same mistake
reappeared one level up: the vehicle was initialised level while the
reference on a curved path already needed a tilt and a body rate, and the
attitude loop spent the first second catching up. That transient set the
whole-window peak tracking error, which is the quantity the keep-out margin is
checked against. The initial attitude, body rate and thrust are now the trim
for the reference at $t=0$.

\begin{remark}
The common feature of all of these is that they were caught by comparison
against an independently computed quantity, not by reading the code. This is
the argument for keeping the analytic layer alive in the implementation even
after it has been superseded.
\end{remark}

\subsection{Errors found by independent review}

An independent review of the engine and of an earlier draft of this paper
found five defects that the tests above had not, each now covered by a
regression test.

\begin{enumerate}[nosep]
  \item The component closure bracketed the first sign change of
        $\mathcal{F}$ on a geometric grid of ratio $1.35$ and called a grid
        with no positive sample infeasible. Near the fold the positive stretch
        is narrower than the grid, and a design with two roots was reported as
        having none. The paper had also described that failure as
        infeasibility. \Cref{sec:numerics} gives the replacement.
  \item The torque authority divided the balanced differential by $N$, a
        factor of four on every axis, which propagated into the gain rule and
        into every layout comparison. \Cref{eq:authority} is now a linear
        program and \cref{app:proofs} the closed form it reduces to.
  \item Telemetry compared the state advanced to $t+\dd t$ with references
        evaluated at $t$, labelled $t$.
  \item The display slab's inertia was assigned to the wrong body axes, and the
        offsets were not moved to the centre of mass, so a boom produced no
        product of inertia. \Cref{sec:components} now builds one tensor from
        one geometry.
  \item After the review, and found by the tracking statistics it demanded:
        the attitude loop used a zero desired angular rate, so on a curved
        path it lagged the rotating thrust direction by a standing error. The
        rate feedforward of \cref{eq:omegad} removed a \SI{0.27}{\metre}
        tracking excursion and a rotor pass at \SI{0.34}{\metre} from the eye.
\end{enumerate}

\section{Design study}
\label{sec:results}

Every number in this section is produced by \code{paper/make\_results.py}
from the engine, together with the configuration that produced it
(\cref{sec:engine}). Unless stated otherwise the flights use a window of
\SI{\rTwindow}{\second} at a step of \SI{\rDt}{\milli\second}, statistics
taken after a settling interval of \SI{\rSettle}{\second}, gust seed
\rSeed, and the governor of \cref{mod:governor} at
$v_{\max}=\SI{2.5}{\metre\per\second}$,
$a_{\max}=\SI{2.5}{\metre\per\second\squared}$,
$\varepsilon=\SI{\rTrackMargin}{\milli\metre}$, heading rate cap
\SI{1.5}{\radian\per\second}.

\subsection{The requirement, worked backwards}

Four human-facing requirements bound the design.

\begin{enumerate}[label=(R\arabic*),nosep]
  \item \emph{Readability.} Two quantities, and they pull in opposite
        directions. Pixel density in pixels per degree,
        $1920/(2\arctan(w/2d))$ for a panel of width $w$ at distance $d$,
        measures sharpness and rises with distance. Text size is an angular
        height: ANSI/HFES 100~\cite{hfes2007} asks for a character height of
        at least \SI{16}{\arcminute}, with \SIrange{20}{22}{\arcminute}
        preferred, and it falls with distance. A desktop meets the second by
        scaling its interface, which divides the usable workspace by the
        scale factor.
  \item \emph{Noise.} A quiet office is \SI{45}{\decibel A}, ordinary speech
        \SI{60}{\decibel A}. Working beside a source for half an hour requires
        something in between.
  \item \emph{Downwash.} Loose paper begins to move at roughly
        \SI{5}{\metre\per\second}. By \cref{eq:dl} the wake velocity is
        $2\vind = 2\sqrt{\mathrm{DL}/2\rho}$, so this bounds disk loading.
  \item \emph{Size and safety.} The machine must fit beside a desk, and by
        \cref{eq:reach} its blades extend $\mathcal{R}$ from its centre.
\end{enumerate}

All four pull the same way, toward large slowly turning lightly loaded rotors,
and (R2) to (R4) conflict with the span limit through $\mathcal{R}=L+R$.

\subsection{The binding constraint is acoustic}

Applying \cref{eq:mmaxexplicit} to the fixed load of the design below
(\SI{\rFixed}{\gram} of panel, electronics and gimbal), the energetic wall is
comfortable at short sessions: at \SI{20}{\minute} the critical effective load
is \SI{\rWallTwenty}{\kilo\gram} against a requirement of
\SI{\rMzeroTwenty}{\kilo\gram}, a utilisation of \SI{\rWallUse}{\percent}.
Half an hour of flight costs about a quarter of a kilogram of cells. Energy
is not what stops this product.

Noise is. \Cref{obs:sixthpower} says the only lever with authority is tip
speed, and \cref{cor:designrule} says solidity is how one buys it. Optimising
under (R1) to (R4) drives the design to four rotors of
\SI{\rDiam}{\milli\metre} diameter with wide, low-aspect-ratio blades turning
at \SI{\rRpm}{rpm}, a tip speed of \SI{\rVtip}{\metre\per\second} against
roughly \SI{68}{\metre\per\second} for a commercial camera drone.

\subsection{Which display}

A resolution is not a payload. \Cref{tab:displays} evaluates the candidates at
the standoff the safety analysis below will impose, \SI{\rStandoff}{\metre},
each with the aircraft re-optimised around it by differential evolution over
rotor diameter, blade loading, aspect ratio and blade count under a span limit
of \SI{\rOptSpan}{\metre}, a free-field level of \SI{\rOptSpl}{\decibel A} at
\SI{1}{\metre}, a wake of \SI{\rOptWash}{\metre\per\second} and a gross mass
of \SI{\rOptMass}{\kilo\gram}. Every $1920$-wide panel is sharp at that
distance; none delivers a full desktop there, for the reason in (R1). Areal
density is what discriminates between the aircraft: a desktop monitor panel is
more than twice as heavy per unit area as a laptop panel, because it carries
thicker glass, a heavier backlight and a metal chassis.

\begin{table}[t]
\centering
\caption{Display options at \SI{\rStandoff}{\metre}, each with the aircraft
re-optimised under identical limits. ``Logical'' is the desktop width in
logical pixels once the interface is scaled to a \SI{16}{\arcminute} character
height. Areal density in \si{\kilo\gram\per\metre\squared}, n/a for a whole
device. Levels in \si{\decibel A}, free field at \SI{1}{\metre} and at the
user's ear in a $5\times4\times\SI{2.4}{\metre}$ room with $\bar\alpha=0.20$
by \cref{prop:room}. ``Over limits'' means the closure exists but the best
design found violates at least one of the four limits.}
\label{tab:displays}
\footnotesize\setlength{\tabcolsep}{4pt}
\begin{tabular}{@{}lrrrrrrrl@{}}
\toprule
Display & px/deg & Logical & Areal & Gross & Power & $L_p(\SI{1}{\metre})$ & At ear & Verdict\\
\midrule
\inputrows{results/tab_displays_rows}
\end{tabular}
\end{table}

\Cref{fig:displays} shows the same comparison graphically.

\begin{figure}[tbp]
\centering
\providecommand{\eDsafeCross}{1.06}
\providecommand{\eDtext}{0.43}
\providecommand{\eDfull}{0.65}
\providecommand{\eLogicalAtCross}{784}
\providecommand{\eArmL}{0.3827}
\providecommand{\eRotorR}{0.2460}
\providecommand{\eSpan}{1.257}
\providecommand{\eReach}{0.629}
\providecommand{\eScreenW}{0.345}
\providecommand{\eScreenH}{0.194}
\providecommand{\eRcp}{0.130}
\providecommand{\eEye}{1.500}
\providecommand{\eZrotor}{1.370}
\providecommand{\eZrotorBelow}{1.630}
\providecommand{\eStandoff}{1.20}
\providecommand{\eDsafe}{0.45}
\providecommand{\eDeltaZ}{0.130}
\providecommand{\eDisplayOrder}{13.3 in laptop panel,15.6 in laptop panel,iPad Pro 13 panel,iPad Pro 13 whole,15.6 in lid assembly,21.5 in monitor panel,24 in monitor panel,13.3 in ultrabook,15.6 in laptop whole}
\providecommand{\eOptSpl}{58}
\providecommand{\eOptMass}{6.0}
\providecommand{\eTurnClear}{0.61}
\providecommand{\eTurnTrack}{64}
\providecommand{\eTrackMargin}{100}
\providecommand{\eTurnPmean}{111.9}
\providecommand{\eHoverBat}{288}
\providecommand{\eReservePct}{12}
\providecommand{\eTethPsrc}{109.6}
\providecommand{\eTethArea}{0.228}
\providecommand{\eTethV}{48}
\providecommand{\eTethL}{5}
\providecommand{\eTethDrop}{5}
%
\begin{tikzpicture}
\begin{groupplot}[
  group style={group size=2 by 1, horizontal sep=0.35cm,
               yticklabels at=edge left},
  paperplot, width=0.40\linewidth, height=6.4cm,
  xbar, /pgf/bar width=6.5pt, /pgf/bar shift=0pt,
  y dir=reverse, symbolic y coords/.expand once={\eDisplayOrder},
  ytick/.expand once={\eDisplayOrder}, enlarge y limits=0.09,
  y tick label style={font=\footnotesize},
  legend style={at={(0.98,0.98)}, anchor=north east}]
\nextgroupplot[xlabel={gross mass $m$ [\si{\kilogram}]}, xmin=0, xmax=7]
  \addplot[fill=cA!70, draw=cA] table[x=gross, y=tag] {results/fig/e_displays_ok.dat};
  \addplot[fill=cD!55, draw=cD] table[x=gross, y=tag] {results/fig/e_displays_bad.dat};
  \draw[cD, dashed, line width=0.8pt]
    (axis cs:\eOptMass,{13.3 in laptop panel}) -- (axis cs:\eOptMass,{15.6 in laptop whole});
  \addlegendimage{line legend, cD, dashed, line width=0.8pt}
  \legend{closes, over limits, limit}
\nextgroupplot[xlabel={$L_p(\SI{1}{\metre})$ [dBA]}, xmin=45, xmax=70]
  \addplot[fill=cA!70, draw=cA] table[x=spl, y=tag] {results/fig/e_displays_ok.dat};
  \addplot[fill=cD!55, draw=cD] table[x=spl, y=tag] {results/fig/e_displays_bad.dat};
  \draw[cD, dashed, line width=0.8pt]
    (axis cs:\eOptSpl,{13.3 in laptop panel}) -- (axis cs:\eOptSpl,{15.6 in laptop whole});
\end{groupplot}
\end{tikzpicture}

\caption{The display candidates of \cref{tab:displays}, each with its aircraft
re-optimised. Left: gross mass, against the \SI{\eOptMass}{\kilo\gram}
limit. Right: free-field level at \SI{1}{\metre}, against the
\SI{\eOptSpl}{dBA} limit. Every candidate that fails exceeds the noise limit,
and only one of them exceeds the mass limit: a heavier display needs a more
heavily loaded rotor, and by \cref{obs:sixthpower} that is paid for in
decibels long before it is paid for in kilograms.}
\label{fig:displays}
\end{figure}

\begin{finding}
The full $1920\times1080$ panel is available, provided it is a laptop panel
and not a monitor panel. Flying a whole laptop is not: the panel is a sixth
of a 15.6-inch machine's mass and the rest is keyboard, battery, cooling and
chassis, none of which is visible while hovering. A tablet is different in
kind because a tablet is almost entirely screen, which is why the iPad flown
whole still closes while the laptop flown whole does not. What no panel of
this size delivers at \SI{\rStandoff}{\metre} is a full desktop's worth of
text: see \cref{sec:readability}.
\end{finding}

\subsection{Geometry relative to the user}

The standoff a user specifies is measured to the \emph{display}. The rotors
extend $\mathcal{R}=L+R$ back toward them from the machine's centre.

\begin{proposition}[Blade clearance]
\label{prop:clearance}
Let the display be at horizontal standoff $d$ from the user's eye, carried a
distance $b\ge0$ ahead of the rotor centroid on a boom, with the rotor plane a
height $\Delta z$ above the eye. Then the nearest blade tip lies at
\begin{equation}
  \ell_0 = \sqrt{\big(\max\{d+b-\mathcal{R},0\}\big)^2 + \Delta z^2} ,
  \label{eq:clearance}
\end{equation}
and the requirement $\ell_0\ge d_{\mathrm{safe}}$ is met either by
\begin{equation}
  d \ge \mathcal{R} + \sqrt{d_{\mathrm{safe}}^2 - \Delta z^2} - b
  \qquad\text{or}\qquad
  b \ge \mathcal{R} + \sqrt{d_{\mathrm{safe}}^2-\Delta z^2} - d .
  \label{eq:clearancefix}
\end{equation}
\end{proposition}

At a nominal $d=\SI{0.70}{\metre}$ with $\mathcal{R}=\SI{\rReach}{\metre}$
and $|\Delta z|=\SI{0.13}{\metre}$, \cref{eq:clearance} gives
$\ell_0=\SI{0.148}{\metre}$ against a safe distance of \SI{0.45}{\metre}: a
violation by a factor of three, invisible to any check evaluated at the
vehicle centre.

\Cref{fig:clearance} draws the two standoffs to scale.

\begin{figure}[tbp]
\centering
\providecommand{\eDsafeCross}{1.06}
\providecommand{\eDtext}{0.43}
\providecommand{\eDfull}{0.65}
\providecommand{\eLogicalAtCross}{784}
\providecommand{\eArmL}{0.3827}
\providecommand{\eRotorR}{0.2460}
\providecommand{\eSpan}{1.257}
\providecommand{\eReach}{0.629}
\providecommand{\eScreenW}{0.345}
\providecommand{\eScreenH}{0.194}
\providecommand{\eRcp}{0.130}
\providecommand{\eEye}{1.500}
\providecommand{\eZrotor}{1.370}
\providecommand{\eZrotorBelow}{1.630}
\providecommand{\eStandoff}{1.20}
\providecommand{\eDsafe}{0.45}
\providecommand{\eDeltaZ}{0.130}
\providecommand{\eDisplayOrder}{13.3 in laptop panel,15.6 in laptop panel,iPad Pro 13 panel,iPad Pro 13 whole,15.6 in lid assembly,21.5 in monitor panel,24 in monitor panel,13.3 in ultrabook,15.6 in laptop whole}
\providecommand{\eOptSpl}{58}
\providecommand{\eOptMass}{6.0}
\providecommand{\eTurnClear}{0.61}
\providecommand{\eTurnTrack}{64}
\providecommand{\eTrackMargin}{100}
\providecommand{\eTurnPmean}{111.9}
\providecommand{\eHoverBat}{288}
\providecommand{\eReservePct}{12}
\providecommand{\eTethPsrc}{109.6}
\providecommand{\eTethArea}{0.228}
\providecommand{\eTethV}{48}
\providecommand{\eTethL}{5}
\providecommand{\eTethDrop}{5}
%
\begin{tikzpicture}[x=3.3cm, y=3.3cm]
\foreach \dd/\sx/\col/\ellv/\ttl in {
    0.70/0.00/cD/0.148/{$d=\SI{0.70}{\metre}$, no boom: violated},
    1.20/2.05/cC/0.586/{$d=\SI{1.20}{\metre}$, no boom: the design}} {
  \begin{scope}[shift={(\sx,0)}]
    \pgfmathsetmacro\tip{\dd-\eReach}
    \pgfmathsetmacro\zlo{\eEye-\eScreenH/2}
    \pgfmathsetmacro\zhi{\eEye+\eScreenH/2}
    % the safe sphere about the eye
    \fill[cD!6] (0,\eEye) circle (\eDsafe);
    \draw[cD, dashed, line width=0.6pt] (0,\eEye) circle (\eDsafe);
    \node[lbl, text=cD, anchor=north] at (0,\eEye-\eDsafe-0.01) {$d_{\mathrm{safe}}$ about the eye};
    % the user
    \draw[line width=0.8pt, fill=white] (-0.10,\eEye+0.02) circle (0.10);
    \draw[line width=0.8pt] (-0.10,\eEye-0.08) -- (-0.10,1.12);
    \fill (0,\eEye) circle (0.9pt);
    \node[lbl, anchor=south] at (-0.10,\eEye+0.12) {eye};
    % rotor plane seen from the side: arms, two disks edge on, hubs
    \draw[line width=1.0pt] (\dd-\eArmL,\eZrotor) -- (\dd+\eArmL,\eZrotor);
    \draw[cB, line width=2.6pt] (\dd-\eArmL-\eRotorR,\eZrotor) -- (\dd-\eArmL+\eRotorR,\eZrotor);
    \draw[cB, line width=2.6pt] (\dd+\eArmL-\eRotorR,\eZrotor) -- (\dd+\eArmL+\eRotorR,\eZrotor);
    \fill (\dd-\eArmL,\eZrotor) circle (0.7pt) (\dd+\eArmL,\eZrotor) circle (0.7pt);
    % mast and display, screen above the rotors
    \draw[cG, line width=1.0pt] (\dd,\eZrotor) -- (\dd,\zlo);
    \draw[cA, line width=2.4pt] (\dd,\zlo) -- (\dd,\zhi);
    \node[lbl, text=cA, anchor=west] at (\dd+0.03,\zhi-0.02) {display};
    % the nearest tip and the clearance
    \draw[\col, line width=1.1pt] (0,\eEye) -- (\tip,\eZrotor);
    \fill[\col] (\tip,\eZrotor) circle (1.2pt);
    \node[lbl, text=\col!80!black, anchor=south west, inner sep=1pt]
      at (\tip+0.04,\eZrotor+0.035) {$\ell_0=\SI{\ellv}{\metre}$};
    % dimensions
    \draw[dim] (0,1.86) -- (\dd,1.86) node[midway, above, lbl] {$d$};
    \draw[construct] (\dd,\zhi) -- (\dd,1.89);
    \draw[construct] (0,\eEye) -- (0,1.89);
    \draw[dim] (\tip,1.24) -- (\dd,1.24) node[midway, below, lbl] {$\mathcal{R}=L+R$};
    \draw[construct] (\tip,\eZrotor) -- (\tip,1.21);
    \draw[construct] (\dd,\eZrotor) -- (\dd,1.21);
    \draw[dim] (-0.33,\eZrotor) -- (-0.33,\eEye) node[midway, left, lbl] {$\Delta z$};
    \draw[construct] (-0.36,\eZrotor) -- (\tip,\eZrotor);
    \draw[construct] (-0.36,\eEye) -- (-0.21,\eEye);
    \node[font=\small, anchor=south] at (\dd/2+0.25,2.00) {\ttl};
  \end{scope}
}
\end{tikzpicture}

\caption{\Cref{prop:clearance} in side view, to scale. The display is held at
eye height at standoff $d$; with the screen above the rotors the rotor plane
lies $|\Delta z|=\SI{\eDeltaZ}{\metre}$ below the eye; the nearest blade tip is
the reach $\mathcal{R}=L+R$ back toward the user from the rotor centroid. At
$d=\SI{0.70}{\metre}$ the tip is inside the sphere of radius $d_{\mathrm{safe}}$
about the eye, although the display is well outside it; a check made at the
vehicle centre would pass this design. At $d=\SI{1.20}{\metre}$ the tip is
outside.}
\label{fig:clearance}
\end{figure}

\Cref{tab:clearance} compares the two remedies, each option closed, its gains
capped by \cref{mod:gains} for its own authority (attitude and position gains
together, since the cascade needs its time-scale separation), and flown
through the turnaround mission with the governor active. For every option but
the last standoff the ideal display position lies inside the keep-out radius
$\varrho = d_{\mathrm{safe}}+\mathcal{R}+\varepsilon$, by up to
\SI{\rCldNearInside}{\centi\metre}; those missions start on the sphere and the
governor holds the reference there, so their ``turn'' clearance measures the
keep-out sphere and not the option. The sphere is conservative for a boom,
whose reach is set by the rear rotors while the front rotors are what
approach the user.

\begin{table}[t]
\centering
\caption{Two ways to recover blade clearance. The boom keeps the display close
and pays in pitch inertia and authority; standing back is free on every axis
and quieter, and pays in text size. ``Turn'' is the minimum distance from the
eye to the nearest rotor disk during the turnaround mission, against the
\SI{0.45}{\metre} required. $J_{yy}$ in \si{\kilo\gram\metre\squared},
$\alpha^{\max}_{\mathrm{pitch}}$ in \si{\radian\per\second\squared}, $K_R$ in
\si{\per\second\squared}, screen angle is the panel's apparent width, logical
width as in \cref{tab:displays}.}
\label{tab:clearance}
\footnotesize\setlength{\tabcolsep}{4pt}
\begin{tabular}{@{}lrrrrrrrrr@{}}
\toprule
Option & Tip gap & Gross & Power & $J_{yy}$ & $\alpha^{\max}_{\mathrm{pitch}}$
       & $K_R$ & Screen & Logical & Turn\\
\midrule
\inputrows{results/tab_clearance_rows}
\end{tabular}
\end{table}

\begin{finding}
A boom of \SI{0.38}{\metre} multiplies the pitch inertia by \rBoomJratio,
reduces the pitch authority to \rBoomAlphaRatio\ of the no-boom value and the
admissible attitude gain to $K_R=\rClbMidKR$, and costs \SI{\rBoomGross}{\gram}
of gross mass and \SI{\rBoomEar}{\decibel A} at the ear against
\SI{\rStandEar}{\decibel A}. Standing back to \SI{1.20}{\metre} achieves the
same static clearance at no cost in mass, power or authority, is the only
option whose ideal position is outside the keep-out sphere, and keeps its
rotors \SI{\rCldFarClear}{\metre} from the eye through the turnaround. The
price is the one (R1) names, and it is paid below.
\end{finding}

\subsection{What the user can read}
\label{sec:readability}

At \SI{\rStandoff}{\metre} the 15.6-inch panel subtends \ang{\rScreenDeg} at
\SI{\rPpd}{px\per\degree}: sharper than any desktop monitor. But a
\SI{16}{px} body-text character, cap height \SI{11}{px}, subtends only
\SI{\rCapArcmin}{\arcminute} at that distance, against the \SI{16}{\arcminute}
minimum of~\cite{hfes2007}. Meeting the minimum takes an interface scale of
\SI{\rUiScale}{\percent} (\SI{\rUiScalePref}{\percent} for the preferred
\SI{20}{\arcminute}), which leaves a logical desktop of
$\rLogicalW\times\rLogicalH$ pixels. That is a tablet's workspace on a
laptop's panel. At \SI{0.70}{\metre} the same panel gives
\SI{\rCapArcminNear}{\arcminute} and a logical width of \rLogicalWNear, a
laptop's workspace, and its blades \SI{0.15}{\metre} from the eye.

\begin{finding}
The two safe ways to fly a panel this size, standing back or a boom, and the
requirement that the user read a full desktop on it, are not simultaneously
satisfiable at $1920\times1080$ in \SI{345}{\milli\metre}. A safe flying
screen at reading distance either carries a physically larger panel, which
\cref{tab:displays} shows the acoustics will not allow, or presents a smaller
workspace than the resolution suggests. The resolution is not the constraint;
the angular size of a character is.
\end{finding}

\Cref{fig:readability} shows the whole trade over standoff.

\begin{figure}[tbp]
\centering
\providecommand{\eDsafeCross}{1.06}
\providecommand{\eDtext}{0.43}
\providecommand{\eDfull}{0.65}
\providecommand{\eLogicalAtCross}{784}
\providecommand{\eArmL}{0.3827}
\providecommand{\eRotorR}{0.2460}
\providecommand{\eSpan}{1.257}
\providecommand{\eReach}{0.629}
\providecommand{\eScreenW}{0.345}
\providecommand{\eScreenH}{0.194}
\providecommand{\eRcp}{0.130}
\providecommand{\eEye}{1.500}
\providecommand{\eZrotor}{1.370}
\providecommand{\eZrotorBelow}{1.630}
\providecommand{\eStandoff}{1.20}
\providecommand{\eDsafe}{0.45}
\providecommand{\eDeltaZ}{0.130}
\providecommand{\eDisplayOrder}{13.3 in laptop panel,15.6 in laptop panel,iPad Pro 13 panel,iPad Pro 13 whole,15.6 in lid assembly,21.5 in monitor panel,24 in monitor panel,13.3 in ultrabook,15.6 in laptop whole}
\providecommand{\eOptSpl}{58}
\providecommand{\eOptMass}{6.0}
\providecommand{\eTurnClear}{0.61}
\providecommand{\eTurnTrack}{64}
\providecommand{\eTrackMargin}{100}
\providecommand{\eTurnPmean}{111.9}
\providecommand{\eHoverBat}{288}
\providecommand{\eReservePct}{12}
\providecommand{\eTethPsrc}{109.6}
\providecommand{\eTethArea}{0.228}
\providecommand{\eTethV}{48}
\providecommand{\eTethL}{5}
\providecommand{\eTethDrop}{5}
%
\begin{tikzpicture}
\begin{groupplot}[
  group style={group size=1 by 2, vertical sep=0.35cm,
               xlabels at=edge bottom, xticklabels at=edge bottom},
  paperplot, height=4.6cm, xmin=0.35, xmax=2.0,
  xlabel={standoff $d$ from the eye to the display [\si{\metre}]}]
% (a) usable desktop once text is scaled to 16 arcmin
\nextgroupplot[ylabel={logical width [px]}, ymin=0, ymax=2150,
  ytick={0,500,1000,1280,1920}]
  \fill[cC!14] (axis cs:0.35,0) rectangle (axis cs:\eDfull,2150);
  \addplot[cA] table[x=d, y=logical] {results/fig/e_readability.dat};
  \addplot[cG, dashed, domain=0.35:2.0, samples=2] {1280};
  \node[lbl, cG, anchor=south east] at (axis cs:1.98,1280) {full desktop, \num{1280} px};
  \node[lbl, text=cC!55!black, anchor=south west, align=left]
    at (axis cs:0.37,120) {full\\ desktop\\ legible};
  \draw[construct] (axis cs:\eDsafeCross,0) -- (axis cs:\eDsafeCross,2150);
  \draw[cA, densely dotted, line width=0.7pt] (axis cs:\eStandoff,0) -- (axis cs:\eStandoff,2150);
  \node[lbl, text=cA, anchor=south west] at (axis cs:\eStandoff,1500) {design};
  \addplot[only marks, mark=*, mark size=1.8pt, cB]
    coordinates {(\eDsafeCross,\eLogicalAtCross)};
  \node[lbl, text=cB!80!black, anchor=south west]
    at (axis cs:\eDsafeCross,\eLogicalAtCross) {\eLogicalAtCross\ px at $d=\SI{\eDsafeCross}{\metre}$};
% (b) nearest blade tip, prop:clearance
\nextgroupplot[ylabel={blade tip to eye $\ell_0$ [\si{\metre}]}, ymin=0, ymax=1.35]
  \fill[cC!14] (axis cs:\eDsafeCross,0) rectangle (axis cs:2.0,1.35);
  \addplot[cB] table[x=d, y=gap] {results/fig/e_readability.dat};
  \addplot[cD, dashed, domain=0.35:2.0, samples=2] {\eDsafe};
  \node[lbl, text=cD, anchor=south west] at (axis cs:0.37,\eDsafe) {$d_{\mathrm{safe}}=\SI{\eDsafe}{\metre}$};
  \node[lbl, text=cC!55!black, anchor=north east, align=right]
    at (axis cs:1.98,1.30) {blades outside\\ $d_{\mathrm{safe}}$};
  \node[lbl, text=cB!80!black, anchor=south west, align=left]
    at (axis cs:0.37,0.16) {rotors overhang the user: $\ell_0=|\Delta z|$};
  \draw[construct] (axis cs:\eDsafeCross,0) -- (axis cs:\eDsafeCross,1.35);
  \draw[cA, densely dotted, line width=0.7pt] (axis cs:\eStandoff,0) -- (axis cs:\eStandoff,1.35);
\end{groupplot}
\end{tikzpicture}

\caption{Readability against blade clearance, over the standoff $d$ from the
eye to the display. Top: the logical desktop width once the interface is
scaled to a \SI{16}{\arcminute} character height; it is the full \num{1920}
pixels only within \SI{\eDtext}{\metre}, where unscaled text already meets the
minimum, and it falls as $1/d$ beyond. A \num{1280}-pixel desktop needs
$d\le\SI{\eDfull}{\metre}$ (green). Bottom: the nearest blade tip $\ell_0$ of
\cref{eq:clearance}; it stays at $|\Delta z|$ while the rotors overhang the
user and clears $d_{\mathrm{safe}}$ only for $d\ge\SI{\eDsafeCross}{\metre}$
(green). The two green regions do not overlap, and at the first safe standoff
the desktop is already down to \eLogicalAtCross\ pixels.}
\label{fig:readability}
\end{figure}

\subsection{Screen above or below the rotors}

Because the display is vertical and the rotor flow is vertical
(\cref{sec:dynamics}), placing the panel above or below the rotor plane is
nearly free aerodynamically: \SI{\rGrossBelow}{\kilo\gram} against
\SI{\rGross}{\kilo\gram}, with the same noise. What changes is geometry. With
the screen below, the rotor plane sits at \SI{\rZrotorBelow}{\metre}, level
with the user's head; with it above, the rotors drop to
\SI{\rZrotorAbove}{\metre} and the panel itself lies between the blades and
the user's face. The latter is adopted. Tracking is marginally better in still
air, because by the remark following \cref{eq:screenmoment} the drag moment
becomes weathercock-stable, and marginally worse in a steady breeze for the
same reason.

\Cref{fig:placement} compares the two at the design standoff.

\begin{figure}[tbp]
\centering
\providecommand{\eDsafeCross}{1.06}
\providecommand{\eDtext}{0.43}
\providecommand{\eDfull}{0.65}
\providecommand{\eLogicalAtCross}{784}
\providecommand{\eArmL}{0.3827}
\providecommand{\eRotorR}{0.2460}
\providecommand{\eSpan}{1.257}
\providecommand{\eReach}{0.629}
\providecommand{\eScreenW}{0.345}
\providecommand{\eScreenH}{0.194}
\providecommand{\eRcp}{0.130}
\providecommand{\eEye}{1.500}
\providecommand{\eZrotor}{1.370}
\providecommand{\eZrotorBelow}{1.630}
\providecommand{\eStandoff}{1.20}
\providecommand{\eDsafe}{0.45}
\providecommand{\eDeltaZ}{0.130}
\providecommand{\eDisplayOrder}{13.3 in laptop panel,15.6 in laptop panel,iPad Pro 13 panel,iPad Pro 13 whole,15.6 in lid assembly,21.5 in monitor panel,24 in monitor panel,13.3 in ultrabook,15.6 in laptop whole}
\providecommand{\eOptSpl}{58}
\providecommand{\eOptMass}{6.0}
\providecommand{\eTurnClear}{0.61}
\providecommand{\eTurnTrack}{64}
\providecommand{\eTrackMargin}{100}
\providecommand{\eTurnPmean}{111.9}
\providecommand{\eHoverBat}{288}
\providecommand{\eReservePct}{12}
\providecommand{\eTethPsrc}{109.6}
\providecommand{\eTethArea}{0.228}
\providecommand{\eTethV}{48}
\providecommand{\eTethL}{5}
\providecommand{\eTethDrop}{5}
%
% Side views at the design standoff, to scale: the display is held at eye
% height either way; what moves is the rotor plane.
\begin{tikzpicture}[x=3.0cm, y=3.0cm]
\foreach \zr/\sx/\anc/\dy/\ttl in {
    \eZrotorBelow/0.00/south east/0.03/{screen below the rotors},
    \eZrotor/2.30/north east/-0.03/{screen above the rotors: adopted}} {
  \begin{scope}[shift={(\sx,0)}]
    \pgfmathsetmacro\tip{\eStandoff-\eReach}
    \pgfmathsetmacro\zlo{\eEye-\eScreenH/2}
    \pgfmathsetmacro\zhi{\eEye+\eScreenH/2}
    % eye height across the panel
    \draw[construct] (-0.22,\eEye) -- (\eStandoff+\eReach,\eEye);
    \node[lbl, text=cG!80!black, anchor=south west] at (-0.22,\eEye+0.13) {eye height \SI{\eEye}{\metre}};
    \draw[construct, -{Stealth[length=1.2mm]}] (-0.02,\eEye+0.14) -- (-0.02,\eEye+0.01);
    % the user: head and shoulders
    \draw[line width=0.8pt, fill=white] (-0.10,\eEye+0.02) circle (0.10);
    \draw[line width=0.8pt] (-0.10,\eEye-0.08) -- (-0.10,1.10);
    \draw[line width=0.8pt] (-0.24,1.36) -- (0.04,1.36);
    \fill (0,\eEye) circle (0.9pt);
    % rotor plane: arms, two disks edge on, hubs
    \draw[line width=1.0pt] (\eStandoff-\eArmL,\zr) -- (\eStandoff+\eArmL,\zr);
    \draw[cB, line width=2.6pt] (\eStandoff-\eArmL-\eRotorR,\zr) -- (\eStandoff-\eArmL+\eRotorR,\zr);
    \draw[cB, line width=2.6pt] (\eStandoff+\eArmL-\eRotorR,\zr) -- (\eStandoff+\eArmL+\eRotorR,\zr);
    \fill (\eStandoff-\eArmL,\zr) circle (0.7pt) (\eStandoff+\eArmL,\zr) circle (0.7pt);
    \node[lbl, text=cB!80!black, anchor=\anc] at (\eStandoff+\eReach,\zr+\dy)
      {rotor plane \SI{\zr}{\metre}};
    % mast and display at eye height
    \draw[cG, line width=1.0pt] (\eStandoff,\zr) -- (\eStandoff,\eEye);
    \draw[cA, line width=2.4pt] (\eStandoff,\zlo) -- (\eStandoff,\zhi);
    \node[lbl, text=cA, anchor=west] at (\eStandoff+0.03,\zhi-0.03) {display};
    % the nearest tip
    \draw[cC, line width=1.1pt] (0,\eEye) -- (\tip,\zr);
    \fill[cC] (\tip,\zr) circle (1.2pt);
    \node[font=\small, anchor=south] at (\eStandoff/2+0.25,1.86) {\ttl};
  \end{scope}
}
\end{tikzpicture}

\caption{The two placements of the display at the design standoff, to scale,
with the display held at eye height in both. The nearest-tip clearance
(green) is the same either way, because \cref{eq:clearance} depends on
$\Delta z$ only through $\Delta z^2$. What differs is where the rotor plane
sits relative to the head: with the screen below it is at the top of the head,
with the screen above it drops below the chin, and the panel is then the
object between the blades and the face. The heights marked are rotor plane
heights above the floor.}
\label{fig:placement}
\end{figure}

\subsection{The design}

\Cref{tab:design} gives the resulting machine with a buildable payload: a bare
laptop panel on a matched controller board, a Raspberry Pi running a remote
desktop client, a separate flight controller, and an on-sensor inference
camera. \Cref{fig:design} draws it to scale.

\begin{figure}[tbp]
\centering
\providecommand{\eDsafeCross}{1.06}
\providecommand{\eDtext}{0.43}
\providecommand{\eDfull}{0.65}
\providecommand{\eLogicalAtCross}{784}
\providecommand{\eArmL}{0.3827}
\providecommand{\eRotorR}{0.2460}
\providecommand{\eSpan}{1.257}
\providecommand{\eReach}{0.629}
\providecommand{\eScreenW}{0.345}
\providecommand{\eScreenH}{0.194}
\providecommand{\eRcp}{0.130}
\providecommand{\eEye}{1.500}
\providecommand{\eZrotor}{1.370}
\providecommand{\eZrotorBelow}{1.630}
\providecommand{\eStandoff}{1.20}
\providecommand{\eDsafe}{0.45}
\providecommand{\eDeltaZ}{0.130}
\providecommand{\eDisplayOrder}{13.3 in laptop panel,15.6 in laptop panel,iPad Pro 13 panel,iPad Pro 13 whole,15.6 in lid assembly,21.5 in monitor panel,24 in monitor panel,13.3 in ultrabook,15.6 in laptop whole}
\providecommand{\eOptSpl}{58}
\providecommand{\eOptMass}{6.0}
\providecommand{\eTurnClear}{0.61}
\providecommand{\eTurnTrack}{64}
\providecommand{\eTrackMargin}{100}
\providecommand{\eTurnPmean}{111.9}
\providecommand{\eHoverBat}{288}
\providecommand{\eReservePct}{12}
\providecommand{\eTethPsrc}{109.6}
\providecommand{\eTethArea}{0.228}
\providecommand{\eTethV}{48}
\providecommand{\eTethL}{5}
\providecommand{\eTethDrop}{5}
%
\begin{tikzpicture}[x=4.9cm, y=4.9cm]
\pgfmathsetmacro\cc{\eArmL*cos(45)}
\pgfmathtruncatemacro\hubmm{round(2*\cc*1000)}
% ---- top view --------------------------------------------------------------
\begin{scope}
  \foreach \sx/\sy/\ang in {1/1/20, -1/1/-35, -1/-1/60, 1/-1/-10} {
    \coordinate (h) at (\sx*\cc,\sy*\cc);
    \draw[line width=1.6pt, cG!70!black] (0,0) -- (h);
    \draw[cB, fill=cB!7, line width=0.6pt] (h) circle (\eRotorR);
    \draw[cG, line width=0.4pt] (h) circle (\eRotorR+0.012);
    \draw[cB!70!black, line width=1.3pt]
      ($(h)+(\ang:\eRotorR-0.01)$) -- ($(h)+(\ang+180:\eRotorR-0.01)$);
    \fill (h) circle (0.8pt);
  }
  \draw[fill=cG!25] (-0.055,-0.055) rectangle (0.055,0.055);
  % the display seen from above: the panel lies in the b2-b3 plane
  \draw[cA, line width=2.4pt] (0,-\eScreenW/2) -- (0,\eScreenW/2);
  \draw[cA, line width=0.4pt] (0,-\eScreenW/2) -- (0,-0.62)
    node[below right, lbl, text=cA, xshift=-2pt] {display, above the rotor plane};
  % hub spacing across the top
  \draw[construct] (-\cc,\cc) -- (-\cc,\cc+\eRotorR+0.07);
  \draw[construct] (\cc,\cc) -- (\cc,\cc+\eRotorR+0.07);
  \draw[dim] (-\cc,\cc+\eRotorR+0.05) -- (\cc,\cc+\eRotorR+0.05)
    node[midway, above, lbl] {$2L\cos\ang{45}=\SI{\hubmm}{\milli\metre}$};
  % rotor diameter on rotor 1
  \draw[construct] (\cc,\cc+\eRotorR) -- (\cc+\eRotorR+0.08,\cc+\eRotorR);
  \draw[construct] (\cc,\cc-\eRotorR) -- (\cc+\eRotorR+0.08,\cc-\eRotorR);
  \draw[dim] (\cc+\eRotorR+0.06,\cc-\eRotorR) -- (\cc+\eRotorR+0.06,\cc+\eRotorR)
    node[midway, right, lbl] {$D=\SI{\rDiam}{\milli\metre}$};
  % body axes, set apart: b1 toward the user
  \draw[vec, cA] (-0.68,-0.64) -- (-0.50,-0.64) node[right, lbl] {$\vect b_1$};
  \draw[vec, cA] (-0.68,-0.64) -- (-0.68,-0.46) node[above, lbl] {$\vect b_2$};
  \node[lbl, anchor=west] at (0.60,-0.10) {user $\rightarrow$};
  \node[font=\small, anchor=south] at (0,0.70) {top view};
\end{scope}
% ---- side view --------------------------------------------------------------
\begin{scope}[shift={(1.50,0)}]
  \pgfmathsetmacro\zlo{\eRcp-\eScreenH/2}
  \pgfmathsetmacro\zhi{\eRcp+\eScreenH/2}
  \draw[line width=1.6pt, cG!70!black] (-\cc,0) -- (\cc,0);
  \draw[cB, line width=2.6pt] (-\cc-\eRotorR,0) -- (-\cc+\eRotorR,0);
  \draw[cB, line width=2.6pt] (\cc-\eRotorR,0) -- (\cc+\eRotorR,0);
  \draw[fill=cG!25] (-0.055,-0.035) rectangle (0.055,0.02);
  \draw[cG, line width=1.0pt] (0,0.02) -- (0,\zlo);
  \draw[cA, line width=2.4pt] (0,\zlo) -- (0,\zhi);
  \draw[cA, fill=white] (0,\eRcp) circle (0.011);
  \node[lbl, text=cA, anchor=west] at (0.02,\zhi) {display, gimbal pivot};
  \draw[vec, cA] (0.40,-0.20) -- (0.58,-0.20) node[right, lbl] {$\vect b_1$};
  \draw[vec, cA] (0.40,-0.20) -- (0.40,-0.02) node[above, lbl] {$\vect b_3$};
  % dimensions
  \draw[dim] (-0.16,0) -- (-0.16,\eRcp) node[midway, left, lbl] {$r_{cp}$};
  \draw[construct] (-0.19,\eRcp) -- (-0.012,\eRcp);
  \draw[dim] (0.12,\zlo) -- (0.12,\zhi) node[midway, right, lbl] {$h=\SI{194}{\milli\metre}$};
  \draw[dim] (-\cc-\eRotorR,-0.10) -- (\cc+\eRotorR,-0.10)
    node[midway, below, lbl] {$2(L\cos\ang{45}+R)$};
  \draw[construct] (-\cc-\eRotorR,0) -- ++(0,-0.12);
  \draw[construct] (\cc+\eRotorR,0) -- ++(0,-0.12);
  \node[font=\small, anchor=south] at (0,0.70) {side view};
\end{scope}
\end{tikzpicture}

\caption{DESK-1/W to scale. Top view: four rotors of diameter $D$ in the
$\times$ layout, the arm length $L$ set by the tip gap rule of
\cref{sec:components} so that adjacent guards (outer circles) nearly touch,
the display edge on above the hub, and $\vect b_1$ pointing at the user. Side
view: the display on its gimbal a height $r_{cp}$ above the rotor plane,
which is what puts the panel, and not the blades, between the rotors and the
user's face. The largest dimension, tip to tip across the diagonal, is the
span $2(L+R)=\SI{\rSpan}{\metre}$.}
\label{fig:design}
\end{figure}

\begin{table}[t]
\centering
\caption{DESK-1/W with the thin-client payload, screen above the rotors,
standoff \SI{\rStandoff}{\metre}. Sized for hover by \cref{eq:fpclosure}.}
\label{tab:design}
\footnotesize\setlength{\tabcolsep}{4pt}
\begin{tabular}{@{}llrl@{}}
\toprule
& Quantity & Value & Note\\
\midrule
\multirow{4}{*}{Geometry}
& rotors & $4\times\SI{\rDiam}{\milli\metre}$ & 2 blades, \SI{\rChord}{\milli\metre} chord, AR 3.0\\
& solidity & \rSolidity & below the 0.25 cascade limit\\
& span & \SI{\rSpan}{\metre} & reach \SI{\rReach}{\metre}\\
& display & $345\times\SI{194}{\milli\metre}$ & 15.6 in, \SI{\rPpd}{px\per\degree}, $\rLogicalW\times\rLogicalH$ logical\\
\midrule
\multirow{5}{*}{Mass}
& fixed payload & \SI{\rFixed}{\gram} & panel \SI{\rPanel}{\gram}, electronics \SI{\rElectronics}{\gram}, gimbal \SI{\rGimbal}{\gram}\\
& frame & \SI{\rFrame}{\gram} & arms \SI{\rArms}{\gram}, guards \SI{\rGuards}{\gram}\\
& propulsion & \SI{\rProp}{\gram} & motors \SI{\rMotors}{\gram}, props \SI{\rProps}{\gram}\\
& battery & \SI{\rBattery}{\gram} & \rPackS S1P \rPackCell, \SI{\rPackWh}{\watt\hour}\\
& \textbf{gross} & \textbf{\SI{\rGross}{\kilo\gram}} & $\FM=\rFM$ implied\\
\midrule
\multirow{4}{*}{Rotor}
& speed & \SI{\rRpm}{rpm} & \SI{\rRpmMax}{rpm} at $T_{\max}$\\
& tip speed & \SI{\rVtip}{\metre\per\second} & \SI{\rVtipMax}{\metre\per\second} at $T_{\max}$\\
& blade loading & \rCtSigma & stall at 0.14\\
& Reynolds & \rReynolds & $\Cdz = \rCdzero$ corrected\\
\midrule
\multirow{3}{*}{Power}
& hover & \SI{\rPower}{\watt} & of which \SI{\rPaux}{\watt} auxiliary\\
& session, hovering & \SI{\rEndReserve}{\minute} & to the \SI{12}{\percent} reserve; \SI{\rEndEmpty}{\minute} to empty\\
& session, worst mission & \SI{\rEndTurn}{\minute} & turnaround, to the reserve\\
\midrule
\multirow{4}{*}{Human factors}
& $L_p$ free field \SI{1}{\metre} & \SI{\rSplOne}{\decibel A} & \\
& $L_p$ at the ear, in room & \SI{\rSplEar}{\decibel A} & \cref{prop:room}, $r_c=\SI{\rCritDist}{\metre}$\\
& $L_p$ elsewhere in room & \SI{\rSplFar}{\decibel A} & does not fall with distance\\
& wake velocity & \SI{\rWake}{\metre\per\second} & Beaufort 3\\
\midrule
\multirow{4}{*}{Control}
& $J_{xx},J_{yy},J_{zz}$ & $\rJxx,\ \rJyy,\ \rJzz$ & \si{\kilo\gram\metre\squared}, about the centre of mass\\
& authority roll/pitch/yaw & $\rAlphaRoll/\rAlphaPitch/\rAlphaYaw$ & \si{\radian\per\second\squared}, \cref{eq:authority}\\
& gains & $K_R=\rKR$, $K_\omega=\rKw$ & below the bound \rKRsug\ of \cref{mod:gains}\\
& blade tip to face, static & \SI{\rTipGap}{\metre} & against \SI{0.45}{\metre} required\\
\bottomrule
\end{tabular}
\end{table}

\subsection{Flight performance}

\Cref{tab:flight} reports the design of \cref{tab:design}, with its hover-sized
battery, flown through each mission with the governor of \cref{mod:governor}
active. Two clearances are reported. The envelope gap is the distance from the
eye to a sphere of radius $\mathcal{R}$ about the centre of mass, the
conservative quantity that \cref{prop:governor}(iv) bounds. The rotor
clearance is the distance from the eye to the nearest oriented rotor disk,
which is what a blade could actually touch.

The six motions are not defined elsewhere, so they are stated here. Each is a
smooth trajectory of the person's position and heading over the
\SI{\rTwindow}{\second} window, and the ideal vehicle position follows from it
by holding the display at the standoff in front of the eyes. Standing is
postural sway of a few centimetres; pacing is walking back and forth along one
line while facing the display; sit to stand lowers the eyes by
\SI{0.45}{\metre}; walking a loop and jogging follow a circle of radius
\SI{1.44}{\metre}, jogging adding a vertical bob at the stride frequency; and
turning around reverses the heading at each end of a walk.
\Cref{fig:missions} draws them, together with the acceleration each one asks
of the display.

\begin{figure}[tbp]
\centering
% The six motions of the person, and where they put the display. Grey is the
% person, blue is the ideal vehicle position that holds the display at the
% standoff in front of their eyes, dashed arrows join the two.
\begin{tikzpicture}
\begin{groupplot}[
  group style={group size=3 by 2, horizontal sep=1.55cm, vertical sep=1.55cm},
  paperplot, width=0.325\linewidth, height=0.31\linewidth,
  title style={font=\small, yshift=-2pt},
  every axis plot/.append style={line width=0.9pt},
  xlabel style={font=\footnotesize}, ylabel style={font=\footnotesize},
  legend style={font=\scriptsize},
  topview/.style={axis equal image, xmin=-2.3, xmax=2.3, ymin=-2.3, ymax=2.3,
    xlabel={$x$ [\si{\metre}]}, ylabel={$y$ [\si{\metre}]}},
  timeaxis/.style={xmin=0, xmax=16, xlabel={$t$ [\si{\second}]}},
  person/.style={cG!60, line width=2.4pt},
  standoff/.style={construct, -{Stealth[length=1.3mm]}, each nth point=40,
    quiver={u=\thisrow{dx}, v=\thisrow{dy}}},
  stations/.style={only marks, mark=*, mark size=1.1pt, cG!80!black, each nth point=40}]
% row 1: time histories
\nextgroupplot[title={standing: postural sway}, timeaxis,
  ylabel={sway [\si{\centi\metre}]}, ymin=-4, ymax=4,
  legend style={at={(0.97,0.03)}, anchor=south east}]
  \addplot[cG!80!black] table[x=t, y expr=100*\thisrow{x}] {results/fig/e_human_standing.dat};
  \addplot[cB] table[x=t, y expr=100*\thisrow{y}] {results/fig/e_human_standing.dat};
  \legend{$x$, $y$}
\nextgroupplot[title={pacing}, timeaxis, ylabel={$x$ [\si{\metre}]}, ymin=-2, ymax=3,
  legend style={at={(0.03,0.97)}, anchor=north west}]
  \addplot[person] table[x=t, y=x] {results/fig/e_human_pacing.dat};
  \addplot[cA] table[x=t, y=rx] {results/fig/e_human_pacing.dat};
  \legend{person, vehicle}
\nextgroupplot[title={sit to stand}, timeaxis, ylabel={height change [\si{\metre}]},
  ymin=-0.55, ymax=0.1, legend style={at={(0.03,0.03)}, anchor=south west}]
  \addplot[person] table[x=t, y=z] {results/fig/e_human_sit_stand.dat};
  \addplot[cA, dashed] table[x=t, y=rzrel] {results/fig/e_human_sit_stand.dat};
  \legend{person, vehicle}
% row 2: top views and the acceleration each motion demands
\nextgroupplot[title={turning around}, topview]
  \addplot[person] table[x=x, y=y] {results/fig/e_human_turnaround.dat};
  \addplot[cA] table[x=rx, y=ry] {results/fig/e_human_turnaround.dat};
  \addplot[standoff] table[x=x, y=y] {results/fig/e_human_turnaround.dat};
  \addplot[stations] table[x=x, y=y] {results/fig/e_human_turnaround.dat};
\nextgroupplot[title={walking a loop, jogging}, topview]
  \addplot[person] table[x=x, y=y] {results/fig/e_human_walk_loop.dat};
  \addplot[cA] table[x=rx, y=ry] {results/fig/e_human_walk_loop.dat};
  \addplot[standoff] table[x=x, y=y] {results/fig/e_human_walk_loop.dat};
  \addplot[stations] table[x=x, y=y] {results/fig/e_human_walk_loop.dat};
\nextgroupplot[title={acceleration demanded}, timeaxis, ymode=log,
  ymin=0.003, ymax=100, ytick={0.01,0.1,1,10,100},
  ylabel={$\|\ddt{\vect r}_{\mathrm{raw}}\|$ [\si{\metre\per\second\squared}]}]
  \addplot[cB] table[x=t, y=acc] {results/fig/e_human_turnaround.dat};
  \addplot[cA] table[x=t, y=acc] {results/fig/e_human_jogging.dat};
  \addplot[cC] table[x=t, y=acc] {results/fig/e_human_pacing.dat};
  \addplot[cD, dashed, domain=0:16, samples=2] {2.5};
  \node[lbl, text=cB, anchor=south] at (axis cs:2.4,22) {turning};
  \node[lbl, text=cA, anchor=south west] at (axis cs:6.4,5.5) {jogging};
  \node[lbl, text=cC!70!black, anchor=north] at (axis cs:9.4,0.35) {pacing};
\end{groupplot}
\end{tikzpicture}

\caption{The six motions flown in \cref{tab:flight}, over one window. Grey is
the person, blue the ideal vehicle position $\vect r_{\mathrm{raw}}$ that holds
the display at the standoff in front of their eyes, and dashed arrows join the
two at intervals. Pacing is mirrored by the vehicle at a fixed offset; sit to
stand is followed in height; on the loop the vehicle runs on an outer circle.
Turning around is different in kind: at each reversal the display must swing
through a half circle of radius $d$ about the person (\cref{prop:turn}). The
last panel shows the price. The raw reference asks for an acceleration an
order of magnitude above the governor's cap $a_{\max}$ (dashed) at every reversal,
jogging's vertical bob alone exceeds it at the stride frequency, and ordinary
walking stays below it.}
\label{fig:missions}
\end{figure}

\begin{table}[t]
\centering
\caption{Missions flown on the hover-sized battery. $\varkappa$ is the mission
overhead of \cref{prop:contraction}; ``session'' is how long the
\SI{\rBattery}{\gram} pack lasts at that power to the \SI{12}{\percent}
reserve. ``Track'' is the peak tracking error over the whole window, start-up
included, to be compared with $\varepsilon=\SI{\rTrackMargin}{\milli\metre}$;
lag is the r.m.s.\ distance from the ideal display position,
\cref{eq:twoerrors}.}
\label{tab:flight}
\footnotesize\setlength{\tabcolsep}{4pt}
\begin{tabular}{@{}lrrrrrrrrr@{}}
\toprule
Mission & Power & $\varkappa$ & Session & Track & Lag rms & Screen tilt & Sat.\ & Envelope & Rotor\\
\midrule
\inputrows{results/tab_flight_rows}
\end{tabular}
\end{table}

The overhead $\varkappa$ ranges from \rKappaMin\ standing to \rKappaMax\ on
the turnaround. The peak tracking error stays below $\varepsilon$ on every
mission (worst \SI{\rTrackWorst}{\milli\metre}), so by
\cref{prop:governor}(iv) the envelope guarantee holds on these runs; the
nearest any rotor disk comes to the eye is \SI{\rClearWorst}{\metre}, on the
turnaround. Disturbance rejection is not what limits the machine indoors: a
\SI{3}{\metre\per\second} breeze with \SI{1.5}{\metre\per\second} gusts, far
stronger than any room, produces \SI{\rBreezeLag}{\milli\metre} of lag and
\ang{\rBreezeTilt} of screen tilt while the airframe leans
\ang{\rBreezeLean}, which is the gimbal doing what it exists for.
\Cref{fig:turn} shows the hardest mission, the turnaround, as time histories.

\begin{figure}[htbp]
\centering
\providecommand{\eDsafeCross}{1.06}
\providecommand{\eDtext}{0.43}
\providecommand{\eDfull}{0.65}
\providecommand{\eLogicalAtCross}{784}
\providecommand{\eArmL}{0.3827}
\providecommand{\eRotorR}{0.2460}
\providecommand{\eSpan}{1.257}
\providecommand{\eReach}{0.629}
\providecommand{\eScreenW}{0.345}
\providecommand{\eScreenH}{0.194}
\providecommand{\eRcp}{0.130}
\providecommand{\eEye}{1.500}
\providecommand{\eZrotor}{1.370}
\providecommand{\eZrotorBelow}{1.630}
\providecommand{\eStandoff}{1.20}
\providecommand{\eDsafe}{0.45}
\providecommand{\eDeltaZ}{0.130}
\providecommand{\eDisplayOrder}{13.3 in laptop panel,15.6 in laptop panel,iPad Pro 13 panel,iPad Pro 13 whole,15.6 in lid assembly,21.5 in monitor panel,24 in monitor panel,13.3 in ultrabook,15.6 in laptop whole}
\providecommand{\eOptSpl}{58}
\providecommand{\eOptMass}{6.0}
\providecommand{\eTurnClear}{0.61}
\providecommand{\eTurnTrack}{64}
\providecommand{\eTrackMargin}{100}
\providecommand{\eTurnPmean}{111.9}
\providecommand{\eHoverBat}{288}
\providecommand{\eReservePct}{12}
\providecommand{\eTethPsrc}{109.6}
\providecommand{\eTethArea}{0.228}
\providecommand{\eTethV}{48}
\providecommand{\eTethL}{5}
\providecommand{\eTethDrop}{5}
%
\begin{tikzpicture}
\begin{groupplot}[
  group style={group size=1 by 4, vertical sep=0.22cm,
               xlabels at=edge bottom, xticklabels at=edge bottom},
  paperplot, height=3.2cm, xmin=0, xmax=16,
  xlabel={time $t$ [\si{\second}]},
  ylabel style={align=center, font=\footnotesize}]
\nextgroupplot[ylabel={$e_{\mathrm{track}}$\\ {[}\si{\milli\metre}{]}}, ymin=0, ymax=125]
  \fill[cG!12] (axis cs:0,0) rectangle (axis cs:3,125);
  \addplot[cA] table[x=t, y=track] {results/fig/e_turn.dat};
  \addplot[cD, dashed, domain=0:16, samples=2] {\eTrackMargin};
  \node[lbl, text=cD, anchor=south east] at (axis cs:15.9,\eTrackMargin) {$\varepsilon$};
  \node[lbl, text=cG!80!black, anchor=north west] at (axis cs:0.1,90) {settling};
\nextgroupplot[ylabel={$e_{\mathrm{lag}}$\\ {[}\si{\metre}{]}}, ymin=0]
  \fill[cG!12] (axis cs:0,0) rectangle (axis cs:3,4);
  \addplot[cB] table[x=t, y expr=\thisrow{lag}/1000] {results/fig/e_turn.dat};
\nextgroupplot[ylabel={rotor to eye\\ {[}\si{\metre}{]}}, ymin=0.3, ymax=1.5]
  \fill[cG!12] (axis cs:0,0.3) rectangle (axis cs:3,1.5);
  \addplot[cC!70!black] table[x=t, y=clear] {results/fig/e_turn.dat};
  \addplot[cD, dashed, domain=0:16, samples=2] {\eDsafe};
  \node[lbl, text=cD, anchor=south] at (axis cs:4.6,\eDsafe) {$d_{\mathrm{safe}}$};
\nextgroupplot[ylabel={power\\ {[}\si{\watt}{]}}]
  \fill[cG!12] (axis cs:0,90) rectangle (axis cs:3,140);
  \addplot[cE] table[x=t, y=P] {results/fig/e_turn.dat};
\end{groupplot}
\end{tikzpicture}

\caption{The turnaround flown with the governor. (a) The tracking error
against the governed reference stays below the margin $\varepsilon$, which is
the hypothesis of \cref{prop:governor}(iv) checked on this run. (b) The lag
behind the ideal display position grows to about \SI{2}{\metre} at each
reversal: the governor lets the display fall behind rather than swing it round
the head at the acceleration of \cref{fig:missions}. (c) The distance from the
eye to the nearest rotor disk never approaches $d_{\mathrm{safe}}$. (d) The
electrical power, whose window mean after the settling interval (grey) is the
$\bar P$ of the energy integral; the brief spikes are the rotor speeds
responding to the jerk of the reference at the start and end of each swing.}
\label{fig:turn}
\end{figure}

What the table also shows is that a battery sized for hover is not sized for
the mission: on every moving mission the pack lands below its reserve, at
\SI{7}{\percent} on the turnaround. That is the reason the coupled loop exists.

\FloatBarrier
\subsection{The fully coupled loop, following the user}

The design of \cref{tab:design} was sized for hover and then flown. That is
not \cref{prob:codesign}. The coupled problem sizes the battery around the
motion actually performed: guess $\mbat$, close frame and propulsion around it
by \cref{prop:inner}, rebuild the vehicle including its inertia tensor and
rotor constants, fly the mission with the governor and keep-out, integrate the
electrical power, and ask what battery that mission required, by
\cref{eq:batreq}. \Cref{tab:coupled} gives the result for each motion. Every
row converged in \rCoupItLo\ battery updates by the secant iteration
\cref{eq:secant} followed by two verification flights, \rCoupSimLo\ full
6-DOF simulations in all, and every row's final flight met the reserve, the
peak-power requirement, the rotor clearance and the readability limit.

\begin{table}[t]
\centering
\caption{The fully coupled closure \cref{eq:residual} on the final design, one
row per motion. ``Hover sizing'' is the first-principles closure of
\cref{eq:fpclosure} with the mission assumed to be hover, i.e.\ the design of
\cref{tab:design}. ``Updates/sims'' counts battery updates and total
simulations including the verification flights; ``reserve'' is the state of
charge left at $t_f$ on the final flight; ``rotor'' the minimum eye-to-disk
distance on it; ``feasible'' whether that flight met every path constraint.}
\label{tab:coupled}
\footnotesize\setlength{\tabcolsep}{4pt}
\begin{tabular}{@{}lrrrrrrrrl@{}}
\toprule
Motion & Gross & Battery & vs.\ hover & Mean power & $\varkappa$ & Upd./sims & Reserve & Rotor & Feasible\\
\midrule
hover sizing & \SI{\rGross}{\kilo\gram} & \SI{\rBattery}{\gram} & n/a & \SI{\rPower}{\watt} & n/a & n/a & \SI{12.0}{\percent} & n/a & n/a\\
\midrule
\inputrows{results/tab_coupled_rows}
\end{tabular}
\end{table}

\begin{figure}[tbp]
\centering
\providecommand{\eDsafeCross}{1.06}
\providecommand{\eDtext}{0.43}
\providecommand{\eDfull}{0.65}
\providecommand{\eLogicalAtCross}{784}
\providecommand{\eArmL}{0.3827}
\providecommand{\eRotorR}{0.2460}
\providecommand{\eSpan}{1.257}
\providecommand{\eReach}{0.629}
\providecommand{\eScreenW}{0.345}
\providecommand{\eScreenH}{0.194}
\providecommand{\eRcp}{0.130}
\providecommand{\eEye}{1.500}
\providecommand{\eZrotor}{1.370}
\providecommand{\eZrotorBelow}{1.630}
\providecommand{\eStandoff}{1.20}
\providecommand{\eDsafe}{0.45}
\providecommand{\eDeltaZ}{0.130}
\providecommand{\eDisplayOrder}{13.3 in laptop panel,15.6 in laptop panel,iPad Pro 13 panel,iPad Pro 13 whole,15.6 in lid assembly,21.5 in monitor panel,24 in monitor panel,13.3 in ultrabook,15.6 in laptop whole}
\providecommand{\eOptSpl}{58}
\providecommand{\eOptMass}{6.0}
\providecommand{\eTurnClear}{0.61}
\providecommand{\eTurnTrack}{64}
\providecommand{\eTrackMargin}{100}
\providecommand{\eTurnPmean}{111.9}
\providecommand{\eHoverBat}{288}
\providecommand{\eReservePct}{12}
\providecommand{\eTethPsrc}{109.6}
\providecommand{\eTethArea}{0.228}
\providecommand{\eTethV}{48}
\providecommand{\eTethL}{5}
\providecommand{\eTethDrop}{5}
%
\begin{tikzpicture}
\begin{groupplot}[
  group style={group size=2 by 1, horizontal sep=1.6cm},
  halfplot, height=5.2cm, ybar, /pgf/bar width=13pt,
  symbolic x coords={standing,pacing,loop,jogging,turning},
  xtick=data, enlarge x limits=0.14,
  x tick label style={font=\footnotesize, rotate=25, anchor=north east}]
\nextgroupplot[ylabel={battery the mission needs [\si{\gram}]}, ymin=270, ymax=320,
  title={closed around the flown mission}]
  \addplot[fill=cA!70, draw=cA] table[x=mission, y=mbat] {results/fig/e_coupled.dat};
  \draw[cG, dashed, line width=0.8pt]
    ({rel axis cs:0,0} |- {axis cs:standing,\eHoverBat}) --
    ({rel axis cs:1,0} |- {axis cs:standing,\eHoverBat});
  \node[lbl, text=cG!80!black, anchor=north west] at (rel axis cs:0.02,0.97)
    {dashed: hover sizing, \eHoverBat\ g};
\nextgroupplot[ylabel={charge left with \eHoverBat\ g [\si{\percent}]}, ymin=0, ymax=14,
  title={hover-sized battery, flown}]
  \addplot[fill=cB!65, draw=cB] table[x=mission, y=soc] {results/fig/e_coupled.dat};
  \draw[cD, dashed, line width=0.8pt]
    ({rel axis cs:0,0} |- {axis cs:standing,\eReservePct}) --
    ({rel axis cs:1,0} |- {axis cs:standing,\eReservePct})
    node[pos=0.98, below left, lbl, text=cD] {reserve \eReservePct\,\%};
\end{groupplot}
\end{tikzpicture}

\caption{The coupled closure against the hover sizing. Left: the battery each
motion needs when \cref{eq:residual} is solved around the flown mission; the
dashed line is the \SI{\eHoverBat}{\gram} of the hover sizing. Right: the
charge left at $t_f$ when that hover-sized pack is flown through each motion
instead. Every moving mission lands below the \SI{\eReservePct}{\percent}
reserve. The two panels are the same fact from either side: the right is the
violation, the left is its cure.}
\label{fig:coupled}
\end{figure}

\Cref{fig:coupled} shows the table's battery column against what the hover
sizing would have done. Three observations follow.

\emph{The standing row reproduces the hover sizing} to within
\SI{\rCoupStandBat}{\gram} of battery. This is a consistency check, not a
validation: the simulation and the closure share the rotor model of
\cref{sec:blade}, and agreement between them says that the substitution of
\cref{mod:battery} at $\varkappa=1$ is implemented correctly, nothing more.

\emph{Following a person is cheap; turning with them is not.} Ordinary motion
costs between \SI{\rCoupOrdLo}{\gram} and \SI{\rCoupOrdHi}{\gram} of battery.
Repeated turnarounds cost \SI{\rCoupTurnBat}{\gram} of battery and
\SI{\rCoupTurnGross}{\gram} of gross mass, almost all of it the mission
overhead $\varkappa=\rCoupTurnKappa$ propagated through the fixed point: by
\cref{eq:kappaclosure} the overhead multiplies $\beta$, and here it moves the
design \SI{\rCoupTurnUp}{\percent} further up the closure curve. Sizing the
battery for the most demanding motion is a worst-case operating policy, and
it is the one adopted here. For a defined session with known phases the
energy is the time-weighted sum over the phases, with the peak-power
requirement enforced separately by \cref{eq:batreq}; a session of occasional
turns needs less than \cref{tab:coupled}'s last row, and one of continuous
turning needs exactly it.

\emph{The inertia coupling is weak for this layout.} The cycle \cref{eq:cycle2}
runs through $J$, but the battery sits near the centre of mass, so $J_{yy}$
moves only from \num{\rCoupJlo} to \SI{\rCoupJhi}{\kilo\gram\metre\squared}
between the lightest and heaviest converged designs. The coupling that matters
is the energy one; the dynamic one is present and small.

\paragraph{Dynamic clearance.} The static clearance of \cref{tab:design},
\SI{\rTipGap}{\metre}, is not the clearance flown. During the turnaround the
governed reference is held outside $\varrho = d_{\mathrm{safe}}+\mathcal{R}+
\varepsilon$ of the head by \cref{eq:barrier}, the vehicle tracks it to within
\SI{\rTurnTrack}{\milli\metre}, and the nearest rotor disk passes the eye at
\SI{\rTurnClear}{\metre} (envelope gap \SI{\rTurnEnv}{\metre}). An earlier
version of the governor, with a keep-out on position only and a velocity test
that ignored the head's own motion, let the vehicle inside the safe distance
by about \SI{4}{\centi\metre} on this mission; the margin $\varepsilon$ and
the moving-boundary barrier are what closed that gap, and
\cref{prop:governor}(iv) states exactly the hypothesis under which they do.

\FloatBarrier
\subsection{The tether}

Taking the battery off the aircraft removes the mission time from the closure.
What replaces it is a conductor whose mass follows from Ohm's law: at bus
voltage $V$ delivering power $P$ over length $L$ with fractional drop
$\delta$,
\begin{equation}
  I = \frac{P}{V},
  \quad
  R_{\mathrm{cond}} = \frac{2\rho_e L}{A_c},
  \quad
  IR_{\mathrm{cond}}\le\delta V
  \;\Longrightarrow\;
  A_c \ge \frac{2\rho_e L P}{\delta V^2},
  \quad
  m_c = 2\rho_m L A_c ,
  \label{eq:tether}
\end{equation}
with $A_c$ taken as the largest of the resistive requirement, the continuous
current density limit $I/J_{\max}$ with $J_{\max}=\SI{\rTethJmax}{\ampere\per\milli\metre\squared}$,
and a handling floor of \SI{\rTethAreaMin}{\milli\metre\squared}.
\Cref{fig:tether} shows how the three requirements trade over power.

\begin{figure}[tbp]
\centering
\providecommand{\eDsafeCross}{1.06}
\providecommand{\eDtext}{0.43}
\providecommand{\eDfull}{0.65}
\providecommand{\eLogicalAtCross}{784}
\providecommand{\eArmL}{0.3827}
\providecommand{\eRotorR}{0.2460}
\providecommand{\eSpan}{1.257}
\providecommand{\eReach}{0.629}
\providecommand{\eScreenW}{0.345}
\providecommand{\eScreenH}{0.194}
\providecommand{\eRcp}{0.130}
\providecommand{\eEye}{1.500}
\providecommand{\eZrotor}{1.370}
\providecommand{\eZrotorBelow}{1.630}
\providecommand{\eStandoff}{1.20}
\providecommand{\eDsafe}{0.45}
\providecommand{\eDeltaZ}{0.130}
\providecommand{\eDisplayOrder}{13.3 in laptop panel,15.6 in laptop panel,iPad Pro 13 panel,iPad Pro 13 whole,15.6 in lid assembly,21.5 in monitor panel,24 in monitor panel,13.3 in ultrabook,15.6 in laptop whole}
\providecommand{\eOptSpl}{58}
\providecommand{\eOptMass}{6.0}
\providecommand{\eTurnClear}{0.61}
\providecommand{\eTurnTrack}{64}
\providecommand{\eTrackMargin}{100}
\providecommand{\eTurnPmean}{111.9}
\providecommand{\eHoverBat}{288}
\providecommand{\eReservePct}{12}
\providecommand{\eTethPsrc}{109.6}
\providecommand{\eTethArea}{0.228}
\providecommand{\eTethV}{48}
\providecommand{\eTethL}{5}
\providecommand{\eTethDrop}{5}
%
\begin{tikzpicture}
\begin{axis}[paperplot, height=5.4cm, xmin=0, xmax=400, ymin=0, ymax=0.9,
  xlabel={power drawn at the source $P$ [\si{\watt}]},
  ylabel={conductor area per leg [\si{\milli\metre\squared}]},
  legend style={at={(0.02,0.98)}, anchor=north west}]
  \addplot[cA] table[x=P, y=drop] {results/fig/e_tether.dat};
  \addplot[cB] table[x=P, y=amp] {results/fig/e_tether.dat};
  \addplot[cG, dashed] table[x=P, y=hand] {results/fig/e_tether.dat};
  \addplot[black, line width=2.2pt, opacity=0.25] table[x=P, y=req] {results/fig/e_tether.dat};
  \addplot[only marks, mark=*, mark size=2.2pt, cD]
    coordinates {(\eTethPsrc,\eTethArea)};
  \draw[cD, line width=0.5pt] (axis cs:\eTethPsrc,\eTethArea) -- (axis cs:190,0.05);
  \node[lbl, text=cD, anchor=west, fill=white, inner sep=1.5pt] at (axis cs:190,0.05)
    {design: \SI{\eTethArea}{\milli\metre\squared} at \SI{\eTethPsrc}{\watt}};
  \legend{voltage drop $\le\SI{\eTethDrop}{\percent}$,
          current density $\le J_{\max}$,
          handling floor,
          required: the largest}
\end{axis}
\end{tikzpicture}

\caption{Conductor area per leg for a \SI{\eTethL}{\metre} lead at
\SI{\eTethV}{\volt}, against the power drawn at the source. Each requirement of
\cref{eq:tether} is linear in $P$ or constant, and the conductor is sized by
their maximum (grey band): the handling floor at low power, the current
density limit above about \SI{60}{\watt}. At this length and voltage the
voltage-drop requirement, the only one that grows with $L^2$, never binds.
The design point lies on the ampacity line, so the carried conductor mass is
proportional to power, which is the $m^{3/2}$ term that keeps the fold.}
\label{fig:tether}
\end{figure}

\begin{observation}
\Cref{eq:tether} contains no mission time, so endurance ceases to be a design
variable and becomes an operating choice. It does not delete the fold. A
carried conductor sized by resistance has mass proportional to $P\propto
m^{3/2}$, which is the same term as the battery's with a coefficient smaller by
orders of magnitude at these lengths and voltages: the fold moves far away, it
does not vanish. The $m^{3/2}$ term is absent only when the carried cable mass
stops depending on power, because the tether is supported from above or
because a gauge floor binds. For the case below the lumped closure is
gauge-limited (\rTethGaugeLimited), which is why its $\beta_{\mathrm{tether}}$
is zero.
\end{observation}

Redesigning the aircraft for a \SI{\rTethLen}{\metre} lead at
\SI{\rTethV}{\volt}, with the same component models as the battery closure,
gives \SI{\rTethGross}{\kilo\gram} at \SI{\rTethPower}{\watt} and
\SI{\rTethEar}{\decibel A} at the ear. This is not the battery design with the
battery removed: propulsion and structure are resized around the lighter
aircraft, and it still carries a \SI{\rTethReserve}{\gram} landing reserve,
sized by the peak power it must deliver (\rTethResLimit-limited) rather than by
energy. The current is \SI{\rTethCurrent}{\ampere}; the conductor is set by
\rTethSizedBy\ at \SI{\rTethArea}{\milli\metre\squared} per leg and
\SI{\rTethMass}{\gram} for the lead. \SI{\rTethV}{\volt} remains below the
\SI{60}{\volt} DC safety-extra-low-voltage threshold, and the design
\rTethFeasible\ the convergence and current-density checks the engine applies.

\section{Limitations}
\label{sec:limitations}

Every result in \cref{sec:results} is the output of a model. This section
states where the model is known to be inadequate, ordered by how much it would
change the conclusions.

\subsection{Recirculation}

\Cref{mod:disk} assumes quiescent inflow. By the exchange-time calculation of
\cref{eq:exchange} the machine passes a room's volume of air through its disks
every \SI{21}{\second}, so within a minute of a thirty-minute session the
assumption is false. The correction is not a coefficient; it requires solving
the room flow with a moving boundary. Two consequences are unquantified: the
degradation of rotor performance in its own recirculated wake, and the
persistent room-scale draught of order \SI{0.1}{\metre\per\second}. This is the
largest open item.

\subsection{Rotor aerodynamics and the flight power model}

The rotor model is momentum theory plus a single profile term plus three
near-field corrections. There is no blade element momentum theory, so the
radial distribution of inflow and loading is not resolved and the effects of
twist and taper enter only through lumped constants. There is no wake model, no
blade-vortex interaction, and no unsteady inflow: by \cref{as:quasisteady} the
inflow responds instantaneously, whereas its true time constant is of order
$R/\vind\approx\SI{0.1}{\second}$, comparable with the attitude loop. The
descent regime is excluded entirely, as noted after \cref{eq:axial}.

The flight simulation compounds this. Its power law \cref{eq:edot} is
calibrated at hover and has no inflow term, so it returns the same power for
a rotor at a given speed whether the vehicle is hovering or translating. The
axial correction derived in \cref{sec:momentum} is therefore not in the
mission energies of \cref{tab:flight,tab:coupled}. Those tables close the
energy balance around a near-hover approximation of the power; that is a
computational coupling, and it is exercised honestly, but it is not a
validation of manoeuvring energy. The figures are appropriate for indoor
following at walking speed and vertical rates of a few tenths of a metre per
second, and are not to be read as predictions for take-off, descent or
aggressive flight. Coupling an axial and oblique inflow model consistently to
both thrust and power is the next step, and it must be done to both at once:
a climb term in the power alone, against a hover-calibrated thrust, would be
inconsistent.

\subsection{Acoustics}

Three separate weaknesses. The source term \cref{eq:splmodel} is a broadband
scaling law fitted to four aircraft of one manufacturer; the exponents are
theoretically motivated but the anchor is empirical and the aircraft are all
conventional two-blade designs, so extrapolation to $\sigma=\rSolidity$
paddle blades is an extrapolation. The room model \cref{eq:roomlevel} uses a
single frequency-averaged absorption coefficient and assumes a diffuse field,
which fails below the Schroeder frequency \cref{eq:schroeder}, i.e.\ around
\SI{200}{\hertz}, which is exactly where the first blade harmonics lie. And
tonal content is computed but not summed into the reported level, so a design
with a strong \SI{\rBpf}{\hertz} fundamental and audible beating between rotors
could be unacceptable at a level the model calls fine. A measured prototype is
the only way to settle this.

\subsection{Uncertainty}

Every result is a single point estimate. The quantities that carry the most
uncertainty are the figure of merit of a low-Reynolds paddle rotor, the motor
mass constant, the screen drag coefficients, the electrical efficiency, the
cell specific energy at the discharge rate used, and above all the acoustic
anchor, whose calibration spread alone is \SI{1.4}{\decibel}. No propagation
of these uncertainties through the closure was performed, so the paper does
not show that the design ranking of \cref{tab:displays,tab:clearance} is
robust to them, nor that the candidate satisfies its constraints across their
plausible ranges. The margins reported are margins of a nominal model. A
proper treatment samples the parameter ranges, re-closes the design at each
sample, and reports the fraction of samples for which each constraint holds.

\subsection{Structures}

\Cref{mod:arm} covers bending stress, tip deflection and the first bending
mode of an idealised cantilever. It does not cover buckling, joints, bonded
interfaces, fatigue, or crash loads. For a machine intended to operate beside a
person, crash and containment analysis is not optional in practice, and none is
performed here. Guard rings are modelled only as a mass; their ability to
actually arrest a blade is asserted, not computed.

\subsection{Failure modes}

The controller of \cref{sec:control} has no failure mode: the allocation
\cref{eq:allocprob} assumes every rotor is available, and after the loss of one
the quadrotor described here falls. That is a statement about this controller
and not about quadrotors. Mueller and D'Andrea~\cite{mueller2014} show that a
quadrotor can be controlled after the loss of one, two or even three
propellers in a relaxed-hover mode in which the vehicle spins continuously
about a tilted axis. That capability is not implemented here, and a
continuously spinning \SI{1.7}{\kilo\gram} machine at arm's length from a
person is not thereby shown to be acceptable.

A hexacopter is not the answer either, or not simply. With the engine's
alternating-spin regular layout, removing one rotor and its opposite leaves an
allocation matrix of rank \rHexRank\ (\cref{sec:verification}): the four
remaining rotors cannot produce thrust, roll, pitch and yaw independently, so
the standard allocation loses a wrench direction, and lift arithmetic alone
would call for $\lambda\ge1.5$ to hover on four rotors of six, not the
$\lambda\ge2$ an earlier draft stated. Re-optimised under the same limits as
the quadrotor, the hexacopter (six rotors of \SI{\rHexD}{\milli\metre}) comes
out at \SI{\rHexGross}{\kilo\gram}, \SI{\rHexPower}{\watt} and
\SI{\rHexEar}{\decibel A} at the ear against \SI{\rQuadGross}{\kilo\gram},
\SI{\rQuadPower}{\watt} and \SI{\rQuadEar}{\decibel A} for the quadrotor
(\cref{tab:hex}), and its authority is lower on every axis:
$\rHexRoll/\rHexPitch/\rHexYaw$ against
$\rQuadRoll/\rQuadPitch/\rQuadYaw\ \si{\radian\per\second\squared}$. Its
weakest axis in angular acceleration is roll, not pitch: the roll torque
authority is the smaller of the two (\cref{tab:hex}, last column), because
with two rotors on the pitch axis the yaw constraint couples into the roll
differential, and the display slab makes $J_{xx}$ exceed $J_{yy}$. The row sums
$\sum_i|\mat M_{j+1,i}|$, which an earlier draft used and which predict pitch
as the weak axis, mislead here (\cref{app:proofs}). The failure case therefore
needs a different answer than more rotors, and none is offered here.

\begin{table}[t]
\centering
\caption{Quadrotor and hexacopter re-optimised under the same limits as
\cref{tab:displays}. Authorities by \cref{eq:authority}; the last column is
the ratio of roll to pitch torque authority.}
\label{tab:hex}
\footnotesize\setlength{\tabcolsep}{4pt}
\begin{tabular}{@{}lrrrrrrrrr@{}}
\toprule
Layout & $D$ & Gross & Power & At ear & Span & $\alpha_{\mathrm{roll}}$ & $\alpha_{\mathrm{pitch}}$ & $\alpha_{\mathrm{yaw}}$ & $\tau_{\mathrm{roll}}/\tau_{\mathrm{pitch}}$\\
\midrule
\inputrows{results/tab_hex_rows}
\end{tabular}
\end{table}

\subsection{Dynamic clearance}

\Cref{prop:governor} is a conditional guarantee, and each of its conditions
is a limitation. It is enforced at update instants on the governed reference,
not in continuous time and not on the vehicle. It transfers to the vehicle only
through the hypothesis $\sup_t\|\vect r-\vect r_g\|\le\varepsilon$, which
\cref{tab:flight} checks for the missions flown and which nothing proves for
missions not flown. It lapses whenever the caps and the keep-out cannot be
satisfied together, which the engine reports as a violation rather than
hides; on the reported runs this did not occur, but only because a mission
whose ideal position lies inside the sphere is started on the sphere rather
than inside it, a choice that removes the conflict from the start-up and not
from the product. And it protects a sphere of radius $\mathcal{R}$ about the
centre of mass; the rotor clearance reported alongside it is the physical
quantity, and the difference between the two is the conservatism of the
envelope, not a margin the design has earned.

\subsection{Estimation and autonomy}

\Cref{as:estimation} gives the controller exact knowledge of the state and of
the user's position. Real vision-based human tracking has latency, dropouts,
and occlusion, and the reference governor of \cref{mod:governor} is precisely
where those errors would surface: a keep-out sphere centred on an estimated
head position is only as safe as the estimate, and a latency of
\SI{100}{\milli\second} at a walking speed of \SI{1.4}{\metre\per\second} is
\SI{14}{\centi\metre} of sphere in the wrong place, more than the whole
tracking margin $\varepsilon$. No estimator is modelled and no sensor noise is
injected.

\subsection{Mission energy and the operating policy}

The windowed extrapolation $E_{\mathrm{mission}}=\bar P t_f$ of
\cref{sec:numerics} is exact only for a mission periodic within the window. For
a session with distinct phases it is an approximation of unquantified error,
and the battery of \cref{tab:coupled} is sized for the worst single phase
repeated for the whole session, which is a policy rather than a law. Reported
endurance figures also assume constant $\eb$; real cells lose capacity at high
discharge and low temperature, and the internal-resistance voltage sag that
would reduce available thrust at low state of charge is not modelled.

\subsection{Numerics}

The convergence check that matters for a sampled-data simulation, holding the
controller rate fixed and refining the integration substeps, was not run
(\cref{sec:numerics}). The differential evolution of \cref{eq:outerprog} is
stochastic and offers no optimality certificate: the reported designs are the
best found, not the best that exists, and the exterior penalty of
\cref{eq:penalty} means a returned point may violate a constraint slightly,
which \cref{tab:displays} reports as such. Runs are seeded for
reproducibility, which makes results repeatable without making them global.

\subsection{Scope}

Nothing was built. The paper computes what a set of models implies, and the
models were calibrated against published data for aircraft that are not this
aircraft. The most valuable next step is not a better model but a measured
rotor: a single \SI{\rDiam}{\milli\metre} two-blade paddle propeller on a
thrust stand in an anechoic space would settle \cref{tab:acoustic},
\cref{prop:fmblade} and \cref{eq:recorrection} simultaneously, and those three
carry most of the uncertainty in the result.


\appendix
\section{Nomenclature}
\label{app:nomenclature}

Dimensionless quantities carry the unit $1$.

\begin{longtable}{@{}llp{7.4cm}@{}}
\caption{Symbols.}\label{tab:symbols}\\
\toprule
Symbol & Unit & Meaning \\
\midrule
\endfirsthead
\toprule
Symbol & Unit & Meaning \\
\midrule
\endhead
\multicolumn{3}{@{}l}{\emph{Geometry and mass}}\\
$m$              & \si{\kilo\gram} & gross mass\\
$\mfix$          & \si{\kilo\gram} & fixed payload and avionics\\
$\mstr$          & \si{\kilo\gram} & structural mass\\
$\mprop$         & \si{\kilo\gram} & motors, speed controllers, propellers\\
$\mbat$          & \si{\kilo\gram} & battery pack\\
$S(m)$           & \si{\kilo\gram} & $\mstr(m)+\mprop(m)$, the mass that scales with $m$\\
$N$              & 1               & number of rotors\\
$N_b$            & 1               & blades per rotor\\
$D,R$            & \si{\metre}     & rotor diameter, radius\\
$A$              & \si{\metre\squared} & total actuator disk area, $N\pi D^2/4$\\
$L$              & \si{\metre}     & arm length, hub to rotor centre\\
$\mathcal{R}$    & \si{\metre}     & reach, $L+R$, centre to nearest blade tip\\
$\ell$           & \si{\metre}     & distance from the user's eye to the nearest rotor disk\\
$c$              & \si{\metre}     & blade chord\\
$\sigma$         & 1               & solidity, $N_b c/(\pi R)$\\
$\mathrm{AR}$    & 1               & blade aspect ratio, $R/c$\\
$J$              & \si{\kilo\gram\metre\squared} & inertia tensor about the centre of mass\\
$I$              & \si{\metre\tothe{4}} & second moment of area\\
\midrule
\multicolumn{3}{@{}l}{\emph{Aerodynamics}}\\
$T$              & \si{\newton}    & total thrust\\
$\vind$          & \si{\metre\per\second} & induced velocity at the disk\\
$w$              & \si{\metre\per\second} & far-wake velocity, $2\vind$\\
$\mathrm{DL}$    & \si{\pascal}    & disk loading, $T/A$\\
$\Vtip$          & \si{\metre\per\second} & tip speed, $\Omega R$\\
$\Ct$            & 1               & thrust coefficient\\
$\Ct/\sigma$     & 1               & blade loading coefficient\\
$\Cdz$           & 1               & blade profile drag coefficient\\
$\kappa$         & 1               & induced power factor\\
$\kappa_{\mathrm{int}}$ & 1        & rotor interference factor\\
$\FM$            & 1               & figure of merit\\
$\rho$           & \si{\kilo\gram\per\metre\cubed} & air density\\
$C_{D,n},C_{D,t}$& 1               & screen drag coefficients, normal and tangential\\
\midrule
\multicolumn{3}{@{}l}{\emph{Power and energy}}\\
$\Pideal$        & \si{\watt}      & ideal induced power, $T\vind$\\
$\Pind,\Pprof$   & \si{\watt}      & induced and profile power\\
$\Pshaft,\Pelec$ & \si{\watt}      & shaft and electrical power\\
$\Paux$          & \si{\watt}      & non-propulsive electrical load\\
$\eta$           & 1               & motor times speed-controller efficiency\\
$\eb$            & \si{\joule\per\kilo\gram} & battery specific energy available to the mission (\cref{mod:battery})\\
$p_b$            & \si{\watt\per\kilo\gram}  & usable battery specific power\\
$E$              & \si{\joule}     & remaining battery energy\\
$t_f$            & \si{\second}    & mission duration\\
$\varsigma$      & 1               & state-of-charge reserve at landing\\
$\Theta_p(m)$    & \si{\second}    & longest mission a mass $m$ closes, \cref{eq:theta}\\
$\hat m$         & \si{\kilo\gram} & mass at which $\Theta_p$ is largest\\
$\vartheta$      & 1               & propulsive share of total power at $\hat m$, \cref{eq:wallslope}\\
\midrule
\multicolumn{3}{@{}l}{\emph{Closure}}\\
$\alpha$         & 1               & thrust coupling, $1-f_s-\gamma_m$\\
$\beta$          & \si{\kilo\gram^{-1/2}} & energy coupling\\
$\Mzero$         & \si{\kilo\gram} & effective fixed load\\
$\Mmax$          & \si{\kilo\gram} & critical fixed load at the fold\\
$m^*$            & \si{\kilo\gram} & fold mass\\
$f_s$            & 1               & structural mass fraction\\
$\gamma_m$       & 1               & propulsion mass fraction, $\lambda g/S_m$\\
$\gamma_b$       & 1               & battery mass fraction\\
$\lambda$        & 1               & thrust margin, $T_{\max}/mg$\\
$S_m$            & \si{\newton\per\kilo\gram} & propulsion specific thrust\\
$p$              & 1               & disk-area scaling exponent, $A\sim m^p$\\
$q$              & 1               & power exponent, $(3-p)/2$\\
$\mathcal{F}$    & \si{\kilo\gram} & component closure residual, \cref{eq:fpclosure}\\
$\mathcal{G},\mathcal{H}$ & \si{\kilo\gram} & inner and outer maps of the co-design fixed point\\
\midrule
\multicolumn{3}{@{}l}{\emph{Dynamics and control}}\\
$\vect r,\vect v$& \si{\metre},\si{\metre\per\second} & position, velocity\\
$R,q$            & 1               & attitude as rotation matrix, quaternion\\
$\vect\omega$    & \si{\radian\per\second} & body angular velocity\\
$\vect\omega_d,\dt{\vect\omega}_d$ & \si{\radian\per\second},\si{\radian\per\second\squared} & desired body rate and its derivative\\
$\vect j_d$      & \si{\metre\per\second\cubed} & reference jerk\\
$\vect h_\omega$ & \si{\radian\per\second} & rate of the desired thrust direction, \cref{eq:omegad}\\
$\Omega_i$       & \si{\radian\per\second} & speed of rotor $i$\\
$k_T,k_Q,k_P$    & 1               & rotor thrust, torque, power coefficients\\
$\mat M$         & 1               & allocation matrix\\
$\Delta$         & \si{\radian\squared\per\second\squared} & balanced squared-speed differential, \cref{eq:authority}\\
$\tau_m$         & \si{\second}    & rotor speed time constant\\
$\vect e_R,\vect e_\omega$ & 1     & attitude and rate errors\\
$K_R,K_\omega$   & \si{\per\second\squared},\si{\per\second} & attitude gains\\
$K_p,K_d,K_i$    & 1               & position gains\\
$\tau_j^{\max},\alpha_j^{\max}$ & \si{\newton\metre},\si{\radian\per\second\squared} & torque and angular acceleration authority\\
$\vect r_g,\vect v_g$ & \si{\metre},\si{\metre\per\second} & governed reference\\
$\vect c$        & \si{\metre}     & user's head position\\
$\varrho$        & \si{\metre}     & keep-out radius, $d_{\mathrm{safe}}+\mathcal{R}+\varepsilon$\\
$\varepsilon$    & \si{\metre}     & tracking margin inside the keep-out radius\\
$e_{\mathrm{track}},e_{\mathrm{lag}}$ & \si{\metre} & error against the governed and the raw reference\\
$\varkappa$      & 1               & mission power overhead, $\bar P/P_{\mathrm{hover}}$\\
\midrule
\multicolumn{3}{@{}l}{\emph{Acoustics}}\\
$L_p,L_W$        & \si{\decibel}   & sound pressure, sound power level\\
$C$              & \si{\decibel}   & acoustic anchor\\
$Q$              & 1               & directivity factor\\
$\mathcal{R}_{\mathrm{room}}$ & \si{\metre\squared} & room constant\\
$\bar\alpha$     & 1               & average absorption coefficient\\
$r_c$            & \si{\metre}     & critical distance\\
$f_{\mathrm{BPF}}$ & \si{\hertz}   & blade passage frequency\\
\bottomrule
\end{longtable}

\section{Supplementary derivations}
\label{app:proofs}

\subsection{The cubic form of the closure}

At $q=3/2$, substituting $u=\sqrt m>0$ into \cref{eq:closure} gives
\begin{equation}
  \beta u^3 - \alpha u^2 + \Mzero = 0 .
  \label{eq:cubic}
\end{equation}

\begin{proposition}
\Cref{eq:cubic} has exactly one negative real root for all
$\alpha,\beta,\Mzero>0$. Its two remaining roots are real and positive if and
only if $\Mzero\le\Mmax$, in agreement with \cref{thm:trichotomy}.
\end{proposition}

\begin{proof}
Write $\Pi(u) = \beta u^3-\alpha u^2+\Mzero$. By Vieta, the product of the
three roots is $-\Mzero/\beta<0$, so either one or three roots are negative.
Since $\Pi(0)=\Mzero>0$ and $\Pi(u)\to+\infty$ as $u\to+\infty$, and
$\Pi'(u)=3\beta u^2-2\alpha u = u(3\beta u - 2\alpha)$ has roots $u=0$ and
$u=2\alpha/(3\beta)$, the function decreases on $(0,2\alpha/(3\beta))$ and
increases thereafter, so it has at most two positive roots. Hence exactly one
root is negative. The local minimum value is
\[
  \Pi\!\left(\frac{2\alpha}{3\beta}\right)
  = \beta\frac{8\alpha^3}{27\beta^3} - \alpha\frac{4\alpha^2}{9\beta^2} + \Mzero
  = \frac{8\alpha^3-12\alpha^3}{27\beta^2} + \Mzero
  = \Mzero - \frac{4\alpha^3}{27\beta^2}
  = \Mzero - \Mmax ,
\]
using \cref{eq:closedform}. Two positive real roots exist iff this is
non-positive, i.e.\ iff $\Mzero\le\Mmax$. Note also that
$u^* = 2\alpha/(3\beta)$ satisfies $(u^*)^2 = 4\alpha^2/(9\beta^2) = m^*$, as
required for consistency.
\end{proof}

\Cref{fig:cubic} draws $\Pi$ in the three cases.

\begin{figure}[tbp]
\centering
% Pi(u) = beta u^3 - alpha u^2 + M0 for the worked example, u = sqrt(m).
% u* = 2 alpha / (3 beta) = 1.8688; Pi(u*) = M0 - M0max.
% Roots: M0 = 0.45: -0.7471, 1.0689, 2.4813; M0 = M0max: -0.9344, 1.8688 (double);
% M0 = 0.95: -1.0429 and a complex pair.
\begin{tikzpicture}
\begin{axis}[paperplot, height=6.2cm,
  xlabel={$u=\sqrt m$ [\si{\kilo\gram^{1/2}}]},
  ylabel={$\Pi(u)=\beta u^3-\alpha u^2+\Mzero$},
  xmin=-1.25, xmax=2.95, ymin=-0.52, ymax=1.25, domain=-1.25:2.95,
  samples=200, legend style={at={(0.5,1.03)}, anchor=south}, legend columns=3]
  \fill[cG!10] (axis cs:-1.25,-0.52) rectangle (axis cs:0,1.25);
  \node[lbl, cG] at (axis cs:-0.62,1.12) {$u<0$: not a mass};
  \draw[black, line width=0.5pt] (axis cs:-1.25,0) -- (axis cs:2.95,0);
  \addplot[cC] {0.2271*x^3 - 0.6366*x^2 + 0.45};
  \addlegendentry{$\Mzero=0.45<\Mmax$}
  \addplot[cB] {0.2271*x^3 - 0.6366*x^2 + 0.7411};
  \addlegendentry{$\Mzero=\Mmax$}
  \addplot[cD] {0.2271*x^3 - 0.6366*x^2 + 0.95};
  \addlegendentry{$\Mzero=0.95>\Mmax$}
  \draw[construct] (axis cs:1.8688,-0.52) -- (axis cs:1.8688,1.25);
  \node[lbl, anchor=north west] at (axis cs:1.89,1.2) {$u^*=\tfrac{2\alpha}{3\beta}$};
  \addplot[only marks, mark=*, mark size=1.7pt, cC, forget plot]
    coordinates {(-0.7471,0) (1.0689,0) (2.4813,0)};
  \addplot[only marks, mark=*, mark size=1.9pt, cB, forget plot]
    coordinates {(-0.9344,0) (1.8688,0)};
  \addplot[only marks, mark=*, mark size=1.7pt, cD, forget plot]
    coordinates {(-1.0429,0)};
  \addplot[only marks, mark=square*, mark size=1.6pt, cC, forget plot]
    coordinates {(1.8688,-0.2911)};
  \addplot[only marks, mark=square*, mark size=1.6pt, cD, forget plot]
    coordinates {(1.8688,0.2089)};
  \node[lbl, anchor=north, cC!70!black] at (axis cs:1.8688,-0.32)
    {$\Pi(u^*)=\Mzero-\Mmax<0$};
  \node[lbl, anchor=south, cD] at (axis cs:1.8688,0.24) {$\Pi(u^*)>0$};
  \node[lbl, anchor=south east] at (axis cs:1.06,0.02) {$\sqrt{m_-}$};
  \node[lbl, anchor=south west] at (axis cs:2.49,0.02) {$\sqrt{m_+}$};
  \node[lbl, anchor=north] at (axis cs:-0.75,-0.04) {$u_3<0$};
\end{axis}
\end{tikzpicture}

\caption{The cubic \cref{eq:cubic} in $u=\sqrt m$ for the worked example of
\cref{sec:closure}. $\Pi(0)=\Mzero>0$ is a local maximum and $\Pi$ has its
local minimum at $u^*=2\alpha/(3\beta)$, where $(u^*)^2=m^*$ and
$\Pi(u^*)=\Mzero-\Mmax$. One root is always negative and is not a mass. The
two positive roots $\sqrt{m_\pm}$ exist exactly when the minimum is not above
the axis, and they merge at $u^*$ when $\Mzero=\Mmax$.}
\label{fig:cubic}
\end{figure}

\subsection{Derivative of the closure solution with respect to endurance}

Differentiating $F(m(t_f))=\Mzero(t_f)$ implicitly, with
$\beta = \beta_0 t_f$ and $\Mzero = \mfix + \Paux t_f/\eb$,
\begin{equation}
  F'(m)\,\frac{\dd m}{\dd t_f}
  = \frac{\partial \Mzero}{\partial t_f} + m^{3/2}\frac{\partial\beta}{\partial t_f}
  = \frac{\Paux}{\eb} + \beta_0 m^{3/2} ,
\end{equation}
so
\begin{equation}
  \frac{\dd m}{\dd t_f}
  = \frac{\Paux/\eb + \beta_0 m^{3/2}}{\alpha - \tfrac32\beta\sqrt m} .
  \label{eq:dmdt}
\end{equation}
The denominator is $F'(m)$, which vanishes at the fold; \cref{eq:dmdt} therefore
diverges there, quantifying the remark following \cref{thm:bifurcation}. Both
terms in the numerator are positive. The denominator is positive on the light
branch and negative on the heavy branch (\cref{lem:concave}), so the light root
increases with endurance and the heavy root decreases: as $t_f$ grows the two
branches approach each other and merge at the fold.

\subsection{Rotor speed at maximum thrust}

For a fixed-pitch rotor at fixed collective, $T=k_T\Omega^2$, so the speed
required for the sizing thrust $T_{\max}=\lambda mg/N$ relative to hover
$T_h = mg/N$ is
\begin{equation}
  \frac{\Omega_{\max}}{\Omega_h} = \sqrt{\frac{T_{\max}}{T_h}} = \sqrt{\lambda} .
  \label{eq:omegaratio}
\end{equation}
Hence hover occurs at $1/\sqrt\lambda$ of full rotor speed: at $\lambda=1.8$
this is $74.5\%$, which is the figure quoted in \cref{sec:results} and which
also fixes the voltage headroom the pack must provide.

\subsection{Torque authority of the symmetric quadrotor}

Holding total thrust at $mg$ and the other two moments at zero, the largest
moment about body axis $j$ is the value of the linear program
\begin{equation}
  \max\ \tau_j
  \quad\text{s.t.}\quad
  \mat M\vect u = (mg,\ \tau_j\vect e_j),
  \qquad
  0\le u_i\le\Omega_{\max}^2 ,
  \label{eq:authlp}
\end{equation}
and its counterpart for $-\tau_j$; the authority is the smaller of the two.
The engine solves \cref{eq:authlp} directly for every layout. For the $\times$
quadrotor with alternating spins it has a closed form.

Let $a_i = \mat M_{j+1,i}$. For each of the three moment rows $|a_i|=a$ is the
same for every rotor ($k_TL/\sqrt2$ for roll and pitch, $k_Q$ for yaw), and the
sign vector $s_i=\operatorname{sign}a_i$ is orthogonal to the thrust row and to
the other two moment rows. Write $\vect u = \Omega_h^2\vect 1 + \vect d$. The
thrust constraint is $\sum_i d_i=0$, and $\tau_j = a\sum_i s_i d_i = 2aD$ with
$D=\sum_{s_i=+1}d_i=-\sum_{s_i=-1}d_i$. Each of the two increments is at most
$\Omega_{\max}^2-\Omega_h^2$ and each of the two decrements at most $\Omega_h^2$,
so $D\le2\Delta$ with
\begin{equation}
  \Delta = \min\{\Omega_h^2,\ \Omega_{\max}^2-\Omega_h^2\} .
\end{equation}
The bound is attained by $\vect d=\Delta\vect s$, which satisfies every
constraint by the orthogonality, hence
\begin{equation}
  \tau_j^{\max} = 4a\Delta = \Delta\sum_{i=1}^{4}\big|\mat M_{j+1,i}\big| ,
  \label{eq:authderiv}
\end{equation}
which is \cref{eq:authority}. There is no factor $1/N$: the thrust sum is held by
pairing increments with decrements, not by dividing the moment. For
$\lambda\le2$, $\Delta=(\lambda-1)\Omega_h^2$. \Cref{fig:differential} shows
the two cases of the minimum.

\begin{figure}[tbp]
\centering
% The balanced differential that attains the roll authority of the x
% quadrotor. Rotor i sits at psi_i = pi/4 + i pi/2, i = 0..3 (numbered 1..4),
% so the roll row y_i k_T has signs (+,+,-,-). Squared speeds normalised by
% Omega_h^2; Omega_max^2 = lambda Omega_h^2, Delta = min(1, lambda - 1).
\begin{tikzpicture}[font=\small]
  % --- the quadrotor, top view, body x to the right, y up ------------------
  \begin{scope}[shift={(0,0)}, scale=0.82]
    \draw[line width=0.8pt, cG] (-1.1,-1.1) -- (1.1,1.1);
    \draw[line width=0.8pt, cG] (-1.1,1.1) -- (1.1,-1.1);
    \foreach \ang/\n/\s/\c in {45/1/+/cA, 135/2/+/cA, 225/3/-/cD, 315/4/-/cD}{
      \draw[\c, line width=0.9pt, fill=\c!12] (\ang:1.556) circle (0.52);
      \node[font=\footnotesize] at (\ang:1.556) {$\n$};
      \node[\c, font=\normalsize] at (\ang:2.42) {$\s$};
    }
    \draw[vec] (0,0) -- (0.85,0) node[right, lbl] {$\vect b_1$};
    \draw[vec] (0,0) -- (0,0.85) node[above, lbl] {$\vect b_2$};
  \end{scope}
  \node[lbl, align=center] at (0,2.45)
    {$s_i=\operatorname{sign}\mat M_{2,i}=(+,+,-,-)$\\$\vect s\perp$ thrust, pitch, yaw rows};
  \node[lbl, align=center, cA] at (0,-2.3)
    {$y_i>0$ pair up, $y_i<0$ pair down:\\$\tau_1=\sum_i y_iT_i>0$};
  % --- bars: lambda = 1.8 --------------------------------------------------
  \begin{scope}[shift={(3.0,-1.7)}, yscale=1.35, xscale=0.8]
    \draw[->] (-0.3,0) -- (4.3,0);
    \draw[->] (-0.3,0) -- (-0.3,2.85) node[above, lbl] {$u_i/\Omega_h^2$};
    \foreach \y in {0,1,2}{\draw (-0.35,\y) -- (-0.25,\y) node[left=2pt, lbl] {$\y$};}
    \draw[cD, dashed] (-0.3,1.8) -- (4.3,1.8) node[right, lbl, cD] {$\Omega_{\max}^2$};
    \draw[cG, densely dotted] (-0.3,1) -- (4.3,1) node[right, lbl, cG] {$\Omega_h^2$};
    \foreach \x/\h/\c in {0.5/1.8/cA, 1.5/1.8/cA, 2.5/0.2/cD, 3.5/0.2/cD}{
      \fill[\c!55] (\x-0.3,0) rectangle (\x+0.3,\h);
      \draw[\c] (\x-0.3,0) rectangle (\x+0.3,\h);
    }
    \foreach \x/\n in {0.5/1,1.5/2,2.5/3,3.5/4}{\node[below, lbl] at (\x,0) {$\n$};}
    \draw[dim] (1.95,1) -- (1.95,1.8) node[midway, right, lbl, black] {$\Delta$};
    \draw[dim] (3.95,0.2) -- (3.95,1) node[midway, right, lbl, black] {$\Delta$};
    \node[lbl, align=center] at (2,-0.72)
      {$\lambda=1.8$: $\Delta=\Omega_{\max}^2-\Omega_h^2$,\\increments reach the limit};
  \end{scope}
  % --- bars: lambda = 2.5 --------------------------------------------------
  \begin{scope}[shift={(8.6,-1.7)}, yscale=1.35, xscale=0.8]
    \draw[->] (-0.3,0) -- (4.3,0);
    \draw[->] (-0.3,0) -- (-0.3,2.85) node[above, lbl] {$u_i/\Omega_h^2$};
    \foreach \y in {0,1,2}{\draw (-0.35,\y) -- (-0.25,\y) node[left=2pt, lbl] {$\y$};}
    \draw[cD, dashed] (-0.3,2.5) -- (4.3,2.5) node[right, lbl, cD] {$\Omega_{\max}^2$};
    \draw[cG, densely dotted] (-0.3,1) -- (4.3,1) node[right, lbl, cG] {$\Omega_h^2$};
    \foreach \x/\h/\c in {0.5/2/cA, 1.5/2/cA}{
      \fill[\c!55] (\x-0.3,0) rectangle (\x+0.3,\h);
      \draw[\c] (\x-0.3,0) rectangle (\x+0.3,\h);
    }
    \foreach \x in {2.5,3.5}{\draw[cD, line width=1.4pt] (\x-0.3,0) -- (\x+0.3,0);}
    \foreach \x/\n in {0.5/1,1.5/2,2.5/3,3.5/4}{\node[below, lbl] at (\x,0) {$\n$};}
    \draw[dim] (1.95,1) -- (1.95,2) node[midway, right, lbl, black] {$\Delta$};
    \node[lbl, align=center] at (2,-0.72)
      {$\lambda=2.5$: $\Delta=\Omega_h^2$,\\decrements reach zero};
  \end{scope}
\end{tikzpicture}

\caption{The balanced differential that attains the torque authority of the
$\times$ quadrotor, drawn for roll. The sign vector $\vect s$ of the roll row
is orthogonal to the thrust, pitch and yaw rows, so
$\vect u=\Omega_h^2\vect 1+\Delta\vect s$ changes the roll moment and nothing
else. $\Delta$ is the smaller of the room to speed rotors up and the room to
slow them down: at $\lambda=1.8$ the increments reach $\Omega_{\max}^2$ first,
at $\lambda=2.5$ the decrements reach zero first. Either way each of the four
rotors contributes $\Delta$, which is why \cref{eq:authderiv} carries no factor
$1/N$.}
\label{fig:differential}
\end{figure}

The closed form uses the equal-magnitude and orthogonality properties of this
layout and should not be extended blindly. For the hexacopter with
$\psi_0=\pi/6$ the roll row has $\sum_i|\sin\psi_i| = 4$ and the pitch row
$\sum_i|\cos\psi_i|=2\sqrt3$, because two rotors sit on the pitch axis and
contribute nothing to pitch; the magnitudes are unequal and the LP must be
solved. The same holds for any layout with a failed rotor.

\subsection{Sound power to pressure with the reference constants}

The reference intensity implied by $p_0=\SI{20}{\micro\pascal}$ is
$I_0 = p_0^2/(\rho_0 c_0)$. At $\rho_0=\SI{1.225}{\kilo\gram\per\metre\cubed}$
and $c_0=\SI{343}{\metre\per\second}$,
$\rho_0 c_0 = \SI{420}{\pascal\second\per\metre}$ and
$I_0 = (2\times10^{-5})^2/420 = \SI{9.52e-13}{\watt\per\metre\squared}$, against
the nominal $W_0/\SI{1}{\metre\squared} = \SI{1e-12}{\watt\per\metre\squared}$.
The discrepancy is $10\log_{10}(1/0.952)=\SI{0.21}{\decibel}$, which is why
\cref{eq:lwlp} is stated as an equality: the reference constants are chosen to
make it one to within a fifth of a decibel at standard conditions.


\bibliographystyle{plain}
\bibliography{refs}

\end{document}
